\documentclass[11pt]{article}
\usepackage[a4paper,margin=24mm]{geometry}
\usepackage{fontspec}
\usepackage[AutoFallBack=true]{xeCJK}
\setmainfont{Noto Serif}
\setCJKmainfont{Noto Serif CJK SC}
\setCJKfallbackfamilyfont{rm}{SimSun-ExtB}
\xeCJKDeclareCharClass{CJK}{"25A1, "25CB}
\usepackage{xcolor}
\usepackage[hidelinks]{hyperref}
\setlength{\parindent}{0pt}
\setlength{\parskip}{0.4em}
\newcommand{\zhsmall}{\footnotesize\color{black!70}}
\newcommand{\unitlabel}[1]{\marginpar{\raggedleft\tiny\color{gray}#1}}
\newenvironment{unitzh}{\par\CJKfontspec{Noto Serif CJK SC}[AutoFallBack=true]\large\linespread{1.25}\selectfont}{\par}
\newenvironment{uniten}{\par\small\color{black!85}}{\par\medskip}
\sloppy
\begin{document}
\begin{titlepage}\centering
{\fontsize{34}{40}\selectfont 九章算術}\par\vspace{1em}
{\LARGE The Nine Chapters on the Mathematical Art}\par\vspace{0.6em}
{\large with the commentary of Liu Hui (魏 劉徽注) and the annotations of Li Chunfeng and others (唐 李淳風等注釋)}\par\vspace{2em}
{\normalsize Chinese text with modern punctuation, and a complete English translation}\par\vfill
\begin{minipage}{0.86\textwidth}\small
Chinese text: Sibu congkan photo-reprint of the Weiboxie edition (四部叢刊, 景上海涵芬樓藏微波榭刊本), as transcribed on Chinese Wikisource (the pages and revisions used are listed at the end).
Punctuation, division into units and English translation: AI-produced by Claude Opus 5.5 (Anthropic), 24--25 October 2026.
The punctuated Chinese was checked by program against the Wikisource text: every character is reproduced, in order,
and small print (commentary) stays small print. In the translation, commentary is given in square brackets with its
author. The labels in the margin identify the units of the machine-readable file units.jsonl.
\end{minipage}
\end{titlepage}
\tableofcontents\clearpage
\section*{序}\addcontentsline{toc}{section}{序}

\begin{unitzh}\unitlabel{jiuzhang-001.1}劉徽《九章算術注》原序\end{unitzh}
\begin{uniten}Original Preface to Liu Hui's Commentary on the Nine Chapters on the Mathematical Art\end{uniten}

\begin{unitzh}\unitlabel{jiuzhang-001.2}昔在包犧氏始畫八卦，以通神明之德，以𩔖萬物之情，作九九之術，以合六爻之變。曁於黃帝神而化之，引而伸之，於是建歷紀，協律呂，用稽道原，然後兩儀四象精微之氣可得而效焉。記稱隸首作數，其詳未之聞也。按周公制禮而有九數，九數之流，則《九章》是矣。往者暴秦焚書，經術散壞。自時厥後，漢北平侯張蒼、大司農中丞耿壽昌皆以善算命世。蒼等因舊文之遺殘，各稱刪補。故校其目則與古或異，而所論者多近語也。\end{unitzh}
\begin{uniten}Of old, Baoxi (包犧氏) first drew the eight trigrams (bagua), in order to communicate with the virtue of the spirit-like and luminous and to classify the dispositions of the myriad things, and made the method of the nine-nines (jiujiu) [Translator's note: i.e. the multiplication table], in order to match the changes of the six lines (liu yao). When it came to the Yellow Emperor (Huangdi), he transformed it in a spirit-like way and drew it out and extended it. Thereupon he established the calendar and its periods (li ji), harmonized the pitch-pipes (lülü), and by these examined the origin of the Way. After that, the subtle and minute qi of the two modes (liangyi) and the four images (sixiang) could be obtained and imitated. The records state that Lishou (隸首) created numbers, but its details we have not heard. Now, when the Duke of Zhou established the rites, there were the nine numbers (jiushu); the stream of the nine numbers is precisely the Nine Chapters. In the past the tyrannical Qin burned the books, and the classical arts were scattered and ruined. From that time on, Zhang Cang (張蒼), Marquis of Beiping of the Han, and Geng Shouchang (耿壽昌), Assistant to the Grand Minister of Agriculture (da sinong zhongcheng), were both renowned in their age for their skill in calculation. Cang and the others, on the basis of the remnants left of the old texts, each, it is said, made deletions and supplements. Therefore, when one collates its headings, they sometimes differ from the ancient ones, and what is discussed is mostly in recent language.\end{uniten}

\begin{unitzh}\unitlabel{jiuzhang-001.3}徽幼習《九章》，長再詳覽。觀陰陽之割裂，總算術之根源，探𦣱之暇，遂悟其意。是以敢竭頑魯，采其所見，爲之作注。事類相推，各有攸歸，故枝條雖分而同本𠏉者，知發其一端而已。又所析理以辭，解體用圖，庶亦約而能周，通而不黷，覽之者思過半矣。\end{unitzh}
\begin{uniten}I, Hui, studied the Nine Chapters when young, and when grown I read through it again in detail. I observed the division of yin and yang and brought together the roots of the mathematical arts; in leisure moments of probing the abstruse, I then came to understand its meaning. Therefore I have ventured to exhaust my dull and stupid [faculties], gathered what I have perceived, and made a commentary on it. Matters of a kind are inferred from one another, and each has that to which it belongs. Therefore, as for the branches which, though divided, share the same root and trunk, one knows that they issue from a single end, that is all. Moreover, I have analysed the principles with words and dissected the bodies with diagrams. I hope that it may also be concise and yet comprehensive, clear and yet not tedious, and that those who read it will have gone more than half the way in their thinking.\end{uniten}

\begin{unitzh}\unitlabel{jiuzhang-001.4}且算在六藝，古者以賓興賢能，教習國子。雖曰九數，其能窮纖入微，探測無方。至於以法相傳，亦猶規矩度量可得而共，非特難爲也。當今好之者寡，故世雖多通才逹學，而未必能綜於此耳。\end{unitzh}
\begin{uniten}Moreover, calculation is among the six arts (liuyi); in antiquity it was used to treat the worthy and able as guests and promote them, and to teach and train the sons of the state. Although it is called the nine numbers, it can exhaust the fine and enter into the minute, probing and measuring without limit. As for passing it on by means of rules (fa), it is just like the compass and square and the measures of length and capacity, which can be obtained and shared; it is not especially difficult to do. Nowadays those who are fond of it are few. Therefore, although the age has many men of comprehensive talent and penetrating learning, they are not necessarily versed in this.\end{uniten}

\begin{unitzh}\unitlabel{jiuzhang-001.5}《周官》大司徒職：夏至日中立八尺之表，其景尺有五寸，謂之地中。說云：「南戴日下萬五千里。」夫云爾者，以術推之。按《九章》立四表望遠及因木望山之術，皆端旁互見，無有超邈若斯之類。然則蒼等爲術，猶未足以博盡羣數也。\end{unitzh}
\begin{uniten}The duties of the Grand Minister of Education (da situ) in the Offices of Zhou (Zhou guan): at noon on the summer solstice set up a gnomon (biao) of 8 chi (尺); its shadow is 1 chi and 5 cun (寸); this is called the centre of the earth. The explanation says: "The point directly beneath the sun in the south is 15000 li (里) away." Those who say so inferred it by a method. Now, in the Nine Chapters, the methods of setting up four gnomons to sight something distant and of sighting a mountain by means of a tree all have end and side mutually visible; there is none so far-reaching as things of this kind. So then, the methods that Cang and the others made were still not sufficient to cover all the numbers broadly and exhaustively.\end{uniten}

\begin{unitzh}\unitlabel{jiuzhang-001.6}徽㝷九數有重差之名，原其指趣，乃所以施於此也。凡望極高、測絕深而兼知其遠者，必用重差、句股，則必以重差爲率，故曰重差也。\end{unitzh}
\begin{uniten}I, Hui, examining the nine numbers, found that there is the name "double difference" (chongcha); tracing its purport back to its origin, it is precisely what is applied to this. Whenever one sights the extremely high or measures the extremely deep and at the same time gets to know its distance, one must use double differences and base and altitude (gougu). Then one must take the double difference as the rate (lü); therefore it is called "double difference".\end{uniten}

\begin{unitzh}\unitlabel{jiuzhang-001.7}立兩表於洛陽之城，令高八尺，南北各盡平地，同日度其正中之時。以景差爲法，表高乘表閒爲實，實如法而一，所得加表高，卽日去地也。以南表之景乘表閒爲實，實如法而一，卽爲從南表至南戴日下也。以南戴日下及日去地爲句股，爲之求弦，卽日去人也。以徑寸之筩南望日，日滿筩空，則定筩之長短以爲股率，以筩徑爲句率，日去人之數爲大股，大股之句卽日徑也。\end{unitzh}
\begin{uniten}Set up two gnomons in the city of Luoyang, making them 8 chi high, [one] south and [one] north, each on completely level ground, and on the same day measure them at the time of exact noon. Take the difference of the shadows as divisor (fa), and the height of the gnomon multiplied by the distance between the gnomons as dividend (shi). Divide the dividend by the divisor, and add the height of the gnomon to what is obtained: this is the distance of the sun from the earth. Take the shadow of the southern gnomon multiplied by the distance between the gnomons as dividend, and divide the dividend by the divisor: this is [the distance] from the southern gnomon to the point directly beneath the sun in the south. Take [the distance to] the point directly beneath the sun in the south and the distance of the sun from the earth as base and altitude, and for them find the hypotenuse (xian): this is the distance of the sun from the person. With a tube 1 cun in diameter sight the sun towards the south. When the sun fills the hollow of the tube, fix the length of the tube and take it as the altitude rate. Take the diameter of the tube as the base rate, and the number of the distance of the sun from the person as the large altitude. The base of the large altitude is the diameter of the sun.\end{uniten}

\begin{unitzh}\unitlabel{jiuzhang-001.8}雖天圓穹之象猶曰可度，又况泰山之高與\end{unitzh}
\begin{uniten}Even the image of the round vault of heaven can still be said to be measurable; how much more so the height of Mount Tai and [continues]\end{uniten}

\section*{目録}\addcontentsline{toc}{section}{目録}

\begin{unitzh}\unitlabel{jiuzhang-002.1}江海之廣哉！徽以爲今之史籍且畧舉天地之物，考論厥數，載之於志，以闡世術之美。輒造重差，并爲注解，以究古人之意，綴於句股之下。度高者重表，測深者累矩，孤離者三望，離而又旁求者四望。觸類而長之，則雖幽遐詭伏，靡所不入。博物君子，詳而覽焉。\end{unitzh}
\begin{uniten}[continued] … the breadth of the rivers and seas! I, Hui, consider that the historical records of today also briefly cite the things of heaven and earth, examine and discuss their numbers, and set them down in the treatises, so as to make plain the excellence of the arts of the age. I have thereupon ventured to compose the Double Differences (chongcha), and have also made notes and explanations for it, in order to fathom the intent of the ancients, and have attached it below "Base and Altitude" (gougu). For measuring a height, double gnomons (chongbiao); for gauging a depth, stacked set-squares (leiju); for what is isolated and apart, three sightings (sanwang); for what is apart and must moreover be sought from the side, four sightings (siwang). If one extends these by analogy to things of the same kind, then even what is dark and distant, strange and hidden: there is nowhere they do not enter. Gentlemen of broad learning, examine it in detail!\end{uniten}

\begin{unitzh}\unitlabel{jiuzhang-002.2}目錄{\zhsmall 附}\end{unitzh}
\begin{uniten}Table of Contents [Note: appended]\end{uniten}

\begin{unitzh}\unitlabel{jiuzhang-002.3}方田第一{\zhsmall 凡三十八問，補圖二}\end{unitzh}
\begin{uniten}Rectangular Fields (fangtian), Chapter 1 [Note: 38 problems in all; 2 diagrams supplied]\end{uniten}

\begin{unitzh}\unitlabel{jiuzhang-002.4}粟米第二{\zhsmall 凡四十六問}\end{unitzh}
\begin{uniten}Millet and Rice (sumi), Chapter 2 [Note: 46 problems in all]\end{uniten}

\begin{unitzh}\unitlabel{jiuzhang-002.5}衰分第三{\zhsmall 凡二十問}\end{unitzh}
\begin{uniten}Graded Distribution (cuifen), Chapter 3 [Note: 20 problems in all]\end{uniten}

\begin{unitzh}\unitlabel{jiuzhang-002.6}少廣第四{\zhsmall 凡二十四問，補圖一}\end{unitzh}
\begin{uniten}Small Breadth (shaoguang), Chapter 4 [Note: 24 problems in all; 1 diagram supplied]\end{uniten}

\begin{unitzh}\unitlabel{jiuzhang-002.7}商功第五{\zhsmall 凡二十八問}\end{unitzh}
\begin{uniten}Assessing Works (shanggong), Chapter 5 [Note: 28 problems in all]\end{uniten}

\begin{unitzh}\unitlabel{jiuzhang-002.8}均輸第六{\zhsmall 凡二十八問}\end{unitzh}
\begin{uniten}Equitable Transport (junshu), Chapter 6 [Note: 28 problems in all]\end{uniten}

\begin{unitzh}\unitlabel{jiuzhang-002.9}盈不足第七{\zhsmall 凡二十問}\end{unitzh}
\begin{uniten}Excess and Deficit (yingbuzu), Chapter 7 [Note: 20 problems in all]\end{uniten}

\begin{unitzh}\unitlabel{jiuzhang-002.10}方程第八{\zhsmall 凡十八問}\end{unitzh}
\begin{uniten}Rectangular Arrays (fangcheng), Chapter 8 [Note: 18 problems in all]\end{uniten}

\begin{unitzh}\unitlabel{jiuzhang-002.11}句股第九{\zhsmall 凡二十四問，補圖十}\end{unitzh}
\begin{uniten}Base and Altitude (gougu), Chapter 9 [Note: 24 problems in all; 10 diagrams supplied]\end{uniten}

\begin{unitzh}\unitlabel{jiuzhang-002.12}音義第十\end{unitzh}
\begin{uniten}Pronunciations and Meanings (yinyi), Chapter 10\end{uniten}

\section*{卷第一}\addcontentsline{toc}{section}{卷第一}

\begin{unitzh}\unitlabel{jiuzhang-003.1}九章算術卷第一{\zhsmall 算經十書之二} 魏劉徽注 唐朝議大夫行太史令上輕車都尉臣李淳風等奉勅注釋\end{unitzh}
\begin{uniten}The Nine Chapters on the Mathematical Art, Juan 1 [Note: The second of the Ten Mathematical Classics (Suanjing shishu).] Wei, Liu Hui: commentary. Tang, Grand Master for Court Discussion (chaoyi dafu), serving as Grand Astrologer (taishi ling), Senior Commandant of Light Chariots (shang qingche duwei), your servant Li Chunfeng and others: annotation by imperial command.\end{uniten}

\begin{unitzh}\unitlabel{jiuzhang-003.2}方田{\zhsmall 以御田疇界域}\end{unitzh}
\begin{uniten}Rectangular fields (fangtian) [Liu Hui: For dealing with the boundaries and extents of fields.]\end{uniten}

\begin{unitzh}\unitlabel{jiuzhang-003.3}今有田廣十五步，從十六步。問：爲田幾何？荅曰：一畝。\end{unitzh}
\begin{uniten}Suppose there is a field 15 bu (步) in breadth (guang) and 16 bu in length (cong). Question: how much field does it make? Answer: 1 mu (畝).\end{uniten}

\begin{unitzh}\unitlabel{jiuzhang-003.4}又有田廣十二步，從十四步。問：爲田幾何？荅曰：一百六十八步。{\zhsmall 圖：從十四，廣十二。}\end{unitzh}
\begin{uniten}Again there is a field 12 bu in breadth and 14 bu in length. Question: how much field does it make? Answer: 168 bu. [Liu Hui: Diagram: length 14, breadth 12.]\end{uniten}

\begin{unitzh}\unitlabel{jiuzhang-003.5}方田 術曰：廣從步數相乘得積步。{\zhsmall 此積謂田羃。凡廣從相乘謂之羃。 臣淳風等謹按：經云「廣從相乘得積步」，注云「廣從相乘謂之羃」。觀斯注意，積羃義同。以理推之，固當不爾。何則？羃是方靣單布之名，積乃衆數聚居之稱。循名責實，二者全殊；雖欲同之，竊恐不可。今以凡言羃者，據廣從之一方；其言積者，舉衆步之都數。經云「相乘得積步」，卽是都數之明文；注云「謂之爲羃」，全乖積步之本意。此注前云「積謂田羃」，於理得通；復云「謂之爲羃」，䌓而不當。今者注釋，存善去非，略爲科𥳑，遺諸後學。}\end{unitzh}
\begin{uniten}Rectangular fields. Method: Multiply the numbers of bu of the breadth and the length together to obtain the area in bu (jibu). [Liu Hui: This area (ji) means the area (mi) of the field. In general, breadth and length multiplied together are called area (mi).] [Li Chunfeng et al.: Your servant Chunfeng and others respectfully note: The Classic says "multiply breadth and length together to obtain the area in bu"; the commentary says "breadth and length multiplied together are called area (mi)". Looking at the intent of this commentary, area (ji) and area (mi) have the same meaning. Reasoning from principle, it certainly ought not to be so. Why? Mi is the name for a square surface spread out in a single layer; ji is the term for many numbers gathered and dwelling together. Following the names and examining the realities, the two are entirely different; although one may wish to make them the same, we humbly fear that it is not possible. Now we take it that whenever one speaks of mi, it is based on the single square of breadth and length; when one speaks of ji, it takes up the total number of the many bu. The Classic says "multiply together to obtain the area in bu": this is precisely an explicit statement of the total number. The commentary says "it is called mi": this entirely departs from the original meaning of "area in bu". This commentary first says "the area (ji) means the area (mi) of the field", which in principle holds; it then again says "it is called mi", which is redundant and inappropriate. In the present annotation we keep what is good and remove what is wrong, briefly sorting and selecting, and leave it to later students.]\end{uniten}

\begin{unitzh}\unitlabel{jiuzhang-003.6}以畝法二百四十步除之，卽畝數。百畝爲一頃。{\zhsmall 臣淳風等謹按：此爲篇端，故特舉頃、畝二法。餘術不復言者，從此可知。一畝之田，廣十五步，從而疏之，令爲十五行，卽每行廣一步而從十六步；又横而截之，令爲十六行，卽毎行廣一步而從十五步。此卽從疏横截之步，各自爲方，凡有二百四十步，爲一畝之地，步數正同。以此言之，卽廣從相乘得積步，驗矣。二百四十步者，畝法也；百畝者，頃法也。故以除之，卽得。}\end{unitzh}
\begin{uniten}Divide it by the divisor (fa) for mu, 240 bu; this gives the number of mu. 100 mu make 1 qing (頃). [Li Chunfeng et al.: Your servant Chunfeng and others respectfully note: This is the beginning of the chapter; therefore the two divisors for qing and mu are specially presented. Where the other methods do not mention them again, it can be understood from this. A field of 1 mu is 15 bu in breadth; split it lengthwise, making 15 rows; then each row is 1 bu in breadth and 16 bu in length. Again cut it crosswise, making 16 rows; then each row is 1 bu in breadth and 15 bu in length. These bu, split lengthwise and cut crosswise, each form a square by themselves; altogether there are 240 bu, which make the ground of 1 mu; the numbers of bu are exactly the same. Speaking from this, that multiplying breadth and length together gives the area in bu is verified. 240 bu is the divisor for mu; 100 mu is the divisor for qing. Therefore one divides by them and obtains it.]\end{uniten}

\begin{unitzh}\unitlabel{jiuzhang-003.7}今有田廣一里，從一里。問：爲田幾何？荅曰：三頃七十五畝。\end{unitzh}
\begin{uniten}Suppose there is a field 1 li (里) in breadth and 1 li in length. Question: how much field does it make? Answer: 3 qing 75 mu.\end{uniten}

\begin{unitzh}\unitlabel{jiuzhang-003.8}又有田廣二里，從三里。問：爲田幾何？荅曰：二十二頃五十畝。\end{unitzh}
\begin{uniten}Again there is a field 2 li in breadth and 3 li in length. Question: how much field does it make? Answer: 22 qing 50 mu.\end{uniten}

\begin{unitzh}\unitlabel{jiuzhang-003.9}里田 術曰：廣從里數相乘得積里，以三百七十五乘之，卽畝數。{\zhsmall 按：此術廣從里數相乘得積里。故方里之中有三頃七十五畝，故以乘之，卽得畝數也。}\end{unitzh}
\begin{uniten}Fields in li (litian). Method: Multiply the numbers of li of the breadth and the length together to obtain the area in li; multiply this by 375, which gives the number of mu. [Liu Hui: Note: In this method the numbers of li of the breadth and the length are multiplied together to obtain the area in li. So within a square li there are 3 qing 75 mu; therefore one multiplies by this and obtains the number of mu.]\end{uniten}

\begin{unitzh}\unitlabel{jiuzhang-003.10}今有十八分之十二。問：約之得幾何？荅曰：三分之二。\end{unitzh}
\begin{uniten}Suppose there is 12/18. Question: reducing it, how much does one obtain? Answer: 2/3.\end{uniten}

\begin{unitzh}\unitlabel{jiuzhang-003.11}又有九十一分之四十九。問：約之得幾何？荅曰：十三分之七。\end{unitzh}
\begin{uniten}Again there is 49/91. Question: reducing it, how much does one obtain? Answer: 7/13.\end{uniten}

\begin{unitzh}\unitlabel{jiuzhang-003.12}約分{\zhsmall 按：約分者，物之數量，不可悉全，必以分言之；分之爲數，䌓則難用。設有四分之二者，繁而言之，亦可爲八分之四；約而言之，則二分之一也。雖則異辭，至於爲數，亦同歸爾。法實相推，動有參差，故爲術者先治諸分。}術曰：可半者半之；不可半者，副置分母、子之數，以少減多，更相減損，求其等也。以等數約之。{\zhsmall 等數約之，卽除也。其所以相減者，皆等數之重疊，故以等數約之。}\end{unitzh}
\begin{uniten}Reducing fractions (yuefen) [Liu Hui: Note: As for reducing fractions: the numbers and quantities of things cannot all be whole, and must be expressed in parts (fen); parts, as numbers, are hard to use when they are complicated. Suppose there is 2/4: expressed in a complicated way, it can also be 4/8; expressed in a reduced way, it is 1/2. Although the wordings differ, as numbers they come to the same. When the divisor and the dividend (shi) are reckoned against each other, there are frequently irregularities; therefore those who make methods first put the fractions in order.] Method: If it can be halved, halve it. If it cannot be halved, set out separately the numbers of the denominator (fenmu) and the numerator (zi); subtract the smaller from the larger, subtracting from each other in turn, to seek their equal. Reduce by the equal number (dengshu). [Liu Hui: Reducing by the equal number is dividing. What is subtracted from each other is in every case repeated layers of the equal number; therefore one reduces by the equal number.]\end{uniten}

\begin{unitzh}\unitlabel{jiuzhang-003.13}今有三分之一、五分之二。問：合之得幾何？荅曰：十五分之十一。\end{unitzh}
\begin{uniten}Suppose there are 1/3 and 2/5. Question: combining them, how much does one obtain? Answer: 11/15.\end{uniten}

\begin{unitzh}\unitlabel{jiuzhang-003.14}又有三分之二、七分之四、九分之五。問：合之得幾何？荅曰：得一、六十三分之五十。\end{unitzh}
\begin{uniten}Again there are 2/3, 4/7 and 5/9. Question: combining them, how much does one obtain? Answer: one obtains 1 50/63.\end{uniten}

\begin{unitzh}\unitlabel{jiuzhang-003.15}又有二分之一、三分之二、四分之三、五分之四。問：合之得幾何？荅曰：得二、六十分之四十三。\end{unitzh}
\begin{uniten}Again there are 1/2, 2/3, 3/4 and 4/5. Question: combining them, how much does one obtain? Answer: one obtains 2 43/60.\end{uniten}

\begin{unitzh}\unitlabel{jiuzhang-003.16}合分{\zhsmall 臣淳風等謹按：合分者，數非一端，分無定凖，諸分子雜互，羣母參差，麄細既殊，理難從一，故齊其衆分，同其羣母，令可相并，故曰合分。}術曰：母互乘子，并以爲實；母相乘爲法。{\zhsmall 母互乘子：約而言之者其分麄，繁而言之者其分細；雖則麄細有殊，然其實一也。衆分錯雜，非細不㑹。乗而散之，所以通之；通之則可并也。凡母互乘子謂之齊，羣母相乘謂之同。同者，相與通同共一母也；齊者，子與母齊，勢不可失本數也。方以類聚，物以羣分。數同類者無遠，數異類者無近。遠而通體者，雖異位而相從也；近而殊形者，雖同列而相違也。然則齊同之術要矣。錯綜度數，動之則諧，其猶佩觽解結，無往而不理焉。乘以散之，約以聚之，齊同以通之，此其算之綱紀乎！其一術者，可令母除爲率，率乘子爲齊。}實如法而一。不滿法者，以法命之。{\zhsmall 今欲求其實，故齊其子，又同其母，令如母而一。其餘以等數約之，卽得。所謂同法爲母，實餘爲子，皆從此例。}其母同者，直相從之。\end{unitzh}
\begin{uniten}Combining fractions (hefen) [Li Chunfeng et al.: Your servant Chunfeng and others respectfully note: As for combining fractions: the numbers are not of one kind, and the fractions have no fixed standard; the various numerators are mixed together, the many denominators are uneven; since coarse and fine are different, it is hard in principle to bring them to one. Therefore one homogenizes (qi) the many fractions and equalizes (tong) the many denominators, so that they can be added together; therefore it is called "combining fractions".] Method: The denominators multiply the numerators crosswise; add these to make the dividend. The denominators multiplied together make the divisor. [Liu Hui: The denominators multiply the numerators crosswise: when it is expressed in reduced terms, its parts are coarse; when it is expressed in complicated terms, its parts are fine; although coarse and fine differ, in reality they are one. When the many fractions are mixed together, they do not meet unless made fine. Multiplying to disperse them is the means of making them communicate; once they communicate, they can be added. In general, the denominators multiplying the numerators crosswise is called homogenizing; the many denominators multiplied together is called equalizing. Equalizing means that they communicate with one another and become the same, sharing one denominator; homogenizing means that the numerators are made homogeneous with the denominators, since by the nature of the situation the original numbers must not be lost. Things of a sort gather by kind; things divide by groups. Numbers of the same kind are never far apart; numbers of different kinds are never close. Those that are far apart but share one body follow each other although in different positions; those that are close but of different forms oppose each other although in the same row. That being so, the method of homogenizing and equalizing is essential. Measures and numbers intricately interwoven harmonize as soon as one sets them in motion; it is like wearing a knot-opening awl (xi) to undo knots: there is nowhere it goes that it does not put things in order. Multiplying in order to disperse them, reducing in order to gather them, homogenizing and equalizing in order to make them communicate: this is surely the guiding principle of computation! Another method: one may let the denominators divide to make the rates (lü), and the rates multiply the numerators to make the homogenized terms.] Divide the dividend by the divisor. Whatever does not fill the divisor, name it by the divisor. [Liu Hui: Now one wishes to find their dividend; therefore one homogenizes their numerators and also equalizes their denominators, and has it divided by the denominator. The remainder is reduced by the equal number, and then one obtains it. What is called "the equalized divisor makes the denominator, the remainder of the dividend makes the numerator" all follows this pattern.] If their denominators are the same, simply combine them.\end{uniten}

\begin{unitzh}\unitlabel{jiuzhang-003.17}今有九分之八，減其五分之一。問：餘幾何？荅曰：四十五分之三十一。\end{unitzh}
\begin{uniten}Suppose there is 8/9; subtract 1/5 from it. Question: how much remains? Answer: 31/45.\end{uniten}

\begin{unitzh}\unitlabel{jiuzhang-003.18}又有四分之三，減其三分之一。問：餘幾何？荅曰：十二分之五。\end{unitzh}
\begin{uniten}Again there is 3/4; subtract 1/3 from it. Question: how much remains? Answer: 5/12.\end{uniten}

\begin{unitzh}\unitlabel{jiuzhang-003.19}減分{\zhsmall 臣淳風等謹按：諸分子母數各不同，以少減多，欲知餘幾，減餘爲實，故曰減分。}術曰：母互乘子，以少減多，餘爲實；母相乘爲法；實如法而一。{\zhsmall 母互乘子者，知以齊其子也；以少減多者，齊故可相減也；母相乘爲法者，同其母也。母同子齊，故如母而一，卽得。}\end{unitzh}
\begin{uniten}Subtracting fractions (jianfen) [Li Chunfeng et al.: Your servant Chunfeng and others respectfully note: The numbers of the numerators and denominators of the various fractions each differ; one subtracts the smaller from the larger, wishing to know how much remains; the remainder of the subtraction is made the dividend; therefore it is called "subtracting fractions".] Method: The denominators multiply the numerators crosswise; subtract the smaller from the larger; the remainder makes the dividend. The denominators multiplied together make the divisor. Divide the dividend by the divisor. [Liu Hui: The denominators multiplying the numerators crosswise: one knows that this serves to homogenize the numerators. Subtracting the smaller from the larger: because they are homogenized, they can be subtracted from each other. The denominators multiplied together making the divisor: this equalizes the denominators. The denominators being equal and the numerators homogenized, one therefore divides by the denominator and obtains it.]\end{uniten}

\begin{unitzh}\unitlabel{jiuzhang-003.20}今有八分之五、二十五分之十六。問：孰多？多幾何？荅曰：二十五分之十六多，多二百分之三。\end{unitzh}
\begin{uniten}Suppose there are 5/8 and 16/25. Question: which is more? By how much is it more? Answer: 16/25 is more; it is more by 3/200.\end{uniten}

\begin{unitzh}\unitlabel{jiuzhang-003.21}又有九分之八、七分之六。問：孰多？多幾何？荅曰：九分之八多，多六十三分之二。\end{unitzh}
\begin{uniten}Again there are 8/9 and 6/7. Question: which is more? By how much is it more? Answer: 8/9 is more; it is more by 2/63.\end{uniten}

\begin{unitzh}\unitlabel{jiuzhang-003.22}又有二十一分之八、五十分之十七。問：孰多？多幾何？荅曰：二十一分之八多，多一千五十分之四十三。\end{unitzh}
\begin{uniten}Again there are 8/21 and 17/50. Question: which is more? By how much is it more? Answer: 8/21 is more; it is more by 43/1050.\end{uniten}

\begin{unitzh}\unitlabel{jiuzhang-003.23}課分{\zhsmall 臣淳風等謹按：分各異名，理不齊一，校其相多之數，故曰課分也。}術曰：母互乘子，以少減多，餘爲實；母相乗爲法；實如法而一，卽相多也。{\zhsmall 臣淳風等謹按：此術母互乘子，以少分減多分，多與減分義同；唯相多之數，意共減分有異：減分知求其餘數有幾，課分知以其餘數相多也。}\end{unitzh}
\begin{uniten}Comparing fractions (kefen) [Li Chunfeng et al.: Your servant Chunfeng and others respectfully note: The fractions each have different names, and in principle they are not uniform; one checks the amount by which one exceeds the other; therefore it is called "comparing fractions".] Method: The denominators multiply the numerators crosswise; subtract the smaller from the larger; the remainder makes the dividend. The denominators multiplied together make the divisor. Divide the dividend by the divisor; this is the amount by which one exceeds the other. [Li Chunfeng et al.: Your servant Chunfeng and others respectfully note: In this method the denominators multiply the numerators crosswise and the smaller fraction is subtracted from the larger fraction; it is largely the same in meaning as subtracting fractions. Only the amount by which one exceeds the other differs in intent from subtracting fractions: in subtracting fractions one knows to seek how much the remaining number is; in comparing fractions one knows to take the remaining number as the amount by which one exceeds the other.]\end{uniten}

\begin{unitzh}\unitlabel{jiuzhang-004.1}今有三分之一、三分之二、四分之三，問：減多益少，各幾何而平？荅曰：減四分之三者二，三分之二者一，并以益三分之一，而各平於十二分之七。\end{unitzh}
\begin{uniten}Suppose there are 1/3, 2/3 and 3/4. Question: diminishing the greater and increasing the lesser, by how much in each case so that they are equal (ping)? Answer: Diminish the 3/4 by 2 and the 2/3 by 1; combine these to increase the 1/3; and each is equal at 7/12.\end{uniten}

\begin{unitzh}\unitlabel{jiuzhang-004.2}又有二分之一、三分之二、四分之三，問：減多益少，各幾何而平？荅曰：減三分之二者一，四分之三者四，并以益二分之一，而各平於三十六分之二十三。\end{unitzh}
\begin{uniten}Again, there are 1/2, 2/3 and 3/4. Question: diminishing the greater and increasing the lesser, by how much in each case so that they are equal? Answer: Diminish the 2/3 by 1 and the 3/4 by 4; combine these to increase the 1/2; and each is equal at 23/36.\end{uniten}

\begin{unitzh}\unitlabel{jiuzhang-004.3}平分{\zhsmall 臣淳風等謹按：平分者，諸分參差，欲令齊等，減彼之多，增此之少，故曰平分也。}\end{unitzh}
\begin{uniten}Equalizing fractions (pingfen) [Li Chunfeng et al.: "Equalizing fractions": the various fractions are uneven, and one wishes to make them equal, diminishing the excess of that one and augmenting the shortfall of this one; hence it is called "equalizing fractions".]\end{uniten}

\begin{unitzh}\unitlabel{jiuzhang-004.4}術曰：母互乘子，{\zhsmall 齊其子也。}副并爲平實。{\zhsmall 臣淳風等謹按：母互乘子，副并爲平實者，定此平實，立限衆子所當損益，如限爲平。}母相乘爲法。{\zhsmall 母相乘爲法者，亦齊其子，又同其母。}以列數乘未并者，各自爲列實。亦以列數乘法。{\zhsmall 此當副并列數爲平實，若然則重有分，故反以列數乘同齊。臣淳風等謹按：問云所平之分多少不定，或三或二，列位無常。平三知置位三重，平二知置位二重。凡此之例，一凖平分，不可預定多少，故直云列數而已。}以平實減列實，餘約之爲所減。并所減以益於少。以法命平實，各得其平。\end{unitzh}
\begin{uniten}Method: The denominators cross-multiply the numerators, [Liu Hui: This makes the numerators uniform (qi).] and these are added in an auxiliary position (fubing) to make the equalizing dividend (pingshi). [Li Chunfeng et al.: "The denominators cross-multiply the numerators, and these are added in an auxiliary position to make the equalizing dividend": this fixes the equalizing dividend and sets up a limit for how much the various numerators should lose or gain; when they accord with the limit, they are equal.] The denominators multiplied together make the divisor (fa). [Liu Hui: "The denominators multiplied together make the divisor": this also makes the numerators uniform, and moreover makes the denominators the same (tong 同).] With the number of entries (lieshu) multiply those not yet added; each by itself makes an entry dividend (lieshi). Also multiply the divisor by the number of entries. [Liu Hui: Here the auxiliary sum and the number of entries ought to make the equalizing dividend; if so, there would be double fractions; therefore, conversely, the number of entries is used to multiply what has been made the same and made uniform.] [Translator's note: in "the auxiliary sum and the number of entries ought to make", the Chinese seems to lack a verb such as "divide".] [Li Chunfeng et al.: The questions say that the fractions to be equalized are not fixed in number, sometimes three, sometimes two; the positions of the entries are not constant. In equalizing three, one knows to set out the positions in three tiers; in equalizing two, one knows to set out the positions in two tiers. All cases of this kind conform uniformly to equalizing fractions; how many cannot be fixed in advance, therefore it simply says "the number of entries", and that is all.] Diminish the entry dividends by the equalizing dividend; reduce (yue) the remainders to make the amounts diminished. Combine the amounts diminished to increase the lesser. Name (ming) the equalizing dividend by the divisor; each obtains its equal value.\end{uniten}

\begin{unitzh}\unitlabel{jiuzhang-004.5}今有七人分八錢三分錢之一，問：人得幾何？荅曰：人得一錢二十一分錢之四。\end{unitzh}
\begin{uniten}Suppose there are 7 people dividing 8 1/3 qian (錢, cash). Question: how much does each person get? Answer: Each person gets 1 4/21 qian.\end{uniten}

\begin{unitzh}\unitlabel{jiuzhang-004.6}又有三人三分人之一，分六錢三分錢之一、四分錢之三，問：人得幾何？荅曰：人得二錢八分錢之一。\end{unitzh}
\begin{uniten}Again, there are 3 1/3 people dividing 6 1/3 qian and 3/4 qian. Question: how much does each person get? Answer: Each person gets 2 1/8 qian.\end{uniten}

\begin{unitzh}\unitlabel{jiuzhang-004.7}經分{\zhsmall 臣淳風等謹按：經分者，自合分已下，皆與諸分相齊，此乃直求一人之分，以人數分所分，故曰經分也。}\end{unitzh}
\begin{uniten}Dividing fractions (jingfen) [Li Chunfeng et al.: "Dividing fractions": from "Combining fractions" (hefen) onwards, all make the various fractions uniform with one another; this one directly seeks the share of one person, dividing what is to be divided by the number of people; hence it is called "dividing fractions".]\end{uniten}

\begin{unitzh}\unitlabel{jiuzhang-004.8}術曰：以人數爲法，錢數爲實，實如法而一。有分者通之。{\zhsmall 母互乘子者，齊其子；母相乘者，同其母。以母通之者，分母乘全內子。乘散全則爲積分，積分則與分子相通，故可令相從。凡數相與者謂之率。率者，自相與通。有分則可散，分重疊則約也。等除法實，相與率也。故散分者，必令兩分母相乘法實也。}重有分者，同而通之。{\zhsmall 又以法分母乘實，實分母乘法。此謂法實俱有分，故令分母各乘全分內子，又令分母互乘上下。}\end{unitzh}
\begin{uniten}Method: Take the number of people as the divisor and the number of qian as the dividend (shi); divide the dividend by the divisor. Where there are fractions, convert (tong 通) them. [Liu Hui: "The denominators cross-multiply the numerators" makes the numerators uniform; "the denominators multiply each other" makes the denominators the same. "Converting them by the denominator" means: the denominator multiplies the integer and the numerator is put in (na). Multiplying disperses the integer, which then becomes accumulated parts (jifen); the accumulated parts are then interconvertible with the numerator, so they can be made to join each other. Whenever numbers stand in relation to each other, they are called rates (lü). Rates are interconvertible among themselves. If there are fractions, they can be dispersed (san); if parts are layered upon each other, they are reduced. When the divisor and the dividend are divided equally, they are still rates in relation to each other. Therefore, in dispersing fractions, one must make the two denominators multiply the divisor and the dividend.] Where there are double fractions, make them the same and convert them. [Liu Hui: Also multiply the dividend by the denominator of the divisor, and the divisor by the denominator of the dividend. This means that the divisor and the dividend both have fractions; therefore one makes each denominator multiply the integer and put in the numerator of its fraction, and moreover makes the denominators cross-multiply above and below.]\end{uniten}

\begin{unitzh}\unitlabel{jiuzhang-004.9}今有田廣七分步之四，從五分步之三，問：爲田幾何？荅曰：三十五分步之十二。\end{unitzh}
\begin{uniten}Suppose there is a field whose width (guang) is 4/7 bu (步) and length (zong) 3/5 bu. Question: how much field does it make? Answer: 12/35 bu.\end{uniten}

\begin{unitzh}\unitlabel{jiuzhang-004.10}又有田廣九分步之七，從十一分步之九，問：爲田幾何？荅曰：十一分步之七。\end{unitzh}
\begin{uniten}Again, there is a field whose width is 7/9 bu and length 9/11 bu. Question: how much field does it make? Answer: 7/11 bu.\end{uniten}

\begin{unitzh}\unitlabel{jiuzhang-004.11}又有田廣五分步之四，從九分步之五，問：爲田幾何？荅曰：九分步之四。\end{unitzh}
\begin{uniten}Again, there is a field whose width is 4/5 bu and length 5/9 bu. Question: how much field does it make? Answer: 4/9 bu.\end{uniten}

\begin{unitzh}\unitlabel{jiuzhang-004.12}乘分{\zhsmall 臣淳風等謹按：乘分者，分母相乘爲法，子相乘爲實，故曰乘分。}\end{unitzh}
\begin{uniten}Multiplying fractions (chengfen) [Li Chunfeng et al.: "Multiplying fractions": the denominators multiplied together make the divisor, and the numerators multiplied together make the dividend; hence it is called "multiplying fractions".]\end{uniten}

\begin{unitzh}\unitlabel{jiuzhang-004.13}術曰：母相乘爲法，子相乘爲實，實如法而一。{\zhsmall 凡實不滿法者，乃有母子之名。若有分，以乘其實而長之，則亦滿法，乃爲全耳。又以子有所乘，故母當報除。報除者，實如法而一也。今子相乘，則母各當報除，因令分母相乘而連除也。此田有廣從，難以廣諭。設有問者曰：馬二十匹，直金十二斤，今賣馬二十匹，三十五人分之，人得幾何？荅曰：三十五分斤之十二。其爲之也，當如經分術，以十二斤金爲實，三十五人爲法。設更言馬五匹，直金三斤，今賣四匹，七人分之，人得幾何？荅曰：人得三十五分斤之十二。其爲之也，當齊其金、人之數，皆合初問，入於經分矣。然則分子相乘爲實者，猶齊其金也；母相乘爲法者，猶齊其人也。同其母爲二十馬，無事於同，但欲求齊而已。又馬五匹直金三斤，完全之率；分而言之，則爲一匹直金五分斤之三。七人賣四馬，一人賣七分馬之四。分子與人交互相生，所從言之異，而計數則三術同歸也。}\end{unitzh}
\begin{uniten}Method: The denominators multiplied together make the divisor; the numerators multiplied together make the dividend; divide the dividend by the divisor. [Liu Hui: Whenever the dividend does not fill the divisor, there are the names "denominator" (mu) and "numerator" (zi). If there is a fraction, and one multiplies its dividend so as to lengthen it, then it too fills the divisor and only then becomes an integer. Moreover, since the numerator has been multiplied by something, the denominator ought to divide in return (baochu). "Dividing in return" is dividing the dividend by the divisor. Now the numerators are multiplied together, so the denominators ought each to divide in return; hence one makes the denominators multiply each other and divides jointly (lianchu). This field has width and length, and it is hard to use it for a broad illustration. Suppose someone asks: "20 horses are worth 12 jin (斤) of gold; now 20 horses are sold and 35 people divide it; how much does each person get?" Answer: 12/35 jin. In doing it, one should follow the method of dividing fractions, taking the 12 jin of gold as the dividend and the 35 people as the divisor. Suppose it is put another way: "5 horses are worth 3 jin of gold; now 4 horses are sold and 7 people divide it; how much does each person get?" Answer: Each person gets 12/35 jin. In doing it, one should make uniform the numbers of the gold and of the people; they all agree with the first question and enter into dividing fractions. Thus "the numerators multiplied together make the dividend" is like making the gold uniform; "the denominators multiplied together make the divisor" is like making the people uniform. Making their denominators the same gives 20 horses; there is no business with making the same, one only wishes to seek uniformity, that is all. Moreover, "5 horses are worth 3 jin of gold" is a rate in whole numbers; stated in parts, it becomes "1 horse is worth 3/5 jin of gold". 7 people sell 4 horses: 1 person sells 4/7 horse. The numerators and the people generate each other interchangeably; what they are stated from differs, but in reckoning the three methods come to the same end.]\end{uniten}

\begin{unitzh}\unitlabel{jiuzhang-004.14}今有田廣三步三分步之一，從五步五分步之二，問：爲田幾何？荅曰：十八步。\end{unitzh}
\begin{uniten}Suppose there is a field whose width is 3 1/3 bu and length 5 2/5 bu. Question: how much field does it make? Answer: 18 bu.\end{uniten}

\begin{unitzh}\unitlabel{jiuzhang-004.15}又有田廣七步四分步之三，從十五步九分步之五，問：爲田幾何？荅曰：一百二十步九分步之五。\end{unitzh}
\begin{uniten}Again, there is a field whose width is 7 3/4 bu and length 15 5/9 bu. Question: how much field does it make? Answer: 120 5/9 bu.\end{uniten}

\begin{unitzh}\unitlabel{jiuzhang-004.16}又有田廣十八步七分步之五，從二十三步十一分步之六，問：爲田幾何？荅曰：一畝二百步十一分步之七。\end{unitzh}
\begin{uniten}Again, there is a field whose width is 18 5/7 bu and length 23 6/11 bu. Question: how much field does it make? Answer: 1 mu (畝) 200 7/11 bu.\end{uniten}

\begin{unitzh}\unitlabel{jiuzhang-004.17}大廣田{\zhsmall 臣淳風等謹按：大廣田者，初術直有全步而無餘分，次術空有餘分而無全步，此術先見全步，復有餘分，可以廣兼三術，故曰大廣。}\end{unitzh}
\begin{uniten}Great broad fields (daguangtian) [Li Chunfeng et al.: "Great broad fields": the first method has only whole bu and no remaining fractions; the next method has merely remaining fractions and no whole bu; this method first shows whole bu and then also has remaining fractions; it can broadly combine the three methods; hence it is called "great broad".]\end{uniten}

\begin{unitzh}\unitlabel{jiuzhang-004.18}術曰：分母各乘其全，分子從之，{\zhsmall 分母各乘其全，分子從之者，通全步內分子。如此則母子皆爲實矣。}相乘爲實。分母相乘爲法，{\zhsmall 猶乘分也。}實如法而一。{\zhsmall 今爲術，廣從俱有分，當各自通其分。命母入者，還須出之，故令分母相乘爲法而連除之。}\end{unitzh}
\begin{uniten}Method: The denominators each multiply their integers, and the numerators are joined (cong) to them; [Liu Hui: "The denominators each multiply their integers, and the numerators are joined to them" means converting the whole bu and putting in the numerators. In this way the denominators and numerators all become dividends.] multiply these together to make the dividend. The denominators multiplied together make the divisor; [Liu Hui: This is like multiplying fractions.] divide the dividend by the divisor. [Liu Hui: Now, in this method, width and length both have fractions, and each should convert its own fraction. What was put in by naming the denominator must be taken out again; therefore one makes the denominators multiply each other to make the divisor and divides by it jointly.]\end{uniten}

\begin{unitzh}\unitlabel{jiuzhang-004.19}今有圭田廣十二步，正從二十一步，問：爲田幾何？荅曰：一百二十六步。\end{unitzh}
\begin{uniten}Suppose there is a jade-tablet field (guitian) whose width is 12 bu and upright length (zhengzong) 21 bu. Question: how much field does it make? Answer: 126 bu.\end{uniten}

\begin{unitzh}\unitlabel{jiuzhang-004.20}又有圭田廣五步二分步之一，從八步三分步之二，問：爲田幾何？荅曰：二十三步六分步之五。\end{unitzh}
\begin{uniten}Again, there is a jade-tablet field whose width is 5 1/2 bu and length 8 2/3 bu. Question: how much field does it make? Answer: 23 5/6 bu.\end{uniten}

\begin{unitzh}\unitlabel{jiuzhang-004.21}術曰：半廣以乘正從。{\zhsmall 半廣者，以盈補虚爲直田也。亦可半正從以乘廣。按：半廣乘從，以取中平之數，故廣從相乘爲積步，畝法除之，卽得也。}\end{unitzh}
\begin{uniten}Method: Halve the width and multiply the upright length by it. [Liu Hui: "Halve the width": using the excess to fill the void, it makes a straight field (zhitian). One may also halve the upright length and multiply the width by it. On examination: halving the width and multiplying the length by it is to take the middle-level number; therefore width and length multiplied together make the area in bu; divide by the mu divisor (mufa), and it is obtained.]\end{uniten}

\begin{unitzh}\unitlabel{jiuzhang-005.1}今有邪田，一頭廣三十步，一頭廣四十二步，正從六十四步。問爲田幾何？荅曰：九畝一百四十四步。\end{unitzh}
\begin{uniten}Suppose there is a slanted field (xie tian 邪田), one end 30 bu (步) wide, the other end 42 bu wide, and 64 bu in straight length (zheng zong 正從). Question: how large is the field? Answer: 9 mu (畝) 144 bu.\end{uniten}

\begin{unitzh}\unitlabel{jiuzhang-005.2}又有邪田，正廣六十五步，一畔從一百步，一畔從七十二步。問爲田幾何？荅曰：二十三畝七十步。\end{unitzh}
\begin{uniten}Again there is a slanted field, 65 bu in straight width (zheng guang 正廣), one side (pan 畔) 100 bu long and the other side 72 bu long. Question: how large is the field? Answer: 23 mu 70 bu.\end{uniten}

\begin{unitzh}\unitlabel{jiuzhang-005.3}術曰：并兩邪而半之，以乘正從若廣；又可半正從若廣，以乘并；畝法而一。{\zhsmall 并而半之者，以盈補虚也。}\end{unitzh}
\begin{uniten}Method: Add the two slanted [sides] and halve [the sum]; with it multiply the straight length or the width. Or one may halve the straight length or the width and with it multiply the sum. Divide by the mu divisor (mu fa 畝法). [Liu Hui: Adding and halving means filling the deficient with the excess.]\end{uniten}

\begin{unitzh}\unitlabel{jiuzhang-005.4}今有箕田，舌廣二十步，踵廣五步，正從三十步。問爲田幾何？荅曰：一畝一百三十五步。\end{unitzh}
\begin{uniten}Suppose there is a dustpan field (ji tian 箕田) whose tongue (she 舌) is 20 bu wide and whose heel (zhong 踵) is 5 bu wide, the straight length being 30 bu. Question: how large is the field? Answer: 1 mu 135 bu.\end{uniten}

\begin{unitzh}\unitlabel{jiuzhang-005.5}又有箕田，舌廣一百一十七步，踵廣五十步，正從一百三十五步。問爲田幾何？荅曰：四十六畝二百三十二步半。\end{unitzh}
\begin{uniten}Again there is a dustpan field whose tongue is 117 bu wide and whose heel is 50 bu wide, the straight length being 135 bu. Question: how large is the field? Answer: 46 mu 232 1/2 bu.\end{uniten}

\begin{unitzh}\unitlabel{jiuzhang-005.6}術曰：并踵舌而半之，以乘正從，畝法而一。{\zhsmall 中分箕田則爲兩邪田，故其術相似。又可并踵舌，半正從，以乘之。}\end{unitzh}
\begin{uniten}Method: Add the heel and the tongue and halve [the sum]; with it multiply the straight length; divide by the mu divisor. [Liu Hui: Dividing the dustpan field down the middle makes two slanted fields; therefore its method is similar. One may also add the heel and the tongue, halve the straight length, and multiply them together.]\end{uniten}

\begin{unitzh}\unitlabel{jiuzhang-005.7}今有圓田，周三十步，徑十步。{\zhsmall 臣淳風等謹按：術意以周三徑一爲率，周三十步合徑十步；今依密率，合徑九步十一分步之六。}問爲田幾何？荅曰：七十五步。{\zhsmall 此於徽術，當爲田七十一步一百五十七分步之一百三。臣淳風等謹依密率，爲田七十一步二十二分步之一十三。}\end{unitzh}
\begin{uniten}Suppose there is a circular field (yuan tian 圓田) with circumference (zhou 周) 30 bu and diameter (jing 徑) 10 bu. [Li Chunfeng et al.: Your servants Chunfeng et al. respectfully note: the intent of the method takes circumference 3, diameter 1 as the rates (lü 率); a circumference of 30 bu corresponds to a diameter of 10 bu. Now, according to the precise rate (mi lü 密率), it corresponds to a diameter of 9 6/11 bu.] Question: how large is the field? Answer: 75 bu. [Liu Hui: By Hui's method this should make a field of 71 103/157 bu.] [Li Chunfeng et al.: Your servants Chunfeng et al., respectfully following the precise rate, make the field 71 13/22 bu.]\end{uniten}

\begin{unitzh}\unitlabel{jiuzhang-005.8}又有圓田，周一百八十一步，徑六十步三分步之一。{\zhsmall 臣淳風等謹按：周三徑一，周一百八十一步，徑六十步三分步之一；依密率，徑五十七步二十二分步之一十三。}問爲田幾何？荅曰：十一畝九十步十二分步之一。{\zhsmall 此於徽術，當爲田十畝二百八步三百一十四分步之一百一十三。臣淳風等謹依密率，當爲田十畝二百五步八十八分步之八十七。}\end{unitzh}
\begin{uniten}Again there is a circular field with circumference 181 bu and diameter 60 1/3 bu. [Li Chunfeng et al.: Your servants Chunfeng et al. respectfully note: with circumference 3, diameter 1, a circumference of 181 bu [has] a diameter of 60 1/3 bu; according to the precise rate, the diameter is 57 13/22 bu.] Question: how large is the field? Answer: 11 mu 90 1/12 bu. [Liu Hui: By Hui's method this should make a field of 10 mu 208 113/314 bu.] [Li Chunfeng et al.: Your servants Chunfeng et al., respectfully following the precise rate, [find that] it should make a field of 10 mu 205 87/88 bu.]\end{uniten}

\begin{unitzh}\unitlabel{jiuzhang-005.9}術曰：半周半徑相乘得積步。{\zhsmall 按：半周爲從，半徑爲廣，故廣從相乘爲積步也。假令圓徑二尺，圓中容六觚之一靣，與圓徑之半，其數均等。合徑率一而外周率三也。又按：爲圖，以六觚之一靣乘半徑，二因而六之，得十二觚之羃。若又割之，次以十二觚之一靣乘一觚之半徑，四因而六之，則得二十四觚之羃。割之彌細，所失彌少。割之又割，以至於不可割，則與圓周合體而無所失矣。觚靣之外，又有餘徑，以靣乘徑，則羃出觚表。若夫觚之細者，與圓合體，則表無餘徑。表無餘徑，則羃不外出矣。以一靣乘半徑，觚而裁之，每輒自倍。故以半周乘半徑而爲圓羃。此以周、徑謂至然之數，非周三徑一之率也。周三者，從其六觚之環耳。以推圓規多少之較，乃弓之與弦也。然世傳此法，莫肯精覈；學者踵古，習其謬失。不有明據，辯之斯難。凡物類形象，不圓則方。方圓之率，誠著於近，則雖遠可知也。由此言之，其用博矣。謹按圓驗，更造密率。恐空設法，數昧而難譬，故置諸撿括，謹詳其記注焉。}\end{unitzh}
\begin{uniten}Method: Multiply half the circumference and the radius (ban jing 半徑, "half the diameter") together to obtain the area in bu (ji bu 積步). [Liu Hui: Note: half the circumference is the length (zong 從) and the radius is the width (guang 廣); therefore width and length multiplied together give the area in bu. Suppose the diameter of a circle is 2 chi (尺): one side of the 6-gon (liu gu 六觚) contained in the circle and half the diameter of the circle are equal in number. [This] matches the rate 1 for the diameter and the rate 3 for the outer perimeter. Further, note: making a diagram, multiply one side of the 6-gon by the radius, multiply by 2 and thereupon by 6, and one obtains the area (mi 羃) of the 12-gon. If one cuts it again, next multiply one side of the 12-gon by the radius of one gon, multiply by 4 and thereupon by 6, and one then obtains the area of the 24-gon. The finer one cuts, the less is lost. Cut it and cut it again until it cannot be cut, and then it coincides with the circumference of the circle and nothing is lost. Outside the sides of the gon there is still a remaining [part of the] radius (yu jing 餘徑); if the side is multiplied by the radius, the area goes beyond the outside of the gon. But when the gon is fine and coincides with the circle, then outside there is no remaining radius; when outside there is no remaining radius, the area does not go outside. Multiply one side by the radius and cut it up gon by gon: each [product] always doubles itself. Therefore multiply half the circumference by the radius to make the area of the circle. Here "circumference" and "diameter" refer to the numbers that are exactly so (zhi ran zhi shu 至然之數), not to the rates of circumference 3, diameter 1. "Circumference 3" merely follows from the ring of the 6-gon. If one uses it to infer the difference in magnitude with respect to the circle (yuan gui 圓規), it is [that of] the bow and its string. Yet this method has been handed down through the ages and no one has been willing to examine it closely; scholars follow in the footsteps of the ancients and practise their errors. Without clear evidence, it is hard to dispute it. All kinds and shapes of things, if not round, are square. If the rates of square and circle are truly made clear in what is near, then even what is far can be known. Speaking from this, its uses are broad. Respectfully, on the basis of verification with the circle, I have further created a precise rate. Fearing that if the method were set out bare, the numbers would be obscure and hard to illustrate, I have therefore set them out in a summary (jian kuo 撿括) and respectfully recorded and annotated them in detail.]\end{uniten}

\begin{unitzh}\unitlabel{jiuzhang-005.10}{\zhsmall 割六觚以爲十二觚術曰：置圓徑二尺，半之爲一尺，卽圓裏六觚之靣也。令半徑一尺爲弦，半靣五寸爲句，爲之求股。以句羃二十五寸減弦羃，餘七十五寸，開方除之，下至秒忽。又一退法，求其微數。微數無名者以爲分子，以下爲分母，約作五分忽之二。故得股八寸六分六釐二秒五忽五分忽之二。以減半徑，餘一寸三分三釐九毫九秒四忽五分忽之三，謂之小句。觚之半靣又謂之小股。爲之求弦，其羃二千六百七十九億四千九百一十九萬三千四百四十五忽，餘分弃之。開方除之，卽十二觚之一靣也。割十二觚以爲二十四觚術曰：亦令半徑爲弦，半靣爲句，爲之求股。置上小弦羃，四而一，得六百六十九億八千七百二十九萬八千三百六十一忽，餘分棄之，卽句羃也。以減弦羃，其餘開方除之，得股九寸六分五釐九毫二秒五忽五分忽之四。以減半徑，餘三分四釐七秒四忽五分忽之一，謂之小句。觚之半靣又謂之小股。爲之求小弦，其羃六百八十一億四千八百三十四萬九千四百六十六忽，餘分棄之。開方除之，卽二十四觚之一靣也。}\end{unitzh}
\begin{uniten}[Liu Hui: Cutting the 6-gon to make a 12-gon. Method: Set up the diameter of the circle, 2 chi; halving it gives 1 chi, which is the side of the 6-gon inside the circle. Let the radius, 1 chi, be the hypotenuse (xian 弦) and the half-side, 5 cun (寸), the base (gou 句), and for them find the altitude (gu 股). Subtract the area of the base, 25 cun, from the area of the hypotenuse; the remainder is 75 cun; extract the square root, going down to miao (秒) and hu (忽). Again move the divisor back once (yi tui fa 一退法) to find its minute number (wei shu 微數). The minute number that has no name is taken as the numerator, and the [number] below as the denominator; reduced, it makes 2/5 hu. Thus one obtains the altitude, 8 cun 6 fen (分) 6 li (釐) 2 miao 5 hu 2/5 hu. Subtracting it from the radius leaves 1 cun 3 fen 3 li 9 hao (毫) 9 miao 4 hu 3/5 hu, which is called the small base (xiao gou 小句). The half-side of the gon is also called the small altitude (xiao gu 小股). For them find the hypotenuse: its area is 267,949,193,445 hu, the remaining fraction being discarded. Extract the square root: this is one side of the 12-gon. Cutting the 12-gon to make a 24-gon. Method: Again let the radius be the hypotenuse and the half-side the base, and for them find the altitude. Set up the area of the small hypotenuse above and divide by 4, obtaining 66,987,298,361 hu, the remaining fraction being discarded: this is the area of the base. Subtract it from the area of the hypotenuse and extract the square root of the remainder, obtaining the altitude, 9 cun 6 fen 5 li 9 hao 2 miao 5 hu 4/5 hu. Subtracting it from the radius leaves 3 fen 4 li 7 miao 4 hu 1/5 hu, which is called the small base. The half-side of the gon is also called the small altitude. For them find the small hypotenuse (xiao xian 小弦): its area is 68,148,349,466 hu, the remaining fraction being discarded. Extract the square root: this is one side of the 24-gon.]\end{uniten}

\begin{unitzh}\unitlabel{jiuzhang-005.11}{\zhsmall 割二十四觚以爲四十八觚術曰：亦令半徑爲弦，半靣爲句，爲之求股。置上小弦羃，四而一，得一百七十億三千七百八萬七千三百六十六忽，餘分棄之，卽句羃也。以減弦羃，其餘開方除之，得股九寸九分一釐四毫四秒四忽五分忽之四。以減半徑，餘八釐五毫五秒五忽五分忽之一，謂之小句。觚之半靣又謂之小股。爲之求小弦，其羃一百七十一億一十二十七萬八千八百一十三忽，餘分棄之。開方除之，得小弦一寸三分八毫六忽，餘分棄之，卽四十八觚之一靣。以半徑一尺乘之，又以二十四乘之，得羃三萬一千三百九十三億四千四百萬忽。以百億除之，得羃三百一十三寸六百二十五分寸之五百八十四，卽九十六觚之羃也。}\end{unitzh}
\begin{uniten}[Liu Hui: Cutting the 24-gon to make a 48-gon. Method: Again let the radius be the hypotenuse and the half-side the base, and for them find the altitude. Set up the area of the small hypotenuse above and divide by 4, obtaining 17,037,087,366 hu, the remaining fraction being discarded: this is the area of the base. Subtract it from the area of the hypotenuse and extract the square root of the remainder, obtaining the altitude, 9 cun 9 fen 1 li 4 hao 4 miao 4 hu 4/5 hu. Subtracting it from the radius leaves 8 li 5 hao 5 miao 5 hu 1/5 hu, which is called the small base. The half-side of the gon is also called the small altitude. For them find the small hypotenuse: its area is 171 hundred-million (yi 億) 10 20 7 ten-thousand (wan 萬) 8,813 hu [Translator's note: the group 一十二十七 before 萬 is garbled as printed], the remaining fraction being discarded. Extract the square root, obtaining the small hypotenuse, 1 cun 3 fen 8 hao 6 hu, the remaining fraction being discarded: this is one side of the 48-gon. Multiply it by the radius, 1 chi, and again multiply by 24, obtaining an area of 3,139,344,000,000 hu. Divide by 10,000,000,000, obtaining an area of 313 584/625 cun: this is the area of the 96-gon.]\end{uniten}

\begin{unitzh}\unitlabel{jiuzhang-005.12}{\zhsmall 割四十八觚以爲九十六觚術曰：亦令半徑爲弦，半靣爲句，爲之求股。置次上弦羃，四而一，得四十二億七千七百五十六萬九千七百三忽，餘分棄之，則句羃也。以減弦羃，其餘開方除之，得股九寸九分七釐八毫五秒八忽十分忽之九。以減半徑，餘二釐一毫四秒一忽十分忽之一，謂之小句。觚之半靣又謂之小股。爲之求小弦，其羃四十二億八千二百一十五萬四千一十二忽，餘分棄之。開方除之，得小弦六分五釐四毫三秒八忽，餘分棄之，卽九十六觚之一靣。以半徑一尺乘之，又以四十八乗之，得羃三萬一千四百一十億二千四百萬忽。以百億除之，得羃三百一十四寸六百二十五分寸之六十四，卽一百九十二觚之羃也。}\end{unitzh}
\begin{uniten}[Liu Hui: Cutting the 48-gon to make a 96-gon. Method: Again let the radius be the hypotenuse and the half-side the base, and for them find the altitude. Set up the area of the hypotenuse next above and divide by 4, obtaining 4,277,569,703 hu, the remaining fraction being discarded: this then is the area of the base. Subtract it from the area of the hypotenuse and extract the square root of the remainder, obtaining the altitude, 9 cun 9 fen 7 li 8 hao 5 miao 8 hu 9/10 hu. Subtracting it from the radius leaves 2 li 1 hao 4 miao 1 hu 1/10 hu, which is called the small base. The half-side of the gon is also called the small altitude. For them find the small hypotenuse: its area is 4,282,154,012 hu, the remaining fraction being discarded. Extract the square root, obtaining the small hypotenuse, 6 fen 5 li 4 hao 3 miao 8 hu, the remaining fraction being discarded: this is one side of the 96-gon. Multiply it by the radius, 1 chi, and again multiply by 48, obtaining an area of 3,141,024,000,000 hu. Divide by 10,000,000,000, obtaining an area of 314 64/625 cun: this is the area of the 192-gon.]\end{uniten}

\begin{unitzh}\unitlabel{jiuzhang-005.13}{\zhsmall 以九十六觚之羃減之，餘六百二十五分寸之一百五，謂之差羃。倍之，爲分寸之二百一十，卽九十六觚之外弧田九十六所，謂以弦乘矢之凡羃也。加此羃於九十六觚之羃，得三百一十四寸六百二十五分寸之一百六十九，則出於圓之表矣。故還就一百九十二觚之全羃三百一十四寸以爲圓羃之定率，而棄其餘分。以半徑一尺除圓羃，倍之，得六尺二寸八分，卽周數。令徑自乘爲方羃四百寸，與圓羃相折，圓羃得一百五十七爲率，方羃得二百爲率。方羃二百，其中容圓羃一百五十七也。圓率猶爲微少。按弧田圖，令方中容圓，圓中容方，内方合外方之半。然則圓羃一百五十七，其中容方羃一百也。又令徑二尺與周六尺二寸八分相約，周得一百五十七，徑得五十，則其相與之率也。周率猶爲微少也。}\end{unitzh}
\begin{uniten}[Liu Hui: Subtract from it the area of the 96-gon: the remainder, 105/625 cun, is called the difference area (cha mi 差羃). Doubling it makes 210 [of these] parts of a cun, which is the arc fields (hu tian 弧田) in 96 places outside the 96-gon, called the total area (fan mi 凡羃) of chord (xian 弦) multiplied by arrow (shi 矢). Adding this area to the area of the 96-gon gives 314 169/625 cun, which then goes beyond the outside of the circle. Therefore one returns to the full area of the 192-gon, 314 cun, taking it as the fixed rate (ding lü 定率) for the area of the circle, and discards its remaining fraction. Divide the area of the circle by the radius, 1 chi, and double it, obtaining 6 chi 2 cun 8 fen: this is the circumference number (zhou shu 周數). Let the diameter be multiplied by itself to make the area of the square (fang mi 方羃), 400 cun; reducing it against the area of the circle (xiang zhe 相折), the area of the circle gets 157 as its rate and the area of the square gets 200 as its rate. A square area of 200 contains within it a circle area of 157. The circle rate is still slightly small. According to the diagram of the arc field, let a square contain a circle and the circle contain a square: the inner square equals half of the outer square. So then a circle area of 157 contains within it a square area of 100. Further, let the diameter, 2 chi, and the circumference, 6 chi 2 cun 8 fen, be reduced together (xiang yue 相約): the circumference gets 157 and the diameter gets 50; these then are their mutual rates. The circumference rate is still slightly small.]\end{uniten}

\begin{unitzh}\unitlabel{jiuzhang-005.14}{\zhsmall 晉武庫中漢時王莽作銅斛，其銘曰：「律嘉量斛，內方尺而圓其外，庣旁九釐五毫，羃一百六十二寸，深一尺，積一千六百二十寸，容十斗。」以此術求之，得羃一百六十一寸有奇，其數相近矣。此術微少。而斛差羃六百二十五分寸之一百五，以十二觚之羃爲率消息，當取此分寸之三十六，以增於一百九十二觚之羃，以爲圓羃三百一十四寸二十五分寸之四。置徑自乘之方羃四百寸，令與圓羃通相約，圓羃三千九百二十七，方羃得五千，是爲率。方羃五千中容圓羃三千九百二十七；圓羃三千九百二十七中容方羃二千五百也。以半徑一尺除圓羃三百一十四寸二十五分寸之四，倍之，得六尺二寸八分二十五分寸之八，卽周數也。全徑二尺與周數通相約，徑得一千二百五十，周得三千九百二十七，卽其相與之率。若此者，蓋盡其纖微矣。舉而用之，上法仍約耳。當求一千五百三十六觚之一靣，得三千七十二觚之羃，而裁其微分，數亦宜然，重其驗耳。}\end{unitzh}
\begin{uniten}[Liu Hui: In the armoury of the Jin there is a bronze hu (斛) made by Wang Mang in the time of the Han. Its inscription says: "Standard excellent-measure hu (lü jia liang hu 律嘉量斛): inside, a square chi, and its outside round; the margin at the side (tiao pang 庣旁) is 9 li 5 hao; the area is 162 cun; the depth is 1 chi; the volume is 1,620 cun; it holds 10 dou (斗)." Finding it by this method, one obtains an area of 161 cun and a fraction (you qi 有奇); the numbers are close. This method is slightly small. And the difference area of the hu is 105/625 cun; taking the area of the 12-gon as the rate and adjusting (xiaoxi 消息), one should take 36 of these parts of a cun and add them to the area of the 192-gon, making the area of the circle 314 4/25 cun. Set up the area of the square from the diameter multiplied by itself, 400 cun, and let it and the area of the circle be made common and reduced together (tong xiang yue 通相約): the area of the circle [is] 3,927 and the area of the square gets 5,000; these are the rates. A square area of 5,000 contains within it a circle area of 3,927; a circle area of 3,927 contains within it a square area of 2,500. Divide the area of the circle, 314 4/25 cun, by the radius, 1 chi, and double it, obtaining 6 chi 2 cun 8 fen 8/25 cun: this is the circumference number. Make the full diameter, 2 chi, and the circumference number common and reduce them together: the diameter gets 1,250 and the circumference gets 3,927; these are their mutual rates. Something like this presumably exhausts its finest minutiae. Taking it up and using it, the method above is still concise. One should find one side of the 1536-gon, obtain the area of the 3072-gon, and trim its minute fractions; the number should also come out so. [This is] only to reconfirm its verification.]\end{uniten}

\begin{unitzh}\unitlabel{jiuzhang-005.15}{\zhsmall 臣淳風等謹按：舊術求圓，皆以周三徑一爲率。若用之求圓周之數，則周少徑多；用之求其六觚之田，乃與此率合會耳。何則？假令六觚之田，觚閒各一尺爲靣，自然從角至角，其徑二尺可知。此則周六徑二，與周三徑一已合。恐此猶以難曉，今更引物爲喻。設令刻物作圭形}\end{unitzh}
\begin{uniten}[Li Chunfeng et al.: Your servants Chunfeng et al. respectfully note: the old method for finding circles always took circumference 3, diameter 1 as the rates. If one uses it to find the number of the circumference of a circle, then the circumference is too little and the diameter too much; if one uses it to find a 6-gon field, then it agrees with this rate. Why is that? Suppose a 6-gon field, each [interval] between the gons being 1 chi as a side; naturally, from corner (jiao 角) to corner, its diameter is known to be 2 chi. This then is circumference 6, diameter 2, which already agrees with circumference 3, diameter 1. Fearing that this may still be hard to understand, we now further adduce an object as an illustration. Suppose one carves an object into the shape of a gui (圭)] [continues]\end{uniten}

\begin{unitzh}\unitlabel{jiuzhang-006.1}{\zhsmall 者六枚，枚别三靣，皆長一尺。攢此六物，悉使鋭頭向裏，則成六觚之周，角徑亦皆一尺。更從觚角外畔圍繞爲規，則六觚之徑盡達規矣。當靣徑短，不至外規。若以六觚言之，則爲周六尺，徑二尺，靣皆一尺。靣徑股不至外畔，定無二尺可知。故周三徑一之率，於圓周乃是徑多周少。徑一周三，理非精密。蓋術從𥳑要，舉大綱略而言之。劉徽特以爲疎，遂乃改張其率。但周徑相乘，數難契合。徽雖出斯一法，終不能究其纖毫也。祖沖之以其不精，就中更推其數。今者修撰，攈摭諸家，考其是非，冲之爲密，故顯之於徽術之下，冀學者之所裁焉。}\end{unitzh}
\begin{uniten}[continued] [Li Chunfeng et al.: …six such pieces, each piece with three sides, all 1 chi (尺) long. Gather these six objects together, making all their sharp tips point inward; then they form the perimeter of a hexagon (liu gu), and the corner diameters are also all 1 chi. Further, going around the outer edge of the polygon's corners, make a circle (gui); then the diameters of the hexagon all reach the circle. The diameters at the sides are short and do not reach the outer circle. If one speaks of it in terms of the hexagon, it makes a perimeter of 6 chi and a diameter of 2 chi, the sides all being 1 chi. The side diameter, as the altitude (gu), does not reach the outer edge; it can be known that it is certainly not 2 chi. Therefore the rate (lü) of circumference 3 to diameter 1, for the circumference of a circle, means that the diameter is too much and the circumference too little. Diameter 1, circumference 3: in principle this is not precise. For the method follows simplicity and essentials, giving the main outline and speaking of it roughly. Liu Hui particularly regarded it as crude, and thereupon changed its rate. But when circumference and diameter are multiplied with each other, the numbers are hard to bring into agreement. Although Hui produced this one method, in the end he could not exhaust its finest details. Zu Chongzhi, because it was not precise, further worked out the numbers therein. Now, in compiling, we have gathered the various authors and examined their rights and wrongs: Chongzhi's is precise. Therefore we display it below Hui's method, hoping that students will judge it.]\end{uniten}

\begin{unitzh}\unitlabel{jiuzhang-006.2}又術曰：周徑相乘，四而一。{\zhsmall 此周與上觚同耳。周徑相乘，各當以半，而今周徑兩全，故兩母相乘爲四，以報除之。於徽術，以五十乘周，一百五十七而一，卽徑也；以一百五十七乘徑，五十而一，卽周也。新術徑率猶當微少。據周以求徑，則失之長；據徑以求周，則失之短。諸據見徑以求羃者，皆失之於微少；據周以求羃者，皆失之於微多。}{\zhsmall 臣淳風等謹依密率，以七乘周，二十二而一，卽徑；以二十二乘徑，七而一，卽周。依術求之，卽得。}\end{unitzh}
\begin{uniten}Another method: Multiply the circumference and the diameter together; divide by 4. [Liu Hui: This circumference is the same as that of the polygon above. When circumference and diameter are multiplied together, each should be taken as half; but now circumference and diameter are both whole, so the two denominators (mu) multiplied together make 4, by which to divide in return (bao chu). In Hui's method, multiply the circumference by 50 and divide by 157: this is the diameter; multiply the diameter by 157 and divide by 50: this is the circumference. In the new method the diameter rate ought still to be slightly smaller. Taking the circumference to find the diameter, it errs on the long side; taking the diameter to find the circumference, it errs on the short side. All cases of taking the given diameter to find the area (mi) err on being slightly too small; those of taking the circumference to find the area all err on being slightly too large.] [Li Chunfeng et al.: Respectfully following the precise rate (mi lü): multiply the circumference by 7 and divide by 22: this is the diameter; multiply the diameter by 22 and divide by 7: this is the circumference. Seek it according to the method, and it is obtained.]\end{uniten}

\begin{unitzh}\unitlabel{jiuzhang-006.3}又術曰：徑自相乘，三之，四而一。{\zhsmall 按圓徑自乘爲外方，三之，四而一者，是爲圓居外方四分之三也。若令六觚之一靣乘半徑，其羃卽外方四分之一也。因而三之，卽亦居外方四分之三也。是爲圓裏十二觚之羃耳。取以爲圓，失之於微少。於徽新術，當徑自乘，又以一百五十七乘之，二百而一。}{\zhsmall 臣淳風等謹按密率：令徑自乘，以十一乘之，十四而一，卽圓羃也。}\end{unitzh}
\begin{uniten}Another method: Multiply the diameter by itself, triple it, and divide by 4. [Liu Hui: Note: the circle's diameter multiplied by itself makes the outer square; tripling it and dividing by 4 means that the circle occupies 3/4 of the outer square. If one lets one side of the hexagon multiply the half diameter, its area is 1/4 of the outer square. Tripling it accordingly, it likewise occupies 3/4 of the outer square. This is just the area of the dodecagon (shi'er gu) inside the circle. Taking it as the circle errs on being slightly too small. In Hui's new method, one should multiply the diameter by itself, then multiply it by 157, and divide by 200.] [Li Chunfeng et al.: We respectfully note the precise rate: multiply the diameter by itself, multiply it by 11, and divide by 14: this is the area of the circle.]\end{uniten}

\begin{unitzh}\unitlabel{jiuzhang-006.4}又術曰：周自相乘，十二而一。{\zhsmall 六觚之周，其於圓徑，三與一也。故六觚之周自相乘爲羃，若圓徑自乘者九方。九方凡爲十二觚者十有二，故曰十二而一，卽十二觚之羃也。今此令周自乘，非但若爲圓徑自乘者九方而已。然則十二而一，所得又非十二觚之類也。若欲以爲圓羃，失之於多矣。以六觚之周，十二而一可也。於徽新術，直令圓周自乘，又以二十五乘之，三百一十四而一，得圓羃。其率三百一十四者，周自乘之羃也。置周數六尺二寸八分，令自乘，得羃三十九萬四千三百八十四分；又置圓羃三萬一千四百分。皆以一千二百五十六約之，得此率。}{\zhsmall 臣淳風等謹按：方靣自乘，卽得其積。圓周求其羃，股率乃通。但此術所求，用三一爲率。圓田正法，半周及半徑以相乘。今乃用全周自乘，故須以十二爲母。何者？據全周而求半周，則須以二爲法；就全周而求半徑，復假六以除之。是二六相乘，除周自乘之數。依密率，以七乘之，八十八而一。}\end{unitzh}
\begin{uniten}Another method: Multiply the circumference by itself; divide by 12. [Liu Hui: The perimeter of the hexagon is to the circle's diameter as 3 to 1. Therefore the perimeter of the hexagon multiplied by itself makes an area like 9 squares of the circle's diameter multiplied by itself. The 9 squares in all amount to 12 dodecagons; therefore it says "divide by 12", which gives the area of the dodecagon. Now here the circumference is multiplied by itself; it is not merely like 9 squares of the circle's diameter multiplied by itself. That being so, what is obtained by dividing by 12 is moreover not of the kind of the dodecagon. If one wishes to take it as the area of the circle, it errs on being too much. With the perimeter of the hexagon, dividing by 12 is acceptable. In Hui's new method, directly multiply the circle's circumference by itself, then multiply it by 25 and divide by 314 to obtain the circle's area. Its rate 314 is the area of the circumference multiplied by itself. Set out the circumference number 6 chi 2 cun (寸) 8 fen (分) and multiply it by itself, obtaining an area of 394384 fen; again set out the circle's area, 31400 fen. Reduce both by 1256 to obtain these rates.] [Li Chunfeng et al.: We respectfully note: a square's side multiplied by itself gives its area at once. For the circumference of a circle, to find its area, only with the altitude rate (gu lü) does it work out. [Translator's note: 股率 "altitude rate" is obscure here; the text may be corrupt.] But what this method seeks uses 3 to 1 as the rate. The proper rule for the circular field is to multiply half the circumference and half the diameter together. Now it uses the whole circumference multiplied by itself, so 12 must be taken as the denominator. Why? Taking the whole circumference to find half the circumference, 2 must be taken as the divisor (fa); from the whole circumference to find half the diameter, one again borrows 6 to divide it. Thus 2 and 6 multiplied together divide the number of the circumference multiplied by itself. According to the precise rate, multiply it by 7 and divide by 88.]\end{uniten}

\begin{unitzh}\unitlabel{jiuzhang-006.5}今有宛田，下周三十步，徑十六步。問：爲田幾何？荅曰：一百二十步。\end{unitzh}
\begin{uniten}Suppose there is a domed field (wan tian) with a lower circumference of 30 bu (步) and a diameter of 16 bu. Question: How much field does it make? Answer: 120 bu.\end{uniten}

\begin{unitzh}\unitlabel{jiuzhang-006.6}又有宛田，下周九十九步，徑五十一步。問：爲田幾何？荅曰：五畝六十二步四分步之一。\end{unitzh}
\begin{uniten}Again there is a domed field with a lower circumference of 99 bu and a diameter of 51 bu. Question: How much field does it make? Answer: 5 mu (畝) 62 1/4 bu.\end{uniten}

\begin{unitzh}\unitlabel{jiuzhang-006.7}術曰：以徑乘周，四而一。{\zhsmall 此術不驗，故推方錐以見其形。假令方錐下方六尺，高四尺。四尺爲股，下方之半三尺爲句，正靣邪爲弦，弦五尺也。令句弦相乘，四因之，得六十尺，卽方錐四靣見者之羃。若令其中容圓錐，圓錐見羃與方錐見羃，其率猶方羃之與圓羃也。按方錐下六尺，則方周二十四尺，以五尺乘而半之，則亦方錐之見羃。故求圓錐之數，折徑以乘下周之半，卽圓錐之羃也。今宛田上徑圓穹，而與圓錐同術，則羃失之於少矣。然其術難用，故略舉大較，施之大廣田也。求圓錐之羃，猶求圓田之羃也。今用兩全相乘，故以四爲法除之，亦如圓田矣。開立圓術說圓方諸率甚備，可以驗此。}\end{unitzh}
\begin{uniten}Method: Multiply the circumference by the diameter; divide by 4. [Liu Hui: This method does not stand up to verification; therefore one reasons by way of a square cone (fang zhui) to show its shape. Suppose a square cone has a lower square of 6 chi and a height of 4 chi. The 4 chi is the altitude (gu); half of the lower square, 3 chi, is the base (gou); the slant of the front face is the hypotenuse (xian); the hypotenuse is 5 chi. Multiply the base and the hypotenuse together and multiply by 4, obtaining 60 chi: this is the area of the four visible faces of the square cone. If one lets a round cone (yuan zhui) be contained within it, the visible area of the round cone and the visible area of the square cone have a rate like that of the square's area to the circle's area. Consider: the square cone's bottom is 6 chi, so the perimeter of the square is 24 chi; multiplying it by 5 chi and halving gives, likewise, the visible area of the square cone. Therefore, to find the number for the round cone, fold the diameter in half and with it multiply half of the lower circumference: this is the area of the round cone. Now the domed field's upper diameter is round and vaulted, and yet it uses the same method as the round cone; then the area errs on being too small. However, its method is difficult to use; therefore it roughly gives a general approximation, applying it to large, broad fields. Finding the area of the round cone is like finding the area of a circular field. Now the two wholes are multiplied together, so 4 is taken as the divisor (fa) to divide it, also as for the circular field. The method of extracting the solid circle (kai li yuan) explains the various rates of circles and squares very fully; it can be used to check this.]\end{uniten}

\begin{unitzh}\unitlabel{jiuzhang-006.8}今有弧田，弦三十步，矢十五步。問：爲田幾何？荅曰：一畝九十七步半。\end{unitzh}
\begin{uniten}Suppose there is an arc field (hu tian) with a chord (xian) of 30 bu and an arrow (shi) of 15 bu. Question: How much field does it make? Answer: 1 mu 97 1/2 bu.\end{uniten}

\begin{unitzh}\unitlabel{jiuzhang-006.9}又有弧田，弦七十八步二分步之一，矢十三步九分步之七。問：爲田幾何？荅曰：二畝一百五十五步八十一分步之五十六。\end{unitzh}
\begin{uniten}Again there is an arc field with a chord of 78 1/2 bu and an arrow of 13 7/9 bu. Question: How much field does it make? Answer: 2 mu 155 56/81 bu.\end{uniten}

\begin{unitzh}\unitlabel{jiuzhang-006.10}術曰：以弦乘矢，矢又自乘，并之，二而一。{\zhsmall 方中之圓，圓裏十二觚之羃，合外方之羃四分之三也。方中合外方之半，則朱實合外方四分之一也。弧田，半圓之羃也，故依半圓之體而爲之術。以弦乘矢而半之，則爲黄羃；矢自乘而半之，爲二青羃。青黃相連爲弧體。弧體法當應規，令觚靣不至外畔，失之於少矣。圓田舊術以周三徑一爲率，俱得十二觚之羃，亦失之於少也。與此相似，指驗半圓之弧耳。若不滿半圓者，益復疎闊。宜依句股鋸圓材之術，以弧弦爲鋸道長，以矢爲句深，而求其徑。旣知圓徑，則弧可割分也。割之者，半弧田之弦以爲股，其矢爲句，爲之求弦，卽小弧之弦也。以半小弧之弦爲句，半圓徑爲弧，爲之求股，以減半徑，其餘卽小弦之矢也。割之又割，使至極細。但舉弦矢相乘之數，則必近密率矣。然於算數差繁，必欲有所㝷究也。若但度田，取其大數，舊術爲約耳。}\end{unitzh}
\begin{uniten}Method: Multiply the arrow by the chord, and multiply the arrow by itself as well; add them together and divide by 2. [Liu Hui: The circle within the square: the area of the dodecagon inside the circle amounts to 3/4 of the area of the outer square. The square in the middle amounts to half of the outer square, so the red region (zhu shi) amounts to 1/4 of the outer square. The arc field is the area of half a circle; therefore the method was made for it according to the form of the half circle. Multiplying the arrow by the chord and halving it gives the yellow area (huang mi); multiplying the arrow by itself and halving it gives the two blue-green areas (qing mi). The blue-green and the yellow joined together make the body of the arc (hu ti). The body of the arc ought by rule to fit the circle; with the sides of the polygon falling short of the outer edge, it errs on being too small. The old method for the circular field, taking circumference 3 to diameter 1 as the rate, likewise obtains the area of the dodecagon, and also errs on being too small. This is similar; the verification points only to the arc of a half circle. If it is less than a half circle, it is even more crude and loose. One should follow the base-and-altitude (gougu) method for sawing round timber: take the chord of the arc as the length of the saw-path and the arrow as the base-depth (gou shen), and find its diameter. Once the circle's diameter is known, the arc can be cut and divided. To cut it: halve the chord of the arc field and take it as the altitude, and its arrow as the base; for these find the hypotenuse: this is the chord of the small arc. Take half the chord of the small arc as the base and half the circle's diameter as the arc [Translator's note: 弧 "arc" here is evidently corrupt; the computation requires the hypotenuse]; for these find the altitude and subtract it from the half diameter: the remainder is the arrow of the small chord. Cut it and cut it again, until it reaches the utmost fineness. If one just takes the numbers of the chords and arrows multiplied together, it must come close to the precise rate. However, in computation this is rather cumbersome; [it is for when] one is set on investigating something. If one merely measures fields and takes their rough numbers, the old method is the concise one.]\end{uniten}

\begin{unitzh}\unitlabel{jiuzhang-007.1}今有環田，中周九十二步，外周一百二十二步，徑五步。{\zhsmall 此欲令與周三徑一之率相應，故言徑五步也。據中外周，以徽術言之，當徑四步一百五十七分步之一百二十二也。臣淳風等謹按：依密率，合徑四步二十二分步之十七。}問爲田幾何？荅曰：二畝五十五步。{\zhsmall 於徽術，當爲田二畝三十一步一百五十七分步之二十三。臣淳風等謹依密率，爲田二畝三十步二十二分步之十五。}\end{unitzh}
\begin{uniten}Suppose there is an annular field (huantian) with inner circumference (zhongzhou) 92 bu (步), outer circumference (waizhou) 122 bu, and diameter (jing) 5 bu. [Translator's note: the jing of an annular field is the width of the ring.] [Liu Hui: This is meant to make it agree with the rate (lü) of circumference 3 to diameter 1; therefore it says "diameter 5 bu". Taking the inner and outer circumferences and speaking by Hui's method (Hui shu), the diameter should be 4 122/157 bu.] [Li Chunfeng et al.: Your servants Chunfeng et al. respectfully note: according to the close rate (milü), the diameter comes to 4 17/22 bu.] Question: how large is the field?
Answer: 2 mu (畝) 55 bu. [Liu Hui: By Hui's method, the field should be 2 mu 31 23/157 bu.] [Li Chunfeng et al.: Your servants Chunfeng et al., respectfully following the close rate: the field is 2 mu 30 15/22 bu.]\end{uniten}

\begin{unitzh}\unitlabel{jiuzhang-007.2}又有環田，中周六十二步四分步之三，外周一百一十三步二分步之一，徑十二步三分步之二。{\zhsmall 此田環而不通匝，故徑十二步三分步之二。若據上周求徑者，此徑失之於多，過周三徑一之率，蓋爲疎矣。於徽術，當徑八步六百二十八分步之五十一。臣淳風等謹按：依周三徑一考之，合徑八步二十四分步之一十一；依密率，合徑八步一百七十六分步之一十三。}問爲田幾何？荅曰：四畝一百五十六步四分步之一。{\zhsmall 於徽術，當爲田二畝二百三十二步五千二十四分步之七百八十七也。依周三徑一，爲田三畝二十五步六十四分步之二十五。臣淳風等謹按：密率，爲田二畝二百三十一步一千四百八分步之七百一十七也。}\end{unitzh}
\begin{uniten}Again, suppose there is an annular field with inner circumference 62 3/4 bu, outer circumference 113 1/2 bu, and diameter 12 2/3 bu. [Liu Hui: This field is ring-shaped but does not go all the way round; therefore the diameter is 12 2/3 bu. If the diameter is sought from the circumferences above, this diameter errs on the side of too much, going beyond the rate of circumference 3 to diameter 1; presumably it is loose (shu). By Hui's method, the diameter should be 8 51/628 bu.] [Li Chunfeng et al.: Your servants Chunfeng et al. respectfully note: examined by circumference 3, diameter 1, the diameter comes to 8 11/24 bu; by the close rate, the diameter comes to 8 13/176 bu.] Question: how large is the field?
Answer: 4 mu 156 1/4 bu. [Liu Hui: By Hui's method, the field should be 2 mu 232 787/5024 bu. By circumference 3, diameter 1, the field is 3 mu 25 25/64 bu.] [Li Chunfeng et al.: Your servants Chunfeng et al. respectfully note: by the close rate, the field is 2 mu 231 717/1408 bu.]\end{uniten}

\begin{unitzh}\unitlabel{jiuzhang-007.3}術曰：并中外周而半之，以徑乘之，爲積步。{\zhsmall 此田截齊中外之周，周則爲長；并而半之者，亦以盈補虚也。此可令中外周各自爲圓田，以中圓減外圓，餘則環實也。按此術：置中外周步數於上，分母子於下。母乘子者，爲中外周俱有餘分，故以互乘齊其子，母相乘同其母。子齊母同，故通全步，內分子。并而半之者，以盈補虚，得中平之周。周則爲從，徑則爲廣，故廣從相乘而得其積。旣合分母，還須分母出之，故令周徑分母相乘而連除之，卽得積步。不盡，以等數除之而命分。以畝法除積步，得畝數也。}\end{unitzh}
\begin{uniten}Method: Add the inner and outer circumferences and halve; multiply by the diameter; this gives the area in bu (jibu). [Liu Hui: In this field, cut and even up the inner and outer circumferences; the circumference then becomes the length (chang). Adding and halving is again using the full to fill the empty. Here one may let the inner and outer circumferences each form a circular field; subtract the inner circle from the outer circle, and the remainder is the content (shi) of the ring. According to this method, place the numbers of bu of the inner and outer circumferences above, and the denominators and numerators below. The denominators multiply the numerators because the inner and outer circumferences both have fractional remainders; therefore one cross-multiplies (hucheng) to homogenize (qi) their numerators, and multiplies the denominators together to equalize (tong) their denominators. The numerators being homogenized and the denominators equalized, one therefore converts (tong) the whole bu and incorporates (na) the numerators. Adding and halving uses the full to fill the empty and gives the averaged (zhongping) circumference. The circumference then serves as the length (zong) and the diameter as the width (guang); therefore width and length multiplied together give its area. Since the denominators have been combined in, they must in turn be taken out by the denominators; therefore multiply together the denominators of the circumference and the diameter and divide by them jointly, which gives the area in bu. If it does not divide exactly, divide by the equal number (dengshu) and name the fraction. Divide the area in bu by the mu divisor (mufa) to obtain the number of mu.]\end{uniten}

\begin{unitzh}\unitlabel{jiuzhang-007.4}密率術曰：置中外周步數，分母子各居其下，母互乘子，分母相乘，通全步，內分子，并而半之。又可以中周減外周，餘半之，以益中周。徑亦通分內子，以乘周爲實，分母相乘爲法，除之爲積步。餘積步之分，等數約之。以畝法除之，卽畝數也。\end{unitzh}
\begin{uniten}Close-rate method: Set out the numbers of bu of the inner and outer circumferences, each with its denominator and numerator placed below; the denominators cross-multiply the numerators; multiply the denominators together; convert the whole bu and incorporate the numerators; add and halve. Alternatively, one may subtract the inner circumference from the outer circumference, halve the remainder, and add it to the inner circumference. The diameter likewise is converted into fractional parts with its numerator incorporated; multiply the circumference by it to form the dividend (shi); multiply the denominators together to form the divisor (fa); divide, which gives the area in bu. The fraction of an area bu left over is reduced by the equal number. Divide by the mu divisor: this is the number of mu.\end{uniten}

\begin{unitzh}\unitlabel{jiuzhang-007.5}九章算術卷一\end{unitzh}
\begin{uniten}The Nine Chapters on the Mathematical Art, Chapter 1\end{uniten}

\begin{unitzh}\unitlabel{jiuzhang-007.6}九章算術卷一訂訛補圖{\zhsmall 算經十書之二} 休寧戴震東原\end{unitzh}
\begin{uniten}The Nine Chapters on the Mathematical Art, Chapter 1: Corrections of Errors and Supplementary Diagrams [Note: second of the Ten Mathematical Classics]
Dai Zhen, Dongyuan, of Xiuning\end{uniten}

\begin{unitzh}\unitlabel{jiuzhang-007.7}臣淳風等謹按：母互乘子，副并爲平實者，定此平實立限，衆子所當損益，如限爲平。{\zhsmall 據首問：第二數母三、第三數母四互乘第一數子一，得十二；第一數母三、第三數母四互乘第二數子二，得二十四；第一、第二數母各三互乘第三數子三，得二十七；并之，共六十三，爲平實。母三三相乘，又與四乘，得三十六，爲法。列數凡三，卽以三乘十二，得三十六；乘二十四，得七十二；乘二十七，得八十一，爲列實。亦以三乘法三十六，得一百八。平實六十三減列實三十六，少二十七；減七十二，餘九；減八十一，餘十八。約之，九爲一，則十八爲二，而二十七爲三，平實六十三爲七，法一百八爲十二，命爲十二分之七。設以十二作三數，三分之一則四也，三分之二則八也，四分之三則九也。定平實七立限，八減一、九減二皆七，所減之一、二益於四亦七，損多益少，適如其限，故云「定此平實立限」，又云「如限爲平」。原本「立」訛作「主」，「如」訛作「知」，遂不可通。}\end{unitzh}
\begin{uniten}Your servants Chunfeng et al. respectfully note: The denominators cross-multiply the numerators, and these are added in an auxiliary position (fu) to form the level dividend (pingshi): this fixes the level dividend and sets up the limit (xian); what the various numerators should lose or gain is such that, matching the limit, they become level. [Dai Zhen: According to the first problem: the denominator 3 of the second number and the denominator 4 of the third number cross-multiply the numerator 1 of the first number, giving 12; the denominator 3 of the first number and the denominator 4 of the third number cross-multiply the numerator 2 of the second number, giving 24; the denominators of the first and second numbers, 3 each, cross-multiply the numerator 3 of the third number, giving 27; adding them, the total 63 is the level dividend. The denominators 3 and 3 multiplied together and further multiplied by 4 give 36, which is the divisor. The listed numbers are 3 in all; then multiply 12 by 3, giving 36; 24, giving 72; 27, giving 81: these are the listed dividends (lieshi). Also multiply the divisor 36 by 3, giving 108. The level dividend 63, subtracted from the listed dividend 36, falls short by 27; subtracted from 72, it leaves 9; subtracted from 81, it leaves 18. Reducing them: 9 becomes 1, so 18 becomes 2 and 27 becomes 3; the level dividend 63 becomes 7 and the divisor 108 becomes 12; name it 7/12. Suppose the three numbers are made with 12: 1/3 is then 4, 2/3 is then 8, 3/4 is then 9. Fixing the level dividend 7 sets up the limit: 8 less 1 and 9 less 2 are both 7, and the 1 and 2 subtracted, added to 4, also make 7. Decreasing the larger and increasing the smaller exactly matches the limit; hence it says "this fixes the level dividend and sets up the limit", and also says "matching the limit, they become level". In the original text 立 "set up" is miswritten as 主, and 如 "match" as 知, so that it could not be understood.]\end{uniten}

\begin{unitzh}\unitlabel{jiuzhang-007.8}按：半周爲從，半徑爲廣，故廣從相乘爲積步也。\end{unitzh}
\begin{uniten}Note: half the circumference serves as the length and half the diameter as the width; therefore width and length multiplied together give the area in bu.\end{uniten}

\begin{unitzh}\unitlabel{jiuzhang-007.9}假令圓徑二尺，圓中容六觚之一面{\zhsmall 「六觚」原本訛作「六弧」。考六角形，其平而亦有六；八角形，其平面亦有八；古人謂之六觚、八觚。若截圓形爲六，古人謂之弧背，其弧卽圓周，不得云「圓中容六弧之一面」。後或言「弧」，或言「觚」，義各不同。原本「觚」皆訛作「弧」，遂𫎇混不可通。}\end{unitzh}
\begin{uniten}Suppose the diameter of a circle is 2 chi (尺); one side of the hexagon (liu gu) inscribed in the circle [Dai Zhen: In the original text 六觚 "hexagon" is miswritten as 六弧 "six arcs". Consider: the six-angled figure (liujiao xing) likewise has six flat faces; the eight-angled figure likewise has eight flat faces; the ancients called them liu gu and ba gu. If a circle is cut into six, the ancients called this the arc-back (hubei); its arc is the circumference of the circle, and one cannot say "one side of the six arcs inscribed in the circle". Further on, it sometimes says 弧 and sometimes 觚, and the meanings are each different. In the original text 觚 is everywhere miswritten as 弧, so that the two are confounded and it cannot be understood.] [Translator's note: 平而 in 其平而亦有六 is evidently for 平面 "flat faces".]\end{uniten}

\begin{unitzh}\unitlabel{jiuzhang-007.10}圓內容六觚之圖 □□□□□□□□□□□□□□□□□ □□□□□□□□□□□□□□□□□ {\zhsmall 劉徽以圓田用周三徑一之率，周三者，從其六觚之環耳，以推圓規多少之較，乃弓之與弦也。六觚之一面與圓徑之半，其數均等，叠雨圓觀之，疎密顯然矣。}\end{unitzh}
\begin{uniten}Diagram of a hexagon inscribed in a circle
[illegible] ×17; [illegible] ×17
[Dai Zhen: Liu Hui takes it that, where the circular field uses the rate of circumference 3 to diameter 1, the "circumference 3" merely follows the ring of its hexagon; inferring from this the difference of more or less with respect to the circle (yuangui), it is that of the bow and its string. One side of the hexagon and half the diameter of the circle are equal in number; superimposing the two circles and viewing them, the looseness and closeness are evident.] [Translator's note: 雨 "rain" in 叠雨圓 is evidently for 兩 "two".]\end{uniten}

\begin{unitzh}\unitlabel{jiuzhang-007.11}以九十六觚之幂減之，餘六百二十五分寸之以百五，謂之差幂；倍之，爲分寸之二百一十。{\zhsmall 「爲分寸」者𫎇上省文，謂六百二十五分寸之二百一十也。}\end{unitzh}
\begin{uniten}Subtract the area (mi) of the 96-gon from it; the remainder, 105/625 cun (寸), is called the difference area (chami); doubling it makes 210 parts of a cun. [Dai Zhen: "Makes parts of a cun" is an abbreviated wording carried over from above; it means 210/625 cun.] [Translator's note: 105 is written 以百五, evidently for 一百五.]\end{uniten}

\begin{unitzh}\unitlabel{jiuzhang-007.12}以此術求之，得幂一百六十一寸有奇，其數相近矣。此術微少，而斛差幂六百二十五分寸之一百五，以十二觚之幂爲率消息，當取此分寸之三十六。{\zhsmall 「取此分寸」亦𫎇上省文，謂六百二十五分寸之三十六也。}\end{unitzh}
\begin{uniten}Seeking it by this method, one obtains an area of 161 cun and something over; the numbers are close. This method is slightly small, and the hu difference area is 105/625 cun; taking the area of the 12-gon as the rate for adjustment (xiaoxi), one should take 36 of these parts of a cun. [Dai Zhen: "Take of these parts of a cun" is likewise an abbreviated wording carried over from above; it means 36/625 cun.] [Translator's note: 斛 "hu" is evidently for 觚 "gu".]\end{uniten}

\begin{unitzh}\unitlabel{jiuzhang-007.13}臣淳風等謹按：依密率，以七乘周，二十二而一，卽徑；以二十二乘徑，七而一，卽周。依術求之，卽得。{\zhsmall 「徑七周二十二」，乃祖氏之約率，非密率也。淳風等以爲密率，失其實矣。徽率與祖氏之約率，率}□{\zhsmall 較於}□{\zhsmall 徽率}□□□□□□□□□□□□□ □□□□□□□□□□□□□□□□□□\end{unitzh}
\begin{uniten}Your servants Chunfeng et al. respectfully note: by the close rate, multiply the circumference by 7 and divide by 22: this is the diameter; multiply the diameter by 22 and divide by 7: this is the circumference. Seeking it by the method, one obtains it. [Dai Zhen: Diameter 7, circumference 22 is Master Zu's approximate rate (yuelü), not the close rate. Chunfeng et al. took it to be the close rate, which misses the truth of it. Hui's rate and Master Zu's approximate rate, rate] [illegible] [compared with] [illegible] [Hui's rate] [illegible] ×31 [continues]\end{uniten}

\section*{卷第二}\addcontentsline{toc}{section}{卷第二}

\begin{unitzh}\unitlabel{jiuzhang-008.1}九章算術卷第二{\zhsmall 算經十書之二} 魏劉徽注 唐朝議大夫行太史令上輕車都尉臣李淳風等奉勅注釋\end{unitzh}
\begin{uniten}The Nine Chapters on the Mathematical Art, juan 2 [Note: The second of the Ten Mathematical Classics.] Commentary by Liu Hui of the Wei. Of the Tang: Grand Master for Court Discussion, acting Grand Astrologer, Senior Commandant of Light Chariots, your servant Li Chunfeng et al., annotated by imperial command.\end{uniten}

\begin{unitzh}\unitlabel{jiuzhang-008.2}粟米{\zhsmall 以御交質變易}\end{unitzh}
\begin{uniten}Millet and Grain (sumi) [Liu Hui: For dealing with exchanges and conversions.]\end{uniten}

\begin{unitzh}\unitlabel{jiuzhang-008.3}粟米之法{\zhsmall 凡此諸率相與大通，其特相求，各如本率。可約者約之，别術然也。}粟率五十，糲米三十，粺米二十七，糳米二十四，御米二十一，小䵂十三半，大䵂五十四，糲飯七十五，粺飯五十四，糳飯四十八，御飯四十二，菽、荅、麻、麥各四十五，稻六十，䜴六十三，飱九十，熟菽一百三半，糵一百七十五。\end{unitzh}
\begin{uniten}The rule of millet and grain [Liu Hui: All these rates (lü) communicate broadly with one another; when particular ones are sought from one another, each follows its original rate. Those that can be reduced (yue) are reduced; this is how the separate methods are.] Millet (su) rate 50; coarse grain (limi) 30; fine grain (baimi) 27; polished grain (zuomi) 24; imperial grain (yumi) 21; small groats (xiaozhi) 13 1/2; large groats (dazhi) 54; cooked coarse grain (lifan) 75; cooked fine grain (baifan) 54; cooked polished grain (zuofan) 48; cooked imperial grain (yufan) 42; soybeans (shu), small beans (da), hemp seed (ma) and wheat (mai), each 45; paddy rice (dao) 60; fermented beans (chi) 63; soaked cooked grain (sun) 90; cooked soybeans (shushu) 103 1/2; malt (nie) 175.\end{uniten}

\begin{unitzh}\unitlabel{jiuzhang-008.4}今有{\zhsmall 此都術也。凡九數以爲篇名，可以廣施諸率，所謂告往而知來，舉一隅而三隅反者也。誠能分詭數之紛雜，通彼此之否塞，因物成率，審辨名分，平其偏頗，齊其參差，則終無不歸於此術也。}\end{unitzh}
\begin{uniten}“Suppose there is” (jinyou) [Liu Hui: This is the general method (dushu). All the nine computations (jiushu), taken as the titles of the chapters, can apply the rates broadly; this is what is called “told what has gone, one knows what is to come” and “one corner being raised, one comes back with three corners”. If one can truly sort out the confused mixture of deceptive numbers, clear the blockages between this and that, form rates according to the things, carefully distinguish names and parts, level their unevenness and even out their irregularities, then in the end there is nothing that does not come back to this method.]\end{uniten}

\begin{unitzh}\unitlabel{jiuzhang-008.5}術曰：以所有數乘所求率爲實，以所有率爲法。{\zhsmall 少者多之始，一者數之母，故爲率者必等之於一。據粟率五、糲率三，是粟五而爲一，糲米三而爲一也。欲化粟爲米者，糲當先本是一。一者，謂以五約之，令五而爲一也。訖，乃以三乘之，令一而爲三。如是，則率等於一，以五爲三矣。然先除後乘，或有餘分，故術反之。又究言之，知粟五升爲糲米三升；分言之，知粟一斗爲糲米五分斗之三。以五爲母，三爲子。以粟求糲米者，以子乘，其母報除也。然則所求之率常爲母也。臣淳風等謹按：宜云「所求之率常爲子，所有之率常爲母」，今乃云「所求之率常爲母」，知脫錯也。}實如法而一。\end{unitzh}
\begin{uniten}Method: With the number one has (suoyou shu), multiply the rate of what is sought (suoqiu lü), making the dividend (shi); take the rate of what one has (suoyou lü) as the divisor (fa). [Liu Hui: The few is the beginning of the many; 1 is the mother of numbers; therefore whoever makes rates must equate them to 1. Take the millet rate 5 and the coarse rate 3: this means that 5 of millet makes 1, and 3 of coarse grain makes 1. For one who wishes to convert millet into grain, the coarse grain [Translator's note: 糲 as printed, although the reduction by 5 below applies to the millet.] should first be based on 1. “1” means reducing it by 5, making 5 into 1. That done, multiply it by 3, making 1 into 3. In this way the rates are equated to 1, and 5 is made into 3. However, if one divides first and multiplies afterwards, there may be leftover fractions; therefore the method reverses this. Moreover, speaking of it in full, one knows that 5 sheng (升) of millet make 3 sheng of coarse grain; speaking of it in parts, one knows that 1 dou (斗) of millet makes 3/5 dou of coarse grain. Take 5 as the denominator (mu) and 3 as the numerator (zi). In seeking coarse grain from millet, one multiplies by the numerator, and its denominator divides in return. Thus the rate of what is sought is always the denominator.] [Li Chunfeng et al.: Your servants respectfully note: it should say “the rate of what is sought is always the numerator; the rate of what one has is always the denominator”. Now it says instead “the rate of what is sought is always the denominator”; we know that something has been dropped and miswritten.] Divide the dividend by the divisor.\end{uniten}

\begin{unitzh}\unitlabel{jiuzhang-008.6}今有粟一斗，欲爲糲米。問得幾何？荅曰：爲糲米六升。術曰：以粟求糲米，三之，五而一。{\zhsmall 臣淳風等謹按：都術以所求率乘有所數，以所有率爲法。此術以粟求米，故粟爲所有數；三是米率，故三爲所求率；五是粟率，故五爲所有率。粟率五十，米率三十，退位求之，故唯云三、五也。}\end{unitzh}
\begin{uniten}Suppose there is 1 dou of millet, which one wishes to make into coarse grain. Question: how much is obtained? Answer: It makes 6 sheng of coarse grain. Method: To seek coarse grain from millet, multiply by 3 and divide by 5. [Li Chunfeng et al.: Your servants respectfully note: the general method uses the rate of what is sought to multiply the number one has [Translator's note: written 有所數, inverted.], and takes the rate of what one has as the divisor. In this method grain is sought from millet; therefore the millet is the number one has. 3 is the grain rate; therefore 3 is the rate of what is sought. 5 is the millet rate; therefore 5 is the rate of what one has. The millet rate is 50 and the grain rate is 30; one seeks with them moved back one place (tuiwei); therefore it says only 3 and 5.]\end{uniten}

\begin{unitzh}\unitlabel{jiuzhang-008.7}今有粟二斗一升，欲爲粺米。問得幾何？荅曰：爲粺米一斗一升五十分升之十七。術曰：以粟求粺米，二十七之，五十而一。{\zhsmall 臣淳風等謹按：粺米之率二十有七，故直以二十七之，五十而一也。}\end{unitzh}
\begin{uniten}Suppose there is 2 dou 1 sheng of millet, which one wishes to make into fine grain. Question: how much is obtained? Answer: It makes 1 dou 1 17/50 sheng of fine grain. Method: To seek fine grain from millet, multiply by 27 and divide by 50. [Li Chunfeng et al.: Your servants respectfully note: the rate of fine grain is 27; therefore one directly multiplies by 27 and divides by 50.]\end{uniten}

\begin{unitzh}\unitlabel{jiuzhang-008.8}今有粟四斗五升，欲爲糳米。問得幾何？荅曰：爲糳米二斗一升五分升之三。術曰：以粟求糳米，十二之，二十五而一。{\zhsmall 臣淳風等謹按：糳米之率二十有四，以爲率太䌓，故因而半之。半所求之率，以乘所有之數；所求之率旣減半，所有之率亦減半，是故十二乘之，二十五而一也。}\end{unitzh}
\begin{uniten}Suppose there is 4 dou 5 sheng of millet, which one wishes to make into polished grain. Question: how much is obtained? Answer: It makes 2 dou 1 3/5 sheng of polished grain. Method: To seek polished grain from millet, multiply by 12 and divide by 25. [Li Chunfeng et al.: Your servants respectfully note: the rate of polished grain is 24; taking this as the rate is too cumbersome, so one accordingly halves it. Halve the rate of what is sought and with it multiply the number one has; since the rate of what is sought has been reduced by half, the rate of what one has is also reduced by half. Therefore one multiplies by 12 and divides by 25.]\end{uniten}

\begin{unitzh}\unitlabel{jiuzhang-008.9}今有粟七斗九升，欲爲御米。問得幾何？荅曰：爲御米三斗三升五十分升之九。術曰：以粟求御米，二十一之，五十而一。\end{unitzh}
\begin{uniten}Suppose there is 7 dou 9 sheng of millet, which one wishes to make into imperial grain. Question: how much is obtained? Answer: It makes 3 dou 3 9/50 sheng of imperial grain. Method: To seek imperial grain from millet, multiply by 21 and divide by 50.\end{uniten}

\begin{unitzh}\unitlabel{jiuzhang-008.10}今有粟一斗，欲爲小䵂。問得幾何？荅曰：爲小䵂二升一十分升之七。術曰：以粟求小䵂，二十七之，百而一。{\zhsmall 臣淳風等謹按：小䵂之率十三有半。半者二爲母，以二通之，得二十七，爲所求率；又以母二通其粟率，得一百，爲所有率。凡本率有分者，須卽乘除也。他皆倣此。}\end{unitzh}
\begin{uniten}Suppose there is 1 dou of millet, which one wishes to make into small groats. Question: how much is obtained? Answer: It makes 2 7/10 sheng of small groats. Method: To seek small groats from millet, multiply by 27 and divide by 100. [Li Chunfeng et al.: Your servants respectfully note: the rate of small groats is 13 1/2. For the half, 2 is the denominator; converting (tong) it by 2 gives 27, which is the rate of what is sought. Again, converting the millet rate by the denominator 2 gives 100, which is the rate of what one has. Whenever an original rate has a fraction, one must accordingly multiply and divide. All the others follow this.]\end{uniten}

\begin{unitzh}\unitlabel{jiuzhang-008.11}今有粟九斗八升，欲爲大䵂。問得幾何？荅曰：爲大䵂一十斗五升二十五分升之二十一。術曰：以粟求大䵂，二十七之，二十五而一。{\zhsmall 臣淳風等謹按：大䵂之率五十有四，因其可半，故二十七之，亦如粟求糳米，半其二率。}\end{unitzh}
\begin{uniten}Suppose there is 9 dou 8 sheng of millet, which one wishes to make into large groats. Question: how much is obtained? Answer: It makes 10 dou 5 21/25 sheng of large groats. Method: To seek large groats from millet, multiply by 27 and divide by 25. [Li Chunfeng et al.: Your servants respectfully note: the rate of large groats is 54; because it can be halved, one therefore multiplies by 27. This is also as in seeking polished grain from millet: halve its two rates.]\end{uniten}

\begin{unitzh}\unitlabel{jiuzhang-008.12}今有粟二斗三升，欲爲糲飯。問得幾何？荅曰：爲糲飯三斗四升半。術曰：以粟求糲飯，三之，二而一。{\zhsmall 臣淳風等謹按：糲飯之率七十有五，粟求糲飯，合以此數乘之。今以等數二十有五約其二率，所求之率得三，所有之率得二，故以三乘、二除。}\end{unitzh}
\begin{uniten}Suppose there is 2 dou 3 sheng of millet, which one wishes to make into cooked coarse grain. Question: how much is obtained? Answer: It makes 3 dou 4 1/2 sheng of cooked coarse grain. Method: To seek cooked coarse grain from millet, multiply by 3 and divide by 2. [Li Chunfeng et al.: Your servants respectfully note: the rate of cooked coarse grain is 75; in seeking cooked coarse grain from millet, one ought to multiply by this number. Now, reducing its two rates by the equal number (dengshu) 25, the rate of what is sought becomes 3 and the rate of what one has becomes 2; therefore one multiplies by 3 and divides by 2.]\end{uniten}

\begin{unitzh}\unitlabel{jiuzhang-008.13}今有粟三斗六升，欲爲粺飯。問得幾何？荅曰：爲粺飯三斗八升二十五分升之二十二。術曰：以粟求粺飯，二十七之，二十五而一。{\zhsmall 臣淳風等謹按：此術與大䵂多同。}\end{unitzh}
\begin{uniten}Suppose there is 3 dou 6 sheng of millet, which one wishes to make into cooked fine grain. Question: how much is obtained? Answer: It makes 3 dou 8 22/25 sheng of cooked fine grain. Method: To seek cooked fine grain from millet, multiply by 27 and divide by 25. [Li Chunfeng et al.: Your servants respectfully note: this method is largely the same as that for large groats.]\end{uniten}

\begin{unitzh}\unitlabel{jiuzhang-008.14}今有粟八斗六升，欲爲糳飯。問得幾何？荅曰：爲糳飯八斗二升二十五分升之一十四。術曰：以粟求糳飯，二十四之，二十五而一。{\zhsmall 臣淳風等謹按：糳飯率四十八，此亦半二率而乘除。}\end{unitzh}
\begin{uniten}Suppose there is 8 dou 6 sheng of millet, which one wishes to make into cooked polished grain. Question: how much is obtained? Answer: It makes 8 dou 2 14/25 sheng of cooked polished grain. Method: To seek cooked polished grain from millet, multiply by 24 and divide by 25. [Li Chunfeng et al.: Your servants respectfully note: the rate of cooked polished grain is 48; here too one halves the two rates and then multiplies and divides.]\end{uniten}

\begin{unitzh}\unitlabel{jiuzhang-008.15}今有粟九斗八升，欲爲御飯。問得幾何？荅曰：爲御飯八斗二升二十五分升之八。術曰：以粟求御飯，二十一之，二十五而一。{\zhsmall 臣淳風等謹按：此術半率，亦與糳飯多同。}\end{unitzh}
\begin{uniten}Suppose there is 9 dou 8 sheng of millet, which one wishes to make into cooked imperial grain. Question: how much is obtained? Answer: It makes 8 dou 2 8/25 sheng of cooked imperial grain. Method: To seek cooked imperial grain from millet, multiply by 21 and divide by 25. [Li Chunfeng et al.: Your servants respectfully note: in this method the halving of the rates is also largely the same as for cooked polished grain.]\end{uniten}

\begin{unitzh}\unitlabel{jiuzhang-008.16}今有粟三斗少半升，欲爲菽。問得幾何？荅曰：爲菽二斗七升一十分升之三。\end{unitzh}
\begin{uniten}Suppose there is 3 dou 1/3 sheng of millet, which one wishes to make into soybeans. Question: how much is obtained? Answer: It makes 2 dou 7 3/10 sheng of soybeans.\end{uniten}

\begin{unitzh}\unitlabel{jiuzhang-008.17}今有粟四斗一升太半升，欲爲荅。問得幾何？荅曰：爲荅三斗七升半。\end{unitzh}
\begin{uniten}Suppose there is 4 dou 1 2/3 sheng of millet, which one wishes to make into small beans. Question: how much is obtained? Answer: It makes 3 dou 7 1/2 sheng of small beans.\end{uniten}

\begin{unitzh}\unitlabel{jiuzhang-008.18}今有粟五斗太半升，欲爲麻。問得幾何？荅曰：爲麻四斗五升五分升之三。\end{unitzh}
\begin{uniten}Suppose there is 5 dou 2/3 sheng of millet, which one wishes to make into hemp seed. Question: how much is obtained? Answer: It makes 4 dou 5 3/5 sheng of hemp seed.\end{uniten}

\begin{unitzh}\unitlabel{jiuzhang-008.19}今有粟一十斗八升五分升之二，欲爲麥。問得幾何？荅曰：爲麥九斗七升二十五分升之一十四。\end{unitzh}
\begin{uniten}Suppose there is 10 dou 8 2/5 sheng of millet, which one wishes to make into wheat. Question: how much is obtained? Answer: It makes 9 dou 7 14/25 sheng of wheat.\end{uniten}

\begin{unitzh}\unitlabel{jiuzhang-008.20}術曰：以粟求菽、荅、麻、麥，皆九之，十而一。{\zhsmall 臣淳風等謹按：四術率並四十五，皆是爲粟所求，俱合以此率乘其本粟。術欲從省，先以等數五約之，所求之率得九，所有之率得十，故九乘十除，義由於此。}\end{unitzh}
\begin{uniten}Method: To seek soybeans, small beans, hemp seed or wheat from millet, in every case multiply by 9 and divide by 10. [Li Chunfeng et al.: Your servants respectfully note: the rates of the four methods are all 45; all are what is sought from the millet, and all ought to multiply their original millet by this rate. The method wishes to follow economy: first reduce them by the equal number 5; the rate of what is sought becomes 9 and the rate of what one has becomes 10. Hence one multiplies by 9 and divides by 10; the reasoning derives from this.]\end{uniten}

\begin{unitzh}\unitlabel{jiuzhang-009.1}今有粟七斗五升七分升之四，欲爲稻，問得幾何？荅曰：爲稻九斗三十五分升之二十四。\end{unitzh}
\begin{uniten}Suppose there are 7 dou (斗) 5 4/7 sheng (升) of millet (su), which one wishes to make into paddy (dao). Question: how much does one get? Answer: it makes 9 dou 24/35 sheng of paddy.\end{uniten}

\begin{unitzh}\unitlabel{jiuzhang-009.2}術曰：以粟求稻，六之，五而一。{\zhsmall 臣淳風等謹按：稻率六十，亦約二率而乘除。}\end{unitzh}
\begin{uniten}Method: To seek paddy from millet, multiply by 6 and divide by 5. [Li Chunfeng et al.: The rate (lü) of paddy is 60; here too one reduces (yue) the two rates and then multiplies and divides.]\end{uniten}

\begin{unitzh}\unitlabel{jiuzhang-009.3}今有粟七斗八升，欲爲䜴，問得幾何？荅曰：爲䜴九斗八升二十五分升之七。\end{unitzh}
\begin{uniten}Suppose there are 7 dou 8 sheng of millet, which one wishes to make into fermented beans (chi). Question: how much does one get? Answer: it makes 9 dou 8 7/25 sheng of fermented beans.\end{uniten}

\begin{unitzh}\unitlabel{jiuzhang-009.4}術曰：以粟求䜴，六十三之，五十而一。\end{unitzh}
\begin{uniten}Method: To seek fermented beans from millet, multiply by 63 and divide by 50.\end{uniten}

\begin{unitzh}\unitlabel{jiuzhang-009.5}今有粟五斗五升，欲爲飱，問得幾何？荅曰：爲飱九斗九升。\end{unitzh}
\begin{uniten}Suppose there are 5 dou 5 sheng of millet, which one wishes to make into soaked cooked grain (sun). Question: how much does one get? Answer: it makes 9 dou 9 sheng of soaked cooked grain.\end{uniten}

\begin{unitzh}\unitlabel{jiuzhang-009.6}術曰：以粟求飱，九之，五而一。{\zhsmall 臣淳風等謹按：飱率九十，退位與求稻多同。}\end{unitzh}
\begin{uniten}Method: To seek soaked cooked grain from millet, multiply by 9 and divide by 5. [Li Chunfeng et al.: The rate of soaked cooked grain is 90; moving back a place (tuiwei), it is largely the same as seeking paddy.]\end{uniten}

\begin{unitzh}\unitlabel{jiuzhang-009.7}今有粟四斗，欲爲熟菽，問得幾何？荅曰：爲熟菽八斗二升五分升之四。\end{unitzh}
\begin{uniten}Suppose there are 4 dou of millet, which one wishes to make into cooked beans (shushu). Question: how much does one get? Answer: it makes 8 dou 2 4/5 sheng of cooked beans.\end{uniten}

\begin{unitzh}\unitlabel{jiuzhang-009.8}術曰：以粟求熟菽，二百七之，百而一。{\zhsmall 臣淳風等謹按：熟菽之率一百三半。半者，其母二，故以母二通之。所求之率旣被二乘，所有之率隨而俱長，故以二百七之，百而一。}\end{unitzh}
\begin{uniten}Method: To seek cooked beans from millet, multiply by 207 and divide by 100. [Li Chunfeng et al.: The rate of cooked beans is 103 1/2. As for the half, its denominator (mu) is 2; therefore one converts (tong) it by the denominator 2. Since the sought rate (suoqiu lü) has been multiplied by 2, the given rate (suoyou lü) grows along with it; therefore one multiplies by 207 and divides by 100.]\end{uniten}

\begin{unitzh}\unitlabel{jiuzhang-009.9}今有粟二斗，欲爲糵，問得幾何？荅曰：爲糵七斗。\end{unitzh}
\begin{uniten}Suppose there are 2 dou of millet, which one wishes to make into malt (nie). Question: how much does one get? Answer: it makes 7 dou of malt.\end{uniten}

\begin{unitzh}\unitlabel{jiuzhang-009.10}術曰：以粟求糵，七之，二而一。{\zhsmall 臣淳風等謹按：糵率一百七十有五，合以此數乘其本粟。術欲從省，先以等數二十五約之，所求之率得七，所有之率得二，故七乘二除。}\end{unitzh}
\begin{uniten}Method: To seek malt from millet, multiply by 7 and divide by 2. [Li Chunfeng et al.: The rate of malt is 175; one should multiply the original millet by this number. The method wishes to economize: first reduce them by the equal number (dengshu) 25; the sought rate gives 7 and the given rate gives 2; therefore multiply by 7 and divide by 2.]\end{uniten}

\begin{unitzh}\unitlabel{jiuzhang-009.11}今有糲米十五斗五升五分升之二，欲爲粟，問得幾何？荅曰：爲粟二十五斗九升。\end{unitzh}
\begin{uniten}Suppose there are 15 dou 5 2/5 sheng of coarse hulled grain (limi), which one wishes to make into millet. Question: how much does one get? Answer: it makes 25 dou 9 sheng of millet.\end{uniten}

\begin{unitzh}\unitlabel{jiuzhang-009.12}術曰：以糲米求粟，五之，三而一。{\zhsmall 臣淳風等謹按：上術以粟求米，故粟爲所有數，三爲所求率，五爲所有率；今此以米求粟，故米爲所有數，五爲所求率，三爲所有率。准都術求之，各合其數。以下所有反求多同，皆准此。}\end{unitzh}
\begin{uniten}Method: To seek millet from coarse hulled grain, multiply by 5 and divide by 3. [Li Chunfeng et al.: The method above seeks hulled grain (mi) from millet, so millet is the given number (suoyou shu), 3 is the sought rate and 5 is the given rate; now here one seeks millet from hulled grain, so hulled grain is the given number, 5 is the sought rate and 3 is the given rate. Seeking them according to the general method (dushu), each agrees with its numbers. All the reverse seekings (fanqiu) below are largely the same; all follow this.]\end{uniten}

\begin{unitzh}\unitlabel{jiuzhang-009.13}今有粺米二斗，欲爲粟，問得幾何？荅曰：爲粟三斗七升二十七分升之一。\end{unitzh}
\begin{uniten}Suppose there are 2 dou of fine hulled grain (baimi), which one wishes to make into millet. Question: how much does one get? Answer: it makes 3 dou 7 1/27 sheng of millet.\end{uniten}

\begin{unitzh}\unitlabel{jiuzhang-009.14}術曰：以粺米求粟，五十之，二十七而一。\end{unitzh}
\begin{uniten}Method: To seek millet from fine hulled grain, multiply by 50 and divide by 27.\end{uniten}

\begin{unitzh}\unitlabel{jiuzhang-009.15}今有糳米三斗少半升，欲爲粟，問得幾何？荅曰：爲粟六斗三升三十六分升之七。\end{unitzh}
\begin{uniten}Suppose there are 3 dou 1/3 sheng of refined hulled grain (zuomi), which one wishes to make into millet. Question: how much does one get? Answer: it makes 6 dou 3 7/36 sheng of millet.\end{uniten}

\begin{unitzh}\unitlabel{jiuzhang-009.16}術日：以糳米求粟，二十五之，十三而一。\end{unitzh}
\begin{uniten}Method: To seek millet from refined hulled grain, multiply by 25 and divide by 13. [Translator's note: the text writes 術日 for 術曰; the answer above corresponds to division by 12.]\end{uniten}

\begin{unitzh}\unitlabel{jiuzhang-009.17}今有御米十四斗，欲爲粟，問得幾何？荅曰：爲粟三十三斗三升少半升。\end{unitzh}
\begin{uniten}Suppose there are 14 dou of imperial hulled grain (yumi), which one wishes to make into millet. Question: how much does one get? Answer: it makes 33 dou 3 1/3 sheng of millet.\end{uniten}

\begin{unitzh}\unitlabel{jiuzhang-009.18}術曰：以御米求粟，五十之，二十一而一。\end{unitzh}
\begin{uniten}Method: To seek millet from imperial hulled grain, multiply by 50 and divide by 21.\end{uniten}

\begin{unitzh}\unitlabel{jiuzhang-009.19}今有稻一十二斗六升一十五分升之一十四，欲爲粟，問得幾何？荅曰：爲粟一十斗五升九分升之七。\end{unitzh}
\begin{uniten}Suppose there are 12 dou 6 14/15 sheng of paddy, which one wishes to make into millet. Question: how much does one get? Answer: it makes 10 dou 5 7/9 sheng of millet.\end{uniten}

\begin{unitzh}\unitlabel{jiuzhang-009.20}術曰：以稻求粟，五之，六而一。\end{unitzh}
\begin{uniten}Method: To seek millet from paddy, multiply by 5 and divide by 6.\end{uniten}

\begin{unitzh}\unitlabel{jiuzhang-009.21}今有糲米一十九斗二升七分升之一，欲爲粺米，問得幾何？荅曰：爲粺米一十七斗二升一十四分升之一十三。\end{unitzh}
\begin{uniten}Suppose there are 19 dou 2 1/7 sheng of coarse hulled grain, which one wishes to make into fine hulled grain. Question: how much does one get? Answer: it makes 17 dou 2 13/14 sheng of fine hulled grain.\end{uniten}

\begin{unitzh}\unitlabel{jiuzhang-009.22}術曰：以糲米求粺米，九之，十而一。{\zhsmall 臣淳風等謹按：粺率二十七，合以此數乘糲米。術欲從省，先以等數三約之，所求之率得九，所有之率得十，故九乘而十除。}\end{unitzh}
\begin{uniten}Method: To seek fine hulled grain from coarse hulled grain, multiply by 9 and divide by 10. [Li Chunfeng et al.: The rate of fine hulled grain is 27; one should multiply the coarse hulled grain by this number. The method wishes to economize: first reduce them by the equal number 3; the sought rate gives 9 and the given rate gives 10; therefore multiply by 9 and divide by 10.]\end{uniten}

\begin{unitzh}\unitlabel{jiuzhang-009.23}今有糲米六斗四升五分升之三，欲爲糲飯，問得幾何？荅曰：爲糲飯一十六斗一升半。\end{unitzh}
\begin{uniten}Suppose there are 6 dou 4 3/5 sheng of coarse hulled grain, which one wishes to make into cooked coarse grain (lifan). Question: how much does one get? Answer: it makes 16 dou 1 1/2 sheng of cooked coarse grain.\end{uniten}

\begin{unitzh}\unitlabel{jiuzhang-009.24}術曰：以糲米求糲飯，五之，二而一。{\zhsmall 臣淳風等謹按：糲飯之率七十有五，宜以本糲米乘此率。術欲從省，先以等數十五約之，所求之率得五，所有之率得二，故五乘二除，義由於此。}\end{unitzh}
\begin{uniten}Method: To seek cooked coarse grain from coarse hulled grain, multiply by 5 and divide by 2. [Li Chunfeng et al.: The rate of cooked coarse grain is 75; one ought to multiply this rate by the original coarse hulled grain. The method wishes to economize: first reduce them by the equal number 15; the sought rate gives 5 and the given rate gives 2; therefore multiply by 5 and divide by 2. The principle derives from this.]\end{uniten}

\begin{unitzh}\unitlabel{jiuzhang-009.25}今有糲飯七斗六升七分升之四，欲爲飱，問得幾何？荅曰：爲飱九斗一升三十五分升之三十一。\end{unitzh}
\begin{uniten}Suppose there are 7 dou 6 4/7 sheng of cooked coarse grain, which one wishes to make into soaked cooked grain. Question: how much does one get? Answer: it makes 9 dou 1 31/35 sheng of soaked cooked grain.\end{uniten}

\begin{unitzh}\unitlabel{jiuzhang-009.26}術曰：以糲飯求飱，六之，五而一。{\zhsmall 臣淳風等謹按：飱率九十，爲糲飯所求，宜以糲飯乘此率。術欲從省，先以等數十五約之，所求之率得六，所有之率得五，以此故六乘五除也。}\end{unitzh}
\begin{uniten}Method: To seek soaked cooked grain from cooked coarse grain, multiply by 6 and divide by 5. [Li Chunfeng et al.: The rate of soaked cooked grain is 90, which is what is sought from the cooked coarse grain; one ought to multiply this rate by the cooked coarse grain. The method wishes to economize: first reduce them by the equal number 15; the sought rate gives 6 and the given rate gives 5; for this reason multiply by 6 and divide by 5.]\end{uniten}

\begin{unitzh}\unitlabel{jiuzhang-009.27}今有菽一斗，欲爲熟菽，問得幾何？荅曰：爲熟菽二斗三升。\end{unitzh}
\begin{uniten}Suppose there is 1 dou of beans (shu), which one wishes to make into cooked beans. Question: how much does one get? Answer: it makes 2 dou 3 sheng of cooked beans.\end{uniten}

\begin{unitzh}\unitlabel{jiuzhang-009.28}術曰：以菽求熟菽，二十三之，十而一。{\zhsmall 臣淳風等謹按：熟菽之率一百三半，因其有半，各以母二通之，宜以熟菽數乘此率。術欲從者，先以等數九約之，所求之率得一十一半，所有之率得五也。}\end{unitzh}
\begin{uniten}Method: To seek cooked beans from beans, multiply by 23 and divide by 10. [Li Chunfeng et al.: The rate of cooked beans is 103 1/2; because it has a half, convert each by the denominator 2; one ought to multiply this rate by the number of cooked beans. As for what the method wishes to follow, first reduce them by the equal number 9: the sought rate gives 11 1/2 and the given rate gives 5.] [Translator's note: in 術欲從者, 者 stands where the parallel notes have 省, "economize".]\end{uniten}

\begin{unitzh}\unitlabel{jiuzhang-009.29}今有菽二斗，欲爲䜴，問得幾何？荅曰：爲䜴二斗八升。\end{unitzh}
\begin{uniten}Suppose there are 2 dou of beans, which one wishes to make into fermented beans. Question: how much does one get? Answer: it makes 2 dou 8 sheng of fermented beans.\end{uniten}

\begin{unitzh}\unitlabel{jiuzhang-009.30}術曰：以菽求䜴，七之，五而一。{\zhsmall 臣淳風等謹按：䜴率六十三，爲菽所求，宜以菽乘此率。術欲從省，先以等數九約之，所求之率得七，而所有之率得五也。}\end{unitzh}
\begin{uniten}Method: To seek fermented beans from beans, multiply by 7 and divide by 5. [Li Chunfeng et al.: The rate of fermented beans is 63, which is what is sought from the beans; one ought to multiply this rate by the beans. The method wishes to economize: first reduce them by the equal number 9; the sought rate gives 7, and the given rate gives 5.]\end{uniten}

\begin{unitzh}\unitlabel{jiuzhang-009.31}今有麥八斗六升七分升之三，欲爲小䵂，問得幾何？荅曰：爲小䵂二斗五升一十四分升之一十三。\end{unitzh}
\begin{uniten}Suppose there are 8 dou 6 3/7 sheng of wheat (mai), which one wishes to make into small groats (xiaozhi). Question: how much does one get? Answer: it makes 2 dou 5 13/14 sheng of small groats.\end{uniten}

\begin{unitzh}\unitlabel{jiuzhang-009.32}術曰：以麥求小䵂，三之，十而一。{\zhsmall 臣淳風等謹按：小䵂之率十三半，宜以母二通之，以乘本麥之數。術欲從省，先以等數九約之，所求之率得三，所有之率得十也。}\end{unitzh}
\begin{uniten}Method: To seek small groats from wheat, multiply by 3 and divide by 10. [Li Chunfeng et al.: The rate of small groats is 13 1/2; one ought to convert it by the denominator 2 and with it multiply the number of the original wheat. The method wishes to economize: first reduce them by the equal number 9; the sought rate gives 3 and the given rate gives 10.]\end{uniten}

\begin{unitzh}\unitlabel{jiuzhang-009.33}今有麥一斗，欲爲大䵂，問得幾何？荅曰：爲大䵂一斗二升。\end{unitzh}
\begin{uniten}Suppose there is 1 dou of wheat, which one wishes to make into large groats (dazhi). Question: how much does one get? Answer: it makes 1 dou 2 sheng of large groats.\end{uniten}

\begin{unitzh}\unitlabel{jiuzhang-009.34}術曰：以麥求大䵂，六之，五而一。{\zhsmall 臣淳風等謹按：大䵂之率五十有四，合以大䵂數乘此率。術欲從省，先以等數九約之，所求之率得六，所有之率得五也。}\end{unitzh}
\begin{uniten}Method: To seek large groats from wheat, multiply by 6 and divide by 5. [Li Chunfeng et al.: The rate of large groats is 54; one should multiply this rate by the number of large groats. The method wishes to economize: first reduce them by the equal number 9; the sought rate gives 6 and the given rate gives 5.]\end{uniten}

\begin{unitzh}\unitlabel{jiuzhang-009.35}今有出錢一百六十，買瓴甓十八枚，{\zhsmall 瓴甓，甎也。}問枚幾何？荅曰：一枚八錢九分錢之八。\end{unitzh}
\begin{uniten}Suppose there are 160 qian (錢, cash) paid out to buy 18 pieces of floor tiles (lingpi) [Liu Hui: lingpi are bricks]. Question: how much per piece? Answer: 1 piece is 8 8/9 qian.\end{uniten}

\begin{unitzh}\unitlabel{jiuzhang-009.36}今有出錢一萬三千五百，買竹二千三百五十箇，問箇幾何？荅曰：一箇五錢四十七分錢之三十五。\end{unitzh}
\begin{uniten}Suppose there are 13500 qian paid out to buy 2350 stems of bamboo. Question: how much per stem? Answer: 1 stem is 5 35/47 qian.\end{uniten}

\begin{unitzh}\unitlabel{jiuzhang-009.37}經率{\zhsmall 臣淳風等謹按：今有之義，以所求率乘所有數，合以瓴甓一枚乘錢一百六十爲實。但以一乘不長，故不復乘，是以徑將所買之率與所出之錢爲法、實也。又按：此今有之義，出錢爲所有數，一枚爲所求率，所買爲所有率，而今有之，卽得所求數。一乘不長，故不復乘，是以徑將所買之率爲法，以所出之錢爲實，實如法得一枚錢。不盡者，等數而命分。}\end{unitzh}
\begin{uniten}Direct rate (jinglü) [Li Chunfeng et al.: According to the principle of jinyou ("Suppose there is"), one multiplies the given number by the sought rate: one should multiply the 160 qian by the 1 piece of floor tile to make the dividend (shi). But multiplying by 1 does not make it grow, so one does not multiply further; therefore one directly takes the rate of what is bought and the cash paid out as divisor (fa) and dividend. Further, we note: in this principle of jinyou, the cash paid out is the given number, the 1 piece is the sought rate, what is bought is the given rate, and applying jinyou to them one at once obtains the sought number (suoqiu shu). Multiplying by 1 does not make it grow, so one does not multiply further; therefore one directly takes the rate of what is bought as divisor and the cash paid out as dividend; dividing the dividend by the divisor, one obtains the cash for 1 piece. Where it does not divide out, take the equal number and name the fraction (mingfen).]\end{uniten}

\begin{unitzh}\unitlabel{jiuzhang-009.38}術曰：以所買率爲法，所出錢數爲實，實如法得一。\end{unitzh}
\begin{uniten}Method: Take the rate of what is bought as divisor and the number of qian paid out as dividend; dividing the dividend by the divisor, one obtains 1.\end{uniten}

\begin{unitzh}\unitlabel{jiuzhang-010.1}今有出錢五千七百八十五買漆一斛六斗七升太半升，欲斗率之。問斗幾何？荅曰：一斗三百四十五錢五百三分錢之一十五。\end{unitzh}
\begin{uniten}Suppose there is an outlay of 5785 qian (錢) to buy 1 hu (斛) 6 dou (斗) 7 2/3 sheng (升) of lacquer, and one wishes to rate (lü) it by the dou. What is it per dou? Answer: 1 dou, 345 15/503 qian.\end{uniten}

\begin{unitzh}\unitlabel{jiuzhang-010.2}今有出錢七百二十買縑一匹二丈一尺，欲丈率之。問丈幾何？荅曰：一丈一百一十八錢六十一分錢之二。\end{unitzh}
\begin{uniten}Suppose there is an outlay of 720 qian to buy 1 pi (匹) 2 zhang (丈) 1 chi (尺) of fine silk, and one wishes to rate it by the zhang. What is it per zhang? Answer: 1 zhang, 118 2/61 qian.\end{uniten}

\begin{unitzh}\unitlabel{jiuzhang-010.3}今有出錢二千三百七十買布九匹二丈七尺，欲匹率之。問匹幾何？荅曰：一匹二百四十四錢一百二十九分錢之一百二十四。\end{unitzh}
\begin{uniten}Suppose there is an outlay of 2370 qian to buy 9 pi 2 zhang 7 chi of cloth, and one wishes to rate it by the pi. What is it per pi? Answer: 1 pi, 244 124/129 qian.\end{uniten}

\begin{unitzh}\unitlabel{jiuzhang-010.4}今有出錢一萬三千六百七十買絲一石二鈞一十七斤，欲石率之。問石幾何？荅曰：一石八千三百二十六錢一百九十七分錢之一百七十八。\end{unitzh}
\begin{uniten}Suppose there is an outlay of 13670 qian to buy 1 dan (石) 2 jun (鈞) 17 jin (斤) of silk, and one wishes to rate it by the dan. What is it per dan? Answer: 1 dan, 8326 178/197 qian.\end{uniten}

\begin{unitzh}\unitlabel{jiuzhang-010.5}經率{\zhsmall 此術猶經分。臣淳風等謹按：今有之義，一斗爲所求率，出錢爲所有數，故以一斗乘錢數。有分者通之，又以分母乘之爲實。所買通分內子爲所有率，故以爲法。實如法而一，得錢數。不盡而命分者，因法爲母，實餘爲子。實見不滿，故以命之。}\end{unitzh}
\begin{uniten}Direct rate (jinglü) [Liu Hui: This method is like direct division (jingfen).] [Li Chunfeng et al.: According to the meaning of 'suppose there is' (jinyou), 1 dou is the rate of what is sought (suoqiu lü) and the cash paid out is the number of what one has (suoyou shu); therefore multiply the number of cash by 1 dou. Where there are fractions, convert them (tong), and further multiply it by the denominator to make the dividend (shi). What is bought, converted into parts with the numerator added in (tongfen nazi), is the rate of what one has (suoyou lü); therefore take it as the divisor (fa). Divide the dividend by the divisor, obtaining the number of cash. Where it is not exhausted and one names the fraction (mingfen), accordingly take the divisor as the denominator and the remainder of the dividend as the numerator. The dividend is seen not to be full, therefore one names it thus.]\end{uniten}

\begin{unitzh}\unitlabel{jiuzhang-010.6}術曰：以所求率乘錢數爲實，以所買率爲法，實如法得一。\end{unitzh}
\begin{uniten}Method: Multiply the number of cash by the rate of what is sought to make the dividend; take the rate of what is bought (suomai lü) as the divisor; divide the dividend by the divisor.\end{uniten}

\begin{unitzh}\unitlabel{jiuzhang-010.7}今有出錢五百七十六買竹七十八箇，欲其大小率之。問各幾何？荅曰：其四十八箇，箇七錢；其三十箇，箇八錢。\end{unitzh}
\begin{uniten}Suppose there is an outlay of 576 qian to buy 78 pieces of bamboo, and one wishes to rate them by large and small. What is each? Answer: 48 pieces of them, 7 qian per piece; 30 pieces of them, 8 qian per piece.\end{uniten}

\begin{unitzh}\unitlabel{jiuzhang-010.8}今有出錢一千一百二十買絲一石二鈞十八斤，欲其貴賤斤率之。問各幾何？荅曰：其二鈞八斤，斤五錢；其一石一十斤，斤六錢。\end{unitzh}
\begin{uniten}Suppose there is an outlay of 1120 qian to buy 1 dan 2 jun 18 jin of silk, and one wishes to rate it, dear and cheap, by the jin. What is each? Answer: 2 jun 8 jin of it, 5 qian per jin; 1 dan 10 jin of it, 6 qian per jin.\end{uniten}

\begin{unitzh}\unitlabel{jiuzhang-010.9}今有出錢一萬三千九百七十買絲一石二鈞二十八斤三兩五銖，欲其貴賤石率之。問各幾何？荅曰：其一鈞九兩一十二銖，石八千五十一錢；其一石一鈞二十七斤九兩一十七銖，石八千五十二錢。\end{unitzh}
\begin{uniten}Suppose there is an outlay of 13970 qian to buy 1 dan 2 jun 28 jin 3 liang (兩) 5 zhu (銖) of silk, and one wishes to rate it, dear and cheap, by the dan. What is each? Answer: 1 jun 9 liang 12 zhu of it, 8051 qian per dan; 1 dan 1 jun 27 jin 9 liang 17 zhu of it, 8052 qian per dan.\end{uniten}

\begin{unitzh}\unitlabel{jiuzhang-010.10}今有出錢一萬三千九百七十買絲一石二鈞二十八斤三兩五銖，欲其貴賤鈞率之。問各幾何？荅曰：其七斤一十兩九銖，鈞二千一十二錢；其一石二鈞二十斤八兩二十銖，鈞二千一十三錢。\end{unitzh}
\begin{uniten}Suppose there is an outlay of 13970 qian to buy 1 dan 2 jun 28 jin 3 liang 5 zhu of silk, and one wishes to rate it, dear and cheap, by the jun. What is each? Answer: 7 jin 10 liang 9 zhu of it, 2012 qian per jun; 1 dan 2 jun 20 jin 8 liang 20 zhu of it, 2013 qian per jun.\end{uniten}

\begin{unitzh}\unitlabel{jiuzhang-010.11}今有出錢一萬三千九百七十買絲一石二鈞二十八斤三兩五銖，欲其貴賤斤率之。問各幾何？荅曰：其一石二鈞七斤十兩四銖，斤六十七錢；其二十斤九兩一銖，斤六十八錢。\end{unitzh}
\begin{uniten}Suppose there is an outlay of 13970 qian to buy 1 dan 2 jun 28 jin 3 liang 5 zhu of silk, and one wishes to rate it, dear and cheap, by the jin. What is each? Answer: 1 dan 2 jun 7 jin 10 liang 4 zhu of it, 67 qian per jin; 20 jin 9 liang 1 zhu of it, 68 qian per jin.\end{uniten}

\begin{unitzh}\unitlabel{jiuzhang-010.12}今有出錢一萬三千九百七十買絲一石二鈞二十八斤三兩五銖，欲其貴賤兩率之。問各幾何？荅曰：其一石一鈞一十七斤一十四兩一銖，兩四錢；其一鈞一十斤五兩四銖，兩五錢。\end{unitzh}
\begin{uniten}Suppose there is an outlay of 13970 qian to buy 1 dan 2 jun 28 jin 3 liang 5 zhu of silk, and one wishes to rate it, dear and cheap, by the liang. What is each? Answer: 1 dan 1 jun 17 jin 14 liang 1 zhu of it, 4 qian per liang; 1 jun 10 jin 5 liang 4 zhu of it, 5 qian per liang.\end{uniten}

\begin{unitzh}\unitlabel{jiuzhang-010.13}其率{\zhsmall 如欲令差分。按：出錢五百七十六買竹七十八箇，以除錢，得七，實餘三十，是爲三十箇復可増一錢。然則實餘之數卽是貴者之數，故曰實貴也。本以七十八箇爲法，今以貴者減之，則其餘悉是賤者之數，故曰法賤也。其求石、鈞、斤、兩，以積銖各除法實，各得其積數，餘各爲銖者，謂石、鈞、斤、兩積銖除實，又以石、鈞、斤、兩積銖除法，餘各爲銖，卽合所問。}\end{unitzh}
\begin{uniten}Its rate (qilü) [Liu Hui: It is as if one wished to make a differentiated division (chafen). Consider: 576 qian are paid out to buy 78 pieces of bamboo; dividing the cash by this gives 7, and the dividend has 30 left over; this means that 30 pieces can again be increased by 1 qian. Thus the number left over in the dividend is precisely the number of the dear ones; therefore it says 'the dividend is the dear'. Originally 78 pieces were taken as the divisor; now subtract the dear ones from it, and then what remains is entirely the number of the cheap ones; therefore it says 'the divisor is the cheap'. As for seeking dan, jun, jin and liang by dividing the divisor and the dividend each by the accumulated zhu (jizhu), each obtaining its accumulated number, with the remainders each being zhu: this means dividing the dividend by the accumulated zhu of the dan, jun, jin or liang, and also dividing the divisor by the accumulated zhu of the dan, jun, jin or liang; the remainders each being zhu, this agrees with what is asked.]\end{uniten}

\begin{unitzh}\unitlabel{jiuzhang-010.14}術曰：各置所買石、鈞、斤、兩以爲法，以所率乘錢數爲實，實如法而一。不滿法者，反以實減法，法賤實貴。\end{unitzh}
\begin{uniten}Method: In each case set out the dan, jun, jin or liang bought as the divisor; multiply the number of cash by what is rated (suolü) to make the dividend; divide the dividend by the divisor. As for what does not fill the divisor, conversely subtract the dividend from the divisor: the divisor is the cheap, the dividend is the dear.\end{uniten}

\begin{unitzh}\unitlabel{jiuzhang-010.15}今有出錢一萬三千九百七十買絲一石二鈞二十八斤三兩五銖，欲其貴賤銖率之。問各幾何？荅曰：其一鈞二十斤六兩十一銖，五銖一錢；其一石一鈞七斤一十二兩一十八銖，六銖一錢。\end{unitzh}
\begin{uniten}Suppose there is an outlay of 13970 qian to buy 1 dan 2 jun 28 jin 3 liang 5 zhu of silk, and one wishes to rate it, dear and cheap, by the zhu. What is each? Answer: 1 jun 20 jin 6 liang 11 zhu of it, 5 zhu for 1 qian; 1 dan 1 jun 7 jin 12 liang 18 zhu of it, 6 zhu for 1 qian.\end{uniten}

\begin{unitzh}\unitlabel{jiuzhang-010.16}今有出錢六百二十買羽二千一百翭{\zhsmall 翭，羽本也。數羽稱其本，猶數草木稱其根株。}，欲其貴賤率之。問各幾何？荅曰：其一千一百四十翭，三翭一錢；其九百六十翭，四翭一錢。\end{unitzh}
\begin{uniten}Suppose there is an outlay of 620 qian to buy 2100 hou (翭) of feathers [Liu Hui: Hou is the base of a feather. In counting feathers one names them by their bases, just as in counting plants and trees one names them by their roots and trunks.], and one wishes to rate them dear and cheap. What is each? Answer: 1140 hou of them, 3 hou for 1 qian; 960 hou of them, 4 hou for 1 qian.\end{uniten}

\begin{unitzh}\unitlabel{jiuzhang-010.17}今有出錢九百八十買矢簳五千八百二十枚，欲其貴賤率之。問各幾何？荅曰：其三百枚，五枚一錢；其五千五百二十枚，六枚一錢。\end{unitzh}
\begin{uniten}Suppose there is an outlay of 980 qian to buy 5820 arrow shafts, and one wishes to rate them dear and cheap. What is each? Answer: 300 of them, 5 for 1 qian; 5520 of them, 6 for 1 qian.\end{uniten}

\begin{unitzh}\unitlabel{jiuzhang-010.18}反其率{\zhsmall 臣淳風等謹按：其率者，錢多物少；反其率者，錢少物多。多少相反，故曰反其率也。其率者，以物數爲法，錢數爲實；反之者，以錢數爲法，物數爲實。不滿法，知實餘也，當以餘物化爲錢矣。法爲凡錢，而今以化錢減之，故以實減法。法少，知經分之所得，故曰法少；實多者，餘分之所益，故曰實多。乘實宜以多，乘法宜以少，故曰各以其所得多少之數乘法實，卽物數。其求石、鈞、斤、兩，以積銖各除法實，各得其數，餘各爲銖者，謂之石、鈞、斤、兩積銖除實，石、鈞、斤、兩積銖除法，餘各爲銖，卽合所問。}\end{unitzh}
\begin{uniten}Inverting its rate (fan qilü) [Li Chunfeng et al.: In 'its rate' the cash is more and the goods fewer; in 'inverting its rate' the cash is fewer and the goods more. More and fewer are mutually reversed, therefore it is called 'inverting its rate'. In 'its rate', the number of goods is taken as the divisor and the number of cash as the dividend; in inverting it, the number of cash is taken as the divisor and the number of goods as the dividend. 'Does not fill the divisor': know that this is the remainder of the dividend; one should convert the remaining goods into cash. The divisor is the total cash, and now the converted cash is subtracted from it; therefore subtract the dividend from the divisor. 'The divisor is the fewer': know that it is what is obtained by direct division, therefore it says 'the divisor is the fewer'; 'the dividend is the more': it is what is added by the remaining parts, therefore it says 'the dividend is the more'. In multiplying the dividend one should use the more; in multiplying the divisor one should use the fewer; therefore it says 'in each case multiply the divisor and the dividend by their numbers obtained, more and fewer; these are the numbers of the goods'. As for seeking dan, jun, jin and liang by dividing the divisor and the dividend each by the accumulated zhu, each obtaining its number, with the remainders each being zhu: this is called dividing the dividend by the accumulated zhu of the dan, jun, jin or liang, and dividing the divisor by the accumulated zhu of the dan, jun, jin or liang; the remainders each being zhu, this agrees with what is asked.]\end{uniten}

\begin{unitzh}\unitlabel{jiuzhang-010.19}術曰：以錢數爲法，所率爲實，實如法而一。不滿法者，反以實減法，法少實多。二物各以所得多少之數乘法實，卽物數。\end{unitzh}
\begin{uniten}Method: Take the number of cash as the divisor and what is rated as the dividend; divide the dividend by the divisor. As for what does not fill the divisor, conversely subtract the dividend from the divisor: the divisor is the fewer, the dividend is the more. For the two goods, in each case multiply the divisor and the dividend by the numbers obtained, more and fewer; these are the numbers of the goods.\end{uniten}

\begin{unitzh}\unitlabel{jiuzhang-010.20}{\zhsmall 按：其率出錢六百二十買羽二千一百翭，反之，當二百四十錢，一錢四翭；其三百八十錢，一錢三翭。是錢有二價，物有貴賤，故以羽乘錢，反二率也。臣淳風等謹按：其率者，以物數爲法，錢爲實；反之者，以錢數爲法，物爲實。不滿法者，實餘也，當以餘物化爲錢矣。法爲凡錢，而今以化錢減之，故曰反以實減法也。法少者，知經分之所得，故曰法少；實多者，知餘分之所益，故曰實多。宜以多乘法，少乘實，故曰各以所得多少數乘法實，卽物數也。}\end{unitzh}
\begin{uniten}[Liu Hui: Consider: by 'its rate', 620 qian are paid out to buy 2100 hou of feathers; inverting it, there should be 240 qian at 1 qian for 4 hou, and of them 380 qian at 1 qian for 3 hou. This means that the cash has two prices and the goods are dear and cheap; therefore one multiplies the cash by the feathers, inverting the two rates.] [Li Chunfeng et al.: In 'its rate', the number of goods is taken as the divisor and the cash as the dividend; in inverting it, the number of cash is taken as the divisor and the goods as the dividend. What does not fill the divisor is the remainder of the dividend; one should convert the remaining goods into cash. The divisor is the total cash, and now the converted cash is subtracted from it; therefore it says 'conversely subtract the dividend from the divisor'. 'The divisor is the fewer': know that it is what is obtained by direct division, therefore it says 'the divisor is the fewer'; 'the dividend is the more': know that it is what is added by the remaining parts, therefore it says 'the dividend is the more'. One should multiply the divisor by the more and the dividend by the fewer; therefore it says 'in each case multiply the divisor and the dividend by the numbers obtained, more and fewer; these are the numbers of the goods'.] [Translator's note: 'multiply the divisor by the more and the dividend by the fewer' reverses the pairing in the method and in the preceding annotation.]\end{uniten}

\begin{unitzh}\unitlabel{jiuzhang-010.21}九章算術卷二終\end{unitzh}
\begin{uniten}End of juan 2 of the Nine Chapters on the Mathematical Art.\end{uniten}

\section*{卷第三}\addcontentsline{toc}{section}{卷第三}

\begin{unitzh}\unitlabel{jiuzhang-011.1}《九章算術》卷二訂訛補圖{\zhsmall 算經十書之二} 休寧戴震東原\end{unitzh}
\begin{uniten}The Nine Chapters on the Mathematical Art, Chapter 2: Corrections of Errors and Supplementary Diagrams [Note: The Ten Mathematical Classics, no. 2]. Dai Zhen, Dongyuan, of Xiuning.\end{uniten}

\begin{unitzh}\unitlabel{jiuzhang-011.2}粟率五十，糲米三十，粺米二十七，糳米二十四，御米二十一。{\zhsmall 《詩·大雅》鄭箋云：「米之率：糲十，粺九，糳八，侍御七。」疏云：「《九章》粟米之法：粟率五十，糲米三十，粺二十七，糳二十四，御二十一。言粟五升爲糲米三升，已下則米漸細，故數益少。」}\end{unitzh}
\begin{uniten}Rate of millet (su) 50, coarse hulled grain (limi) 30, polished grain (baimi) 27, fine-polished grain (zuomi) 24, imperial grain (yumi) 21. [Dai Zhen: Zheng's commentary (jian) on the "Greater Odes" (Daya) of the Book of Odes says: "The rates of hulled grain (mi): coarse 10, polished 9, fine-polished 8, for imperial attendance (shiyu) 7." The subcommentary (shu) says: "The method of millet and grain in the Nine Chapters: rate of millet 50, coarse hulled grain 30, polished 27, fine-polished 24, imperial 21. That is, 5 sheng (升) of millet make 3 sheng of coarse hulled grain; from there downward the grain becomes gradually finer, and therefore the numbers become ever smaller."]\end{uniten}

\begin{unitzh}\unitlabel{jiuzhang-011.3}按其率，出錢六百二十，買羽二千一百翭，反之，當二百四十錢，一錢四翭；其三百八十錢，一錢三翭。{\zhsmall 已上舛誤不可通。參考上注，當云：「按其率，錢多物少；反之，錢少物多。出錢六百二十，買羽二千一百翭，當以除羽，得三，實餘二百四十，是謂三翭，復可增一翭。然則實餘之數卽是多者之錢，故曰實多。本以六百二十錢爲法，今以多者減之，則其餘三百八十悉是少者之錢，故曰法少也。二百四十錢，一錢四翭，乘得九百六十；其三百八十錢，一錢三翭，乘得一千一百四十。」}\end{unitzh}
\begin{uniten}Note: by its rates (qi lü), paying 620 cash (qian) and buying 2100 hou (翭) of feathers; inverting it, it should be 240 cash, 1 cash for 4 hou; the 380 cash, 1 cash for 3 hou. [Dai Zhen: The above is erroneous and cannot be made sense of. Consulting the commentary above, it should say: "Note: by its rates, the cash is more and the things fewer; inverting it, the cash is fewer and the things more. Paying 620 cash and buying 2100 hou of feathers, one should divide the feathers by it, getting 3, with 240 remaining in the dividend (shi); this means that at 3 hou one can further add 1 hou. Thus the number remaining in the dividend is the cash of the 'more' ones; therefore it says 'the dividend is the more'. Originally 620 cash was taken as the divisor (fa); now subtracting the 'more' ones from it, the remaining 380 is all the cash of the 'fewer' ones; therefore it says 'the divisor is the fewer'. 240 cash, 1 cash for 4 hou, multiplied gives 960; the 380 cash, 1 cash for 3 hou, multiplied gives 1140."]\end{uniten}

\begin{unitzh}\unitlabel{jiuzhang-011.4}《九章算術》卷三{\zhsmall 算經十書之二} 魏劉徽注 唐朝議大夫、行太史令、上輕車都尉臣李淳風等奉勅注釋\end{unitzh}
\begin{uniten}The Nine Chapters on the Mathematical Art, Chapter 3 [Note: The Ten Mathematical Classics, no. 2]. Commentary by Liu Hui of the Wei. Annotated and explained by imperial order by Your servant Li Chunfeng et al. of the Tang, Grand Master for Court Discussion (chaoyi dafu), acting Grand Astrologer (taishi ling), Senior Commandant of Light Chariots (shang qingche duwei).\end{uniten}

\begin{unitzh}\unitlabel{jiuzhang-011.5}衰分{\zhsmall 以御貴賤禀稅}\end{unitzh}
\begin{uniten}Proportional Distribution (cuifen) [Liu Hui: For dealing with the noble and the base, rations and taxes.]\end{uniten}

\begin{unitzh}\unitlabel{jiuzhang-011.6}衰分{\zhsmall 衰分，差也。}術曰：各置列衰，{\zhsmall 列衰，相與率也。重疊則可約。}副并爲法，以所分乘未并者，各自爲實。{\zhsmall 法集而衰别，數本一也。今以所分乘上别，以下集除之，一乘一除，適足相消，故所分猶存，且各應率而别也。於今有術，列衰各爲所求率，副并爲所有率，所分爲所有數。又以經分言之，假令甲家三人，乙家二人，丙家一人，并六人，共分十二，爲人得二也。欲復作逐家者，則當列置人數，以一人所得乘之。今此術先乗而後除也。}實如法而一。不滿法者，以法命之。\end{unitzh}
\begin{uniten}Proportional distribution [Liu Hui: Cuifen is grading (cha).] Method: Set up each of the arrayed proportions (liecui); [Liu Hui: The arrayed proportions are rates relative to one another (xiangyu lü). If they are reduplicated (chongdie), they can be reduced.] add them up on the side as the divisor (fa); multiply the unadded ones by what is to be distributed, each to be its own dividend (shi). [Liu Hui: The divisor is collected while the proportions are separate, but the number is fundamentally one. Now we multiply the separate ones above by what is to be distributed and divide by the collected one below; one multiplication and one division exactly cancel each other, so that what is to be distributed still remains, and moreover it is separated, each in accordance with its rate. In terms of the "suppose there is" method (jinyou shu), the arrayed proportions are each the sought rate (suoqiu lü), their side sum is the given rate (suoyou lü), and what is to be distributed is the given number (suoyou shu). Again, speaking in terms of ordinary division (jingfen): suppose household A has 3 persons, household B 2 persons, household C 1 person, together 6 persons, who share 12 in common; each person gets 2. If one wishes to work it out again household by household, one should set out the numbers of persons and multiply them by what one person gets. Now this method multiplies first and divides afterwards.] Divide the dividend by the divisor. What does not fill the divisor, name it with the divisor.\end{uniten}

\begin{unitzh}\unitlabel{jiuzhang-011.7}今有大夫、不更、簪褭、上造、公士，凡五人，共獵得五鹿，欲以爵次分之，問各得幾何？荅曰：大夫得一鹿三分鹿之二，不更得一鹿三分鹿之一，簪褭得一鹿，上造得三分鹿之二，公士得三分鹿之一。\end{unitzh}
\begin{uniten}Suppose there are a dafu, a bugeng, a zanniao, a shangzao and a gongshi, 5 persons in all, who together hunt and catch 5 deer. One wishes to divide them according to the order of rank. Question: how much does each get? Answer: The dafu gets 1 2/3 deer; the bugeng gets 1 1/3 deer; the zanniao gets 1 deer; the shangzao gets 2/3 deer; the gongshi gets 1/3 deer.\end{uniten}

\begin{unitzh}\unitlabel{jiuzhang-011.8}術曰：列置爵數，各自爲衰，{\zhsmall 爵數者，謂大夫五，不更四，簪褭三，上造二，公士一也。《墨子·號令篇》以爵級爲賜，然則戰國之初有此名也。今有術，列衰各爲所求率，副并爲所有率，今有鹿數爲所有數，而今有之，卽得。}副并爲法，以五鹿乘未并者，各自爲實，實如法得一鹿。\end{unitzh}
\begin{uniten}Method: Set out the rank numbers, each to be its own proportion (cui); [Liu Hui: The rank numbers means: dafu 5, bugeng 4, zanniao 3, shangzao 2, gongshi 1. The "Orders and Commands" (Haoling) chapter of the Mozi takes rank grades as rewards; thus at the beginning of the Warring States these names existed. By the "suppose there is" method, the arrayed proportions are each the sought rate, their side sum is the given rate, the number of deer now had is the given number; apply "suppose there is" to it, and one gets [the result].] add them up on the side as the divisor; multiply the unadded ones by the 5 deer, each to be its own dividend; each time the dividend matches the divisor, one gets 1 deer.\end{uniten}

\begin{unitzh}\unitlabel{jiuzhang-011.9}今有牛、馬、羊食人苗，苗主責之粟五斗。羊主日：「我羊食半馬。」馬主日：「我馬食半牛。」今欲衰償之，問各出幾何？荅曰：牛主出二斗八升七分升之四，馬主出一斗四升七分升之二，羊主出七升七分升之一。\end{unitzh}
\begin{uniten}Suppose there are an ox, a horse and a sheep that eat someone's sprouts. The owner of the sprouts demands 5 dou (斗) of millet from them. The owner of the sheep says: "My sheep ate half as much as the horse." The owner of the horse says: "My horse ate half as much as the ox." [Translator's note: in both places the text has 日 "day" for 曰 "says".] Now one wishes to repay it proportionally. Question: how much does each pay? Answer: The owner of the ox pays 2 dou 8 4/7 sheng; the owner of the horse pays 1 dou 4 2/7 sheng; the owner of the sheep pays 7 1/7 sheng.\end{uniten}

\begin{unitzh}\unitlabel{jiuzhang-011.10}術曰：置牛四、馬二、羊一，各自爲列衰，副并爲法，以五斗乘未并者，各自爲實，實如法得一斗。{\zhsmall 臣淳風等謹按：此術問意，羊食半馬，馬食半牛，是謂四羊當一牛，二羊當一馬。今術置羊一、馬二、牛四者，通其率以爲列衰。}\end{unitzh}
\begin{uniten}Method: Set up ox 4, horse 2, sheep 1, each to be its own arrayed proportion; add them up on the side as the divisor; multiply the unadded ones by 5 dou, each to be its own dividend; each time the dividend matches the divisor, one gets 1 dou. [Li Chunfeng et al.: We respectfully note: the intent of the question in this method is that the sheep ate half as much as the horse and the horse ate half as much as the ox; this means that 4 sheep are equivalent to 1 ox and 2 sheep are equivalent to 1 horse. Now the method sets up sheep 1, horse 2, ox 4: it unifies (tong) their rates to make the arrayed proportions.]\end{uniten}

\begin{unitzh}\unitlabel{jiuzhang-011.11}今有甲持錢五百六十，乙持錢三百五十，丙持錢一百八十，凡三人俱出關，關稅百錢。欲以錢數多少衰出之，問各幾何？荅曰：甲出五十一錢一百九分錢之四十一，乙出三十二錢一百九分錢之一十二，丙出一十六錢一百九分錢之五十六。\end{unitzh}
\begin{uniten}Suppose there is A holding 560 cash, B holding 350 cash, and C holding 180 cash; in all 3 persons go out through a pass together, and the pass tax is 100 cash. One wishes them to pay it proportionally according to how much cash [each holds]. Question: how much does each [pay]? Answer: A pays 51 41/109 cash; B pays 32 12/109 cash; C pays 16 56/109 cash.\end{uniten}

\begin{unitzh}\unitlabel{jiuzhang-011.12}術曰：各置錢數爲列衰，副并爲法，以百錢乘未并者，各自爲實，實如法得一錢。{\zhsmall 臣淳風等謹按：此術甲、乙、丙持錢數以爲列衰，副并爲所有率，未并者各爲所求率，百錢爲所有數，而今有之，卽得。}\end{unitzh}
\begin{uniten}Method: Set up each of the amounts of cash as the arrayed proportions; add them up on the side as the divisor; multiply the unadded ones by 100 cash, each to be its own dividend; each time the dividend matches the divisor, one gets 1 cash. [Li Chunfeng et al.: We respectfully note: in this method the amounts of cash held by A, B and C are taken as the arrayed proportions; their side sum is the given rate; the unadded ones are each the sought rate; 100 cash is the given number; apply "suppose there is" to it, and one gets [the result].]\end{uniten}

\begin{unitzh}\unitlabel{jiuzhang-011.13}今有女子善織，日自倍，五日織五尺。問日織幾何？荅曰：初日織一寸三十一分寸之十九，次日織三寸三十一分寸之七，次日織六寸三十一分寸之十四，次日織一尺二寸三十一分寸之二十八，次日織二尺五寸三十一分寸之二十五。\end{unitzh}
\begin{uniten}Suppose there is a woman who is good at weaving; each day she doubles her own [output], and in 5 days she weaves 5 chi (尺). Question: how much does she weave each day? Answer: On the first day she weaves 1 19/31 cun (寸); on the next day she weaves 3 7/31 cun; on the next day she weaves 6 14/31 cun; on the next day she weaves 1 chi 2 28/31 cun; on the next day she weaves 2 chi 5 25/31 cun.\end{uniten}

\begin{unitzh}\unitlabel{jiuzhang-011.14}術曰：置一、二、四、八、十六爲列衰，副并爲法，以五尺乘未并者，各自爲實，實如法得一尺。\end{unitzh}
\begin{uniten}Method: Set up 1, 2, 4, 8, 16 as the arrayed proportions; add them up on the side as the divisor; multiply the unadded ones by 5 chi, each to be its own dividend; each time the dividend matches the divisor, one gets 1 chi.\end{uniten}

\begin{unitzh}\unitlabel{jiuzhang-011.15}今有北鄉算八千七百五十八，西鄉算七千二百三十六，南鄉算八千三百五十六，凡三鄉發傜三百七十八人。欲以算數多少衰出之，問各幾何？荅曰：北鄉遣一百三十五人一萬二千一百七十五分人之一萬一千六百二十七，西鄉遣一百一十二人一萬二千一百七十五分人之四千四，南鄉遣一百二十九人一萬二千一百七十五分人之八千七百九。\end{unitzh}
\begin{uniten}Suppose there is a northern township (xiang) with 8758 tax counts (suan), a western township with 7236 tax counts, and a southern township with 8356 tax counts; in all, the 3 townships levy 378 persons for corvée (yao). One wishes them to provide [the persons] proportionally according to how many tax counts [each has]. Question: how many does each [send]? Answer: The northern township sends 135 11627/12175 persons; the western township sends 112 4004/12175 persons; the southern township sends 129 8709/12175 persons.\end{uniten}

\begin{unitzh}\unitlabel{jiuzhang-011.16}術曰：各置算數爲列衰，{\zhsmall 臣淳風等謹按：三鄉算數，約，可半者爲列衰。}副并爲法，以所發傜人數乘未并者，各自爲實，實如法得一人。{\zhsmall 按此術，今有之義也。}\end{unitzh}
\begin{uniten}Method: Set up each of the numbers of tax counts as the arrayed proportions; [Li Chunfeng et al.: We respectfully note: the tax counts of the three townships are reduced; those which can be halved are made the arrayed proportions.] add them up on the side as the divisor; multiply the unadded ones by the number of persons levied for corvée, each to be its own dividend; each time the dividend matches the divisor, one gets 1 person. [Liu Hui: Note: this method has the sense of "suppose there is".]\end{uniten}

\begin{unitzh}\unitlabel{jiuzhang-011.17}今有禀粟，大夫、不更、簪褭、上造、公士凡五人，一十五斗。今有大夫一人後來，亦當禀五斗。倉無粟，欲以衰出之，問各幾何？荅曰：大夫出一斗四分斗之一，不更出一斗，𬖂褭出四分斗之三，上造出四分斗之二，公士出四分斗之一。\end{unitzh}
\begin{uniten}Suppose there is a millet ration: a dafu, a bugeng, a zanniao, a shangzao and a gongshi, 5 persons in all, [receive] 15 dou. Suppose there is 1 dafu who comes later and should also receive 5 dou as ration. The granary has no millet; one wishes them to contribute it proportionally. Question: how much does each [contribute]? Answer: The dafu contributes 1 1/4 dou; the bugeng contributes 1 dou; the zanniao contributes 3/4 dou; the shangzao contributes 2/4 dou; the gongshi contributes 1/4 dou.\end{uniten}

\begin{unitzh}\unitlabel{jiuzhang-011.18}術曰：各置所禀粟斛斗數，爵次均之，以爲列衰，副并而加後來大夫亦五斗，得二十，以爲法，以五斗乘未并者，各自爲實，實如法得一斗。{\zhsmall 禀前五人十五斗者，大夫得五斗，不更得四斗，簪褭得三斗，上造得二斗，公士得一斗。欲令五人各依所得粟多少減與後來大夫，卽與前來大夫同。據前來大夫已得五斗，故言亦也。各以所得斗數爲衰，并得十五，而加後來大夫亦五斗，凡二十爲法也。是爲六人共出五斗，後來大夫亦俱損折。今有術，副并爲所有率，未并者各爲所求率，五斗爲所有數，而今有之，卽得。}\end{unitzh}
\begin{uniten}Method: Set up each of the numbers of hu (斛) and dou of millet received as rations, apportion them by rank order, and take them as the arrayed proportions; add them up on the side and add the 5 dou that the later-coming dafu also [receives], getting 20, and take this as the divisor; multiply the unadded ones by 5 dou, each to be its own dividend; each time the dividend matches the divisor, one gets 1 dou. [Liu Hui: As for the 15 dou issued to the first 5 persons: the dafu got 5 dou, the bugeng got 4 dou, the zanniao got 3 dou, the shangzao got 2 dou, the gongshi got 1 dou. One wishes the 5 persons each to reduce [their share] according to how much millet they got and give it to the later-coming dafu, so that he is then the same as the earlier-coming dafu. Because the earlier-coming dafu has already got 5 dou, it says "also". Each takes the number of dou received as its proportion; added together they give 15, and adding the 5 dou that the later-coming dafu also [receives], in all 20 is the divisor. This makes 6 persons jointly contribute 5 dou, the later-coming dafu also sharing in the reduction. By the "suppose there is" method, the side sum is the given rate, the unadded ones are each the sought rate, and 5 dou is the given number; apply "suppose there is" to it, and one gets [the result].]\end{uniten}

\begin{unitzh}\unitlabel{jiuzhang-012.1}今有禀粟五斛，五人分之，欲令三人得三，二人得二。問各幾何？荅曰：三人，人得一斛一斗五升十三分升之五；二人，人得七斗六升十三分升之十二。\end{unitzh}
\begin{uniten}Suppose there are 5 hu (斛) of ration millet (禀粟), which 5 persons divide; it is desired that 3 persons get 3 and 2 persons get 2. Question: how much is each? Answer: The 3 persons: each person gets 1 hu 1 dou (斗) 5 5/13 sheng (升). The 2 persons: each person gets 7 dou 6 12/13 sheng.\end{uniten}

\begin{unitzh}\unitlabel{jiuzhang-012.2}術曰：置三人人三、二人人二爲列衰，副并爲法。以五斛乘未并者，各自爲實。實如法得一斛。\end{unitzh}
\begin{uniten}Method: Set out 3 per person for the 3 persons and 2 per person for the 2 persons as the arrayed proportions (liecui 列衰); add them together in a duplicate (fubing 副并) as the divisor (fa 法). Multiply the ones not added by 5 hu, each separately as the dividend (shi 實). Each time the dividend matches the divisor, one gets 1 hu.\end{uniten}

\begin{unitzh}\unitlabel{jiuzhang-012.3}返衰\end{unitzh}
\begin{uniten}Inverse proportions (fancui 返衰)\end{uniten}

\begin{unitzh}\unitlabel{jiuzhang-012.4}{\zhsmall 以爵次言之，大夫五，不更四。欲令高爵得多者，當使大夫一人受五分，不更一人受四分。人數爲母，分數爲子。母同則子齊，齊卽衰也。故上衰分宜以五、四爲列焉。今此令高爵出少，則當使大夫五人共出一人分，不更四人共出一人分，故謂之返衰。人數不同，則分數不齊，當令母互乘子。母互乘子，則動者爲不動者衰也。亦可先同其母，各以分母約其子，爲返衰。副并爲法，以所分乘未并者，各自爲實，實如法而一。}\end{unitzh}
\begin{uniten}[Liu Hui: Speaking in terms of the order of ranks, the dafu (大夫) is 5 and the bugeng (不更) is 4. If it is desired that the higher rank get more, one should make 1 dafu receive 5 parts and 1 bugeng receive 4 parts. The number of persons is the denominator (mu 母), the number of parts is the numerator (zi 子). When the denominators are the same, the numerators are homogenized (qi 齊); being homogenized, they are the proportions (cui 衰). Therefore, in the distribution by proportions (cuifen 衰分) above, it is fitting to take 5 and 4 as the array. Now here the higher rank is made to pay less; then one should make 5 dafu jointly pay one person's share and 4 bugeng jointly pay one person's share; therefore it is called inverse proportions. When the numbers of persons are not the same, the numbers of parts are not homogenized, and one should make the denominators multiply the numerators crosswise (hucheng 互乘). When the denominators multiply the numerators crosswise, what is moved becomes the proportion of what is not moved. One may also first make their denominators the same and, each with its denominator, reduce its numerator, to make the inverse proportions. Add them together in a duplicate as the divisor; multiply the ones not added by the amount to be divided, each separately as the dividend; divide the dividend by the divisor.]\end{uniten}

\begin{unitzh}\unitlabel{jiuzhang-012.5}術曰：列置衰而令相乘，動者爲不動者衰。\end{unitzh}
\begin{uniten}Method: Set out the proportions in an array and make them multiply each other; what is moved becomes the proportion of what is not moved.\end{uniten}

\begin{unitzh}\unitlabel{jiuzhang-012.6}今有大夫、不更、簪褭、上造、公士，凡五人，共出百錢。欲令高爵出少，以次漸多。問各幾何？荅曰：大夫出八錢一百三十七分錢之一百四；不更出一十錢一百三十七分錢之一百三十；簪褭出一十四錢一百三十七分錢之八十二；上造出二十一錢一百三十七分錢之一百二十三；公士出四十三錢一百三十七分錢之一百九。\end{unitzh}
\begin{uniten}Suppose there are a dafu, a bugeng, a zanniao (簪褭), a shangzao (上造) and a gongshi (公士), altogether 5 persons, who jointly pay 100 qian (錢). It is desired that the higher rank pay less, becoming gradually more in order. Question: how much is each? Answer: The dafu pays 8 104/137 qian; the bugeng pays 10 130/137 qian; the zanniao pays 14 82/137 qian; the shangzao pays 21 123/137 qian; the gongshi pays 43 109/137 qian.\end{uniten}

\begin{unitzh}\unitlabel{jiuzhang-012.7}術曰：置爵數，各自爲衰，而返衰之。副并爲法。以百錢乘未并者，各自爲實。實如法得一錢。\end{unitzh}
\begin{uniten}Method: Set out the numbers of the ranks, each separately as a proportion, and make them inverse proportions. Add them together in a duplicate as the divisor. Multiply the ones not added by 100 qian, each separately as the dividend. Each time the dividend matches the divisor, one gets 1 qian.\end{uniten}

\begin{unitzh}\unitlabel{jiuzhang-012.8}今有甲持粟三升，乙持糲米三升，丙持糲飯三升。欲令合而分之，問各幾何？荅曰：甲二升一十分升之七；乙四升一十分升之五；丙一升一十分升之八。\end{unitzh}
\begin{uniten}Suppose there is A (jia 甲) holding 3 sheng of millet (su 粟), B (yi 乙) holding 3 sheng of coarse hulled grain (limi 糲米), and C (bing 丙) holding 3 sheng of coarse cooked grain (lifan 糲飯). It is desired to combine and divide them. Question: how much is each? Answer: A, 2 7/10 sheng; B, 4 5/10 sheng; C, 1 8/10 sheng.\end{uniten}

\begin{unitzh}\unitlabel{jiuzhang-012.9}術日：以粟率五十、糲米率三十、糲飯率七十五爲衰，而返衰之。副并爲法。以九升乘未并者，各自爲實。實如法得一升。{\zhsmall 按：此術三人所持升數雖等，論其本率，精麤不同。米率雖少，令最得多；飯率雖多，返使得少。故令返之，使精得多而麤得少。於今有術，副并爲所有率，未并者各爲所求率，九升爲所有數，而今有之，卽得。}\end{unitzh}
\begin{uniten}Method: Take the rate (lü 率) of millet, 50, the rate of coarse hulled grain, 30, and the rate of coarse cooked grain, 75, as the proportions, and make them inverse proportions. Add them together in a duplicate as the divisor. Multiply the ones not added by 9 sheng, each separately as the dividend. Each time the dividend matches the divisor, one gets 1 sheng. [Liu Hui: Note: in this method, although the numbers of sheng held by the three persons are equal, if one considers their basic rates (benlü 本率), the fine and the coarse are not the same. Although the rate of the hulled grain is small, it is made to get the most; although the rate of the cooked grain is large, it is conversely made to get little. Therefore they are made inverse, so that the fine gets more and the coarse gets less. In the "Suppose there is" method (jinyou shu 今有術), the duplicate sum is the rate of what one has (suoyou lü 所有率), the ones not added are each the rate of what one seeks (suoqiu lü 所求率), and 9 sheng is the number of what one has (suoyou shu 所有數); apply "Suppose there is" to them, and one obtains it.]\end{uniten}

\begin{unitzh}\unitlabel{jiuzhang-012.10}今有絲一斤價直二百四十，今有錢一千三百二十八，問得絲幾何？荅曰：五斤八兩一十二銖五分銖之四。\end{unitzh}
\begin{uniten}Suppose there is 1 jin (斤) of silk whose price is 240; now one has 1328 qian. Question: how much silk does one get? Answer: 5 jin 8 liang (兩) 12 4/5 zhu (銖).\end{uniten}

\begin{unitzh}\unitlabel{jiuzhang-012.11}術曰：以一斤價數爲法，以一斤乘今有錢數爲實，實如法得絲數。{\zhsmall 按：此術今有之義。以一斤價爲所有率，一斤爲所求率，今有錢爲所有數，而今有之，卽得。}\end{unitzh}
\begin{uniten}Method: Take the number of the price of 1 jin as the divisor; multiply the amount of cash one now has by 1 jin as the dividend; divide the dividend by the divisor to get the amount of silk. [Liu Hui: Note: this method has the sense of "Suppose there is". Take the price of 1 jin as the rate of what one has, 1 jin as the rate of what one seeks, and the cash one now has as the number of what one has; apply "Suppose there is" to them, and one obtains it.]\end{uniten}

\begin{unitzh}\unitlabel{jiuzhang-012.12}今有絲一斤價直三百四十五，今有絲七兩一十二銖，問得錢幾何？荅曰：一百六十一錢三十二分錢之二十三。\end{unitzh}
\begin{uniten}Suppose there is 1 jin of silk whose price is 345; now one has 7 liang 12 zhu of silk. Question: how much cash does one get? Answer: 161 23/32 qian.\end{uniten}

\begin{unitzh}\unitlabel{jiuzhang-012.13}術曰：以一斤銖數爲法，以一斤價數乘七兩一十二銖爲實，實如法得錢數。{\zhsmall 按：此術亦今有之義。以絲一斤銖數爲所有率，價錢爲所求率，今有絲爲所有數，而今有之，卽得。}\end{unitzh}
\begin{uniten}Method: Take the number of zhu in 1 jin as the divisor; multiply 7 liang 12 zhu by the number of the price of 1 jin as the dividend; divide the dividend by the divisor to get the amount of cash. [Liu Hui: Note: this method also has the sense of "Suppose there is". Take the number of zhu in 1 jin of silk as the rate of what one has, the price in cash as the rate of what one seeks, and the silk one now has as the number of what one has; apply "Suppose there is" to them, and one obtains it.]\end{uniten}

\begin{unitzh}\unitlabel{jiuzhang-012.14}今有縑一丈價直一百二十八，今有縑一匹九尺五寸，問得錢幾何？荅曰：六百三十三錢五分錢之三。\end{unitzh}
\begin{uniten}Suppose there is 1 zhang (丈) of fine silk (jian 縑) whose price is 128; now one has 1 pi (匹) 9 chi (尺) 5 cun (寸) of fine silk. Question: how much cash does one get? Answer: 633 3/5 qian.\end{uniten}

\begin{unitzh}\unitlabel{jiuzhang-012.15}術曰：以一丈寸數爲法，以價錢數乘今有縑寸數爲實，實如法得錢數。{\zhsmall 臣淳風等謹按：此術亦今有之義。以縑一丈寸數爲所有率，價錢爲所求率，今有縑寸數爲所有數，而今有之，卽得。}\end{unitzh}
\begin{uniten}Method: Take the number of cun in 1 zhang as the divisor; multiply the number of cun of fine silk one now has by the amount of the price in cash as the dividend; divide the dividend by the divisor to get the amount of cash. [Li Chunfeng et al.: Your servant Chunfeng and others respectfully note: this method also has the sense of "Suppose there is". Take the number of cun in 1 zhang of fine silk as the rate of what one has, the price in cash as the rate of what one seeks, and the number of cun of fine silk one now has as the number of what one has; apply "Suppose there is" to them, and one obtains it.]\end{uniten}

\begin{unitzh}\unitlabel{jiuzhang-012.16}今有布一匹價直一百二十五，今有布二丈七尺，問得錢幾何？荅曰：八十四錢八分錢之三。\end{unitzh}
\begin{uniten}Suppose there is 1 pi of cloth (bu 布) whose price is 125; now one has 2 zhang 7 chi of cloth. Question: how much cash does one get? Answer: 84 3/8 qian.\end{uniten}

\begin{unitzh}\unitlabel{jiuzhang-012.17}術曰：以一匹尺數爲法，今有布尺數乘價錢爲實，實如法得錢數。{\zhsmall 按：此術亦今有之義。以一匹尺數爲所有率，價錢爲所求率，今有布爲所有數，今有之，卽得。}\end{unitzh}
\begin{uniten}Method: Take the number of chi in 1 pi as the divisor; multiply the price in cash by the number of chi of cloth one now has as the dividend; divide the dividend by the divisor to get the amount of cash. [Liu Hui: Note: this method also has the sense of "Suppose there is". Take the number of chi in 1 pi as the rate of what one has, the price in cash as the rate of what one seeks, and the cloth one now has as the number of what one has; apply "Suppose there is" to them, and one obtains it.]\end{uniten}

\begin{unitzh}\unitlabel{jiuzhang-012.18}今有素一匹一丈價直六百二十五，今有錢五百，問得素幾何？荅曰：得素一匹。\end{unitzh}
\begin{uniten}Suppose there is 1 pi 1 zhang of plain silk (su 素) whose price is 625; now one has 500 qian. Question: how much plain silk does one get? Answer: One gets 1 pi of plain silk.\end{uniten}

\begin{unitzh}\unitlabel{jiuzhang-012.19}術日：以價直爲法，以一匹一丈尺數乘今有錢數爲實，實如法得素數。{\zhsmall 按：此術亦今有之義。以價錢爲所有率，五丈尺數爲所求率，今有錢爲所有數，今有之，卽得。}\end{unitzh}
\begin{uniten}Method: Take the price as the divisor; multiply the amount of cash one now has by the number of chi in 1 pi 1 zhang as the dividend; divide the dividend by the divisor to get the amount of plain silk. [Liu Hui: Note: this method also has the sense of "Suppose there is". Take the price in cash as the rate of what one has, the number of chi in 5 zhang as the rate of what one seeks, and the cash one now has as the number of what one has; apply "Suppose there is" to them, and one obtains it.]\end{uniten}

\begin{unitzh}\unitlabel{jiuzhang-012.20}今有與人絲一十四斤，約得縑一十斤。今與人絲四十五斤八兩，問得縑幾何？荅曰：三十二斤八兩。\end{unitzh}
\begin{uniten}Suppose there is the giving of 14 jin of silk to a person, for which by agreement one gets 10 jin of fine silk. Now one gives a person 45 jin 8 liang of silk. Question: how much fine silk does one get? Answer: 32 jin 8 liang.\end{uniten}

\begin{unitzh}\unitlabel{jiuzhang-012.21}術曰：以一十四斤兩數爲法，以一十斤乘今有絲兩數爲實，實如法得縑數。{\zhsmall 此術亦今有之義。以一十四斤兩數爲所有率，一十斤爲所求率，今有絲爲所有數，今有之，卽得。}\end{unitzh}
\begin{uniten}Method: Take the number of liang in 14 jin as the divisor; multiply the number of liang of silk one now has by 10 jin as the dividend; divide the dividend by the divisor to get the amount of fine silk. [Liu Hui: This method also has the sense of "Suppose there is". Take the number of liang in 14 jin as the rate of what one has, 10 jin as the rate of what one seeks, and the silk one now has as the number of what one has; apply "Suppose there is" to them, and one obtains it.]\end{uniten}

\begin{unitzh}\unitlabel{jiuzhang-012.22}今有絲一斤耗七兩，今有絲二十三斤五兩，問耗幾何？荅曰：一百六十三兩四銖半。\end{unitzh}
\begin{uniten}Suppose there is 1 jin of silk that loses (hao 耗) 7 liang; now one has 23 jin 5 liang of silk. Question: how much is the loss? Answer: 163 liang 4 1/2 zhu.\end{uniten}

\begin{unitzh}\unitlabel{jiuzhang-012.23}術曰：以一斤展十六兩爲法，以七兩乘今有絲兩數爲實，實如法得耗數。{\zhsmall 按：此術亦今有之義。以一斤爲十六兩爲所有率，七兩爲所求率，今有絲爲所有數，而今有之，卽得。}\end{unitzh}
\begin{uniten}Method: Take 1 jin, expanded (zhan 展) into 16 liang, as the divisor; multiply the number of liang of silk one now has by 7 liang as the dividend; divide the dividend by the divisor to get the amount of loss. [Liu Hui: Note: this method also has the sense of "Suppose there is". Take 1 jin, made 16 liang, as the rate of what one has, 7 liang as the rate of what one seeks, and the silk one now has as the number of what one has; apply "Suppose there is" to them, and one obtains it.]\end{uniten}

\begin{unitzh}\unitlabel{jiuzhang-012.24}今有生絲三十斤，乾之，耗三斤十二兩。今有乾絲一十二斤，問生絲幾何？荅曰：一十三斤一十一兩十銖七分銖之二。\end{unitzh}
\begin{uniten}Suppose there are 30 jin of raw silk; on drying it, it loses 3 jin 12 liang. Now one has 12 jin of dried silk. Question: how much is the raw silk? Answer: 13 jin 11 liang 10 2/7 zhu.\end{uniten}

\begin{unitzh}\unitlabel{jiuzhang-012.25}術曰：置生絲兩數，除耗數，餘以爲法。{\zhsmall 餘四百二十兩，卽乾絲率。}三十斤乘乾絲兩數爲實，實如法得生絲數。\end{unitzh}
\begin{uniten}Method: Set out the number of liang of raw silk, subtract the amount of loss, and take the remainder as the divisor. [Liu Hui: The remainder, 420 liang, is the rate of dried silk.] Multiply the number of liang of dried silk by 30 jin as the dividend; divide the dividend by the divisor to get the amount of raw silk.\end{uniten}

\begin{unitzh}\unitlabel{jiuzhang-012.26}{\zhsmall 凡所謂率者，細則俱細，麤則俱麤，兩數相推而已。故品物不同，如上縑、絲之比，得相與乘焉。三十斤凡四百八十兩，今生絲率四百八十兩，乾絲率四百二十兩，則其數相通，可俱爲銖，可俱爲兩，可俱爲斤，無所歸滯也。若然，宜以所有乾絲斤數乘生絲兩數爲實。今斤兩錯互而亦同歸者，使乾絲以兩數爲率，生絲以斤數爲率。譬之異類，亦各有一定之勢。}\end{unitzh}
\begin{uniten}[Liu Hui: In general, what are called rates are such that if fine (xi 細), both are fine, and if coarse (cu 麤), both are coarse; it is merely a matter of two numbers inferred from each other. Therefore, when the kinds of things are different, as in the ratio (bi 比) of fine silk and silk above, they may be multiplied with each other. 30 jin is in all 480 liang; now the rate of raw silk is 480 liang and the rate of dried silk is 420 liang; then their numbers are mutually convertible (xiangtong 相通): they may both be made zhu, may both be made liang, may both be made jin, and there is nothing at which they are obstructed. If so, it would be fitting to multiply the number of liang of raw silk by the number of jin of dried silk one has, as the dividend. Now, that jin and liang are interchanged crosswise and yet arrive at the same is a matter of letting the dried silk take the number of liang as its rate and the raw silk take the number of jin as its rate. It is like different kinds: each also has a fixed configuration (shi 勢).]\end{uniten}

\begin{unitzh}\unitlabel{jiuzhang-012.27}{\zhsmall 臣淳風等謹按：此術置生絲兩數，除耗數，餘卽乾絲之率，於今有術爲所有率；三十斤爲所求率；乾絲兩數爲所有數。凡所謂率者，細則俱細，麤則俱麤。今以斤乘兩者，乾絲卽以兩數爲率，生絲卽以斤數爲率，譬之異物，各有一定之率也。}\end{unitzh}
\begin{uniten}[Li Chunfeng et al.: Your servant Chunfeng and others respectfully note: in this method one sets out the number of liang of raw silk and subtracts the amount of loss; the remainder is the rate of dried silk, which in the "Suppose there is" method is the rate of what one has; 30 jin is the rate of what one seeks; the number of liang of dried silk is the number of what one has. In general, what are called rates are such that if fine, both are fine, and if coarse, both are coarse. Now, as for multiplying liang by jin: the dried silk then takes the number of liang as its rate, and the raw silk then takes the number of jin as its rate; it is like different things, each having a fixed rate.]\end{uniten}

\begin{unitzh}\unitlabel{jiuzhang-013.1}今有田一畝，收粟六升太半升。今有田一頃二十六畝一百五十九步，問收粟幾何？荅曰：八斛四斗四升一十二分升之五。\end{unitzh}
\begin{uniten}Suppose there is a field of 1 mu (畝) that yields 6 2/3 sheng (升) of millet (su 粟). Suppose there is a field of 1 qing (頃) 26 mu 159 bu (步). Question: how much millet does it yield? Answer: 8 hu (斛) 4 dou (斗) 4 5/12 sheng.\end{uniten}

\begin{unitzh}\unitlabel{jiuzhang-013.2}術曰：以畝二百四十步爲法，以六升太半升乘今有田積步爲實，實如法得粟數。{\zhsmall 按此術亦今有之義。以一畝步數爲所有率，六升太半升爲所求率，今有田積步爲所有數，而今有之，卽得。}\end{unitzh}
\begin{uniten}Method: Take the 240 bu of a mu as the divisor (fa). With 6 2/3 sheng multiply the area in bu of the field one now has, making the dividend (shi). Divide the dividend by the divisor to obtain the amount of millet. [Liu Hui: Note: this method too has the sense of Suppose-there-is (jinyou). Take the number of bu in 1 mu as the rate of what one has (suoyou lü), 6 2/3 sheng as the rate of what one seeks (suoqiu lü), and the area in bu of the field one now has as the number one has (suoyou shu). Then apply Suppose-there-is to them, and it is obtained at once.]\end{uniten}

\begin{unitzh}\unitlabel{jiuzhang-013.3}今有取保一歲，價錢二千五百。今先取一千二百，問當作日幾何？荅曰：一百六十九日二十五分日之二十三。\end{unitzh}
\begin{uniten}Suppose there is the hiring of a laborer (qubao 取保) for 1 year, at a price of 2500 qian (錢). Now 1200 is taken in advance. Question: how many days ought to be worked? Answer: 169 23/25 days.\end{uniten}

\begin{unitzh}\unitlabel{jiuzhang-013.4}術曰：以價錢爲法，以一歲三百五十四日乘先取錢數爲實，實如法得日數。{\zhsmall 按此術亦今有之義。以價爲所有率，一歲日數爲所求率，取錢爲所有數，而今有之，卽得。}\end{unitzh}
\begin{uniten}Method: Take the price in qian as the divisor (fa). With the 354 days of 1 year multiply the number of qian taken in advance, making the dividend (shi). Divide the dividend by the divisor to obtain the number of days. [Liu Hui: Note: this method too has the sense of Suppose-there-is. Take the price as the rate of what one has, the number of days in 1 year as the rate of what one seeks, and the qian taken as the number one has. Then apply Suppose-there-is to them, and it is obtained at once.]\end{uniten}

\begin{unitzh}\unitlabel{jiuzhang-013.5}今有貸人千錢，月息三十。今有貸人七百五十錢，九日歸之，問息幾何？荅曰：六錢四分錢之三。\end{unitzh}
\begin{uniten}Suppose there is a loan of 1000 qian to someone, with an interest of 30 per month. Suppose there is a loan of 750 qian to someone, which is returned after 9 days. Question: how much is the interest? Answer: 6 3/4 qian.\end{uniten}

\begin{unitzh}\unitlabel{jiuzhang-013.6}術曰：以月三十日乘千錢爲法。{\zhsmall 以三十日乘千錢爲法者，得三萬，是爲貸人錢三萬，一日息三十也。}以息三十乘今所貸錢數，又以九日乘之，爲實。實如法得一錢。{\zhsmall 以九日乘今所貸錢，爲今一日所有錢，於今有術爲所有數；息三十爲所求率；三萬錢爲所有率。此又可以一月三十日約息三十錢，爲十分一日，以乘今一日所有錢爲實，千錢爲法。爲率者，當等之於一也。故三十日或可乘本，或可約息，皆所以等之也。}\end{unitzh}
\begin{uniten}Method: With the 30 days of a month multiply 1000 qian, making the divisor (fa). [Liu Hui: Taking 30 days times 1000 qian as the divisor gives 30000. This amounts to lending someone 30000 qian with an interest of 30 for one day.] With the interest of 30 multiply the number of qian now lent, and further multiply it by 9 days, making the dividend (shi). Divide the dividend by the divisor, obtaining 1 qian for each time it contains the divisor. [Liu Hui: Multiplying the qian now lent by 9 days makes the qian one now has for one day. In the Suppose-there-is method (jinyou shu 今有術) this is the number one has; the interest of 30 is the rate of what one seeks; the 30000 qian is the rate of what one has. Here one may also reduce (yue 約) the interest of 30 qian by the 30 days of 1 month, making 1/10 for 1 day. With this, multiply the qian one now has for one day, making the dividend, with 1000 qian as the divisor. As for forming rates, one should equalize them to one. Therefore the 30 days may either multiply the principal (ben 本) or reduce the interest; both are means of equalizing them.] [Translator's note: 爲十分一日, rendered "making 1/10 for 1 day", is obscure as it stands.]\end{uniten}

\begin{unitzh}\unitlabel{jiuzhang-013.7}九章算術卷三\end{unitzh}
\begin{uniten}The Nine Chapters on the Mathematical Art, juan (卷) 3.\end{uniten}

\section*{卷第四}\addcontentsline{toc}{section}{卷第四}

\begin{unitzh}\unitlabel{jiuzhang-014.1}《九章算術》卷四{\zhsmall 《算經十書》之二} 魏劉徽注 唐朝議大夫、行太史令、上輕車都尉臣李淳風等奉勅注釋\end{unitzh}
\begin{uniten}The Nine Chapters on the Mathematical Art, Juan 4 [Note: The second of the Ten Mathematical Classics]. Wei: commentary by Liu Hui. Tang: annotated by imperial command by the Grand Master for Court Discussion, acting Grand Astrologer, Supreme Commandant of Light Chariots, Your servant Li Chunfeng et al.\end{uniten}

\begin{unitzh}\unitlabel{jiuzhang-014.2}少廣{\zhsmall 以御積羃方圓}\end{unitzh}
\begin{uniten}Small Breadth (shaoguang) [Liu Hui: By which to deal with areas (jimi), squares and circles.]\end{uniten}

\begin{unitzh}\unitlabel{jiuzhang-014.3}少廣{\zhsmall 臣淳風等謹按：一畝之田，廣一步，長二百四十步。今欲截取其從少，以益其廣，故曰少廣。}\end{unitzh}
\begin{uniten}Small Breadth [Li Chunfeng et al.: Your servants Chunfeng et al. respectfully note: A field of 1 mu (畝) is 1 bu (步) broad (guang) and 240 bu long (chang). Now one wishes to cut off [part of] its length (cong), making it less, in order to add to its breadth; therefore it is called Small Breadth.]\end{uniten}

\begin{unitzh}\unitlabel{jiuzhang-014.4}術曰：置全步及分母子，以最下分母徧乘諸分子及全步，{\zhsmall 臣淳風等謹按：以分母乘全步者，通其分也；以母乘子者，齊其子也。}各以其母除其子，置之於左；命通分者，又以分母徧乘諸分子及已通者，皆通而同之，并之爲法。{\zhsmall 臣淳風等謹按：諸子悉通，故可并之爲法。亦不宜用合分術，列數尤多，若用乘則算數至䌓，故别制此術，從省約。}置所求步數，以全步積分乘之爲實。{\zhsmall 此以田廣爲法，以畝積步爲實。法有分者，當同其母，齊其子，以同乘法實，而并齊於法。今以分母乘全步及子，子如母而一，竝以并全法，則法實俱長，意亦等也。故如法而一，得從步數。}實如法而一，得從步。\end{unitzh}
\begin{uniten}Method: Put down the whole bu (quanbu) and the denominators and numerators of the fractions (fen mu zi). With the bottom-most denominator multiply throughout all the numerators and the whole bu. [Li Chunfeng et al.: Your servants Chunfeng et al. respectfully note: Multiplying the whole bu by the denominator is to convert (tong) it into its parts (fen); multiplying the numerators by the denominator is to homogenize (qi) the numerators.] Divide each numerator by its own denominator, and place [the results] on the left. For those designated (ming) as fractions to be converted (tongfen), again with the denominator multiply throughout all the numerators and those already converted. When all have been converted and made the same, add them to make the divisor (fa). [Li Chunfeng et al.: Your servants Chunfeng et al. respectfully note: All the numerators have been converted, therefore they can be added to make the divisor. It is also not suitable to use the method for combining fractions (hefen shu): the numbers laid out are especially many, and if multiplication is used, the counting-rod numbers become extremely complicated; therefore this method was devised separately, for the sake of economy and brevity.] Put down the number of bu sought, and multiply it by the accumulated parts (jifen) of the whole bu to make the dividend (shi). [Liu Hui: Here the breadth of the field is taken as the divisor, and the area (ji) of 1 mu in bu as the dividend. When the divisor has fractions, one should make their denominators the same (tong) and homogenize their numerators, multiply the divisor and the dividend by the common [denominator] (tong), and add the homogenized [numerators] (qi) into the divisor. Now the whole bu and the numerators are multiplied by the denominator; the numerators are divided by the denominator, and all are added into the whole divisor (quanfa); then divisor and dividend are both enlarged, and the meaning is likewise equal. Therefore dividing by the divisor gives the number of bu of the length.] Divide the dividend by the divisor, obtaining the length in bu.\end{uniten}

\begin{unitzh}\unitlabel{jiuzhang-014.5}今有田廣一步半，求田一畝，問從幾何？荅曰：一百六十步。\end{unitzh}
\begin{uniten}Suppose there is a field whose breadth is 1 1/2 bu. One seeks a field of 1 mu. What is the length? Answer: 160 bu.\end{uniten}

\begin{unitzh}\unitlabel{jiuzhang-014.6}術曰：下有半，是二分之一。以一爲二，半爲一，并之得三，爲法。置田二百四十步，亦以一爲二乘之，爲實。實如法得從步。\end{unitzh}
\begin{uniten}Method: Below there is a half; this is 1/2. Take 1 as 2 and the half as 1; adding them gives 3, which makes the divisor. Put down the field, 240 bu, and likewise multiply it by taking 1 as 2, to make the dividend. Divide the dividend by the divisor, obtaining the length in bu.\end{uniten}

\begin{unitzh}\unitlabel{jiuzhang-014.7}今有田廣一步半、三分步之一，求田一畝，問從幾何？荅曰：一百三十步一十一分步之一十。\end{unitzh}
\begin{uniten}Suppose there is a field whose breadth is 1 1/2 bu and 1/3 bu. One seeks a field of 1 mu. What is the length? Answer: 130 10/11 bu.\end{uniten}

\begin{unitzh}\unitlabel{jiuzhang-014.8}術曰：下有三分，以一爲六，半爲三，三分之一爲二，并之得一十一，爲法。置田二百四十步，亦以一爲六乘之，爲實。實如法得從歩。\end{unitzh}
\begin{uniten}Method: Below there are thirds. Take 1 as 6, the half as 3, 1/3 as 2; adding them gives 11, which makes the divisor. Put down the field, 240 bu, and likewise multiply it by taking 1 as 6, to make the dividend. Divide the dividend by the divisor, obtaining the length in bu.\end{uniten}

\begin{unitzh}\unitlabel{jiuzhang-014.9}今有田廣一步半、三分步之一、四分步之一，求田一畝，問從幾何？荅曰：一百一十五步五分步之一。\end{unitzh}
\begin{uniten}Suppose there is a field whose breadth is 1 1/2 bu, 1/3 bu and 1/4 bu. One seeks a field of 1 mu. What is the length? Answer: 115 1/5 bu.\end{uniten}

\begin{unitzh}\unitlabel{jiuzhang-014.10}術曰：下有四分，以一爲一十二，半爲六，三分之一爲四，四分之一爲三，并之得二十五，以爲法。置田二百四十步，亦以一爲一十二乘之，爲實。實如法而一，得從步。\end{unitzh}
\begin{uniten}Method: Below there are quarters. Take 1 as 12, the half as 6, 1/3 as 4, 1/4 as 3; adding them gives 25, which is taken as the divisor. Put down the field, 240 bu, and likewise multiply it by taking 1 as 12, to make the dividend. Divide the dividend by the divisor, obtaining the length in bu.\end{uniten}

\begin{unitzh}\unitlabel{jiuzhang-014.11}今有田廣一步半、三分步之一、四分步之一、五分步之一，求田一畝，問從幾何？荅曰：一百五步一百三十七分步之一十五。\end{unitzh}
\begin{uniten}Suppose there is a field whose breadth is 1 1/2 bu, 1/3 bu, 1/4 bu and 1/5 bu. One seeks a field of 1 mu. What is the length? Answer: 105 15/137 bu.\end{uniten}

\begin{unitzh}\unitlabel{jiuzhang-014.12}術曰：下有五分，以一爲六十，半爲三十，三分之一爲二十，四分之一爲一十五，五分之一爲一十二，并之得一百三十七，以爲法。置田二百四十步，亦以一爲六十乘之，爲實。實如法得從步。\end{unitzh}
\begin{uniten}Method: Below there are fifths. Take 1 as 60, the half as 30, 1/3 as 20, 1/4 as 15, 1/5 as 12; adding them gives 137, which is taken as the divisor. Put down the field, 240 bu, and likewise multiply it by taking 1 as 60, to make the dividend. Divide the dividend by the divisor, obtaining the length in bu.\end{uniten}

\begin{unitzh}\unitlabel{jiuzhang-014.13}今有田廣一步半、三分步之一、四分步之一、五分步之一、六分步之一，求田一畝，問從幾何？荅曰：九十七步四十九分步之四十七。\end{unitzh}
\begin{uniten}Suppose there is a field whose breadth is 1 1/2 bu, 1/3 bu, 1/4 bu, 1/5 bu and 1/6 bu. One seeks a field of 1 mu. What is the length? Answer: 97 47/49 bu.\end{uniten}

\begin{unitzh}\unitlabel{jiuzhang-014.14}術曰：下有六分，以一爲一百二十，半爲六十，三分之一爲四十，四分之一爲三十，五分之一爲二十四，六分之一爲二十，并之得二百九十四，以爲法。置田二百四十歩，亦以一爲一百二十乘之，爲實。實如法得從步。\end{unitzh}
\begin{uniten}Method: Below there are sixths. Take 1 as 120, the half as 60, 1/3 as 40, 1/4 as 30, 1/5 as 24, 1/6 as 20; adding them gives 294, which is taken as the divisor. Put down the field, 240 bu, and likewise multiply it by taking 1 as 120, to make the dividend. Divide the dividend by the divisor, obtaining the length in bu.\end{uniten}

\begin{unitzh}\unitlabel{jiuzhang-014.15}今有田廣一步半、三分步之一、四分步之一、五分步之一、六分步之一、七分步之一，求田一畝，問從幾何？荅曰：九十二步一百二十一分步之六十八。\end{unitzh}
\begin{uniten}Suppose there is a field whose breadth is 1 1/2 bu, 1/3 bu, 1/4 bu, 1/5 bu, 1/6 bu and 1/7 bu. One seeks a field of 1 mu. What is the length? Answer: 92 68/121 bu.\end{uniten}

\begin{unitzh}\unitlabel{jiuzhang-014.16}術曰：下有七分，以一爲四百二十，半爲二百一十，三分之一爲一百四十，四分之一爲一百五，五分之一爲八十四，六分之一爲七十，七分之一爲六十，并之得一千八十九，以爲法。置田二百四十步，亦以一爲四百二十乘之，爲實。實如法得從步。\end{unitzh}
\begin{uniten}Method: Below there are sevenths. Take 1 as 420, the half as 210, 1/3 as 140, 1/4 as 105, 1/5 as 84, 1/6 as 70, 1/7 as 60; adding them gives 1089, which is taken as the divisor. Put down the field, 240 bu, and likewise multiply it by taking 1 as 420, to make the dividend. Divide the dividend by the divisor, obtaining the length in bu.\end{uniten}

\begin{unitzh}\unitlabel{jiuzhang-014.17}今有田廣一步半、三分步之一、四分步之一、五分步之一、六分步之一、七分步之一、八分步之一，求田一畝，問從幾何？荅曰：八十八步七百六十一分步之二百三十二。\end{unitzh}
\begin{uniten}Suppose there is a field whose breadth is 1 1/2 bu, 1/3 bu, 1/4 bu, 1/5 bu, 1/6 bu, 1/7 bu and 1/8 bu. One seeks a field of 1 mu. What is the length? Answer: 88 232/761 bu.\end{uniten}

\begin{unitzh}\unitlabel{jiuzhang-014.18}術曰：下有八分，以一爲八百四十，半爲四百二十，三分之一爲二百八十，四分之一爲二百一十，五分之一爲一百六十八，六分之一爲一百四十，七分之一爲一百二十，八分之一爲一百五，并之得二千二百八十三，以爲法。置田二百四十步，亦以一爲八百四十乘之，爲實。實如法得從步。\end{unitzh}
\begin{uniten}Method: Below there are eighths. Take 1 as 840, the half as 420, 1/3 as 280, 1/4 as 210, 1/5 as 168, 1/6 as 140, 1/7 as 120, 1/8 as 105; adding them gives 2283, which is taken as the divisor. Put down the field, 240 bu, and likewise multiply it by taking 1 as 840, to make the dividend. Divide the dividend by the divisor, obtaining the length in bu.\end{uniten}

\begin{unitzh}\unitlabel{jiuzhang-014.19}今有田廣一步半、三分步之一、四分步之一、五分步之一、六分步之一、七分步之一、八分步之一、九分步之一，求田一畝，問從幾何？荅曰：八十四步七千一百二十九分步之五千九百六十四。\end{unitzh}
\begin{uniten}Suppose there is a field whose breadth is 1 1/2 bu, 1/3 bu, 1/4 bu, 1/5 bu, 1/6 bu, 1/7 bu, 1/8 bu and 1/9 bu. One seeks a field of 1 mu. What is the length? Answer: 84 5964/7129 bu.\end{uniten}

\begin{unitzh}\unitlabel{jiuzhang-014.20}術曰：下有九分，以一爲二千五百二十，半爲一千二百六十，三分之一爲八百四十，四分之一爲六百三十，五分之一爲五百四，六分之一爲四百二十，七分之一爲三百六十，八分之一爲三百一十五，九分之一爲二百八十，并之得七千一百二十九，以爲法。置田二百四十步，亦以一爲二千五百二十乘之，爲實。實如法得從步。\end{unitzh}
\begin{uniten}Method: Below there are ninths. Take 1 as 2520, the half as 1260, 1/3 as 840, 1/4 as 630, 1/5 as 504, 1/6 as 420, 1/7 as 360, 1/8 as 315, 1/9 as 280; adding them gives 7129, which is taken as the divisor. Put down the field, 240 bu, and likewise multiply it by taking 1 as 2520, to make the dividend. Divide the dividend by the divisor, obtaining the length in bu.\end{uniten}

\begin{unitzh}\unitlabel{jiuzhang-014.21}今有田廣一步半、三分步之一、四分步之一、五分步之一、六分步之一、七分步之一、八分步之一、九分步之一、十分步之一，求田一畝，問從幾何？荅曰：八十一步七千三百八十一分步之六千九百三十九。\end{unitzh}
\begin{uniten}Suppose there is a field whose breadth is 1 1/2 bu, 1/3 bu, 1/4 bu, 1/5 bu, 1/6 bu, 1/7 bu, 1/8 bu, 1/9 bu and 1/10 bu. One seeks a field of 1 mu. What is the length? Answer: 81 6939/7381 bu.\end{uniten}

\begin{unitzh}\unitlabel{jiuzhang-014.22}術曰：下有一十分，以一爲二千五百二十，半爲一千二百六十，三分之一爲八百四十，四分之一爲六百三十，五分之一爲五百四，六分之一爲四百二十，七分之一爲三百六十，八分之一爲三百一十五，九分之一爲二百八十，十分之一爲二百五十二，并之得七千三百八十一，以爲法。置田二百四十步，亦以一爲二千五百二十乘之，爲實。實如法得從步。\end{unitzh}
\begin{uniten}Method: Below there are tenths. Take 1 as 2520, the half as 1260, 1/3 as 840, 1/4 as 630, 1/5 as 504, 1/6 as 420, 1/7 as 360, 1/8 as 315, 1/9 as 280, 1/10 as 252; adding them gives 7381, which is taken as the divisor. Put down the field, 240 bu, and likewise multiply it by taking 1 as 2520, to make the dividend. Divide the dividend by the divisor, obtaining the length in bu.\end{uniten}

\begin{unitzh}\unitlabel{jiuzhang-015.1}今有田廣一步、半、三分步之一、四分步之一、五分步之一、六分步之一、七分步之一、八分步之一、九分步之一、十分步之一、十一分步之一，求田一畝。問從幾何？荅曰：七十九步八萬三千七百一十一分步之三萬九千六百三十一。\end{unitzh}
\begin{uniten}Suppose there is a field whose width (guang) is 1 bu (步), 1/2, 1/3 bu, 1/4 bu, 1/5 bu, 1/6 bu, 1/7 bu, 1/8 bu, 1/9 bu, 1/10 bu, and 1/11 bu; one seeks a field of 1 mu (畝). Question: what is the length (cong)? Answer: 79 39631/83711 bu.\end{uniten}

\begin{unitzh}\unitlabel{jiuzhang-015.2}術曰：下有一十一分，以一爲二萬七千七百二十，半爲一萬三千八百六十，三分之一爲九千二百四十，四分之一爲六千九百三十，五分之一爲五千五百四十四，六分之一爲四千六百二十，七分之一爲三千九百六十，八分之一爲三千四百六十五，九分之一爲三千八十，一十分之一爲二千七百七十二，一十一分之一爲二千五百二十，并之得八萬三千七百一十一，以爲法。置田二百四十步，亦以一爲二萬七千七百二十乘之，爲實。實如法得從步。\end{unitzh}
\begin{uniten}Method: Below there are 11 parts. Take 1 as 27720; then the half is 13860, 1/3 is 9240, 1/4 is 6930, 1/5 is 5544, 1/6 is 4620, 1/7 is 3960, 1/8 is 3465, 1/9 is 3080, 1/10 is 2772, 1/11 is 2520. Add them together to obtain 83711, which is taken as the divisor (fa). Set out the field of 240 bu, and likewise multiply it by taking 1 as 27720, to make the dividend (shi). Divide the dividend by the divisor to obtain the length in bu.\end{uniten}

\begin{unitzh}\unitlabel{jiuzhang-015.3}今有田廣一步、半、三分步之一、四分步之一、五分步之一、六分步之一、七分步之一、八分步之一、九分步之一、十分步之一、十一分步之一、十二分步之一，求田一畝。問從幾何？荅曰：七十七步八萬六千二十一分步之二萬九千一百八十三。\end{unitzh}
\begin{uniten}Suppose there is a field whose width is 1 bu, 1/2, 1/3 bu, 1/4 bu, 1/5 bu, 1/6 bu, 1/7 bu, 1/8 bu, 1/9 bu, 1/10 bu, 1/11 bu, and 1/12 bu; one seeks a field of 1 mu. Question: what is the length? Answer: 77 29183/86021 bu.\end{uniten}

\begin{unitzh}\unitlabel{jiuzhang-015.4}術曰：下有一十二分，以一爲八萬三千一百六十，半爲四萬一千五百八十，三分之一爲二萬七千七百二十，四分之一爲二萬七百九十，五分之一爲一萬六千六百三十二，六分之一爲一萬三千八百六十，七分之一爲一萬一千八百八十，八分之一爲一萬三百九十五，九分之一爲九千二百四十，一十分之一爲八千三百一十六，十一分之一爲七千五百六十，十二分之一爲六千九百三十，并之得二十五萬八千六十三，以爲法。置田二百四十步，亦以一爲八萬三千一百六十乘之，爲實。實如法得從步。\end{unitzh}
\begin{uniten}Method: Below there are 12 parts. Take 1 as 83160; then the half is 41580, 1/3 is 27720, 1/4 is 20790, 1/5 is 16632, 1/6 is 13860, 1/7 is 11880, 1/8 is 10395, 1/9 is 9240, 1/10 is 8316, 1/11 is 7560, 1/12 is 6930. Add them together to obtain 258063, which is taken as the divisor. Set out the field of 240 bu, and likewise multiply it by taking 1 as 83160, to make the dividend. Divide the dividend by the divisor to obtain the length in bu.\end{uniten}

\begin{unitzh}\unitlabel{jiuzhang-015.5}{\zhsmall 臣淳風等謹按：凡爲術之意，約省爲善。宜云：下有一十二分，以一爲二萬七千七百二十，半爲一萬三千八百六十，三分之一爲九千二百四十，四分之一爲六千九百三十，五分之一爲五千五百四十四，六分之一爲四千六百二十，七分之一爲三千九百六十，八分之一爲三千四百六十五，九分之一爲三千八十，十分之一爲二千七百七十二，十一分之一爲二千五百二十，十二分之一爲二千三百一十，并之得八萬六千二十一，以爲法。置田二百四十步，亦以一爲二萬七千七百二十乘之，以爲實。實如法得從步。其術亦得，知不繁也。}\end{unitzh}
\begin{uniten}[Li Chunfeng et al.: In general, the intent in making a method is that conciseness and economy are best. It should say: Below there are 12 parts. Take 1 as 27720; then the half is 13860, 1/3 is 9240, 1/4 is 6930, 1/5 is 5544, 1/6 is 4620, 1/7 is 3960, 1/8 is 3465, 1/9 is 3080, 1/10 is 2772, 1/11 is 2520, 1/12 is 2310. Add them together to obtain 86021, which is taken as the divisor. Set out the field of 240 bu, and likewise multiply it by taking 1 as 27720, to make the dividend. Divide the dividend by the divisor to obtain the length in bu. This method also works, and one knows that it is not cumbersome.]\end{uniten}

\begin{unitzh}\unitlabel{jiuzhang-015.6}今有積五萬五千二百二十五步，問爲方幾何？荅曰：二百三十五步。\end{unitzh}
\begin{uniten}Suppose there is an area of 55225 bu. Question: what is [the side of] the square? Answer: 235 bu.\end{uniten}

\begin{unitzh}\unitlabel{jiuzhang-015.7}又有積二萬五千二百八十一步，問爲方幾何？荅曰：一百五十九步。\end{unitzh}
\begin{uniten}Again, there is an area of 25281 bu. Question: what is [the side of] the square? Answer: 159 bu.\end{uniten}

\begin{unitzh}\unitlabel{jiuzhang-015.8}又有積七萬一千八百二十四步，問爲方幾何？荅曰：二百六十八步。\end{unitzh}
\begin{uniten}Again, there is an area of 71824 bu. Question: what is [the side of] the square? Answer: 268 bu.\end{uniten}

\begin{unitzh}\unitlabel{jiuzhang-015.9}又有積五十六萬四千七百五十二步四分步之一，問爲方幾何？荅曰：七百五十一步半。\end{unitzh}
\begin{uniten}Again, there is an area of 564752 1/4 bu. Question: what is [the side of] the square? Answer: 751 1/2 bu.\end{uniten}

\begin{unitzh}\unitlabel{jiuzhang-015.10}又有積三十九億七千二百一十五萬六百二十五步，問爲方幾何？荅曰：六萬三千二十五步。\end{unitzh}
\begin{uniten}Again, there is an area of 3972150625 bu. Question: what is [the side of] the square? Answer: 63025 bu.\end{uniten}

\begin{unitzh}\unitlabel{jiuzhang-015.11}開方{\zhsmall 求方羃之一面也。}\end{unitzh}
\begin{uniten}Extracting the square root (kaifang) [Liu Hui: This is seeking one side (mian) of a square area (mi).]\end{uniten}

\begin{unitzh}\unitlabel{jiuzhang-015.12}術曰：置積爲實。借一算，步之，超一等。{\zhsmall 言百之面十也，言萬之面百也。}議所得，以一乘所借一算爲法，而以除。{\zhsmall 先得黃甲之面，上下相命，是自乘而除也。}除已，倍法爲定法。{\zhsmall 倍之者，豫張兩面朱羃定袤，以待復除，故曰定法。}其復除，折法而下。{\zhsmall 欲除朱羃者，本當副置所得成方，倍之爲定法，以折議乘而以除。如是當復步之而止，乃得相命，故使就上折下。}復置借算，步之如初，以復議一乘之，{\zhsmall 欲除朱羃之角黃乙之羃，其意如初之所得也。}所得副，以加定法，以除。以所得副從定法。{\zhsmall 再以黄乙之面加定法者，是則張兩青羃之袤。}復除，折下如前。\end{unitzh}
\begin{uniten}Method: Set out the area as the dividend (shi). Borrow 1 counting rod (suan) and step it forward (bu zhi), skipping one rank (chao yi deng). [Liu Hui: This says that the side of 100 is 10, and says that the side of 10000 is 100.] Estimate (yi) what is obtained, and multiply the 1 borrowed rod by it once to make the divisor (fa); then divide by it. [Liu Hui: First one obtains the side of the yellow A (huang jia); the upper and the lower name each other (xiang ming): this is multiplying by itself and removing.] When the division is done, double the divisor to make the fixed divisor (ding fa). [Liu Hui: As for doubling it: one lays out in advance the fixed lengths (ding mao) of the vermilion areas (zhu mi) on the two sides, in order to await the next division; therefore it is called the fixed divisor.] For the next division, shift back (zhe) the divisor and [move it] down. [Liu Hui: To remove the vermilion areas, one ought originally to set out in duplicate (fu zhi) what has been obtained, forming the square [side], double it to make the fixed divisor, shift it back, multiply by the estimate, and divide by it. In this way one would have to step it forward again and stop, and only then could they name each other; therefore one makes it approach the upper and shift down.] Set out the borrowed rod again, step it forward as at first, and multiply it once by the renewed estimate (fu yi); [Liu Hui: One wishes to remove the area of the yellow B (huang yi) at the corner of the vermilion areas; the idea is like that of what was obtained at first.] set out what is obtained in duplicate (fu), add it to the fixed divisor, and divide. Join (cong) the duplicate of what is obtained to the fixed divisor. [Liu Hui: Adding the side of the yellow B to the fixed divisor a second time is, then, to lay out the lengths of the two blue-green areas (qing mi).] Divide again, shifting back down as before.\end{uniten}

\begin{unitzh}\unitlabel{jiuzhang-015.13}若開之不盡者，爲不可開，當以面命之。{\zhsmall 術或有以借算加定法而命分者，雖麤相近，不可用也。凡開積爲方，方之自乘當還復其積分。令不加借算而命分，則常微少；其加借算而命分，則又微多。其數不可得而定。故惟以面命之，爲不失耳。譬猶以三除十，以其餘爲三分之一，而復其數可舉。不以面命之，加定法如前，求其微數。微數無名者以爲分子，其一退以十爲母，其再退以百爲母。退之彌下，其分彌細，則朱羃雖有所乘之數，不足言之也。}\end{unitzh}
\begin{uniten}If, when it is extracted, it is not exhausted, it is "not extractable" (bu ke kai), and one should name (ming) it by the side. [Liu Hui: Among the methods, some name the fraction (ming fen) by adding the borrowed rod to the fixed divisor; although this is roughly close, it cannot be used. In general, when one extracts an area into a square, the square [side] multiplied by itself ought to restore again its area and parts (ji fen). If one names the fraction without adding the borrowed rod, it is always slightly too small; if one names the fraction adding the borrowed rod, it is again slightly too large. Its number cannot be determined. Therefore only naming it by the side is without error. It is like dividing 10 by 3: one takes its remainder as 1/3, and its number can be restored. If one does not name it by the side, add to the fixed divisor as before and seek its fine numbers (wei shu). Fine numbers that have no name are taken as numerators (fenzi): at the first step back (tui), take 10 as the denominator (mu); at the second step back, take 100 as the denominator. The further down one steps back, the finer its parts; then, although the vermilion areas have numbers that are multiplied, they are not worth mentioning.]\end{uniten}

\begin{unitzh}\unitlabel{jiuzhang-015.14}若實有分者，通分內子爲定實，乃開之。訖，開其母，報除。{\zhsmall 臣淳風等謹按：分母可開者，竝通之積先合二母。旣開之後，一母尚存，故開分母，求一母爲法，以報除也。}若母不可開者，又以母再乘定實，乃開之。訖，令如母而一。{\zhsmall 臣淳風等謹按：分母不可開者，本一母也。又以母乘之，乃合二母。旣開之後，亦一母存焉。故令如母而一，得全面也。}\end{unitzh}
\begin{uniten}If the dividend has parts, convert into parts (tong fen) and add in the numerator (na zi) to make the fixed dividend (ding shi), and then extract it. When finished, extract its denominator and divide back (bao chu). [Li Chunfeng et al.: When the denominator can be extracted, the area, converted together (bing tong), first combines two denominators. After it has been extracted, one denominator still remains; therefore one extracts the denominator, seeking one denominator to make the divisor, so as to divide back.] If the denominator cannot be extracted, further multiply the fixed dividend by the denominator a second time, and then extract it. When finished, divide by the denominator. [Li Chunfeng et al.: When the denominator cannot be extracted, there is originally one denominator. Multiplying it further by the denominator then combines two denominators. After it has been extracted, one denominator likewise remains. Therefore dividing by the denominator gives the complete side.]\end{uniten}

\begin{unitzh}\unitlabel{jiuzhang-015.15}{\zhsmall 又按：此術「開方」者，求方羃之一面也。「借一算」者，假借一算，空有列位之名，而無除積之實。方隅得面，是故借算列之於下也。「步之超一等」者，方十自乘，其積有百；方百自乘，其積有萬。故超位至百而言十，至萬而言百也。「議所得，以一乘所借一算爲法，而以除」者，先得黃甲之面，以方爲積者兩相乘，故開方除之，還令兩面上下相命，是自乘而除之也。「除已，倍法爲定法」者，實積未盡，當復更除，故豫張兩面朱羃定袤，以待復除，故曰定除也。「其復除，折法而下」者，欲除朱羃，本當副置所得成方，倍之爲定法，以折議乘之而以除。如是當復步之而止，乃得相命，故使就上折之而下也。「復置借算，步之如初，以復議一乘之，所得副，以加定法，以除」者，欲除朱羃之角黃乙之羃。「以所得副從定法」者，再以黃乙之羃加定法，是則張兩青羃之袤。故如前開之，卽合所問。}\end{unitzh}
\begin{uniten}[Li Chunfeng et al.: Further note: In this method, "extracting the square root" is seeking one side of a square area. "Borrowing 1 rod" is to borrow 1 rod provisionally (jia jie); it has, emptily, the name of a position in the array (lie wei), but not the substance of removing the area. From the corner (yu) of the square one obtains the side; this is why the borrowed rod is set out below. "Step it forward, skipping one rank": a square of 10 multiplied by itself has an area of 100; a square of 100 multiplied by itself has an area of 10000. Therefore, skipping places, at 100 one speaks of 10, and at 10000 one speaks of 100. "Estimate what is obtained, and multiply the 1 borrowed rod by it once to make the divisor; then divide by it": first one obtains the side of the yellow A; to make an area from a square, the two [sides] are multiplied together; therefore, in extracting the square root, one again lets the two sides, upper and lower, name each other: this is multiplying by itself and removing it. "When the division is done, double the divisor to make the fixed divisor": the area of the dividend is not yet exhausted and must be divided yet again; therefore one lays out in advance the fixed lengths of the vermilion areas on the two sides, in order to await the next division; therefore it is called the "fixed division" (ding chu). "For the next division, shift back the divisor and [move it] down": to remove the vermilion areas, one ought originally to set out in duplicate what has been obtained, forming the square [side], double it to make the fixed divisor, shift it back, multiply it by the estimate, and divide by it. In this way one would have to step it forward again and stop, and only then could they name each other; therefore one makes it approach the upper, shifting it back and down. "Set out the borrowed rod again, step it forward as at first, multiply it once by the renewed estimate, set out what is obtained in duplicate, add it to the fixed divisor, and divide": one wishes to remove the area of the yellow B at the corner of the vermilion areas. "Join the duplicate of what is obtained to the fixed divisor": adding the area of the yellow B to the fixed divisor a second time is, then, to lay out the lengths of the two blue-green areas. Therefore, extracting it as before, one accords with what is asked.]\end{uniten}

\begin{unitzh}\unitlabel{jiuzhang-016.1}今有積一千五百一十八步四分步之三，問爲圓周幾何？荅曰：一百三十五步。{\zhsmall 於徽術，當周一百三十八步一十分步之一。臣淳風等謹按：此依密率，爲周一百三十八步五十分步之九。}\end{unitzh}
\begin{uniten}Suppose there is an area of 1518 3/4 bu (步). Question: as a circle, how much is the circumference? Answer: 135 bu. [Liu Hui: By Hui's method (Hui shu), the circumference should be 138 1/10 bu.] [Li Chunfeng et al.: Your servant Chunfeng and others respectfully note: this, according to the accurate rate (milü), gives a circumference of 138 9/50 bu.]\end{uniten}

\begin{unitzh}\unitlabel{jiuzhang-016.2}又有積三百步，問爲圓周幾何？荅曰：六十步。{\zhsmall 於徽術，當周六十一步五十分步之十九。臣淳風等謹按：依密率，爲周六十一步一百分步之四十一。}\end{unitzh}
\begin{uniten}Again, there is an area of 300 bu. Question: as a circle, how much is the circumference? Answer: 60 bu. [Liu Hui: By Hui's method, the circumference should be 61 19/50 bu.] [Li Chunfeng et al.: Your servant Chunfeng and others respectfully note: according to the accurate rate, the circumference is 61 41/100 bu.]\end{uniten}

\begin{unitzh}\unitlabel{jiuzhang-016.3}開圓\end{unitzh}
\begin{uniten}Extracting the circle (kai yuan)\end{uniten}

\begin{unitzh}\unitlabel{jiuzhang-016.4}術曰：置積步數，以十二乘之，以開方除之，卽得周。{\zhsmall 此術以周三徑一爲率，與舊圓田術相返覆也。於徽術，以三百一十四乘積，如二十五而一，所得，開方除之，卽周也。開方除之，卽徑。是爲據見羃以求周，猶失之於微少。其以二百乘積，一百五十七而一，開方除之，卽徑，猶失之於微多。}\end{unitzh}
\begin{uniten}Method: Set out the number of bu of the area, multiply it by 12, and extract the square root: this gives the circumference. [Liu Hui: This method takes circumference 3, diameter 1 as the rates (lü); it and the old method for the circular field are the reverse of each other. By Hui's method, multiply the area by 314 and divide by 25; extract the square root of what is obtained: that is the circumference. Extract the square root: that is the diameter. This is seeking the circumference on the basis of the given area (mi), and it still errs by being minutely too small. As for multiplying the area by 200, dividing by 157 and extracting the square root, which is the diameter, it still errs by being minutely too large.]\end{uniten}

\begin{unitzh}\unitlabel{jiuzhang-016.5}{\zhsmall 臣淳風等謹按：此注於徽術求周之法，其中不用「開方除之卽徑」六字，今本有者，衍賸也。依密率，八十八乘之，七而一。按：周三徑一之率，假令周六徑二，半周半徑相乘，得羃三；周六自乘，得三十六；俱以等數除，羃得一，周之數十二也。其積本周自乘，合以一乘之，十二而一，得積三也。術爲一乘不長，故以十二而一，得此積。今還原，置此積三，以十二乘之，復其本周自乘之數。凡物自乘，開方除之，復其本數。故開方除之，卽周。}\end{unitzh}
\begin{uniten}[Li Chunfeng et al.: Your servant Chunfeng and others respectfully note: in this commentary, as to the method of seeking the circumference by Hui's method, the six characters 開方除之卽徑 ("extract the square root: that is the diameter") are not used in it; where the present text has them, they are superfluous. According to the accurate rate, multiply it by 88 and divide by 7. Note: with the rates of circumference 3, diameter 1, suppose the circumference is 6 and the diameter 2. Half the circumference and half the diameter multiplied together give an area (mi) of 3; the circumference 6 multiplied by itself gives 36; dividing both by the equal number (deng shu), the area gives 1, and the number of the circumference is 12. As for its area: the circumference is originally multiplied by itself; this should be multiplied by 1 and divided by 12, giving the area 3. In the method, since multiplying by 1 does not increase it, one therefore divides by 12 to obtain this area. Now, returning to the origin, set out this area 3 and multiply it by 12; this restores the number of the original circumference multiplied by itself. Whenever a thing is multiplied by itself, extracting the square root restores its original number. Therefore extract the square root, and that is the circumference.]\end{uniten}

\begin{unitzh}\unitlabel{jiuzhang-016.6}今有積一百八十六萬八百六十七尺{\zhsmall 此尺謂立方之尺也。凡物有高深而言積者，曰立方。}問爲立方幾何？荅曰：一百二十三尺。\end{unitzh}
\begin{uniten}Suppose there is a volume of 1,860,867 chi (尺). [Liu Hui: This chi means a chi of a cube (lifang). Whenever a thing has height or depth and one speaks of its volume, it is called "cube" (lifang).] Question: as a cube, how much is it? Answer: 123 chi.\end{uniten}

\begin{unitzh}\unitlabel{jiuzhang-016.7}今有積一千九百五十三尺八分尺之一，問爲立方幾何？荅曰：一十二尺半。\end{unitzh}
\begin{uniten}Suppose there is a volume of 1953 1/8 chi. Question: as a cube, how much is it? Answer: 12 1/2 chi.\end{uniten}

\begin{unitzh}\unitlabel{jiuzhang-016.8}今有積六萬三千四百一尺五百一十二分尺之四百四十七，問爲立方幾何？荅曰：三十九尺八分尺之七。\end{unitzh}
\begin{uniten}Suppose there is a volume of 63,401 447/512 chi. Question: as a cube, how much is it? Answer: 39 7/8 chi.\end{uniten}

\begin{unitzh}\unitlabel{jiuzhang-016.9}又有積一百九十三萬七千五百四十一尺二十七分尺之一十七，問爲立方幾何？荅曰：一百二十四尺太半尺。\end{unitzh}
\begin{uniten}Again, there is a volume of 1,937,541 17/27 chi. Question: as a cube, how much is it? Answer: 124 2/3 chi.\end{uniten}

\begin{unitzh}\unitlabel{jiuzhang-016.10}開立方{\zhsmall 立方適等，求其一面也。}\end{unitzh}
\begin{uniten}Extracting the cube root (kai lifang) [Liu Hui: The cube being exactly equal, one seeks one of its sides (mian).]\end{uniten}

\begin{unitzh}\unitlabel{jiuzhang-016.11}術曰：置積爲實，借一算，步之，超二等。{\zhsmall 言千之面十，言百萬之面百。}議所得，以再乘所借一算爲法，而除之。{\zhsmall 再乘者，亦求爲方羃，以上議命而除之，則立方等也。}除已，三之爲定法。{\zhsmall 爲當復除，故豫張三面，以定方羃爲定法也。}復除，折而下。{\zhsmall 復除者，三面方羃以皆自乘之數，須得折議，定其厚薄爾。開平羃者，方百之面十；開立羃者，方千之面十。據定法已有成方之羃，故復除當以千爲百，折下一等也。}\end{unitzh}
\begin{uniten}Method: Set out the volume as the dividend (shi). Borrow one counting rod (suan) and step it forward, skipping two places. [Liu Hui: That is to say, the side of 1000 is 10; the side of 1,000,000 is 100.] Deliberate (yi) what is obtained; multiply the borrowed rod by it twice to make the divisor (fa), and divide. [Liu Hui: Multiplying twice is again seeking to make a square area (fang mi); naming (ming) it by the upper deliberated number and dividing, the cube is then equal.] When the division is done, triple it to make the fixed divisor (ding fa). [Liu Hui: Because one must divide again, one therefore lays out in advance the three sides, taking the fixed square areas as the fixed divisor.] Divide again, shifting back (zhe) and moving down. [Liu Hui: As for dividing again: since the square areas of the three sides are all numbers multiplied by themselves, one must obtain a shifted deliberation to fix their thickness. In extracting a plane area (ping mi), the side of a square of 100 is 10; in extracting a solid area (li mi), the side of a square of 1000 is 10. Since the fixed divisor already has the area of a completed square, in dividing again one should take 1000 as 100 and shift back down one place.]\end{uniten}

\begin{unitzh}\unitlabel{jiuzhang-016.12}以三乘所得數，置中行。{\zhsmall 設三亷之定長。}復借一算，置下行。{\zhsmall 欲以爲隅方，立方等未有定數，且置一算定其位。}步之，中超一，下超二位。{\zhsmall 上方法長自乘，而一折；中亷法但有長，故降一等；下隅法無面長，故又降一等也。}復置議，以一乘中，{\zhsmall 爲三亷備羃也。}再乘下，{\zhsmall 令隅自乘爲方羃也。}皆副以加定法。以定法除。{\zhsmall 三面、三亷、一隅皆已有羃，以上議命之，而除去三袤之厚也。}\end{unitzh}
\begin{uniten}Multiply the obtained number by 3 and set it in the middle row. [Liu Hui: This sets up the fixed length of the three edges (lian).] Again borrow one counting rod and set it in the lower row. [Liu Hui: One wishes to make it the corner square (yu fang); the cube being equal has as yet no fixed number, so for the time being one sets one rod to fix its position.] Step them forward, the middle skipping one, the lower skipping two places. [Liu Hui: The upper square divisor (fang fa), its length multiplied by itself, shifts back once; the middle edge divisor (lian fa) has only length, so it descends one place; the lower corner divisor (yu fa) has no side or length, so it descends one place further.] Again set out the deliberated number; multiply the middle by it once [Liu Hui: this is to complete the areas of the three edges], multiply the lower by it twice [Liu Hui: this is to make the corner, multiplied by itself, into a square area]; in each case copy (fu) it and add it to the fixed divisor. Divide by the fixed divisor. [Liu Hui: The three sides, the three edges and the one corner all already have areas; naming them by the upper deliberated number, one removes the thickness of the three lengths (mao).]\end{uniten}

\begin{unitzh}\unitlabel{jiuzhang-016.13}除已，倍下，并中，從定法。{\zhsmall 凡再以中、三以下加定法者，三亷各當以兩面之羃連於兩方之面，一隅連於三亷之端，以待復除也。言不盡意，解此要當以棊，乃得明耳。}復除，折下如前。開之不盡者，亦爲不可開。{\zhsmall 術亦有以定法命分者，不如故羃開方，以微數爲分也。}\end{unitzh}
\begin{uniten}When the division is done, double the lower, combine it with the middle, and join them to the fixed divisor. [Liu Hui: As for adding, in all, the middle twice and the lower three times to the fixed divisor: each of the three edges should be joined by the areas of its two sides to the sides of two squares, and the one corner is joined to the ends of the three edges, in readiness for dividing again. Words do not exhaust the meaning; to explain this one must use blocks (qi), and only then can it be made clear.] Divide again, shifting back down as before. If extracting it does not come to an end, this too is "cannot be extracted". [Liu Hui: Among the methods there are also those that name the fraction by the fixed divisor; this is not as good as the former extraction of the square root of an area, taking minute numbers (wei shu) as the fraction.]\end{uniten}

\begin{unitzh}\unitlabel{jiuzhang-016.14}若積有分者，通分內子爲定實，定實乃開之，訖，開其母以報除。{\zhsmall 臣淳風等按：分母可開者，竝通之。積先合三母，旣開之後，一母尙存，故開分母，求一母爲法，以報除也。}若母不可開者，又以母再乘定實，乃開之，訖，令如母而一。{\zhsmall 臣淳風等謹按：分母不可開者，本一母也。又以母再乘之，今合三母，旣開之後，一母猶存，故令如母而一，得全面也。}\end{unitzh}
\begin{uniten}If the volume has a fraction, convert it into parts (tong fen) and add in the numerator (na zi) to make the fixed dividend (ding shi); then extract the root of the fixed dividend; when this is done, extract the root of its denominator, to divide in return (bao chu). [Li Chunfeng et al.: Your servant Chunfeng and others note: when the denominator can be extracted, convert them all. The volume first combines three denominators; after it has been extracted, one denominator still remains; therefore one extracts the root of the denominator, seeking one denominator to make the divisor, so as to divide in return.] If the denominator cannot be extracted, further multiply the fixed dividend twice by the denominator, then extract its root; when this is done, divide by the denominator. [Li Chunfeng et al.: Your servant Chunfeng and others respectfully note: when the denominator cannot be extracted, there is originally one denominator. Further multiplying it twice by the denominator, it now combines three denominators; after it has been extracted, one denominator still remains; therefore divide by the denominator, and one obtains the complete side.]\end{uniten}

\begin{unitzh}\unitlabel{jiuzhang-016.15}{\zhsmall 按：開立方者，立方適等，求其一面之數也。借一算步之，超二等者，立方求積，方再自乘，就積開之，故超二位，言千之面十，言百萬之面百也。議所得，以再乘所借一算爲法，而以除者，求爲方羃，以議命之而除，則立方等也。除已，三之爲定法者，爲積未盡，當復更除，故豫張三面，已定方羃爲定法也。復除，折而下者，三面方羃皆已有自乘之數，須得折議，定其厚薄。據開平方，百之面十；其開立方，卽千之面十。而定法已有成方之羃，故復除之，當以千爲百，折卞一等也。}\end{unitzh}
\begin{uniten}[Li Chunfeng et al.: Note: as for extracting the cube root: the cube being exactly equal, one seeks the number of one of its sides. "Borrow one counting rod and step it forward, skipping two places": for the cube, in seeking the volume the side is multiplied by itself twice; extracting from the volume, one therefore skips two places; that is to say, the side of 1000 is 10, and the side of 1,000,000 is 100. "Deliberate what is obtained; multiply the borrowed rod by it twice to make the divisor, and divide by it": one seeks to make a square area; naming it by the deliberated number and dividing, the cube is then equal. "When the division is done, triple it to make the fixed divisor": because the volume is not yet exhausted and one must divide once more, one therefore lays out in advance the three sides, the already fixed square areas being made the fixed divisor. "Divide again, shifting back and moving down": the square areas of the three sides all already have numbers multiplied by themselves, and one must obtain a shifted deliberation to fix their thickness. According to extracting the square root (kai pingfang), the side of 100 is 10; in extracting the cube root, then, the side of 1000 is 10. And the fixed divisor already has the area of a completed square; therefore, in dividing again, one should take 1000 as 100 and shift back down one place.] [Translator's note: 卞 in 折卞一等 is evidently an error for 下 "down".]\end{uniten}

\begin{unitzh}\unitlabel{jiuzhang-016.16}{\zhsmall 以三乘所得數，置中行者，設三亷之定長也。復借一算，置下行者，欲以爲隅方，立方等未有數，且置一算定其位也。步之，中超一，下超二者，上方法長自乘，而一折；中亷法但有長，故降一等；下隅法無面長，故又降一等也。復置議，以一乘中者，爲三亷借羃也。再乘下者，當令隅自乘爲方羃也。皆副以加定法，以定法除者，三面、三亷、一隅皆已有羃，以上議命之，而除去三袤之厚也。除已，倍下，併中，從定法者，三亷各當以兩面之羃連於兩方之面，一隅連於三亷之端，以待復除也。}\end{unitzh}
\begin{uniten}[Li Chunfeng et al.: "Multiply the obtained number by 3 and set it in the middle row": this sets up the fixed length of the three edges. "Again borrow one counting rod and set it in the lower row": one wishes to make it the corner square; the cube being equal has as yet no number, so for the time being one sets one rod to fix its position. "Step them forward, the middle skipping one, the lower skipping two": the upper square divisor, its length multiplied by itself, shifts back once; the middle edge divisor has only length, so it descends one place; the lower corner divisor has no side or length, so it descends one place further. "Again set out the deliberated number; multiply the middle by it once": this is to borrow areas for the three edges. "Multiply the lower by it twice": one should make the corner, multiplied by itself, into a square area. "In each case copy it and add it to the fixed divisor; divide by the fixed divisor": the three sides, the three edges and the one corner all already have areas; naming them by the upper deliberated number, one removes the thickness of the three lengths. "When the division is done, double the lower, combine it with the middle, and join them to the fixed divisor": each of the three edges should be joined by the areas of its two sides to the sides of two squares, and the one corner is joined to the ends of the three edges, in readiness for dividing again.]\end{uniten}

\begin{unitzh}\unitlabel{jiuzhang-016.17}{\zhsmall 其開之不盡者，折下如前開方，卽合所問。有分者，通分內子，開之，訖，開其母以報除。可開者，並通之。積先合三母，旣開之後，一母尙存，故開分母者，求一母爲法，以報除。若母不可開者，又以母再乘定實，乃開之，訖，令如母而一。分母不可開者，本一母，又以母再乘，今合三母，旣開之後，亦一母尙存，故令如母而一，得全面也。}\end{unitzh}
\begin{uniten}[Li Chunfeng et al.: Where extracting it does not come to an end, shift back down and extract the root as before; then it accords with what is asked. Where there is a fraction, convert it into parts and add in the numerator, extract its root, and when this is done, extract the root of its denominator to divide in return. When it can be extracted, convert them all. The volume first combines three denominators; after it has been extracted, one denominator still remains; therefore, extracting the root of the denominator, one seeks one denominator to make the divisor, to divide in return. If the denominator cannot be extracted, further multiply the fixed dividend twice by the denominator, then extract its root; when this is done, divide by the denominator. When the denominator cannot be extracted, there is originally one denominator; further multiplying twice by the denominator, it now combines three denominators; after it has been extracted, likewise one denominator still remains; therefore divide by the denominator, and one obtains the complete side.]\end{uniten}

\begin{unitzh}\unitlabel{jiuzhang-017.1}今有積四千五百尺{\zhsmall 亦謂立方之尺也}，問爲立圓徑幾何？荅曰：二十尺。{\zhsmall 依密率，立圓徑二十八尺，計積四千一百九十尺二十一分尺之一十。}\end{unitzh}
\begin{uniten}Suppose there is a volume of 4500 chi (尺) [Liu Hui: this also means cubic (lifang) chi]. Question: if it is made into a sphere (liyuan), what is the diameter? Answer: 20 chi. [Liu Hui: According to the accurate rate (milü), a sphere of diameter 28 chi is reckoned to have a volume of 4190 10/21 chi.]\end{uniten}

\begin{unitzh}\unitlabel{jiuzhang-017.2}又有積一萬六千四百四十八億六千六百四十三萬七千五百尺，問爲立圓徑幾何？荅曰：一萬四千三百尺。{\zhsmall 依密率，爲徑一萬四千六百四十三尺四分尺之三。}\end{unitzh}
\begin{uniten}Again, there is a volume of 1,644,866,437,500 chi. Question: if it is made into a sphere, what is the diameter? Answer: 14,300 chi. [Liu Hui: According to the accurate rate, the diameter is 14,643 3/4 chi.]\end{uniten}

\begin{unitzh}\unitlabel{jiuzhang-017.3}開立圓\end{unitzh}
\begin{uniten}Extracting the sphere (kai liyuan)\end{uniten}

\begin{unitzh}\unitlabel{jiuzhang-017.4}術曰：置積尺數，以十六乘之，九而一，所得開立方除之，卽丸徑。\end{unitzh}
\begin{uniten}Method: Set out the number of chi of the volume, multiply it by 16 and divide by 9; extract the cube root of what is obtained; this is the diameter of the ball (wan).\end{uniten}

\begin{unitzh}\unitlabel{jiuzhang-017.5}{\zhsmall 立圓卽丸也。爲術者蓋依周三徑一之率，令圓羃居方羃四分之三，圓囷居立方亦四分之三。更令圓囷爲方率十二，爲丸率九，丸居圓囷又四分之三也。置四分自乘得十六，三分自乘得九，故丸居立方十六分之九也。故以十六乘積，九而一，得立方之積。丸徑與立方等，故開立方而除，得徑也。}\end{unitzh}
\begin{uniten}[Liu Hui: The sphere is the ball. Those who made the method presumably relied on the rate (lü) of circumference 3 to diameter 1, making the area (mi) of the circle 3/4 of the area of the square, and the round granary (yuanqun) likewise 3/4 of the cube (lifang). Further, they made the round granary the rate of the square (fanglü), 12, and the rate of the ball 9: the ball is again 3/4 of the round granary. Set out the 4 parts and multiply by itself, obtaining 16; the 3 parts multiplied by itself give 9; therefore the ball is 9/16 of the cube. Therefore multiply the volume by 16 and divide by 9 to obtain the volume of the cube. The diameter of the ball is equal to the cube; therefore extracting the cube root gives the diameter.]\end{uniten}

\begin{unitzh}\unitlabel{jiuzhang-017.6}{\zhsmall 然此意非也。何以驗之？取立方棊八枚，皆令立方一寸，積之爲立方二寸。規之爲圓囷，徑二寸，高二寸。又復橫規之，則其形有似牟合方蓋矣。八棊皆似陽馬，圓然也。按合蓋者，方率也，丸居其中，卽圓率也。推此言之，謂夫圓囷爲方率，豈不闕哉？以周三徑一爲圓率，則圓羃傷少；令圓囷爲方率，則丸積傷多。互相通補，是以丸與十六之率，偶與實相近，而丸猶傷多耳。觀立方之內，合蓋之外，雖衰殺有漸，而多少不掩。判合總結，方圓相纒，濃纖詭互，不可等正。欲陋形措意，懼失正理。敢不闕疑，以俟能言者。}\end{unitzh}
\begin{uniten}[Liu Hui: But this idea is wrong. How can it be verified? Take 8 cubic blocks (qi), each made a cube of 1 cun (寸), and pile them up into a cube of 2 cun. Cut it with the compasses (gui) into a round granary of diameter 2 cun and height 2 cun. Then cut it again crosswise with the compasses; its shape then resembles a "matched square umbrella" (mouhe fanggai). The 8 blocks all resemble yangma (陽馬), but rounded. Now, the combined umbrella (hegai) is the rate of the square; the ball lies inside it and is the rate of the circle (yuanlü). Reasoning from this, to say that the round granary is the rate of the square — is that not defective? If circumference 3, diameter 1 is taken as the rate of the circle, the area of the circle errs on the small side; if the round granary is taken as the rate of the square, the volume of the ball errs on the large side. These make up for each other; that is why the rate of the ball to 16 happens to come close to the actual value, while the ball is still too large. Looking at what lies inside the cube and outside the combined umbrella: although it diminishes by degrees, its amount cannot be grasped. Split and joined, wholly bound together, square and round entwined, thick and thin interchanging irregularly, it cannot be made equal and regular. I wish to set my mind on this crude form, but fear losing the correct principle. I dare not fail to leave the doubt open, to await one who is able to explain it.]\end{uniten}

\begin{unitzh}\unitlabel{jiuzhang-017.7}{\zhsmall 黃金方寸重十六兩，金丸徑寸重九兩，率生於此，未曾驗也。《周官·考工記》：「㮚氏爲量，改煎金錫則不耗，不耗然後權之，權之然後凖之，凖之然後量之。」言錬金使極精，而後分之，則可以爲率也。}\end{unitzh}
\begin{uniten}[Liu Hui: Yellow gold 1 cun square weighs 16 liang (兩); a gold ball 1 cun in diameter weighs 9 liang. The rate arose from this, but it has never been verified. In the Kaogong ji of the Zhou guan: "The Li (㮚) artisans made the measures. When the metal and tin were smelted over and over, there was no further loss; when there was no further loss, they weighed it; after weighing it, they leveled it; after leveling it, they measured it." This says that the metal is refined to the utmost purity and only then divided; then it can be used to establish a rate.]\end{uniten}

\begin{unitzh}\unitlabel{jiuzhang-017.8}{\zhsmall 令丸徑自乘，三而一，開方除之，卽丸中之立方也。假令丸中立方五尺，五尺爲句，句自乘羃二十五尺，倍之得五十尺，以爲股羃，謂平面方五尺之弦也。以此弦羃爲股，亦以五尺爲句，并句股羃得七十五尺，是爲大弦羃。開方除之，則大弦可知也。大弦則中立方之長邪，邪卽丸徑也。故中立方自乘之羃，於丸徑自乘之羃三分之一也。令大弦還乘其羃，卽丸外立方之積也。大弦羃開之不盡，令其羃七十五再自乘之爲面，命得外立方積四十二萬一千八百七十五尺之面。又令中立方五尺自乘，又以方乘之，得積一百二十五尺。一百二十五尺自乘爲面，命得積一萬五千六百二十五尺之面。皆以六百二十五約之，外立方積六百七十五尺之面，中立方積二十五尺之面也。}\end{unitzh}
\begin{uniten}[Liu Hui: Multiply the diameter of the ball by itself, divide by 3, and extract the square root: this is the cube inside the ball. Suppose the cube inside the ball is 5 chi. Take 5 chi as the base (gou); the base multiplied by itself gives an area of 25 chi; double it to obtain 50 chi, and take this as the area of the altitude (gu): this means the hypotenuse (xian) of a plane square of 5 chi. Take this area of the hypotenuse as the altitude, and again take 5 chi as the base; add the areas of the base and the altitude to obtain 75 chi: this is the area of the great hypotenuse. Extract its square root, and the great hypotenuse can be known. The great hypotenuse is the long diagonal (changxie) of the inner cube, and the diagonal is the diameter of the ball. Therefore the area of the inner cube multiplied by itself is 1/3 of the area of the diameter of the ball multiplied by itself. Let the great hypotenuse in turn multiply its area: this is the volume of the cube outside the ball. The square root of the area of the great hypotenuse cannot be extracted exactly; so multiply its area, 75, by itself twice and take it as a side (mian) [Translator's note: "the side of N" stands for the square root of N]; naming the result, the volume of the outer cube is the side of 421,875 chi. Again, multiply the 5 chi of the inner cube by itself and multiply again by the side (fang), obtaining a volume of 125 chi; multiply 125 chi by itself and take it as a side; naming the result, the volume is the side of 15,625 chi. Reduce both by 625: the volume of the outer cube is the side of 675 chi, and the volume of the inner cube is the side of 25 chi.]\end{uniten}

\begin{unitzh}\unitlabel{jiuzhang-017.9}{\zhsmall 張衡算又謂立方爲質，立圓爲渾。衡言質之與中外之渾：六百七十五尺之面，開方除之，不足一，謂外質積二十六也；內渾二十五之面，謂積五尺也。今徽令質言中渾，渾又言質，則二質相與之率，猶衡二渾相與之率也。衡蓋亦先二質之率，推以言渾之率也。衡又言質六十四之面，渾二十五之面。質復言渾，謂居質八分之五也。又云方八之面，圓六之面，圓渾相推，知其復以圓囷爲方率，渾爲圓率也，失之遠矣。衡說之自然，欲恊其陰陽奇耦之說，而不顧疎密矣。雖有文辭，斯亂道破義病也。置外質積二十六，以九乘之，十六而一，得積一十四尺八分之五，卽質中之渾也。以分母乘全內子，得一百一十七。又置內質積五，以分母乘之，得四十。是爲質居渾一百一十七分之四十，而渾率猶爲傷多也。}\end{unitzh}
\begin{uniten}[Liu Hui: Zhang Heng's computation further calls the cube zhi ("solid") and the sphere hun ("globe"). Heng speaks of the zhi in relation to the inner and outer hun: the side of 675 chi, when its square root is extracted, has a deficit of 1, and is said to be the volume of the outer zhi, 26; the inner hun, the side of 25, is said to be a volume of 5 chi. Now I, Hui, let zhi denote the inner hun, and hun in turn denote zhi; then the rate between the two zhi is like Heng's rate between the two hun. Heng presumably also started from the rate of the two zhi, and inferred from it to state the rate of the hun. Heng further says: zhi, the side of 64; hun, the side of 25. Zhi in turn denoting hun, this means that it occupies 5/8 of the zhi. He also says: square, the side of 8; circle, the side of 6. Inferring circle and hun from each other, we know that he again took the round granary as the rate of the square and the hun as the rate of the circle; this is far off. Heng spoke of it as natural, wishing to make it accord with his doctrine of yin and yang, odd and even, and paid no regard to coarseness or precision. Although it has fine wording, this is a fault that confuses the Way and breaks the meaning. Set out the volume of the outer zhi, 26; multiply it by 9 and divide by 16, obtaining a volume of 14 5/8 chi: this is the hun inside the zhi. Multiply the integer by the denominator and add in the numerator, obtaining 117. Again set out the volume of the inner zhi, 5, and multiply it by the denominator, obtaining 40. This makes the zhi 40/117 of the hun, and the rate of the hun is still too large.]\end{uniten}

\begin{unitzh}\unitlabel{jiuzhang-017.10}{\zhsmall 假令方二尺，方四面并得八尺也，謂之方周。其中令圓徑與方等，亦二尺也。丸半徑以乘圓周之半，卽圓羃也；半方以乘方周之半，卽方羃也。然則方周知方羃之率也，圓周知圓羃之率也。按如衡術，方周率八之面，圓周率五之面也。令方周六十四尺之面，卽圓周四十尺之面也。又令徑一尺，方周四尺，自乘得十六尺之面，是爲圓周率十二之面，而徑率一之面也。衡亦以周三徑一之率爲非是，故更著此法，然增周太多，過其實矣。}\end{unitzh}
\begin{uniten}[Liu Hui: Suppose a square of 2 chi; the four sides of the square added together give 8 chi; this is called the perimeter of the square (fangzhou). Inside it, let the diameter of a circle be equal to the square, also 2 chi. Half the diameter of the ball multiplied by half the circumference of the circle is the area of the circle; half the square multiplied by half the perimeter of the square is the area of the square. Thus the perimeter of the square indicates the rate of the area of the square, and the circumference of the circle indicates the rate of the area of the circle. Now, according to Heng's method, the rate of the perimeter of the square is the side of 8, and the rate of the circumference of the circle the side of 5. If the perimeter of the square is the side of 64 chi, then the circumference of the circle is the side of 40 chi. Again, let the diameter be 1 chi; the perimeter of the square is 4 chi; multiplied by itself it gives the side of 16 chi; this makes the rate of the circumference of the circle the side of 12, and the rate of the diameter the side of 1. Heng too regarded the rate of circumference 3 to diameter 1 as incorrect, and therefore composed this method anew; but he increased the circumference too much, beyond its actual value.]\end{uniten}

\begin{unitzh}\unitlabel{jiuzhang-017.11}{\zhsmall 臣淳風等謹按：祖暅之謂劉徽、張衡二人皆以圓囷爲方率，丸爲圓率，乃設新法。祖暅之開立圓術日：以二十一乘積，十一而一，開立方除之，卽立圓徑。其意何也？取立方棊一枚，令立樞於左後之下隅，從規去其右上之亷；又合而横規之，去其前上之亷、右前之亷。於是立方之棊分而爲四：規內棊一，謂之內棊；規外棊三，謂之外棊規。更合四棊，復橫斷之。以句股言之，令餘高爲句，內棊斷上方爲股，本方之數，其弦也。句股之法，以句羃減弦羃，則餘爲股羃。若令餘高自乘，減本方之羃，餘卽內減棊斷上方之羃也。本方之羃，卽外四棊之斷上羃。然則餘高自乘，卽外三棊之斷上羃矣。不問高卑，勢皆然也。}\end{unitzh}
\begin{uniten}[Li Chunfeng et al.: Your servants Chunfeng and others respectfully note: Zu Gengzhi said that Liu Hui and Zhang Heng both took the round granary as the rate of the square and the ball as the rate of the circle, and so he set up a new method. Zu Gengzhi's method for extracting the sphere says: Multiply the volume by 21, divide by 11, and extract the cube root; this is the diameter of the sphere. What is its idea? Take one cubic block; set the vertical pivot (shu) at the lower corner at the left rear, cut vertically with the compasses and remove its upper right edge (lian); again put it together and cut it crosswise with the compasses, removing its upper front edge and its right front edge. Thereupon the cubic block is divided into four: inside the compass cuts there is one block, called the inner block; outside the compass cuts there are three blocks, called the outer blocks, compass [Translator's note: this 規 "compass" appears superfluous]. Put the four blocks together again, and again cut them crosswise. Speaking in terms of base and altitude: let the remaining height be the base, and the square at the cut top of the inner block the altitude; the number of the original square is the hypotenuse. By the method of base and altitude, subtract the area of the base from the area of the hypotenuse, and the remainder is the area of the altitude. If the remaining height is multiplied by itself and subtracted from the area of the original square, the remainder is the area of the square at the cut top of the inner reduced block [Translator's note: 減 "reduce, subtract" here appears superfluous]. The area of the original square is the area at the cut tops of the outer four blocks. Then the remaining height multiplied by itself is the area at the cut tops of the three outer blocks. Whether high or low, the situation is always so.]\end{uniten}

\begin{unitzh}\unitlabel{jiuzhang-017.12}{\zhsmall 然固有所歸同而塗殊者爾，而乃控遠以演類，借况以析㣲。按陽馬方高數參等者，列而立之，横截去上，則高自乘與斷上羃數亦等焉。夫疊棊成立積，緣羃勢旣同，則積不容異。由此觀之，規之外三棊旁蹙爲一，卽一陽馬也。三分立方，則陽馬居一，內棊居二，可知矣。合八小方成一大方，合八內棊成一合蓋。內棊居小方三分之二，則合蓋居立方亦三分之二，較然驗矣。置三分之二，以圓羃率三乘之，如方羃率四而一，約而定之，以爲丸率。故曰丸居立方三分之一也。等數旣密，心亦昭晰。張衡放舊，貽哂於後；劉徽循故，未暇校新。夫豈難哉？抑未之思也。依密率，立此圓積，本以圓徑再自乘，十一乘之，二十一而一，約此積。今欲求其本積，故以二十一乘之，十一而一。凡物再自乘，開立方除之，復其本數。故立方除之，卽丸徑也。}\end{unitzh}
\begin{uniten}[Li Chunfeng et al.: Yet there are indeed things that come to the same end by different paths; and so one draws on what is distant to extend by analogy, and borrows comparisons to analyze what is subtle. Now, take a yangma whose square and height are all three equal in number; set it out and stand it up, and cut it across horizontally, removing the top; then the height multiplied by itself and the area of the cut top are also equal in number. Now, when blocks are stacked to form solid volumes, if the areas along the heights (shi) are the same, the volumes cannot differ. Seen in this light, the three blocks outside the compass cuts, pressed together from the sides into one, are one yangma. Dividing the cube into three, the yangma occupies one part and the inner block two: this can be known. Eight small cubes put together make one large cube; eight inner blocks put together make one combined umbrella. The inner block occupies 2/3 of the small cube; then the combined umbrella likewise occupies 2/3 of the cube: this is clearly verified. Set out 2/3, multiply it by the rate of the area of the circle, 3, and divide by the rate of the area of the square, 4; reduce and fix it, and take it as the rate of the ball. Therefore it is said: the ball occupies 1/3 of the cube. Since the equal numbers are precise, the mind too is enlightened. Zhang Heng imitated the old and left something for later ages to laugh at; Liu Hui followed the old and had no leisure to examine the new. Was it really so difficult? Rather, they did not think of it. According to the accurate rate, this volume of the sphere is set up by originally multiplying the diameter of the circle by itself twice, multiplying by 11 and dividing by 21, reducing it to this volume. Now one wishes to find its original volume; therefore multiply by 21 and divide by 11. In general, when a thing has been multiplied by itself twice, extracting the cube root restores its original number. Therefore extracting the cube root gives the diameter of the ball.]\end{uniten}

\begin{unitzh}\unitlabel{jiuzhang-017.13}九章算術卷四\end{unitzh}
\begin{uniten}Nine Chapters on the Mathematical Art, juan 4\end{uniten}

\begin{unitzh}\unitlabel{jiuzhang-017.14}九章算術卷四訂訛補圖{\zhsmall 算經十書之二}\end{unitzh}
\begin{uniten}Nine Chapters on the Mathematical Art, juan 4: Corrections of Errors and Supplementary Figures (ding'e butu) [Note: Ten Mathematical Classics (Suanjing shishu), the second]\end{uniten}

\begin{unitzh}\unitlabel{jiuzhang-017.15}□□□□□□□□□□□□□□□□□□□□□□□□□□□□□□□□□□□□\end{unitzh}
\begin{uniten}[illegible: 18 characters] [illegible: 18 characters]\end{uniten}

\begin{unitzh}\unitlabel{jiuzhang-017.16}祖暅之開立圓術曰：以二十一乘積，十一而一，開立方除之，卽立圓徑。其意何也？取立方棊一枚，令立樞於左後之下隅，從規去其右上之廉；又合而橫規之，去其前上之廉、右前之廉。於是立方之棊分而爲四：規內棊一，謂之內棊；規外棊三，謂之外棊。更合四棊，復横斷之。以句股言之，令餘高爲句，內棊斷上方爲股，本方之數，其弦也。句股之法，以句幂減弦幂，則餘爲股幂。若令餘高自乘，減本方之幂，餘卽內減棊斷上方之幂也。本方之幂，卽外四棊之斷上幂。然則餘高自乘，卽外三棊之斷上幂矣。不問高卑，勢皆然也。{\zhsmall 以上借立方棊以論立圓，而所言僅及句股弦與平幂，不足見圓術，當有脫誤。}\end{unitzh}
\begin{uniten}Zu Gengzhi's method for extracting the sphere says: Multiply the volume by 21, divide by 11, and extract the cube root; this is the diameter of the sphere. What is its idea? Take one cubic block; set the vertical pivot at the lower corner at the left rear, cut vertically with the compasses and remove its upper right edge; again put it together and cut it crosswise with the compasses, removing its upper front edge and its right front edge. Thereupon the cubic block is divided into four: inside the compass cuts there is one block, called the inner block; outside the compass cuts there are three blocks, called the outer blocks. Put the four blocks together again, and again cut them crosswise. Speaking in terms of base and altitude: let the remaining height be the base, and the square at the cut top of the inner block the altitude; the number of the original square is the hypotenuse. By the method of base and altitude, subtract the area of the base from the area of the hypotenuse, and the remainder is the area of the altitude. If the remaining height is multiplied by itself and subtracted from the area of the original square, the remainder is the area of the square at the cut top of the inner reduced block. The area of the original square is the area at the cut tops of the outer four blocks. Then the remaining height multiplied by itself is the area at the cut tops of the three outer blocks. Whether high or low, the situation is always so. [Note: The above borrows the cubic block to discuss the sphere, but what it says touches only on base, altitude and hypotenuse and on plane areas; it is insufficient to show the method for the round, and there must be omissions and errors.]\end{uniten}

\begin{unitzh}\unitlabel{jiuzhang-017.17}故曰丸居立方三分之一也。{\zhsmall 此句舛誤。據上言「置三分之二，以三乘之，如四而一」，乃丸居立方二分之一，非三分之一。况以上眀祖氏圓術，其率乃丸居立方二十一分之十一，下云「圓徑再自乘，十一乘之，如二十一而一」是也。若二分之一，於祖氏術不恊矣。又祖氏方幂率十四，圓幂率十一，亦不得用方幂四、圓幂三之疎率以解祖氏說。自「祖暅之開立方圓術曰」至此，似因傳寫旣訛，後人妄加竄改，遂不可通。今考立方與圓囷，猶之平方與平圓也，其率亦方積十四，圓囷積十一，而丸居圓囷三分之二，與十四分之十一通之，分母乘分母得四十二，分子乘分子得二十二，是爲九居立方四十二分之二十二，卽二十一分之十一也。祖氏求圓囷、立圓、平圖三法，本條貫爲一。}\end{unitzh}
\begin{uniten}Therefore it is said: the ball occupies 1/3 of the cube. [Note: This sentence is erroneous. According to what is said above, "set out 2/3, multiply by 3, divide by 4", then the ball occupies 1/2 of the cube, not 1/3. Moreover, the above makes clear Mr. Zu's method for the sphere: its rate is that the ball occupies 11/21 of the cube, and what is said below, "multiply the diameter of the circle by itself twice, multiply by 11, divide by 21", is this. If it were 1/2, it would not accord with Mr. Zu's method. Further, Mr. Zu's rate of the area of the square is 14 and his rate of the area of the circle 11; one may not, either, use the coarse rates of square area 4 and circle area 3 to explain Mr. Zu's account. From "Zu Gengzhi's method for extracting the solid square and round says" to this point, it seems that, after errors had arisen in copying, later people rashly altered it, so that it became unintelligible. Now, on examination, the cube and the round granary are like the plane square and the plane circle; their rate is likewise: the volume of the square 14, the volume of the round granary 11; and the ball occupies 2/3 of the round granary. Combine this with 11/14: denominator times denominator gives 42, numerator times numerator gives 22; this makes the nine [Translator's note: 九 "nine" evidently stands for 丸 "ball"] occupy 22/42 of the cube, that is, 11/21. Mr. Zu's three methods for finding the round granary, the sphere and the plane figure (pingtu) are fundamentally linked into one.]\end{uniten}

\section*{卷第五}\addcontentsline{toc}{section}{卷第五}

\begin{unitzh}\unitlabel{jiuzhang-018.1}《九章算術》卷第五{\zhsmall 《算經十書》之二} 魏劉徽注 唐朝議大夫、行太史令、上輕車都尉臣李淳風等奉勅注釋\end{unitzh}
\begin{uniten}The Nine Chapters on the Mathematical Art, juan 5 [Note: The second of the Ten Mathematical Classics.] Commentary by Liu Hui of the Wei. Of the Tang, Grand Master for Court Discussion, acting Director of the Astronomical Bureau, Senior Commandant of Light Chariots, Your servant Li Chunfeng et al., annotated by imperial command.\end{uniten}

\begin{unitzh}\unitlabel{jiuzhang-018.2}商功{\zhsmall 以御功程積實}\end{unitzh}
\begin{uniten}Shang gong ("Assessing Works") [Liu Hui: For handling labour standards (gongcheng) and volumes (jishi).]\end{uniten}

\begin{unitzh}\unitlabel{jiuzhang-018.3}今有穿地積一萬尺，問：爲堅、壤各幾何？荅曰：爲堅七千五百尺，爲壤一萬二千五百尺。\end{unitzh}
\begin{uniten}Suppose there is excavated ground (chuandi) with a volume of 10000 chi (尺). Question: as hard earth (jian) and as loose earth (rang), how much is it each? Answer: As hard earth, 7500 chi; as loose earth, 12500 chi.\end{uniten}

\begin{unitzh}\unitlabel{jiuzhang-018.4}術曰：穿地四，爲壤五{\zhsmall 壤謂息土}，爲堅三{\zhsmall 堅爲築土}，爲墟四{\zhsmall 墟謂穿坑。此皆其常率}。以穿地求壤，五之；求堅，三之；皆四而一{\zhsmall 今有術也}。以壤求穿，四之；求堅，三之；皆五而一。以堅求穿，四之；求壤，五之；皆三而一。\end{unitzh}
\begin{uniten}Method: Excavated ground 4 makes loose earth 5 [Liu Hui: "Loose earth" means swollen earth.], makes hard earth 3 [Liu Hui: "Hard earth" is rammed earth.], makes pit (xu) 4 [Liu Hui: "Pit" means the dug hole. These are all their constant rates (lü).]. From excavated ground, to find loose earth, multiply by 5; to find hard earth, multiply by 3; in both cases divide by 4 [Liu Hui: This is the "Suppose there is" method.]. From loose earth, to find excavated ground, multiply by 4; to find hard earth, multiply by 3; in both cases divide by 5. From hard earth, to find excavated ground, multiply by 4; to find loose earth, multiply by 5; in both cases divide by 3.\end{uniten}

\begin{unitzh}\unitlabel{jiuzhang-018.5}{\zhsmall 臣淳風等謹按：此術竝今有之義也。重張穿地積一萬尺爲所有數，堅率三、壤率五各爲所求率，墟率四爲所有率，而今有之，卽得。}\end{unitzh}
\begin{uniten}[Li Chunfeng et al.: This method throughout has the sense of "Suppose there is". Lay out the excavated-ground volume of 10000 chi twice as the number one has (suoyou shu); the hard-earth rate 3 and the loose-earth rate 5 are each the rate sought (suoqiu lü); the pit rate 4 is the rate one has (suoyou lü); apply "Suppose there is" to them, and one obtains the results.]\end{uniten}

\begin{unitzh}\unitlabel{jiuzhang-018.6}城垣隄溝壍渠皆同術。術曰：并上下廣而半之{\zhsmall 損廣補狹}，以高若深乘之，又以袤乘之，卽積尺。\end{unitzh}
\begin{uniten}City walls (cheng), walls (yuan), dikes (di), ditches (gou), moats (qian) and canals (qu) all have the same method. Method: Add the upper and lower widths (guang) and halve [Liu Hui: Reduce the wide to make up the narrow.]; multiply by the height or depth, and further multiply by the length (mao): this is the volume in chi.\end{uniten}

\begin{unitzh}\unitlabel{jiuzhang-018.7}{\zhsmall 按此術「并上下廣而半之」者，以盈補虚，得中平之廣。「以高若深乘之」，堅率三、壤率五各爲所求率，墟率四爲所有率，而今有之，得一頭之立羃。「又以袤乘之」者，得立實之積，故爲積尺。}\end{unitzh}
\begin{uniten}[Liu Hui: In this method, "add the upper and lower widths and halve" is to use the surplus to fill the deficiency, obtaining the evened width of the middle. "Multiply by the height or depth": the hard-earth rate 3 and the loose-earth rate 5 are each the rate sought, the pit rate 4 is the rate one has, and applying "Suppose there is" to them one obtains the standing area (limi) of one end. "Further multiply by the length": this obtains the volume of the standing solid (lishi); hence it is the volume in chi.] [Translator's note: the words from "the hard-earth rate 3" to "applying Suppose there is to them" repeat the preceding note of Li Chunfeng et al. and are out of place here.]\end{uniten}

\begin{unitzh}\unitlabel{jiuzhang-018.8}今有城，下廣四丈，上廣二丈，高五丈，袤一百二十六丈五尺。問：積幾何？荅曰：一百八十九萬七千五百尺。\end{unitzh}
\begin{uniten}Suppose there is a city wall with lower width 4 zhang (丈), upper width 2 zhang, height 5 zhang and length 126 zhang 5 chi. What is the volume? Answer: 1897500 chi.\end{uniten}

\begin{unitzh}\unitlabel{jiuzhang-018.9}今有垣，下廣三尺，上廣二尺，高一丈二尺，袤二十二丈五尺八寸。問：積幾何？荅曰：六千七百七十四尺。\end{unitzh}
\begin{uniten}Suppose there is a wall with lower width 3 chi, upper width 2 chi, height 1 zhang 2 chi and length 22 zhang 5 chi 8 cun (寸). What is the volume? Answer: 6774 chi.\end{uniten}

\begin{unitzh}\unitlabel{jiuzhang-018.10}今有隄，下廣二丈，上廣八尺，高四尺，袤一十二丈七尺。問：積幾何？荅曰：七千一百一十二尺。\end{unitzh}
\begin{uniten}Suppose there is a dike with lower width 2 zhang, upper width 8 chi, height 4 chi and length 12 zhang 7 chi. What is the volume? Answer: 7112 chi.\end{uniten}

\begin{unitzh}\unitlabel{jiuzhang-018.11}冬程人功四百四十四尺。問：用徒幾何？荅曰：一十六人一百一十一分人之二。\end{unitzh}
\begin{uniten}The winter standard (cheng) for one person's work is 444 chi. Question: how many labourers (tu) are used? Answer: 16 2/111 persons.\end{uniten}

\begin{unitzh}\unitlabel{jiuzhang-018.12}術曰：以積尺爲實，程功尺數爲法，實如法而一，卽用徒人數。\end{unitzh}
\begin{uniten}Method: Take the volume in chi as the dividend (shi) and the number of chi of the standard work as the divisor (fa); divide the dividend by the divisor: this is the number of labourers used.\end{uniten}

\begin{unitzh}\unitlabel{jiuzhang-018.13}今有溝，上廣一丈五尺，下廣一丈，深五尺，袤七丈。問：積幾何？荅曰：四千三百七十五尺。\end{unitzh}
\begin{uniten}Suppose there is a ditch with upper width 1 zhang 5 chi, lower width 1 zhang, depth 5 chi and length 7 zhang. What is the volume? Answer: 4375 chi.\end{uniten}

\begin{unitzh}\unitlabel{jiuzhang-018.14}春程人功七百六十六尺，并出土功五分之四，定功六百一十二尺五分尺之四。問：用徒幾何？荅曰：七人三千六十四分人之四百二十七。\end{unitzh}
\begin{uniten}The spring standard for one person's work is 766 chi, together with the work of carrying out the earth (chutu gong), 4/5; the fixed work (dinggong) is 612 4/5 chi. Question: how many labourers are used? Answer: 7 427/3064 persons. [Translator's note: "4/5" is as the text reads; at the corresponding point the summer problem below has 1/5, and the method that follows removes 1/5.]\end{uniten}

\begin{unitzh}\unitlabel{jiuzhang-018.15}術曰：置本人功，去其五分之一，餘爲法{\zhsmall 去其五分之一者，謂以四乘、五除也}；以溝積尺爲實；實如法而一，得用徒人數。\end{unitzh}
\begin{uniten}Method: Set out the original work of one person and remove 1/5 of it; the remainder is the divisor [Liu Hui: "Remove 1/5 of it" means multiply by 4 and divide by 5.]. Take the volume of the ditch in chi as the dividend. Divide the dividend by the divisor to obtain the number of labourers used.\end{uniten}

\begin{unitzh}\unitlabel{jiuzhang-018.16}{\zhsmall 按此術「本置人功，去其五分之一」者，謂以四乘之，五而一，除去出土之功，取其定功。乃通分内子以爲法。以分母乘溝積尺爲實者，法裏有分，實裏通之。實如法而一，卽用徒人數。此以一人之積尺除其衆尺，故用徒人數。不盡者，等數約之而命分也。}\end{unitzh}
\begin{uniten}[Liu Hui: In this method, "originally set out the person's work and remove 1/5 of it" means multiplying it by 4 and dividing by 5, removing the work of carrying out the earth and taking its fixed work. Then convert to parts and bring in the numerator (tongfen nazi) to make the divisor. As for multiplying the volume of the ditch in chi by the denominator (fenmu) to make the dividend: the divisor has parts in it, so in the dividend one converts to them. Divide the dividend by the divisor: this is the number of labourers used. Here the total chi are divided by one person's volume in chi; hence it is the number of labourers used. Where it does not divide exactly, reduce by the equal number (dengshu) and name the fraction (mingfen).]\end{uniten}

\begin{unitzh}\unitlabel{jiuzhang-018.17}今有壍，上廣一丈六尺三寸，下廣一丈，深六尺三寸，袤一十三丈二尺一寸。問：積幾何？荅曰：一萬九百四十三尺八寸{\zhsmall 八寸者，謂穿地方尺，深八寸。此積餘有方寸中二分四釐五毫，乘之。貴欲從易，非其常定也}。\end{unitzh}
\begin{uniten}Suppose there is a moat with upper width 1 zhang 6 chi 3 cun, lower width 1 zhang, depth 6 chi 3 cun and length 13 zhang 2 chi 1 cun. What is the volume? Answer: 10943 chi 8 cun [Liu Hui: The 8 cun means ground dug one chi square and 8 cun deep. This volume has a remainder of 2 fen (分) 4 li (釐) 5 hao (毫) within a square cun; multiply it. What is valued is the wish to follow the easy way; this is not a constant, fixed rule.]. [Translator's note: "multiply it" (乘之) does not fit the context; the sense seems to be that the remainder is dropped.]\end{uniten}

\begin{unitzh}\unitlabel{jiuzhang-018.18}夏程人功八百七十一尺，并出土功五分之一，沙礫水石之功作太半，定功二百三十二尺一十五分尺之四。問：用徒幾何？荅曰：四十七人三千四百八十四分人之四百九。\end{unitzh}
\begin{uniten}The summer standard for one person's work is 871 chi, together with the work of carrying out the earth, 1/5, and the work for sand, gravel, water and stone, making up 2/3 (taiban); the fixed work is 232 4/15 chi. Question: how many labourers are used? Answer: 47 409/3484 persons.\end{uniten}

\begin{unitzh}\unitlabel{jiuzhang-018.19}術曰：置本人功，去其出土功五分之一，又去沙礫水石之功太半，餘爲法；以壍積尺爲實；實如法而一，卽用徒人數。\end{unitzh}
\begin{uniten}Method: Set out the original work of one person; remove from it the work of carrying out the earth, 1/5, and further remove the work of sand, gravel, water and stone, 2/3; the remainder is the divisor. Take the volume of the moat in chi as the dividend. Divide the dividend by the divisor: this is the number of labourers used.\end{uniten}

\begin{unitzh}\unitlabel{jiuzhang-018.20}{\zhsmall 按此術「置本人功，去其出土功五分之一」者，謂以四乘、五除。「又去沙礫水石作太半」者，一乘、三除，存其少半，取其定功。乃通分内子以爲法。以分母乘壍積尺爲實，法裏有分，實裏通之。實如法而一，卽用徒人數。不盡者，等數約之而命分也。}\end{unitzh}
\begin{uniten}[Liu Hui: In this method, "set out the original work of one person and remove from it the work of carrying out the earth, 1/5" means multiplying by 4 and dividing by 5. "Further remove the 2/3 that sand, gravel, water and stone make up" means multiplying by 1 and dividing by 3, keeping the 1/3 (shaoban), and taking the fixed work. Then convert to parts and bring in the numerator to make the divisor. Multiply the volume of the moat in chi by the denominator to make the dividend: the divisor has parts in it, so in the dividend one converts to them. Divide the dividend by the divisor: this is the number of labourers used. Where it does not divide exactly, reduce by the equal number and name the fraction.]\end{uniten}

\begin{unitzh}\unitlabel{jiuzhang-018.21}今有穿渠，上廣一丈八尺，下廣三尺六寸，深一丈八尺，袤五萬一千八百二十四尺。問：積幾何？荅曰：一千七萬四千五百八十五尺六寸。\end{unitzh}
\begin{uniten}Suppose there is a canal to be dug with upper width 1 zhang 8 chi, lower width 3 chi 6 cun, depth 1 zhang 8 chi and length 51824 chi. What is the volume? Answer: 10074585 chi 6 cun.\end{uniten}

\begin{unitzh}\unitlabel{jiuzhang-018.22}秋程人功三百尺。問：用徒幾何？荅曰：三萬三千五百八十一，一人功内少一十四尺四寸。\end{unitzh}
\begin{uniten}The autumn standard for one person's work is 300 chi. Question: how many labourers are used? Answer: 33581, and one person within whose work 14 chi 4 cun is lacking.\end{uniten}

\begin{unitzh}\unitlabel{jiuzhang-018.23}一千人先到。問：當受袤幾何？荅曰：一百五十四丈三尺二寸八十一分寸之八。\end{unitzh}
\begin{uniten}1000 persons arrive first. Question: what length should they receive? Answer: 154 zhang 3 chi 2 8/81 cun.\end{uniten}

\begin{unitzh}\unitlabel{jiuzhang-018.24}術曰：以一人功尺數乘先到人數爲實{\zhsmall 以一千人一日功爲實，立幂爲功}；并渠上下廣而半之，以深乘之爲法{\zhsmall 以渠廣深之立幂爲法}；實如法得袤尺。\end{unitzh}
\begin{uniten}Method: Multiply the number of persons arriving first by the number of chi of one person's work to make the dividend [Liu Hui: The one day's work of 1000 persons is taken as the dividend; the standing area is the work.]; add the upper and lower widths of the canal and halve, and multiply by the depth to make the divisor [Liu Hui: The standing area of the canal's width and depth is taken as the divisor.]; divide the dividend by the divisor to obtain the length in chi.\end{uniten}

\begin{unitzh}\unitlabel{jiuzhang-018.25}今有方堢壔{\zhsmall 堢者，堢城也。壔，音丁老切，又音纛，謂以土擁木也}，方一丈六尺，高一丈五尺。問：積幾何？荅曰：三千八百四十尺。\end{unitzh}
\begin{uniten}Suppose there is a square fort-mound (fang baodao) [Liu Hui: Bao (堢) is a fortress wall (baocheng). Dao (壔) is read by the fanqie 丁老, and is also read as 纛; it means earth banked around timber.], 1 zhang 6 chi square and 1 zhang 5 chi high. What is the volume? Answer: 3840 chi.\end{uniten}

\begin{unitzh}\unitlabel{jiuzhang-018.26}術曰：方自乘，以高乘之，卽積尺。\end{unitzh}
\begin{uniten}Method: Multiply the square side by itself and multiply by the height: this is the volume in chi.\end{uniten}

\begin{unitzh}\unitlabel{jiuzhang-018.27}今有圓堢壔，周四丈八尺，高一丈一尺。問：積幾何？荅曰：二千一百一十二尺{\zhsmall 於徽術，當積二千一十七尺一百五十七分尺之一百三十一。臣淳風等謹按：依密率，積二千一十六尺}。\end{unitzh}
\begin{uniten}Suppose there is a round fort-mound (yuan baodao) with circumference 4 zhang 8 chi and height 1 zhang 1 chi. What is the volume? Answer: 2112 chi [Liu Hui: By Hui's method the volume should be 2017 131/157 chi.] [Li Chunfeng et al.: According to the precise rate (milü), the volume is 2016 chi.].\end{uniten}

\begin{unitzh}\unitlabel{jiuzhang-018.28}術曰：周自相乘，以高乘之，十二而一。\end{unitzh}
\begin{uniten}Method: Multiply the circumference by itself, multiply by the height, and divide by 12.\end{uniten}

\begin{unitzh}\unitlabel{jiuzhang-018.29}{\zhsmall 此章諸術亦以周三徑一爲率，皆非也。於徽術，當以周自乘，以高乘之，又以二十五乘之，三百一十四而一。此之圓羃亦如圓田之羃也，求羃亦如圓田，而以高乘羃也。臣淳風等謹按：依密率，以七乘之，八十八而一。}\end{unitzh}
\begin{uniten}[Liu Hui: The methods of this chapter also take circumference 3, diameter 1 as the rate; all are wrong. By Hui's method one should multiply the circumference by itself, multiply by the height, further multiply by 25 and divide by 314. The circular area (mi) here is like the area of the circular field; the area is found as for the circular field, and then the area is multiplied by the height.] [Li Chunfeng et al.: According to the precise rate, multiply by 7 and divide by 88.]\end{uniten}

\begin{unitzh}\unitlabel{jiuzhang-018.30}今有方亭，下方五丈，上方四丈，高五丈。問：積幾何？荅曰：一十萬一千六百六十六尺太半尺。\end{unitzh}
\begin{uniten}Suppose there is a square pavilion (fangting) whose lower square is 5 zhang, upper square 4 zhang and height 5 zhang. What is the volume? Answer: 101666 2/3 chi.\end{uniten}

\begin{unitzh}\unitlabel{jiuzhang-018.31}術曰：上下方相乘，又各自乘，并之，以高乘之，三而一。\end{unitzh}
\begin{uniten}Method: Multiply the upper and lower squares by each other, and also multiply each by itself; add these together, multiply by the height and divide by 3.\end{uniten}

\begin{unitzh}\unitlabel{jiuzhang-018.32}{\zhsmall 此章有壍堵、陽馬，皆合而成立方，蓋說算者乃立棊三品，以效高深之積。假令方亭上方一尺，下方三尺，高一尺，其用棊也，中央立方一，四面壍堵四，四角陽馬四。上下方相乘爲三尺，以高乘之，約積三尺，是爲得中央立方一、四面壍堵各一。上方自乘，亦得中央立方一。下方自乘爲九，以高乘之，得積九尺，是爲中央立方一、四面壍堵各二、四角陽馬各三也。上方自乘，以高乘之，得積一尺，又爲中央立方一。凡三品棊皆一而爲三，故三而一，得積尺。用棊之數：立方三，壍堵、陽馬各十二，凡二十七棊。十二與三更差次之，而成方亭者三，驗矣。爲術又可令方差自乘，以高乘之，三而一，卽四陽馬也；上下方相乘，以高乘之，卽中央立方及四面壍堵也；并之，以爲方亭積數也。}\end{unitzh}
\begin{uniten}[Liu Hui: This chapter has moat-walls (qiandu) and yang-horses (yangma), which all combine to form cubes (lifang). Presumably, then, those who explained the calculations set up blocks (qi) of three kinds to demonstrate the volumes of heights and depths. Suppose a square pavilion with upper square 1 chi, lower square 3 chi and height 1 chi; the blocks it uses are: in the centre, 1 cube; on the four sides, 4 moat-walls; at the four corners, 4 yang-horses. The upper and lower squares multiplied together make 3 chi; multiplied by the height, the volume comes to 3 chi: this is obtaining 1 central cube and 1 moat-wall on each of the four sides. The upper square multiplied by itself also gives 1 central cube. The lower square multiplied by itself makes 9; multiplied by the height, it gives a volume of 9 chi: this is 1 central cube, 2 moat-walls on each of the four sides and 3 yang-horses at each of the four corners. The upper square multiplied by itself and multiplied by the height gives a volume of 1 chi, which again is 1 central cube. Altogether, the three kinds of blocks, each being one, have become three; therefore divide by 3 to obtain the volume in chi. The numbers of blocks used: cubes, 3; moat-walls and yang-horses, 12 each; in all 27 blocks. The twelves and the three, interchanged and set in order, form three square pavilions: this is verified. As a method, one may also take the difference of the squares, multiply it by itself, multiply by the height and divide by 3: this is the four yang-horses. Multiply the upper and lower squares together and multiply by the height: this is the central cube and the moat-walls on the four sides. Add them together to make the volume of the square pavilion.] [Translator's note: qiandu: a right prism whose cross-section is a right triangle; yangma: a pyramid on a rectangular base with one lateral edge perpendicular to the base.]\end{uniten}

\begin{unitzh}\unitlabel{jiuzhang-019.1}今有圓亭，下周三丈，上周二丈，高一丈。問積幾何？荅曰：五百二十七尺九分尺之七。{\zhsmall 於徽術，當積五百四尺四百七十一分尺之一百一十六也。按密率，爲積五百三尺三十三分尺之二十六。}\end{unitzh}
\begin{uniten}Suppose there is a round pavilion (yuanting) whose lower circumference is 3 zhang (丈), upper circumference 2 zhang, and height 1 zhang. What is the volume? Answer: 527 7/9 chi (尺). [Liu Hui: By Hui's method the volume should be 504 116/471 chi. According to the precise rate (mi lü), the volume is 503 26/33 chi.]\end{uniten}

\begin{unitzh}\unitlabel{jiuzhang-019.2}術曰：上下周相乘，又各自乘，并之，以高乘之，三十六而一。{\zhsmall 此術周三徑一之義，合以三除上下周，各爲上下徑，以相乘，又各自乘，并，以高乘之，三而一，爲方亭之積。假令三約上下周俱不盡，還通之，卽各爲上下徑。令上下徑分母相乘，又各自乘，并，以高乘之，爲三方亭之積。此合分母三相乘得九，分母各自乘亦得九，爲法除之，又三而一，得方亭之積。從方亭求圓亭之積，亦猶方羃中求圓羃，乃令圓率三乘之，方率四而一，得圓亭之積。前求方亭之積，乃以三而一；今求圓亭之積，亦各三乘之。二母旣同，故相凖折。惟以方羃四乘分母九，得三十六而連除之。}\end{unitzh}
\begin{uniten}Method: Multiply the upper and lower circumferences together, also multiply each by itself, add them, multiply by the height, and divide by 36. [Liu Hui: This method follows the sense of "circumference 3, diameter 1". One should divide the upper and lower circumferences by 3, making the upper and lower diameters respectively, multiply these together, also multiply each by itself, add, multiply by the height, and divide by 3: this is the volume of the square pavilion (fangting). Suppose that when the upper and lower circumferences are reduced by 3 neither comes out exactly: convert them back (tong), and they are then the upper and lower diameters respectively. Let the upper and lower diameters with their denominators be multiplied together, also each multiplied by itself, added, and multiplied by the height: this is the volume of 3 square pavilions. Here one should multiply the denominators 3 together, giving 9; the denominators each multiplied by itself also give 9; take this as the divisor (fa) and divide, and again divide by 3, to obtain the volume of the square pavilion. Finding the volume of the round pavilion from the square pavilion is just like finding the circle's area (mi) within the square's area: so multiply by the circle rate 3 and divide by the square rate 4 to obtain the volume of the round pavilion. Before, in finding the volume of the square pavilion, one divided by 3; now, in finding the volume of the round pavilion, one also multiplies each by 3. Since the 2 denominators are the same, they therefore offset (zhunzhe) each other. Only multiply the denominator 9 by the square area 4, giving 36, and divide by it in one go (lianchu).]\end{uniten}

\begin{unitzh}\unitlabel{jiuzhang-019.3}{\zhsmall 於徽術，當上下周相乘，又各自乘，并，以高乘之，又二十五乘之，九百四十二而一。此方亭四角圓殺，比於方亭，二百分之一百五十七。爲術之意，先作方亭，三而一，則此據上下徑爲之者，當又以一百五十七乘之，六百而一也。今據周爲之，若於圓堢壔，又以二十五乘之，三百一十四而一，則先得三圓亭矣。故以三百一十四爲九百四十二而一，并除之。}{\zhsmall 臣淳風等謹按：依密率，以七乘之，二百六十四而一。}\end{unitzh}
\begin{uniten}[Liu Hui: By Hui's method, one should multiply the upper and lower circumferences together, also multiply each by itself, add, multiply by the height, further multiply by 25, and divide by 942. This is the square pavilion with its 4 corners rounded off (yuansha); compared with the square pavilion it is 157/200. The intent of the method is first to make the square pavilion and divide by 3; so, where this is done on the basis of the upper and lower diameters, one should further multiply by 157 and divide by 600. Now it is done on the basis of the circumferences: as for the round fort (yuan baodao), further multiply by 25 and divide by 314, and one has then first obtained 3 round pavilions. Therefore 314 is made into "divide by 942", and they are divided together.] [Li Chunfeng et al.: Your servants respectfully note: according to the precise rate, multiply by 7 and divide by 264.]\end{uniten}

\begin{unitzh}\unitlabel{jiuzhang-019.4}今有方錐，下方二丈七尺，高二丈九尺。問積幾何？荅曰：七千四十七尺。\end{unitzh}
\begin{uniten}Suppose there is a square cone (fangzhui) whose lower side is 2 zhang 7 chi and height 2 zhang 9 chi. What is the volume? Answer: 7047 chi.\end{uniten}

\begin{unitzh}\unitlabel{jiuzhang-019.5}術曰：下方自乘，以高乘之，三而一。{\zhsmall 按：此術假令方錐下方二尺，高一尺，卽四陽馬。如術爲之，用十二陽馬成三方錐，故三而一，得陽馬也。}\end{unitzh}
\begin{uniten}Method: Multiply the lower side by itself, multiply by the height, and divide by 3. [Liu Hui: Note: in this method, suppose a square cone has lower side 2 chi and height 1 chi: it is then 4 yangma ('yang horse'). Doing it according to the method, 12 yangma are used to form 3 square cones; therefore dividing by 3 gives the yangma.]\end{uniten}

\begin{unitzh}\unitlabel{jiuzhang-019.6}今有圓錐，下周三丈五尺，高五丈一尺。問積幾何？荅曰：一千七百三十五尺一十二分尺之五。{\zhsmall 於徽術，當積一千六百五十八尺三百一十四分尺之十三。依密率，爲積一千六百五十六尺八十八分尺之四十七。}\end{unitzh}
\begin{uniten}Suppose there is a round cone (yuanzhui) whose lower circumference is 3 zhang 5 chi and height 5 zhang 1 chi. What is the volume? Answer: 1735 5/12 chi. [Liu Hui: By Hui's method the volume should be 1658 13/314 chi. According to the precise rate, the volume is 1656 47/88 chi.]\end{uniten}

\begin{unitzh}\unitlabel{jiuzhang-019.7}術曰：下周自乘，以高乘之，三十六而一。{\zhsmall 按：此術圓錐下周以爲方錐下方。方錐下方令自乘，以高乘之，合三而一，得大錐方之積。大錐方之積合十二圓矣。今求一圓，復合十二除之，故令三乘十二，得三十六而連除。於徽術，當下周自乘，以高乘之，又以二十五乘之，九百四十二而一。圓錐比於方錐，亦二百分之一百五十七。令徑自乘者，亦當以一百五十七乘之，六百而一。其說如圓亭也。}{\zhsmall 臣淳風等謹按：依密率，以七乘之，二百六十四而一。}\end{unitzh}
\begin{uniten}Method: Multiply the lower circumference by itself, multiply by the height, and divide by 36. [Liu Hui: Note: in this method the lower circumference of the round cone is taken as the lower side of a square cone. The lower side of the square cone is multiplied by itself and multiplied by the height, and one should divide by 3, obtaining the volume of the large square cone. The volume of the large square cone should amount to 12 round ones. Now, to find 1 round one, one should further divide by 12; therefore multiply 12 by 3, giving 36, and divide in one go. By Hui's method, one should multiply the lower circumference by itself, multiply by the height, further multiply by 25, and divide by 942. The round cone compared with the square cone is likewise 157/200. Where the diameter is multiplied by itself, one should likewise multiply by 157 and divide by 600. The explanation is as for the round pavilion.] [Li Chunfeng et al.: Your servants respectfully note: according to the precise rate, multiply by 7 and divide by 264.]\end{uniten}

\begin{unitzh}\unitlabel{jiuzhang-019.8}今有壍堵，下廣二丈，袤一十八丈六尺，高二丈五尺。問積幾何？荅曰：四萬六千五百尺。\end{unitzh}
\begin{uniten}Suppose there is a qiandu ('trench wall') whose lower breadth is 2 zhang, length 18 zhang 6 chi, and height 2 zhang 5 chi. What is the volume? Answer: 46500 chi.\end{uniten}

\begin{unitzh}\unitlabel{jiuzhang-019.9}術曰：廣袤相乘，以高乘之，二而一。{\zhsmall 邪解立方得兩壍堵，雖復橢方，亦爲壍堵，故二而一。此則合所規羃。推其物體，蓋爲壍上疊也。其形如城而無上廣，與所規棊形異而同實。未聞所以名之爲壍堵之說也。}\end{unitzh}
\begin{uniten}Method: Multiply the breadth and the length together, multiply by the height, and divide by 2. [Liu Hui: Cutting a cube (lifang) obliquely (xiejie) gives 2 qiandu; even if it is an oblong square (tuofang), it still makes qiandu; therefore divide by 2. This then accords with the planned area. Inferring its physical body, it is probably what is piled up on a trench. Its shape is like a city wall but without upper breadth; it differs in shape from the planned blocks (qi) but is the same in volume. I have not heard an explanation of why it is named qiandu.]\end{uniten}

\begin{unitzh}\unitlabel{jiuzhang-019.10}今有陽馬，廣五尺，袤七尺，高八尺。問積幾何？荅曰：九十三尺少半尺。\end{unitzh}
\begin{uniten}Suppose there is a yangma whose breadth is 5 chi, length 7 chi, and height 8 chi. What is the volume? Answer: 93 chi and a lesser half (shaoban, 1/3) chi.\end{uniten}

\begin{unitzh}\unitlabel{jiuzhang-019.11}術曰：廣袤相乘，以高乘之，三而一。{\zhsmall 按：此術，陽馬之形，方錐一隅也。今謂四柱屋隅爲陽馬。假令廣袤各一尺，高一尺，相乘之，得立方積一尺。邪解立方得兩壍堵；邪解壍堵，其一爲陽馬，一爲鼈臑。陽馬居二，鼈臑居一，不易之率也。合兩鼈臑成一陽馬，合三陽馬而成一立方，故三而一。驗之以棊，其形露矣。悉割陽馬，凡爲六鼈臑。觀其割分，則體勢互通，蓋易了也。其棊或脩短，或廣狹，立方不等者，亦割分以爲六鼈臑。其形不悉相似，然見數同，積實均也。鼈臑殊形，然陽馬異體，則不純合。不純合，則難爲之矣。何則？按：邪解方棊以爲壍堵者，必當以半爲分；邪解壍堵以爲陽馬者，亦必當以半爲分，一從一横耳。設陽馬爲分内，鼈臑爲分外，棊雖或隨脩短廣狹，猶有此分常率。如殊形異體亦同也者，以此而已。}\end{unitzh}
\begin{uniten}Method: Multiply the breadth and the length together, multiply by the height, and divide by 3. [Liu Hui: Note: in this method, the shape of the yangma is one corner of a square cone. Nowadays the corner of a house with 4 pillars is called yangma. Suppose the breadth and the length are each 1 chi and the height 1 chi; multiplying them together gives a cube of volume 1 chi. Cutting the cube obliquely gives 2 qiandu; cutting a qiandu obliquely, one of them is a yangma and one is a bienao ('turtle's forelimb'). The yangma takes 2 and the bienao takes 1: this is an unchanging rate. Combining 2 bienao forms 1 yangma, and combining 3 yangma forms 1 cube; therefore divide by 3. Verify it with blocks, and the shapes become evident. Cut up all the yangma, and altogether there are 6 bienao. Observe how they are cut and divided, and the bodies and their configurations (tishi) communicate with one another; this is easily understood. Where the blocks are long or short, or broad or narrow — cubes that are unequal — they too are cut and divided to make 6 bienao. Their shapes are not all alike, yet the numbers seen are the same and the volumes are equal. The bienao differ in shape, yet the yangma differ in body, and then they do not combine purely. If they do not combine purely, it is difficult to deal with them. Why so? Note: when a square block is cut obliquely to make qiandu, it must be divided in half; when a qiandu is cut obliquely to make a yangma, it must likewise be divided in half — one lengthwise, one crosswise, that is all. Suppose the yangma is the inside of the division and the bienao the outside of the division: although the blocks may vary with being long or short, broad or narrow, there is still this constant rate of division. If different shapes and different bodies are also the same, it is by this alone.]\end{uniten}

\begin{unitzh}\unitlabel{jiuzhang-019.12}{\zhsmall 其使鼈臑廣袤各高二尺，用壍堵、鼈臑之棊各二，皆用赤棊；又使陽馬之廣、袤、高各二尺，用立方之棊一，壍堵、陽馬之棊各二，皆用黑棊。棊之赤黑，接爲壍堵，廣、袤、高各二尺。於是中效其廣，又中分其高，令赤黑壍堵各自適當一方，高二尺，方二尺。每二分鼈臑，則一陽馬也。其餘兩端各積本體，合成一方焉。是爲别種而方者率居三，通其體而方者率居一。雖方隨棊改，而固有常然之勢也。}\end{unitzh}
\begin{uniten}[Liu Hui: Let the breadth and length of the bienao each, and the height, be 2 chi, using 2 each of qiandu and bienao blocks, all of them red blocks; further let the breadth, length and height of the yangma each be 2 chi, using 1 cube block and 2 each of qiandu and yangma blocks, all of them black blocks. The red and black blocks are joined to make a qiandu whose breadth, length and height are each 2 chi. Thereupon divide (xiao) its breadth at the middle, and again divide its height at the middle, so that the red and the black qiandu each, by itself, exactly matches 1 cube, of height 2 chi and side 2 chi. Every 2 parts of bienao are then 1 yangma. As for the rest, the 2 ends each pile up their own bodies and combine to form 1 cube. Thus those formed into cubes by separate kinds take 3 in the rate, and those formed into a cube by joining their bodies take 1 in the rate. Although the cubes change with the blocks, there is inherently a constant configuration.]\end{uniten}

\begin{unitzh}\unitlabel{jiuzhang-019.13}{\zhsmall 按：餘數具而可知者，有一、二分之别，卽一、二之爲率定矣。其於理也，豈虛矣？若爲數而窮之，置餘廣、袤、高之數，各半之，則四分之三又可知也。半之彌少，其餘彌細。至細曰微，微則無形。由是言之，安取餘哉？數而求窮之者，謂以情推，不用籌算。鼈臑之物，不同器用；陽馬之形，或隨脩短廣狹。然不有鼈臑，無以審陽馬之數；不有陽馬，無以知錐亭之數，功實之主也。}\end{unitzh}
\begin{uniten}[Liu Hui: Note: the amounts besides the remainder, which are complete and can be known, have the distinction of 1 and 2 parts; thus 1 and 2 as the rate are settled. As to its principle, how could it be empty? If one counts and exhausts it: set out the amounts of the breadth, length and height of the remainder and halve each, and then 3/4 can again be known. The more it is halved, the smaller it becomes, and the finer its remainder. The utmost fineness is called minute (wei); what is minute has no shape. Speaking from this, how could any remainder be taken? "Counting and seeking to exhaust it" means inferring by reasoning (qing), not calculating with counting rods (chousuan). The bienao as a thing is unlike implements in use; the shape of the yangma may vary with being long or short, broad or narrow. Yet without the bienao there is no means to ascertain the amount of the yangma, and without the yangma there is no means to know the amounts of cones and pavilions: it is the master of the volumes of works.]\end{uniten}

\begin{unitzh}\unitlabel{jiuzhang-020.1}今有鼈臑，下廣五尺，無袤；上袤四尺，無廣；高七尺。問積幾何？荅曰：二十三尺少半尺。\end{unitzh}
\begin{uniten}Suppose there is a turtle's forelimb (bienao) with a lower width of 5 chi (尺) and no length, an upper length of 4 chi and no width, and a height of 7 chi. What is the volume? Answer: 23 1/3 chi.\end{uniten}

\begin{unitzh}\unitlabel{jiuzhang-020.2}術曰：廣袤相乘，以高乘之，六而一。{\zhsmall 按此術：臑者，臂骨也。或曰半陽馬，其形有似鼈肘，故以名云。中破陽馬，得兩鼈臑，鼈臑之見數卽陽馬之半數。數同而實據半，故云六而一卽得。}\end{unitzh}
\begin{uniten}Method: Multiply the width and the length together, multiply by the height, and divide by 6. [Liu Hui: Consider this method: nao (臑) is the arm bone. Some say it is half a corner beam (yangma); its shape resembles a turtle's elbow, hence the name. Splitting a corner beam through the middle gives two turtle's forelimbs; the manifest number of a turtle's forelimb is half the number of the corner beam. The numbers are the same, but the content (shi) amounts to half; therefore it says "divide by 6", and the result is obtained.]\end{uniten}

\begin{unitzh}\unitlabel{jiuzhang-020.3}今有羡除，下廣六尺，上廣一丈，深三尺；末廣八尺，無深；袤七尺。問積幾何？荅曰：八十四尺。\end{unitzh}
\begin{uniten}Suppose there is a tomb passage (xianchu) with a lower width of 6 chi, an upper width of 1 zhang (丈) and a depth of 3 chi; an end width of 8 chi and no depth; and a length of 7 chi. What is the volume? Answer: 84 chi.\end{uniten}

\begin{unitzh}\unitlabel{jiuzhang-020.4}術曰：并三廣，以深乘之，又以袤乘之，六而一。{\zhsmall 按此術：羨除，實隧道也。其所穿地，上平下邪，似兩鼈臑夾一壍堵，卽羨除之形。假令用此棊：上廣三尺，深一尺，下廣一尺；末廣一尺，無深；袤一尺。下廣皆壍堵之廣；上廣者，兩鼈臑與一壍堵相連之廣也。以深袤乘，得積五尺：鼈臑居二，壍堵居三，其於本棊皆以爲六，故六而一。合四陽馬以爲方錐。邪畫方錐之底，亦令爲中方。就中方削而上合，全爲中方錐之半。於是陽馬之棊悉中解矣。中錐離而爲四鼈臑焉，故外錐之半亦爲四鼈臑。雖背正異形，與常所謂鼈臑參不相似，實則同也。所云「夾壍堵」者，中錐之鼈臑也。凡壍堵上袤短者，連陽馬也；下袤短者，與鼈臑連也；下兩袤相等者，亦與鼈臑連也。并三廣，以高袤乘，六而一，皆其積也。今此羨除之廣，卽壍堵之袤也。}\end{unitzh}
\begin{uniten}Method: Add the three widths, multiply by the depth, multiply again by the length, and divide by 6. [Liu Hui: Consider this method: a tomb passage is in fact a tunnel. The ground it cuts through is level above and slanting below, like two turtle's forelimbs flanking one trench wall (qiandu): this is the shape of a tomb passage. Let us use these blocks (qi): upper width 3 chi, depth 1 chi, lower width 1 chi; end width 1 chi, with no depth; length 1 chi. The lower width is in every case the width of the trench wall; the upper width is the width of the two turtle's forelimbs and the one trench wall joined together. Multiplying by the depth and the length gives a volume of 5 chi: the turtle's forelimbs occupy 2 and the trench wall occupies 3; with respect to the original blocks, all are taken sixfold; therefore divide by 6. Put 4 corner beams together to make a square cone (fangzhui). Draw slanting lines on the base of the square cone, also making a central square. Cutting from the central square and joining at the top, the whole has the central square cone as its half. Thereupon the corner-beam blocks are all split through the middle. The central cone comes apart into 4 turtle's forelimbs; therefore the half that is the outer cone also makes 4 turtle's forelimbs. Although back and front differ in shape, and compared with what are ordinarily called turtle's forelimbs they are not alike, in content they are the same. Those said to "flank the trench wall" are the turtle's forelimbs of the central cone. In general, a trench wall whose upper length is short is joined to corner beams; one whose lower length is short is joined with turtle's forelimbs; one whose two lower lengths are equal is also joined with turtle's forelimbs. Add the three widths, multiply by the height and the length, and divide by 6: in every case this is their volume. Now the widths of this tomb passage are the lengths of the trench wall.]\end{uniten}

\begin{unitzh}\unitlabel{jiuzhang-020.5}{\zhsmall 按此本是三廣不等，卽與鼈臑連者。别而言之：中央壍堵廣六尺，高三尺，袤七尺；末廣之兩旁，各一小鼈臑，皆與壍堵等。令小鼈臑居裏，大鼈臑居表，則大鼈臑出橢，皆方錐：下廣三尺，袤六尺，高七尺。分取其半，則爲袤三尺，以高廣乘之，三而一，卽半錐之積也。邪解半錐，得此兩大鼈臑。求其積，亦當六而一，合與常率矣。按陽馬之棊兩邪，棊底方。當其方也，不問旁角而割之，相半可知也。推此上連無成不方，故方錐與陽馬同實。角而割之者，相半之勢。此大小鼈臑可知更相表裏，但體有背正也。}\end{unitzh}
\begin{uniten}[Liu Hui: Consider: this is originally a case in which the three widths are unequal, that is, one joined with turtle's forelimbs. Speaking of the parts separately: the central trench wall has a width of 6 chi, a height of 3 chi and a length of 7 chi; on the two sides of the end width there is one small turtle's forelimb each, both equal to the trench wall. Let the small turtle's forelimbs lie inside and the large turtle's forelimbs outside; then the large turtle's forelimbs come out oblong, all making a square cone with a lower width of 3 chi, a length of 6 chi and a height of 7 chi. Take half of it: then the length is 3 chi; multiply by the height and the width, and divide by 3: this is the volume of the half cone. Splitting the half cone obliquely gives these two large turtle's forelimbs. To find their volume one should likewise divide by 6, which agrees with the usual rate (lü). Consider: the corner-beam block has two slants, and the block's base is square. As to its square, if one cuts it without regard to side or corner, it can be known that the parts are halves of each other. Inferring from this, joining upward nothing is formed that is not square; therefore the square cone and the corner beam have the same content. Cutting it at the corner is the configuration of halving. From this it can be known that the large and small turtle's forelimbs are in turn outside and inside of each other; only their bodies have back and front.]\end{uniten}

\begin{unitzh}\unitlabel{jiuzhang-020.6}今有芻甍，下廣三丈，袤四丈；上袤二丈，無廣；高一丈。問積幾何？荅曰：五千尺。\end{unitzh}
\begin{uniten}Suppose there is a hay ridge (chumeng) with a lower width of 3 zhang and a length of 4 zhang; an upper length of 2 zhang and no width; and a height of 1 zhang. What is the volume? Answer: 5000 chi.\end{uniten}

\begin{unitzh}\unitlabel{jiuzhang-020.7}術曰：倍下袤，上袤從之，以廣乘之，又以高乘之，六而一。{\zhsmall 推明義理者，舊說云：凡積芻甍有上下廣曰童，甍謂其屋蓋之茨也。是故甍之下廣袤與童之上廣袤等。正斬方亭兩邊，合之，卽芻甍之形也。假令下廣二尺，袤三尺；上袤一尺，無廣；高一尺。其用棊也，中央壍堵二，兩端陽馬各二。倍下袤，上袤從之，爲七尺，以高廣乘之，得羃十四尺。陽馬之羃各居一，壍堵之羃各居三，以高乘之，得積十四尺。其於本棊也，皆一而爲六，故六而一，卽得。亦可令上下袤差乘廣，以高乘之，三而一，卽四陽馬也；下廣乘上袤而半之，高乘之，卽二壍堵；并之，以爲甍積也。}\end{unitzh}
\begin{uniten}Method: Double the lower length, add the upper length to it, multiply by the width, multiply again by the height, and divide by 6. [Liu Hui: To elucidate the meaning and principle: the old explanation says that any stacked hay ridge that has an upper and a lower width is called tong (童); meng (甍) refers to the thatch of a house's roof covering. Therefore the lower width and length of the meng are equal to the upper width and length of the tong. Cut off the two sides of a square pavilion (fangting) straight and put them together: that is the shape of the hay ridge. Let the lower width be 2 chi and the length 3 chi; the upper length 1 chi, with no width; and the height 1 chi. As for the blocks used: in the centre, 2 trench walls; at the two ends, 2 corner beams each. Double the lower length and add the upper length to it, making 7 chi; multiply by the height and the width, obtaining an area (mi) of 14 chi. The areas of the corner beams each occupy 1, the areas of the trench walls each occupy 3; multiply by the height, obtaining a volume of 14 chi. With respect to the original blocks, in every case one has become six; therefore divide by 6, and it is obtained. One may also multiply the width by the difference of the upper and lower lengths, multiply by the height, and divide by 3: that is the 4 corner beams; multiply the upper length by the lower width and halve, then multiply by the height: that is the 2 trench walls; add them together to make the volume of the ridge.]\end{uniten}

\begin{unitzh}\unitlabel{jiuzhang-020.8}芻童、曲池、盤池、⿱冝八 -- 𡨋谷，皆同術。術曰：倍上袤，下袤從之；亦倍下袤，上袤從之；各以其廣乘之；并，以高若深乘之，皆六而一。{\zhsmall 按此術：假令芻童上廣一尺，袤二尺；下廣三尺，袤四尺；高一尺。其用棊也，中央立方二，四面壍堵六，四角陽馬四。倍下袤爲八，上袤從之爲十，以下廣乘之，得積三十尺，是爲得中央立方各三，兩邊壍堵各四，兩旁壍堵各六，四角陽馬亦各六。復倍上袤，下袤從之爲八，以高廣乘之，得積八尺，是爲得中央立方亦各三，兩端壍堵各二。并兩旁，三品棊皆一而爲六，故六而一，卽得。爲術又可令上下廣袤差相乘，以高乘之，三而一，亦四陽馬；上下廣袤互相乘，并而半之，以高乘之，卽四面六壍堵與二立方；并之，爲芻童積。又可令上下廣袤互相乘而半之，上下廣袤又各自乘，并，以高乘之，三而一，卽得也。}\end{unitzh}
\begin{uniten}Hay stack (chutong), curved pool (quchi), basin pool (panchi) and dark valley (minggu) all have the same method. [Translator's note: the transcription gives the first character of "dark valley" as ⿱冝八 -- 𡨋, a form of 冥 "dark".] Method: Double the upper length and add the lower length to it; also double the lower length and add the upper length to it; multiply each by its width; add them, multiply by the height or the depth, and in every case divide by 6. [Liu Hui: Consider this method: let a hay stack have an upper width of 1 chi and a length of 2 chi; a lower width of 3 chi and a length of 4 chi; and a height of 1 chi. As for the blocks used: in the centre, 2 cubes (lifang); on the four sides, 6 trench walls; at the four corners, 4 corner beams. Double the lower length, making 8; add the upper length to it, making 10; multiply by the lower width, obtaining a volume of 30 chi: this gives the central cubes 3 each, the trench walls on the two edges 4 each, the trench walls on the two flanks 6 each, and the corner beams at the four corners also 6 each. Again double the upper length and add the lower length to it, making 8; multiply by the height and the width, obtaining a volume of 8 chi: this gives the central cubes also 3 each, and the trench walls at the two ends 2 each. Adding the two flanks, for the three kinds of blocks in every case one has become six; therefore divide by 6, and it is obtained. As a method, one may also multiply together the differences of the upper and lower widths and lengths, multiply by the height, and divide by 3: this is likewise the 4 corner beams; multiply the upper and lower widths and lengths crosswise, add, halve, and multiply by the height: this is the 6 trench walls of the four sides and the 2 cubes; add them together to make the volume of the hay stack. Again, one may multiply the upper and lower widths and lengths crosswise and halve; also multiply the upper and the lower width and length each among themselves; add, multiply by the height, and divide by 3: this gives it.]\end{uniten}

\begin{unitzh}\unitlabel{jiuzhang-020.9}其曲池者，并上中外周而半之，以爲上袤；亦并下中外周而半之，以爲下袤。{\zhsmall 此池環而不通匝，形如盤虵而曲之。亦云「周」者，謂如委榖依垣之周耳。引而伸之，周爲袤。求袤之意，環田也。}\end{unitzh}
\begin{uniten}As for the curved pool, add the upper inner and outer circumferences and halve, taking this as the upper length; likewise add the lower inner and outer circumferences and halve, taking this as the lower length. [Liu Hui: This pool forms a ring but does not close all the way round; its shape is like a coiled snake, curving. That it is also spoken of as "circumference" merely means that it is like the circumference of grain heaped against a wall. Pull it out and stretch it: the circumference becomes the length. The idea of finding the length is that of the ring field (huantian).]\end{uniten}

\begin{unitzh}\unitlabel{jiuzhang-020.10}今有芻童，下廣二丈，袤三丈；上廣三丈，袤四丈；高三丈。問積幾何？荅曰：一萬六千五百尺。\end{unitzh}
\begin{uniten}Suppose there is a hay stack with a lower width of 2 zhang and a length of 3 zhang; an upper width of 3 zhang and a length of 4 zhang; and a height of 3 zhang. What is the volume? Answer: 16500 chi.\end{uniten}

\begin{unitzh}\unitlabel{jiuzhang-020.11}今有曲池，上中周二丈，外周四丈，廣一丈；下中周一丈四尺，外周二丈四尺，廣五尺；深一丈。問積幾何？荅曰：一千八百八十三尺三寸少半寸。\end{unitzh}
\begin{uniten}Suppose there is a curved pool with an upper inner circumference of 2 zhang, an outer circumference of 4 zhang and a width of 1 zhang; a lower inner circumference of 1 zhang 4 chi, an outer circumference of 2 zhang 4 chi and a width of 5 chi; and a depth of 1 zhang. What is the volume? Answer: 1883 chi 3 1/3 cun (寸).\end{uniten}

\begin{unitzh}\unitlabel{jiuzhang-020.12}今有盤池，上廣六丈，袤八丈；下廣四丈，袤六丈；深二丈。問積幾何？荅曰：七萬六百六十六尺太半尺。\end{unitzh}
\begin{uniten}Suppose there is a basin pool with an upper width of 6 zhang and a length of 8 zhang; a lower width of 4 zhang and a length of 6 zhang; and a depth of 2 zhang. What is the volume? Answer: 70666 2/3 chi.\end{uniten}

\begin{unitzh}\unitlabel{jiuzhang-020.13}負土往來七十步，其二十步上下棚除，棚除二當平道五；踟蹰之間，十加一；載輸之間，三十步；定一返一百四十步。土籠積一尺六寸。秋程人功行五十九里半。問人到積尺及用徒各幾何？荅曰：人到二百四尺。用徒三百四十六人一百五十三分人之六十二。\end{unitzh}
\begin{uniten}Carrying earth on the back, the going and returning is 70 bu (步); of this, 20 bu are up and down scaffold steps (pengchu), and on the scaffold steps 2 count as 5 on level road; for the interval of lingering, add 1 for every 10; for the interval of loading and unloading, 30 bu; one round trip is fixed at 140 bu. The volume of an earth basket is 1 chi 6 cun. By the autumn norm (cheng), a person's labour is to walk 59 1/2 li (里). What are the volume in chi that a person delivers and the labourers used? Answer: A person delivers 204 chi. Labourers used: 346 62/153 persons.\end{uniten}

\begin{unitzh}\unitlabel{jiuzhang-020.14}術曰：以一籠積尺乘程行步數，爲實。往來上下棚除二當平道五。{\zhsmall 棚閣除，邪道，有上下之難，故使二當五也。}置定往來步數，十加一，及載輸之間三十步，以爲法。除之，所得卽一人所到尺。{\zhsmall 按此術：棚閣除，邪道，有上下之難，故使二當五。置定往來步數，十加一，及載輸之間三十步，是爲往來求一返，凡用一百四十步。於今有術，爲所有行率；籠積一尺六寸，爲所求到土率；程行五十九里半，爲所有數；而今有之，卽人到尺數。「以所到約積尺，卽用徒人數」者，此一人之積除其衆積尺，故得用徒人數。爲術又可令往來一返所用之步約程行，爲返數；乘籠積，爲一人所到。以此術與今有術相反覆，則乘除之或先後，意各有所在而同歸耳。}以所到約積尺，卽用徒人數。\end{unitzh}
\begin{uniten}Method: Multiply the number of bu of the norm walk by the volume in chi of one basket, as the dividend (shi). Going and returning, up and down the scaffold steps, 2 count as 5 on level road. [Liu Hui: The scaffold-gallery (pengge) steps are a slanting path, with the difficulty of going up and down; therefore 2 are made to count as 5.] Set out the fixed number of bu going and returning, add 1 for every 10, together with the 30 bu for loading and unloading, and take this as the divisor (fa). Divide: what is obtained is the chi that one person delivers. [Liu Hui: Consider this method: the scaffold-gallery steps are a slanting path, with the difficulty of going up and down; therefore 2 are made to count as 5. Setting out the fixed number of bu going and returning, adding 1 for every 10, together with the 30 bu for loading and unloading: this is to find one round trip from the going and returning; in all 140 bu are used. In the "Suppose there is" method (jinyou shu), this is the walking rate one has (suoyou); the basket volume of 1 chi 6 cun is the rate of earth delivered that one seeks (suoqiu); the norm walk of 59 1/2 li is the number one has; applying "Suppose there is" to them gives the number of chi a person delivers. As for "divide the volume in chi by what is delivered; this gives the number of labourers used": this divides the total volume in chi by the volume of one person, and therefore one obtains the number of labourers used. As a method, one may also divide the norm walk by the bu used for one round trip going and returning, as the number of round trips; multiply the basket volume by it, as what one person delivers. When this method and the "Suppose there is" method are reversed against each other, multiplication and division come sometimes earlier, sometimes later; the intentions each have their place, but they return to the same end.] Divide (yue) the volume in chi by what is delivered; this gives the number of labourers used.\end{uniten}

\begin{unitzh}\unitlabel{jiuzhang-021.1}今有⿱冝八 -- 𡨋谷，上廣二丈，袤七丈；下廣八尺，袤四丈；深六丈五尺。問積幾何？荅曰：五萬二千尺。載土往來二百步，載輸之間一里，程行五十八里；六人共車，車載三十四尺七寸。問人到積尺及用徒各幾何？荅曰：人到二百一尺五十分尺之十三；用徒二百五十八人一萬六十三分人之三千七百四十六。\end{unitzh}
\begin{uniten}Suppose there is a dark valley (minggu) [Translator's note: the first character is given in the transcription as ⿱冝八 (冝 above 八), followed by 𡨋 as an alternative; it is read as a form of 冥 'dark'.]: upper width 2 zhang (丈), length 7 zhang; lower width 8 chi (尺), length 4 zhang; depth 6 zhang 5 chi. What is the volume? Answer: 52000 chi. In carrying the earth, going and coming is 200 bu (步); between loading and delivering is 1 li (里); the standard march (chengxing) is 58 li; 6 persons share a cart, and a cart carries 34 chi 7 cun (寸). What are the volume in chi delivered per person and the laborers used, each? Answer: Delivered per person: 201 13/50 chi. Laborers used: 258 3746/10063 persons.\end{uniten}

\begin{unitzh}\unitlabel{jiuzhang-021.2}術曰：以一車積尺乘程行步數，爲實；置今往來步數，加載輸之間一里，以車六人乘之，爲法；除之，所得卽一人所到尺。{\zhsmall 按此術今有之義。以載輸及往來并得五百步，爲所有行率；車載三十四尺七寸，爲所求到土率；程行五十八里，通之爲步，爲所有數；而今有之，所得則一車所到。欲得人到者，當以六人除之，卽得。術有分，故亦更令乘法而并除者，亦用以半尺數以爲一人到土率，六人乘五百步爲行率也。又可以五百步爲行率，令六人約半積尺數爲一人到土率，以載土術入之。入之者，亦可求返數也。要取其會通而已。術恐有分，故令乘法而并除。以所到約積尺，卽用徒人數者，以一人所積尺除其衆積，故得射徒人數也。}以所到約積尺，卽用徒人數。\end{unitzh}
\begin{uniten}Method: Take the volume in chi of one cart to multiply the number of bu of the standard march, as the dividend (shi). Set out the present number of bu of going and coming, add the 1 li between loading and delivering, and multiply it by the 6 persons of a cart, as the divisor (fa). Divide; what is obtained is the chi delivered by one person. [Liu Hui: Note: this method has the sense of 'Suppose there is' (jinyou). Take the loading-and-delivering distance and the going and coming, added together to give 500 bu, as the travel rate one has (suoyou xinglü); the cart's load of 34 chi 7 cun as the sought rate of delivered earth (suoqiu daotu lü); and the standard march of 58 li, converted into bu, as the number one has (suoyou shu); then apply 'Suppose there is' to them, and what is obtained is what one cart delivers. If one wishes to get what a person delivers, one should divide it by the 6 persons, and then one obtains it. Because the method has fractions, it therefore further has [the 6 persons] multiply the divisor and divides all together: this also uses the half [Translator's note: 半 'half', here and in the next sentence, does not fit the computation; the text appears corrupt.] number of chi as the rate of earth delivered by one person, and 6 persons multiplied by 500 bu as the travel rate. Alternatively, one may take 500 bu as the travel rate, have the 6 persons divide the half volume-in-chi number to make the rate of earth delivered by one person, and enter it into the method for carrying earth. Entering it thus, one can also find the number of round trips (fanshu). What matters is only to grasp how they connect (huitong). The method fears there would be fractions; therefore it has [the 6 persons] multiply the divisor and divides all together. 'Divide the volume in chi by what is delivered; this is the number of laborers used': this divides the total volume by the volume in chi of one person, and therefore one can determine the number of laborers.] Divide the volume in chi by what is delivered; this is the number of laborers used.\end{uniten}

\begin{unitzh}\unitlabel{jiuzhang-021.3}今有委粟平地，下周一十二丈，高二丈。問積及爲粟幾何？荅曰：積八千尺，{\zhsmall 於徽術，當積七千六百四十三尺一百五十七分尺之四十九。臣淳風等謹按：依密率，爲積七千六百三十六尺十一分尺之四。}爲粟二千九百六十二斛二十七分斛之二十六。{\zhsmall 於徽術，當粟二千八百三十斛一千四百一十三分斛之一千二百一十。臣淳風等謹按：依密率，爲粟二千八百二十八斛九十九分斛之二十八。}\end{unitzh}
\begin{uniten}Suppose there is heaped millet (weisu) on level ground, lower circumference 12 zhang, height 2 zhang. What are the volume and the amount it makes as millet? Answer: Volume 8000 chi [Liu Hui: By Hui's method (Hui shu), the volume should be 7643 49/157 chi.] [Li Chunfeng et al.: We respectfully note: according to the precise rate (milü), the volume is 7636 4/11 chi.]; as millet, 2962 26/27 hu (斛). [Liu Hui: By Hui's method, the millet should be 2830 1210/1413 hu.] [Li Chunfeng et al.: We respectfully note: according to the precise rate, the millet is 2828 28/99 hu.]\end{uniten}

\begin{unitzh}\unitlabel{jiuzhang-021.4}今有委菽依垣，下周三丈，高七尺。問積及爲菽各幾何？荅曰：積三百五十尺，{\zhsmall 依徽術，當積三百三十四尺四百七十一分尺之一百八十六也。臣淳風等謹按：依密率，爲積三百三十四尺十一分尺之一。}爲菽一百四十四斛二百四十三分斛之八。{\zhsmall 依徽術，當菽一百三十七斛一萬二千七百一十七分斛之十千七百七十一。臣淳風等謹按：依密率，爲菽一百三十七斛八百九十一分斛之四百三十三。}\end{unitzh}
\begin{uniten}Suppose there are heaped beans leaning against a wall, lower circumference 3 zhang, height 7 chi. What are the volume and the amount it makes as beans, each? Answer: Volume 350 chi [Liu Hui: According to Hui's method, the volume should be 334 186/471 chi.] [Li Chunfeng et al.: We respectfully note: according to the precise rate, the volume is 334 1/11 chi.]; as beans, 144 8/243 hu. [Liu Hui: According to Hui's method, the beans should be 137 10771/12717 hu. [Translator's note: the numerator is written 十千七百七十一, literally 'ten thousand seven hundred seventy-one', which exceeds the denominator; the text appears corrupt.]] [Li Chunfeng et al.: We respectfully note: according to the precise rate, the beans are 137 433/891 hu.]\end{uniten}

\begin{unitzh}\unitlabel{jiuzhang-021.5}今有委米依垣内角，下周八尺，高五尺。問積及爲米幾何？荅曰：積三十五尺九分尺之五，{\zhsmall 於徽術，當積三十二尺四百七十一分尺之四百五十七。臣淳風等謹按：依密率，當積三十三尺三十三分尺之三十一。}爲米二十一斛七百二十九分斛之六百九十一。{\zhsmall 於徽術，當米二十斛三萬八千一百五十一分斛之三萬六千九百八十。臣淳風等謹按：依密率，爲米二十斛二千六百七十三分斛之二千五百四十。}\end{unitzh}
\begin{uniten}Suppose there is heaped hulled grain leaning against the inner corner of a wall, lower circumference 8 chi, height 5 chi. What are the volume and the amount it makes as hulled grain? Answer: Volume 35 5/9 chi [Liu Hui: By Hui's method, the volume should be 32 457/471 chi.] [Li Chunfeng et al.: We respectfully note: according to the precise rate, the volume should be 33 31/33 chi.]; as hulled grain, 21 691/729 hu. [Liu Hui: By Hui's method, the hulled grain should be 20 36980/38151 hu.] [Li Chunfeng et al.: We respectfully note: according to the precise rate, the hulled grain is 20 2540/2673 hu.]\end{uniten}

\begin{unitzh}\unitlabel{jiuzhang-021.6}委粟術曰：下周自乘，以高乘之，三十六而一。{\zhsmall 此猶圓錐也。於徽術，亦當下周自乘，以高乘之，又以二十五乘之，九百四十二而一也。}其依垣者，{\zhsmall 居圓錐之半也。}十八而一。{\zhsmall 於徽術，當令此下周自乘，以高乘之，又以二十五乘之，四百七十一而一。依垣之周半於全周，其自乘之羃居全周自乘之羃四分之一，故半全周之法以爲法也。}其依垣内角者，{\zhsmall 角，隅也。居圓錐四分之一也。}九而一。{\zhsmall 於徽術，當令此下周自乘而倍之，以高乘之，又以二十五乘之，四百七十一而一。依隅之周半於依垣，其自乘之羃居依垣自乘之羃四分之一，當半依垣之法以爲法。法不可半，故倍其實。}\end{unitzh}
\begin{uniten}Method for heaped millet: Multiply the lower circumference by itself, multiply by the height, and divide by 36. [Liu Hui: This is like a circular cone (yuanzhui). By Hui's method, one should likewise multiply the lower circumference by itself, multiply by the height, further multiply by 25, and divide by 942.] For that leaning against a wall, [Liu Hui: It occupies half of a circular cone.] divide by 18. [Liu Hui: By Hui's method, one should multiply this lower circumference by itself, multiply by the height, further multiply by 25, and divide by 471. The circumference of that leaning against a wall is half of the full circumference; the area (mi) of its self-multiplication occupies 1/4 of the area of the self-multiplication of the full circumference; therefore one halves the divisor (fa) for the full circumference to make the divisor.] For that leaning against the inner corner of a wall, [Liu Hui: Jiao (角, angle) is yu (隅, corner). It occupies 1/4 of a circular cone.] divide by 9. [Liu Hui: By Hui's method, one should multiply this lower circumference by itself and double it, multiply by the height, further multiply by 25, and divide by 471. The circumference of that leaning in a corner is half of that leaning against a wall; the area of its self-multiplication occupies 1/4 of the area of the self-multiplication for that leaning against a wall; one should halve the divisor for leaning against a wall to make the divisor. The divisor cannot be halved; therefore one doubles its dividend (shi).]\end{uniten}

\begin{unitzh}\unitlabel{jiuzhang-021.7}{\zhsmall 又此術亦用周三徑一之率。假令以三除周得徑，若不盡，通分内子，卽爲徑之積分。令自乘，以高乘之，爲三方錐之積。分母自相乘得九，爲法，又當三而一，約方錐之積。從方錐中求圓錐之積，亦猶方羃求圓羃，乃當三乘之，四而一，得圓羃之積。前求方積乃合三而一，今求圓錐之積復合三乘之，二母旣同，故相凖折。惟以四乘分母九，得三十六，而連除圓錐之積。其圓錐之積與平地聚粟同，故三十六而一。臣淳風等謹依密率，以七乘之，其平地者二百六十四而一，依垣者一百三十二而一，依隅者六十六而一也。}\end{unitzh}
\begin{uniten}[Liu Hui: Moreover, this method also uses the rate (lü) of circumference 3 to diameter 1. Suppose one divides the circumference by 3 to get the diameter; if it does not divide exactly, convert to parts and incorporate the numerator (tongfen nazi), and this is then the diameter in accumulated parts (jifen). Multiply it by itself and multiply by the height: this is the volume of 3 square cones (fangzhui). The denominator multiplied by itself gives 9, as the divisor; further one should divide by 3, reducing it to the volume of a square cone. Seeking the volume of the circular cone from within the square cone is like seeking the circular area (mi) from the square area (mi): one should then multiply by 3 and divide by 4, to obtain the area of the circular mi. Before, in seeking the square volume, one had to divide by 3; now, in seeking the volume of the circular cone, one must in turn multiply by 3; the two denominators (mu) being the same, they therefore offset each other. Only multiply the denominator 9 by 4, getting 36, and with it divide the volume of the circular cone in one continued division. The volume of this circular cone is the same as that of millet piled on level ground; therefore divide by 36.] [Li Chunfeng et al.: We respectfully, according to the precise rate, multiply it by 7: for that on level ground, divide by 264; for that leaning against a wall, divide by 132; for that leaning in a corner, divide by 66.]\end{uniten}

\begin{unitzh}\unitlabel{jiuzhang-021.8}程粟一斛，積二尺七寸；{\zhsmall 二尺七寸者，謂方一尺，深二尺七寸，凡積二千七百寸。}其米一斛，積一尺六寸五分寸之一；{\zhsmall 謂積一千六百二十寸。}其菽、荅、麻、麥一斛，皆二尺四寸十分寸之三。{\zhsmall 謂積二千四百三十寸。此爲以精麤爲率，而不等其槩也。粟率五，米率三，故米一斛於粟一斛五分斛之三。菽、荅、麻、麥亦如本率云。故謂此三量器爲槩，而皆不合於今斛。}\end{unitzh}
\begin{uniten}By the standard (cheng), 1 hu of millet has a volume of 2 chi 7 cun; [Liu Hui: '2 chi 7 cun' means a square of 1 chi with a depth of 2 chi 7 cun, a volume of 2700 cun in all.] 1 hu of hulled grain has a volume of 1 chi 6 1/5 cun; [Liu Hui: This means a volume of 1620 cun.] 1 hu of beans, small beans, hemp or wheat is in each case 2 chi 4 3/10 cun. [Liu Hui: This means a volume of 2430 cun. This takes fineness and coarseness as the rates and does not make their measuring standards (gai) equal. The rate of millet is 5 and the rate of hulled grain is 3; therefore 1 hu of hulled grain is, relative to 1 hu of millet, 3/5 hu. Beans, small beans, hemp and wheat are likewise according to their own rates. Therefore these three measuring vessels are called gai, and none of them agrees with the present hu.]\end{uniten}

\begin{unitzh}\unitlabel{jiuzhang-021.9}{\zhsmall 當今大司農斛，圓徑一尺三寸五分五釐，正深一尺，於徽術爲積一千四百四十一寸，排成餘分，又有十分寸之三。王莽銅斛，於今尺爲深九寸五分五釐，徑一尺三寸六分八釐二毫，以徽術計之，於今斛爲容九斗七升四合有奇。《周官·考工記》：「㮚氏爲量，深一尺，内方一尺而圓外，其實一鬴。」於徽術，此圓積一千五百七十六寸。《左氏傳》曰：「齊舊四量：豆、區、釜、鍾。四升曰豆，各自其四，以登於釜，釜十則鍾。」鍾六斛四斗，釜六斗四升。方一尺，深一尺，其積一千寸。若此方積容四斗二升，則通外圓積成旁，容十斗四合一龠五分龠之三也。以數相乘之，則斛之制：方一尺而圓其外，庣旁一釐七毫，羃一百五十六寸四分寸之一，深一尺，積一千五百六十二寸半，容十斗。王莽銅斛與《漢書·律歷志》所論斛同。}\end{unitzh}
\begin{uniten}[Liu Hui: The hu of the present Grand Minister of Agriculture (Da sinong) has a circular diameter of 1 chi 3 cun 5 fen (分) 5 li (釐) and a depth of exactly 1 chi; by Hui's method its volume is 1441 cun, and, setting out the remaining parts, there is a further 3/10 cun. Wang Mang's bronze hu, in present chi, is 9 cun 5 fen 5 li deep, with a diameter of 1 chi 3 cun 6 fen 8 li 2 hao (毫); reckoned by Hui's method, it holds, in present hu, 9 dou (斗) 7 sheng (升) 4 ge (合) and a little more. The Kaogong ji of the Zhou guan: 'The Lishi (㮚氏) make the measure: depth 1 chi, square 1 chi inside and round outside; its contents are 1 fu (鬴).' By Hui's method, the volume of this circle is 1576 cun. The Zuo shi zhuan says: 'The four old measures of Qi: dou (豆), ou (區), fu (釜), zhong (鍾). 4 sheng are called a dou; each [larger one] being 4 of its own [smaller], they rise to the fu; 10 fu make a zhong.' A zhong is 6 hu 4 dou; a fu is 6 dou 4 sheng. A square of 1 chi with a depth of 1 chi has a volume of 1000 cun. If this square volume holds 4 dou 2 sheng, then, taking in the volume of the outer circle that completes the sides, it holds 10 dou 4 ge 1 3/5 yue (龠). Multiplying the numbers together, the standard of the hu is then: a square of 1 chi with its outside made round, the margin (tiaopang) 1 li 7 hao, the area (mi) 156 1/4 cun, depth 1 chi, volume 1562 1/2 cun, holding 10 dou. Wang Mang's bronze hu is the same as the hu discussed in the Treatise on Pitch-pipes and the Calendar (Lüli zhi) of the Han shu.]\end{uniten}

\begin{unitzh}\unitlabel{jiuzhang-022.1}今有穿地，袤一丈六尺，深一丈，上廣六尺，爲垣積五百七十六尺。問穿地下廣幾何？荅曰：三尺五分尺之三。\end{unitzh}
\begin{uniten}Suppose there is dug earth (chuandi) with length (mao) 1 zhang (丈) 6 chi (尺), depth 1 zhang and upper breadth (guang) 6 chi, which makes a wall (yuan) with a volume of 576 chi. Question: what is the lower breadth of the dug earth? Answer: 3 3/5 chi.\end{uniten}

\begin{unitzh}\unitlabel{jiuzhang-022.2}術曰：置垣積尺，四之爲實。{\zhsmall 實：穿地四，爲堅三。垣，堅也。以堅求穿地，當四之，三而一也。}以深袤相乘，{\zhsmall 爲深袤之立實也。}又以三之爲法。{\zhsmall 以深袤乘之立實除垣積，則阬廣。又三之者，與堅率并除之。}所得倍之，{\zhsmall 阬有兩廣，先并而半之，卽爲廣狹之中平。今先得其中平，故又倍之，知兩廣全也。}減上廣，餘卽下廣。\end{unitzh}
\begin{uniten}Method: Set out the volume of the wall in chi and multiply it by 4 to make the dividend (shi). [Liu Hui: As for the dividend: dug earth 4 makes firm earth (jian) 3; the wall is firm earth. To find the dug earth from the firm earth, one should multiply by 4 and divide by 3.] Multiply the depth and the length together [Liu Hui: this makes the standing content (li shi) of depth and length], and further multiply it by 3 to make the divisor (fa). [Liu Hui: Dividing the volume of the wall by the standing content obtained by multiplying by the depth and length gives the breadth of the pit (keng). Further multiplying it by 3 is to divide it together with the rate (lü) of firm earth.] Double the result. [Liu Hui: The pit has two breadths; adding them first and halving gives the mean (zhongping) of the broad and the narrow. Now one has first obtained this mean; therefore one further doubles it, and knows the two breadths in full.] Subtract the upper breadth; the remainder is the lower breadth.\end{uniten}

\begin{unitzh}\unitlabel{jiuzhang-022.3}{\zhsmall 按此術：穿地四，爲堅三。垣，卽堅也。今以堅求穿地，當四乘之，三而一。深袤相乘者，爲深袤立羃。以深袤立羃除積，卽阬廣。又三之爲法，與堅率并除。所得倍之者，爲阬有兩廣，先并而半之，爲中平之廣。今此得中平之廣，故倍之，還爲兩廣并。故減上廣，餘卽下廣也。}\end{unitzh}
\begin{uniten}[Liu Hui: Examining this method: dug earth 4 makes firm earth 3; the wall is firm earth. Now, to find the dug earth from the firm earth, one should multiply it by 4 and divide by 3. Multiplying the depth and the length together makes the standing area (li mi) of depth and length; dividing the volume by the standing area of depth and length gives the breadth of the pit. Further multiplying it by 3 to make the divisor is to divide together with the rate of firm earth. Doubling the result is because the pit has two breadths: adding them first and halving makes the breadth of the mean. Now here one has obtained the breadth of the mean, so doubling it returns to the sum of the two breadths. Therefore, subtracting the upper breadth, the remainder is the lower breadth.]\end{uniten}

\begin{unitzh}\unitlabel{jiuzhang-022.4}今有倉，廣三丈，袤四丈五尺，容粟一萬斛。問高幾何？荅曰：二丈。\end{unitzh}
\begin{uniten}Suppose there is a granary (cang) with breadth 3 zhang and length 4 zhang 5 chi, holding 10,000 hu (斛) of millet (su). Question: what is the height? Answer: 2 zhang.\end{uniten}

\begin{unitzh}\unitlabel{jiuzhang-022.5}術曰：置粟一萬斛積尺爲實，廣袤相乘爲法，實如法而一，得高尺。{\zhsmall 以廣袤之羃除積，故得高。按此術：本以廣袤相乘，以高乘之，得此積。今還原，置此廣袤相乘爲法除之，故得高也。}\end{unitzh}
\begin{uniten}Method: Set out the volume in chi of the 10,000 hu of millet as the dividend; multiply the breadth and the length together as the divisor; divide the dividend by the divisor, obtaining the height in chi. [Liu Hui: Since the volume is divided by the area (mi) of the breadth and length, one obtains the height. Examining this method: originally one multiplies the breadth and the length together and multiplies by the height, obtaining this volume. Now, returning to the original, one sets out this product of the breadth and length as the divisor and divides by it; therefore one obtains the height.]\end{uniten}

\begin{unitzh}\unitlabel{jiuzhang-022.6}今有圓囷，{\zhsmall 圓囷，廩也，亦云圓囤也。}高一丈三尺三寸少半寸，容米二千斛。問周幾何？荅曰：五丈四尺。{\zhsmall 於徽術，當周五丈五尺二寸二十分寸之九。臣淳風等謹依密率，爲周五丈五尺一百分尺之二十七。}\end{unitzh}
\begin{uniten}Suppose there is a round granary (yuan qun) [Liu Hui: a round qun is a granary (lin); it is also called a round bin (yuan dun)] with height 1 zhang 3 chi 3 1/3 cun (寸), holding 2,000 hu of husked grain (米, mi). Question: what is the circumference? Answer: 5 zhang 4 chi. [Liu Hui: By Hui's method (Hui shu) the circumference should be 5 zhang 5 chi 2 9/20 cun.] [Li Chunfeng et al.: Your servants Chunfeng et al. respectfully follow the precise rate (milü): the circumference is 5 zhang 5 chi 27/100 chi.]\end{uniten}

\begin{unitzh}\unitlabel{jiuzhang-022.7}術曰：置米積尺，{\zhsmall 此卽猶圓堢壔之積。}以十二乘之，令高而一，所得開方除之，卽周。{\zhsmall 於徽術，當置米積尺，以三百一十四乘之爲實，二十五乘囷高爲法，所得開方除之，卽周也。此亦據見羃以求周，失之於微少也。}\end{unitzh}
\begin{uniten}Method: Set out the volume of the husked grain in chi [Liu Hui: this is just like the volume of a round fort (yuan baodao)], multiply it by 12, divide by the height, and extract the square root of the result; this is the circumference. [Liu Hui: By Hui's method, one should set out the volume of the husked grain in chi and multiply it by 314 to make the dividend, multiply the height of the granary by 25 to make the divisor, and extract the square root of the result; this is the circumference. This too seeks the circumference on the basis of the visible area (jian mi), and it errs by being slightly too small.]\end{uniten}

\begin{unitzh}\unitlabel{jiuzhang-022.8}{\zhsmall 晉武庫中有漢時王莽所作銅斛，其篆書字題斛旁云：「律嘉量斛，方一尺而圓其外，庣旁九釐五毫，羃一百六十二寸，深一尺，積一千六百二十寸，容十斗。」及斛底云：「律嘉量斗，方尺而圜其外，庣旁九釐五毫，羃一尺六寸二分，深一寸，積一百六十二寸，容一斗。」合龠皆有文字。升居斛旁，合龠在斛耳上。後有讃文，與今《律歷志》同，亦魏晉所常用。今祖疏王莽銅斛文字、尺寸分數，然不盡得升合勺之文字。}\end{unitzh}
\begin{uniten}[Liu Hui: In the Jin armory there was a bronze hu made by Wang Mang in the Han period. Its seal-script characters, inscribed on the side of the hu, say: 'Standard excellent-measure (lü jialiang) hu: a square of 1 chi with its outside made round; side margin (tiaopang) 9 li (釐) 5 hao (毫); area 162 cun; depth 1 chi; volume 1,620 cun; holds 10 dou (斗).' And on the base of the hu it says: 'Standard excellent-measure dou: a square chi with its outside made round; side margin 9 li 5 hao; area 1 chi 6 cun 2 fen (分); depth 1 cun; volume 162 cun; holds 1 dou.' The ge (合) and yue (龠) all have inscriptions. The sheng (升) is located on the side of the hu; the ge and yue are on the ears of the hu. After these there is a eulogy, the same as in the present Treatise on Pitch-pipes and the Calendar (Lü li zhi); it was also commonly used in Wei and Jin. Now Zu has itemized the inscriptions of Wang Mang's bronze hu and its numbers of chi, cun and fen, but did not obtain all the inscriptions of the sheng, ge and shao (勺).] [Translator's note: the side margin (tiaopang) is the amount by which the circle's radius exceeds half the diagonal of the square.]\end{uniten}

\begin{unitzh}\unitlabel{jiuzhang-022.9}{\zhsmall 按此術：本周自相乘，以高乘之，十二而一，得此積。今還元，置此積，以十二乘之，令高而一，卽復本周自乘之數。凡物自乘，開方除之，復其本。周自乘之數，故開方除之，卽得也。臣淳風等謹依密率，以八十八乘之爲實，七乘囷高爲法，實如法而一，開方除之，卽周也。}\end{unitzh}
\begin{uniten}[Liu Hui: Examining this method: originally the circumference is multiplied by itself, multiplied by the height and divided by 12, giving this volume. Now, returning to the original, one sets out this volume, multiplies it by 12 and divides by the height; this restores the number of the original circumference multiplied by itself. In general, when a thing is multiplied by itself, extracting the square root restores the original. It is the number of the circumference multiplied by itself; therefore, extracting the square root, one obtains it.] [Li Chunfeng et al.: Your servants Chunfeng et al. respectfully follow the precise rate: multiply it by 88 to make the dividend, multiply the height of the granary by 7 to make the divisor, divide the dividend by the divisor, and extract the square root; this is the circumference.]\end{uniten}

\begin{unitzh}\unitlabel{jiuzhang-022.10}九章算術卷五\end{unitzh}
\begin{uniten}The Nine Chapters on the Mathematical Art, Chapter 5\end{uniten}

\begin{unitzh}\unitlabel{jiuzhang-022.11}九章算術卷五訂訛{\zhsmall 算經十書之二}休寧 戴震 東原\end{unitzh}
\begin{uniten}The Nine Chapters on the Mathematical Art, Chapter 5: Corrections of Errors (ding'e) [Note: The Ten Mathematical Classics, no. 2]. Dai Zhen (Dongyuan) of Xiuning\end{uniten}

\begin{unitzh}\unitlabel{jiuzhang-022.12}今有壍，上廣一丈六尺三寸，下廣一丈，深六尺三寸，袤十三丈二尺一寸。問積幾何？荅曰：一萬九百四十三尺八寸。{\zhsmall 此積以立方一尺者計之，餘有方尺而深八寸二分四釐五毫，實八十二萬四千五百分也。}\end{unitzh}
\begin{uniten}Suppose there is a moat (qian) with upper breadth 1 zhang 6 chi 3 cun, lower breadth 1 zhang, depth 6 chi 3 cun and length 13 zhang 2 chi 1 cun. Question: what is the volume? Answer: 10,943 chi 8 cun. [Note: This volume is reckoned in cubes of 1 chi; there remains a square chi with a depth of 8 cun 2 fen 4 li 5 hao, whose actual content is 824,500 fen.]\end{uniten}

\begin{unitzh}\unitlabel{jiuzhang-022.13}今有穿渠，上廣一丈八尺，下廣三尺六寸，深一丈八尺，袤五萬一千八百二十四尺。問積幾何？荅曰：一千七萬四千五百八十五尺六寸。{\zhsmall 此六寸亦方尺而深六寸，實積六百寸也。}\end{unitzh}
\begin{uniten}Suppose there is a canal being dug (chuanqu) with upper breadth 1 zhang 8 chi, lower breadth 3 chi 6 cun, depth 1 zhang 8 chi and length 51,824 chi. Question: what is the volume? Answer: 10,074,585 chi 6 cun. [Note: These 6 cun are likewise a square chi with a depth of 6 cun; the actual volume is 600 cun.]\end{uniten}

\section*{卷第六}\addcontentsline{toc}{section}{卷第六}

\begin{unitzh}\unitlabel{jiuzhang-023.1}《九章算術》卷第六{\zhsmall 《算經十書》之二}魏劉徽注。唐朝議大夫、行太史令、上輕車都尉臣李淳風等奉勅注釋。\end{unitzh}
\begin{uniten}The Nine Chapters on the Mathematical Art, juan 6 [Note: The second of the Ten Mathematical Classics (Suanjing shishu).] Commentary by Liu Hui of Wei. Of the Tang: Grand Master for Court Discussion, Acting Grand Astrologer, Supreme Commandant of Light Chariots, Your servant Li Chunfeng et al., annotated by imperial command.\end{uniten}

\begin{unitzh}\unitlabel{jiuzhang-023.2}均輸{\zhsmall 以御遠近勞費}\end{unitzh}
\begin{uniten}Equitable Transport (junshu) [Liu Hui: To deal with labor and expense, far and near.]\end{uniten}

\begin{unitzh}\unitlabel{jiuzhang-023.3}今有均輸粟：甲縣一萬戸，行道八日；乙縣九千五百戸，行道十日；丙縣一萬二千三百五十戸，行道十三日；丁縣一萬二千二百戸，行道二十日，各到輸所。凡四縣賦當輸二十五萬斛，用車一萬乘。欲以道里遠近、戸數多少衰出之。問粟、車各幾何？荅曰：甲縣粟八萬三千一百斛，車三千三百二十四乘；乙縣粟六萬三千一百七十五斛，車二千五百二十七乘；丙縣粟六萬三千一百七十五斛，車二千五百二十七乘；丁縣粟四萬五百五十斛，車一千六百二十二乘。\end{unitzh}
\begin{uniten}Suppose there is an equitable transport of millet (su). County A has 10000 households and travels 8 days on the road; county B has 9500 households and travels 10 days on the road; county C has 12350 households and travels 13 days on the road; county D has 12200 households and travels 20 days on the road; each goes to the place of delivery. In all, the levy on the four counties is that they should deliver 250000 hu (斛), using 10000 carts. It is desired to have them provide these in proportional parts (cui 衰) according to the distance of the road, far or near, and the number of households, many or few. What are the millet and the carts of each? Answer: County A: millet 83100 hu, 3324 carts; county B: millet 63175 hu, 2527 carts; county C: millet 63175 hu, 2527 carts; county D: millet 40550 hu, 1622 carts.\end{uniten}

\begin{unitzh}\unitlabel{jiuzhang-023.4}均輸{\zhsmall 按：此均輸猶均運也。令戸率出車，以行道日數爲均，發粟爲輸。}術曰：令縣戸數各如其本行道日數而一，以爲衰。{\zhsmall 據甲行道八日，因使八戸共出一車；乙行道十日，因使十戸共出一車。計其在道，則皆戸一日出一車，故可爲均平之率也。}甲衰一百二十五，乙、丙衰各九十五，丁衰六十一，副并爲法。以賦粟車數乘未并者，各自爲實。{\zhsmall 衰分科率。}實如法得一車。{\zhsmall 各置所當出車，以其行道日數乘之，如戸數而一，得率，戸用車二百四十七分日之三十一，故謂之均。求此率以戸，當各計車之衰分也。}\end{unitzh}
\begin{uniten}Equitable transport [Liu Hui: Note that this "equitable transport" (junshu) is like "equitable conveyance" (junyun). Let the households provide carts by a rate (lü), taking the number of days of travel on the road as the "equitable" (jun) and the dispatch of millet as the "transport" (shu).] Method: Let the number of households of each county be divided by its own number of days of travel on the road, and take the results as the proportional parts. [Liu Hui: Based on A's travelling 8 days on the road, one accordingly makes 8 households jointly provide 1 cart; B travelling 10 days on the road, one accordingly makes 10 households jointly provide 1 cart. Reckoning their time on the road, it is in every case 1 cart for 1 day per household; therefore this can be taken as the rate of equalization (junping).] The proportional part of A is 125, those of B and C are 95 each, and that of D is 61; add them together on the side (fu) as the divisor (fa). Multiply the unadded ones by the number of carts for the levied millet, each separately, as the dividend (shi). [Liu Hui: Graded rates (kelü) of proportional distribution (cuifen).] Divide the dividend by the divisor to obtain the carts. [Liu Hui: Set out for each the carts it should provide, multiply them by its number of days of travel on the road, and divide by the number of households: one obtains the rate, each household using a cart for 31/247 of a day [Translator's note: the computation gives 2 31/47 days; 百 is evidently an error for 日]. Therefore it is called "equitable". Seeking this rate by means of the households, one should compute for each the proportional distribution of the carts.]\end{uniten}

\begin{unitzh}\unitlabel{jiuzhang-023.5}{\zhsmall 臣淳風等謹按：縣戸有多少之差，行道有遠近之異，欲其均等，故各令行道日數約戸爲衰。行道多者少其戸，行道少者多其戸，故各令約戸爲衰。以八日約除甲縣，得一百二十五；乙、丙各九十五；丁六十一。於今有術，副并爲所有率，未并者各爲所求率，以賦粟車數爲所有數，而今有之，各得車數。一旬除乙，十三除丙，各得九十五；二旬除丁，得六十一也。}\end{unitzh}
\begin{uniten}[Li Chunfeng et al.: Your servants Chunfeng et al. respectfully note: the households of the counties differ in being many or few, and the travel on the road differs in being far or near. Wishing to make them equal, one therefore lets each county's number of days of travel on the road reduce (yue) its households, making the proportional parts. For those with more travel, their households are made fewer; for those with less travel, their households are made more; therefore one lets each reduce its households, making the proportional parts. Reducing county A by division by 8 days gives 125; B and C, 95 each; D, 61. In the "Suppose there is" method (jinyou shu), the ones added together on the side are taken as the given rate (suoyou lü), the unadded ones each as the sought rate (suoqiu lü), and the number of carts for the levied millet as the given number (suoyou shu); applying "Suppose there is" to them, one obtains the number of carts for each. Dividing B by one ten-day period (xun) and C by 13, each gives 95; dividing D by two ten-day periods gives 61.]\end{uniten}

\begin{unitzh}\unitlabel{jiuzhang-023.6}有分者上下軰之。{\zhsmall 軰，配也。車牛人之數不可分裂，推少就多，均賦之宜。今按：甲分旣少，宜從於乙，滿法除之，有餘從丙；丁分又少，亦宜就丙，除之適盡。加乙、丙各一，上下軰益，以少從多也。}以二十五斛乘車數，卽粟數。\end{unitzh}
\begin{uniten}Where there are fractions, match them (bei) above and below. [Liu Hui: Bei (軰) means "to pair" (pei). The numbers of carts, oxen and men cannot be split up; pushing the smaller to go to the larger is what is fitting for an equitable levy. Now consider: since A's fraction is small, it should go over to B; when it fills the divisor, divide it out, and the remainder goes over to C. D's fraction is also small; it too should go to C, and on dividing, it is exactly used up. Add 1 each to B and C. Matching and augmenting above and below is letting the smaller go over to the larger.] Multiply the number of carts by 25 hu: this is the amount of millet.\end{uniten}

\begin{unitzh}\unitlabel{jiuzhang-023.7}今有均輸卒：甲縣一千二百人，薄塞；乙縣一千五百五十人，行道一日；丙縣一千二百八十人，行道二日；丁縣九百九十人，行道三日；戊縣一千七百五十人，行道五日。凡五縣賦輸卒一月一千二百人，欲以遠近、戸率多少衰出之。問縣各幾何？荅曰：甲縣二百二十九人；乙縣二百八十六人；丙縣二百二十八人；丁縣一百七十一人；戊縣二百八十六人。\end{unitzh}
\begin{uniten}Suppose there is an equitable transport of soldiers (zu). County A has 1200 men and adjoins the frontier; county B has 1550 men and travels 1 day on the road; county C has 1280 men and travels 2 days on the road; county D has 990 men and travels 3 days on the road; county E has 1750 men and travels 5 days on the road. In all, the levy on the five counties is to deliver soldiers for one month, 1200 men. It is desired to have them provide these in proportional parts according to distance, far or near, and the household rate, many or few. What is the number for each county? Answer: County A: 229 men; county B: 286 men; county C: 228 men; county D: 171 men; county E: 286 men.\end{uniten}

\begin{unitzh}\unitlabel{jiuzhang-023.8}術曰：令縣卒各如其居所及行道日數而一，以爲衰。{\zhsmall 按：此亦以日數爲均，居所爲輸。甲無行道日，但以居所三十日爲率。言欲爲均平之率者，當使甲三十人而出一人，乙三十一人而出一人。出一人者，計役則皆一人一日，是以可爲均平之率。}甲衰四，乙衰五，丙衰四，丁衰三，戊衰五，副并爲法。以人數乘未并者，各自爲實。實如法而一。{\zhsmall 各置所當出人數，以其居所及行道日數乘之，如縣人數而一，得戸率，人役五日七分日之五。}{\zhsmall 臣淳風等謹按：爲衰，於今有術，副并爲所有率，未并者各爲所求率，以賦卒人數爲所有數。此術以别考，則意同，以廣異聞，故存之也。}\end{unitzh}
\begin{uniten}Method: Let the soldiers of each county be divided by its number of days at the place of stay (juso) and on the road, and take the results as the proportional parts. [Liu Hui: Note that here too the numbers of days are taken as the "equitable", and the place of stay as the "transport". A has no days of travel on the road; one takes only the 30 days at the place of stay as its rate. This says that to make a rate of equalization one should make A provide 1 man out of 30 men, and B provide 1 man out of 31 men. As for providing 1 man: reckoning the service, it is in every case 1 man for 1 day; therefore this can be taken as the rate of equalization.] The proportional part of A is 4, of B 5, of C 4, of D 3, of E 5; add them together on the side as the divisor. Multiply the unadded ones by the number of men, each separately, as the dividend. Divide the dividend by the divisor. [Liu Hui: Set out for each the number of men it should provide, multiply it by its number of days at the place of stay and on the road, and divide by the number of men of the county: one obtains the household rate, each man serving 5 5/7 days.] [Li Chunfeng et al.: Your servants Chunfeng et al. respectfully note: make the proportional parts; in the "Suppose there is" method, the ones added together on the side are taken as the given rate, the unadded ones each as the sought rate, and the number of levied soldiers as the given number. When this method is examined by another way, its idea is the same; to broaden acquaintance with different approaches, we therefore preserve it.]\end{uniten}

\begin{unitzh}\unitlabel{jiuzhang-023.9}有分者上下軰之。{\zhsmall 軰，配也。今按：丁分最少，宜就戊除；不從乙者，丁近戊故也。滿法除之，有餘從乙。丙分又少，亦就乙除，有餘從甲，除之適盡。從甲、丙二分，其數正等，二者於乙遠近皆同，不以甲從乙者，方以下從上也。}\end{unitzh}
\begin{uniten}Where there are fractions, match them above and below. [Liu Hui: Bei means "to pair". Now consider: D's fraction is the smallest; it should go to E to be divided. That it does not go over to B is because D is near E. When it fills the divisor, divide it out, and the remainder goes over to B. C's fraction is also small; it too goes to B to be divided, and the remainder goes over to A; on dividing, it is exactly used up. As for going over, the two fractions of A and C are exactly equal in amount, and the two are equally far or near with respect to B; not making A go over to B is precisely letting the lower go over to the upper.]\end{uniten}

\begin{unitzh}\unitlabel{jiuzhang-023.10}今有均賦粟：甲縣二萬五百二十戸，粟一斛二十錢，自輸其縣；乙縣一萬二千三百一十二戸，粟一斛十錢，至輸所二百里；丙縣七千一百八十二戸，粟一斛十二錢，至輸所一百五十里；丁縣一萬三千三百三十八戸，粟一斛十七錢，至輸所二百五十里；戊縣五千一百三十戸，粟一斛十三錢，至輸所一百五十里。凡五縣賦輸粟一萬斛。一車載二十五斛，與僦一里一錢。欲以縣戸賦粟，令費勞等。問縣各粟幾何？荅曰：甲縣三千五百七十一斛二千八百七十三分斛之五百一十七；乙縣二千三百八十斛二千八百七十三分斛之二千二百六十；丙縣一千三百八十八斛二千八百七十三分斛之二千二百七十六；丁縣一千七百一十九斛二千八百七十三分斛之一千三百一十三；戊縣九百三十九斛二千八百七十三分斛之二千二百五十三。\end{unitzh}
\begin{uniten}Suppose there is an equitable levy (junfu) of millet. County A has 20520 households; 1 hu of millet is 20 qian (錢); it delivers within its own county. County B has 12312 households; 1 hu of millet is 10 qian; it is 200 li (里) to the place of delivery. County C has 7182 households; 1 hu of millet is 12 qian; it is 150 li to the place of delivery. County D has 13338 households; 1 hu of millet is 17 qian; it is 250 li to the place of delivery. County E has 5130 households; 1 hu of millet is 13 qian; it is 150 li to the place of delivery. In all, the levy on the five counties is to deliver 10000 hu of millet. One cart carries 25 hu, and the hire (jiu) paid is 1 qian per li. It is desired to levy the millet by the households of the counties so that expense and labor are equal. What is the millet of each county? Answer: County A: 3571 517/2873 hu; county B: 2380 2260/2873 hu; county C: 1388 2276/2873 hu; county D: 1719 1313/2873 hu; county E: 939 2253/2873 hu.\end{uniten}

\begin{unitzh}\unitlabel{jiuzhang-023.11}術曰：以一里僦價乘至輸所里，{\zhsmall 此以出錢爲均也。問者曰「一車載二十五斛，與僦一里一錢」，一錢卽一里僦價也。以乘里數者，欲知僦一車到輸所所用錢也。甲自輸其縣，則無取僦價也。}以一車二十五斛除之，{\zhsmall 欲知僦一斛所用錢。}加一斛粟價，則致一斛之費。{\zhsmall 加一斛之價於一斛僦直，卽凡餘粟取僦錢也。甲一斛之費二十，乙、丙各十八，丁二十七，戊十九也。}各以約其戸數，爲衰。{\zhsmall 言使甲二十戸共出一斛，乙、丙十八戸共出一斛，計其所費，則皆戸一錢，故可爲均賦之率也。}甲衰一千二十六，乙衰六百八十四，丙衰三百九十九，丁衰四百九十四，戊衰二百七十，副并爲法。所賦粟乘未并者，各自爲實。實如法而一。{\zhsmall 各置所當出粟，以其十斛之費乘之，如戸數而一，得率，戸出三錢二千八百七十三分錢之一千三百八十一。}\end{unitzh}
\begin{uniten}Method: Multiply the li to the place of delivery by the hire price (jiujia) for 1 li. [Liu Hui: Here the paying out of cash is taken as the "equitable". The question says: "One cart carries 25 hu, and the hire paid is 1 qian per li"; the 1 qian is the hire price for 1 li. Multiplying the number of li by it is in order to know the cash used in hiring 1 cart to reach the place of delivery. A delivers within its own county, so it takes no hire price.] Divide it by the 25 hu of one cart. [Liu Hui: In order to know the cash used in hiring for 1 hu.] Add the price of 1 hu of millet: this is the expense of delivering 1 hu. [Liu Hui: Adding the price of 1 hu to the hire charge (jiuzhi) for 1 hu: this is the case for all the other millet, which takes hire cash. The expense for 1 hu of A is 20; of B and C, 18 each; of D, 27; of E, 19.] With each, reduce its number of households, making the proportional parts. [Liu Hui: This says: make 20 households of A jointly provide 1 hu, and 18 households of B or C jointly provide 1 hu; reckoning what they spend, it is in every case 1 qian per household; therefore it can be taken as the rate of equitable levy.] The proportional part of A is 1026, of B 684, of C 399, of D 494, of E 270; add them together on the side as the divisor. Multiply the unadded ones by the levied millet, each separately, as the dividend. Divide the dividend by the divisor. [Liu Hui: Set out for each the millet it should provide, multiply it by its expense for 10 hu [Translator's note: the stated result requires the expense for 1 hu], and divide by the number of households: one obtains the rate, each household paying 3 1381/2873 qian.]\end{uniten}

\begin{unitzh}\unitlabel{jiuzhang-023.12}{\zhsmall 臣淳風等謹按：此以出錢爲均。問者曰「一車載二十五斛，與僦一里一錢」，一錢卽一里僦價也。以乘里數者，欲知僦一車到輸所用錢。甲自出其縣，則無取僦之價。以一車二十五斛除之者，欲知僦一斛所用錢。加一斛之價於一斛僦直，卽凡餘粟取僦錢。甲一斛之費二十，乙、丙各十八，丁二十七，戊十九。各以約其戸爲衰：甲衰一千二十六，乙衰六百八十四，丙衰三百九十九，丁衰四百九十四，戊衰二百七十。言使甲二十戸共出一斛，乙、丙十八戸共出一斛，計其所費，則皆戸一錢，故可爲均賦之率也。於今有術，副并爲所有率，未并者各爲所求率，賦粟一萬斛爲所有數。此今有、衰分之義也。}\end{unitzh}
\begin{uniten}[Li Chunfeng et al.: Your servants Chunfeng et al. respectfully note: here the paying out of cash is taken as the "equitable". The question says: "One cart carries 25 hu, and the hire paid is 1 qian per li"; the 1 qian is the hire price for 1 li. Multiplying the number of li by it is in order to know the cash used in hiring 1 cart to reach the place of delivery. A provides within its own county, so it takes no hire price. Dividing by the 25 hu of one cart is in order to know the cash used in hiring for 1 hu. Adding the price of 1 hu to the hire charge for 1 hu: this is the case for all the other millet, which takes hire cash. The expense for 1 hu of A is 20; of B and C, 18 each; of D, 27; of E, 19. With each, reduce its households, making the proportional parts: the proportional part of A is 1026, of B 684, of C 399, of D 494, of E 270. This says: make 20 households of A jointly provide 1 hu, and 18 households of B or C jointly provide 1 hu; reckoning what they spend, it is in every case 1 qian per household; therefore it can be taken as the rate of equitable levy. In the "Suppose there is" method, the ones added together on the side are taken as the given rate, the unadded ones each as the sought rate, and the levied millet, 10000 hu, as the given number. This is the meaning of "Suppose there is" and of proportional distribution.]\end{uniten}

\begin{unitzh}\unitlabel{jiuzhang-023.13}{\zhsmall 計經賦之率，旣有戸算之率，亦有遠近貴賤之率。此二率者，各自相與通。通則甲二十，乙十二，丙七，丁十三，戊五。一斛之費爲之錢率。錢率約戸率者，則錢爲母，戸爲子。子不齊，令母互乘爲齊，卽衰也。若其不然，以一斛之費約戸數取衰，竝有分，當通分内子約之，於算甚䌓。此一章皆相與通功共率，畧相依似。以上二率、下一率亦可放此，從其簡易而已。又以分言之，使甲一戸出二十分斛之一，乙一戸出十八分斛之一，各以戸數乘之，亦可得一縣凡所當輸，俱爲衰也。乘之者，乘其子，母報除之。以此觀之，則以一斛之費約戸數者，其意不異矣。然則可置一斛之費而反衰之，約戸以乘戸率爲衰也。合分注曰：「母除爲率，率乘子爲齊。」反衰注曰：「先同其母，各以分母約其子，爲反衰。」以施其率，爲算旣約，且不妨處下也。}\end{unitzh}
\begin{uniten}[Li Chunfeng et al. (continued): Reckoning the rates of the levy in the Classic (jing), there is a rate reckoned by households (husuan), and there is also a rate of distance, far or near, and of price, dear or cheap. These two rates are each made to communicate (tong) among themselves. When made to communicate, A is 20, B 12, C 7, D 13, E 5. The expense of 1 hu is taken as their cash rate. When the cash rate reduces the household rate, the cash is the denominator (mu) and the households are the numerator (zi). The numerators are not homogenized (qi); let the denominators multiply them crosswise (hucheng) to homogenize them: these are the proportional parts. If one does not do so, but reduces the numbers of households by the expense of 1 hu to obtain the proportional parts, they all have fractions, and one must convert into parts, include the numerators (tongfen nazi) and reduce them, which is very cumbersome in the computation. In this one chapter all make their efforts communicate and share rates with one another (tonggong gonglü), and roughly resemble one another. The two rates above and the one rate below can also imitate this, simply following what is simple and easy. Again, speaking in terms of fractions: let one household of A provide 1/20 hu, and one household of B provide 1/18 hu; multiplying each by the number of households, one can likewise obtain all that one county should deliver, and all of these serve as proportional parts. As for multiplying them: one multiplies their numerators, and the denominators divide them in return (baochu). Seen in this way, reducing the number of households by the expense of 1 hu does not differ in meaning. That being so, one may set out the expense of 1 hu and take its inverse proportions (fancui), reducing the households by multiplying the household rates, to make the proportional parts. The commentary on combining fractions (hefen) says: "The denominators divide, making the rates; the rates multiply the numerators, making them homogenized." The commentary on inverse proportion (fancui) says: "First equalize (tong) their denominators; reduce each numerator by its denominator, making the inverse proportions." In applying these rates, the computation is simplified, and moreover it does no harm to the placement below.]\end{uniten}

\begin{unitzh}\unitlabel{jiuzhang-024.1}今有均賦粟：甲縣四萬二千算，粟一斛二十，傭價一日一錢，自輸其縣；乙縣三萬四千二百七十二算，粟一斛十八，傭價一日十錢，到輸所七十里；丙縣一萬九千三百二十八算，粟一斛十六，傭價一日五錢，到輸所一百四十里；丁縣一萬七千七百算，粟一斛十四，傭價一日五錢，到輸所一百七十五里；戊縣二萬三千四十算，粟一斛十二，傭價一日五錢，到輸所二百一十里；已縣一萬九千一百三十六算，粟一斛一十，傭價一日五錢，到輸所二百八十里。凡六縣賦粟六萬斛，皆輸甲縣。六人共車，車載二十五斛，重車日行五十里，空車日行七十里，載輸之間各一日。粟有貴賤，傭各别價，以算出錢，令費勞等。問縣各粟幾何？荅曰：甲縣一萬八千九百四十七斛一百三十三分斛之四十九；乙縣一萬八百二十七斛一百三十三分斛之九；丙縣七千二百一十八斛一百三十三分斛之六；丁縣六千七百六十六斛一百三十三分斛之一百二十二；戊縣九千二十二斛一百三十三分斛之七十四；己縣七千二百一十八斛一百三十三分斛之六。\end{unitzh}
\begin{uniten}Suppose there is an equitable levy (junfu) of millet (su). County A: 42000 tax units (suan); millet is 20 per hu (斛); the hire price is 1 qian (錢) a day; it delivers in its own county. County B: 34272 tax units; millet is 18 per hu; the hire price is 10 qian a day; to the delivery place is 70 li (里). County C: 19328 tax units; millet is 16 per hu; the hire price is 5 qian a day; to the delivery place is 140 li. County D: 17700 tax units; millet is 14 per hu; the hire price is 5 qian a day; to the delivery place is 175 li. County E: 23040 tax units; millet is 12 per hu; the hire price is 5 qian a day; to the delivery place is 210 li. County F: 19136 tax units; millet is 10 per hu; the hire price is 5 qian a day; to the delivery place is 280 li. Altogether the six counties are levied 60000 hu of millet, all to be delivered to County A. 6 persons share a cart; a cart carries 25 hu; a loaded cart travels 50 li a day; an empty cart travels 70 li a day; loading and delivering take 1 day each. Millet is dear or cheap, and hire has a different price in each; money is paid out by tax unit, so that the expense and labour are equal. Question: how much millet for each county?
Answer: County A, 18947 49/133 hu; County B, 10827 9/133 hu; County C, 7218 6/133 hu; County D, 6766 122/133 hu; County E, 9022 74/133 hu; County F, 7218 6/133 hu.\end{uniten}

\begin{unitzh}\unitlabel{jiuzhang-024.2}術曰：以車程行空、重相乘爲法，并空、重以乘道里，各自爲實，實如法得一日。{\zhsmall 臣淳風等謹按：此術重往空還，一輸再行道也。置空行一里用七十分日之一，重行一里用五十分日之一，齊而同之，空、重行一里之路，往返用一百七十五分日之六。定言之者，一百七十五里之路，往返用六日也。故并空、重者，齊其子也；空、重相乘者，同其母也。於今有術，至輸所里爲所有數，六爲所求率，齊一百七十五爲所有率，而今有之，卽各得輸所用日也。}\end{unitzh}
\begin{uniten}Method: Take the cart's standard daily travel (chengxing), empty and loaded, and multiply them together to make the divisor (fa). Add the empty and loaded and with this multiply the road distance (daoli), each separately making the dividend (shi). Divide the dividend by the divisor to get [the number of] days. [Li Chunfeng et al.: Your servants respectfully note: in this method [the cart] goes loaded and returns empty, so one delivery travels the road twice. Set up: travelling 1 li empty uses 1/70 day; travelling 1 li loaded uses 1/50 day. Equalize (qi) and homogenize (tong) them: for a road of 1 li travelled empty and loaded, going and returning uses 6/175 day. Stated in definite terms, a road of 175 li, going and returning, uses 6 days. Hence 'adding the empty and loaded' is equalizing their numerators; 'multiplying the empty and loaded together' is homogenizing their denominators. In the 'Suppose there is' rule (jinyou shu), the li to the delivery place are the given number (suoyou shu), 6 is the sought rate (suoqiu lü), the equalized 175 is the given rate (suoyou lü). Applying jinyou to them, one obtains for each the days used for [the road to] the delivery place.]\end{uniten}

\begin{unitzh}\unitlabel{jiuzhang-024.3}加載輸各一日，{\zhsmall 欲得幾日也。}而以六人乘之，{\zhsmall 欲知致一車用人也。}又以傭價乘之，{\zhsmall 欲知致車人傭直幾錢。}以二十五斛除之，{\zhsmall 欲知致一斛之傭直也。}加一斛粟價，卽致一斛之費。{\zhsmall 加一斛之價於致一斛之傭直，卽凡輸一斛餘粟取傭所用錢。}各以約其算數爲衰。{\zhsmall 今按：甲衰四十二，乙衰二十四，丙衰十六，丁衰十五，戊衰二十，已衰十六。於今有術，副并爲所有率，未并者各自爲所求率，所賦粟爲所有數，此今有衰分之義也。}副并爲法，以所賦粟乘未并者，各自爲實，實如法得一斛。{\zhsmall 各置所當出粟，以其一斛之費乘之，如算數而一，得率，算出九錢一百三十三分錢之三。又載輸之間各一日者，卽二日也。}\end{unitzh}
\begin{uniten}Add 1 day each for loading and delivering, [Liu Hui: One wants to obtain how many days.] and multiply by the 6 persons; [Liu Hui: One wants to know the persons used in conveying one cart.] then multiply by the hire price; [Liu Hui: One wants to know how many qian the hire value of the persons conveying the cart is.] divide by 25 hu; [Liu Hui: One wants to know the hire value of conveying 1 hu.] add the price of 1 hu of millet: this is the cost of conveying 1 hu. [Liu Hui: Add the price of 1 hu to the hire value of conveying 1 hu: this is the money used in all in delivering 1 hu, for the millet and, besides it, for taking hire.] With each [cost], divide its number of tax units to make the proportions (cui). [Liu Hui: Now note: A's proportion is 42, B's proportion 24, C's proportion 16, D's proportion 15, E's proportion 20, F's proportion 16. In the jinyou rule, the side sum (fubing) is the given rate, the unsummed ones each separately are the sought rates, and the levied millet is the given number. This is the meaning of jinyou [and] proportional distribution (cuifen).] Sum them on the side to make the divisor; multiply the unsummed ones by the levied millet, each separately making the dividend; divide the dividend by the divisor to get [the result in] hu. [Liu Hui: Set up for each the millet it should pay out, multiply it by its cost of 1 hu, and divide by the number of tax units: one obtains the rate, each tax unit paying out 9 3/133 qian. Also, 'loading and delivering 1 day each' means 2 days.]\end{uniten}

\begin{unitzh}\unitlabel{jiuzhang-024.4}今有粟七斗，三人分春之，一人爲糲米，一人爲粺米，一人爲糳米，令米數等。問取粟爲米各幾何？荅曰：糲米取粟二斗一百二十一分斗之一十；粺米取粟二斗一百二十一分斗之三十八；糳米取粟二斗一百二十一分斗之七十三；爲米各一斗六百五分斗之一百五十一。\end{unitzh}
\begin{uniten}Suppose there is 7 dou (斗) of millet; 3 persons share in husking it [Translator's note: 春 'spring' is written for 舂 'to husk'.]: one person makes coarse rice (limi), one person makes polished rice (baimi), one person makes fine rice (zuomi), so that the amounts of rice are equal. Question: how much millet is taken and how much rice is made by each?
Answer: Coarse rice takes 2 10/121 dou of millet; polished rice takes 2 38/121 dou of millet; fine rice takes 2 73/121 dou of millet; the rice made is 1 151/605 dou each.\end{uniten}

\begin{unitzh}\unitlabel{jiuzhang-024.5}術曰：列置糲米三十，粺米二十七，糳米二十四，而反衰之。{\zhsmall 此先約三率：糲爲十，粺爲九，糳爲八。欲令米等者，其取粟糲率十分之一，粺率九分之一，糳率八分之一，當齊其子，故曰反衰也。臣淳風等謹按：米有精麤之異，有多少之差。據率，粺、糳少而糲多；用粟，則粺、糳多而糲少。米若依本率之分，粟當倍率，故今反衰之，使精取多而麤得少。}副并爲法，以七斗乘未并者，各自爲取粟實，實如法得一斗。{\zhsmall 於今有術，副并爲所有率，未并者各爲所求率，粟七斗爲所有數，而今有之，故各得取粟也。}若求米等者，以本率各乘定所取粟爲實，以粟率五十爲法，實如法得一斗。{\zhsmall 若徑求爲米等數者，置糲米三用粟五，粺米二十七用粟五十，糳米十二用粟二十五，齊其粟，同其米，并齊爲法，以七斗乘同爲實，所得卽爲米斗數。}\end{unitzh}
\begin{uniten}Method: Lay out coarse rice 30, polished rice 27, fine rice 24, and take them in inverse proportion (fancui). [Liu Hui: Here one first reduces the three rates: coarse becomes 10, polished 9, fine 8. If one wants the rice to be equal, then for the millet taken the coarse rate is 1/10, the polished rate 1/9, the fine rate 1/8; one should equalize their numerators; hence it is called 'inverse proportion'.] [Li Chunfeng et al.: Your servants respectfully note: rice has the difference of refined and coarse, and the disparity of more and less. According to the rates, polished and fine are less and coarse is more; in the millet used, polished and fine are more and coarse is less. If the rice is to follow the division by the original rates, the millet should go against (bei 倍) the rates; hence now one takes them in inverse proportion, so that the refined takes more and the coarse gets less.] Sum them on the side to make the divisor; multiply the unsummed ones by 7 dou, each separately making the dividend for the millet taken; divide the dividend by the divisor to get [the result in] dou. [Liu Hui: In the jinyou rule, the side sum is the given rate, the unsummed ones are each the sought rates, the millet 7 dou is the given number, and one applies jinyou to them; hence one obtains for each the millet taken.] If one seeks the equal [amount of] rice, multiply each determined [amount of] millet taken by its original rate to make the dividend, take the millet rate 50 as the divisor, and divide the dividend by the divisor to get [the result in] dou. [Liu Hui: If one seeks directly the equal amount of rice made: set up coarse rice 3 using millet 5, polished rice 27 using millet 50, fine rice 12 using millet 25; equalize their millet and homogenize their rice; sum the equalized [terms] to make the divisor; multiply the homogenized [term] by 7 dou to make the dividend; what is obtained is the number of dou of rice.]\end{uniten}

\begin{unitzh}\unitlabel{jiuzhang-024.6}今有人當稟粟二斛，倉無粟，欲與米一、菽二以當所稟粟。問各幾何？荅曰：米五斗一升七分升之三；菽一斛二升七分升之六。\end{unitzh}
\begin{uniten}Suppose there is a person who should receive a ration (bing) of 2 hu of millet. The granary has no millet, and one wants to give rice 1 [and] beans (shu) 2 in place of the millet to be received. Question: how much of each?
Answer: Rice, 5 dou 1 3/7 sheng (升); beans, 1 hu 2 6/7 sheng.\end{uniten}

\begin{unitzh}\unitlabel{jiuzhang-024.7}術曰：置米一、菽二，求爲粟之數，并之得三九分之八，以爲法。亦置米一、菽二，而以粟二斛乘之，各自爲實，實如法得一斛。{\zhsmall 臣淳風等謹按：置粟率五乘米一，米率三除之，得一三分之二，卽是米一之粟也；粟率十以乘菽二，菽率九除之，得二九分之二，卽是菽二之粟也。并全得三，齊子并之得二十四，同母得二十七，約之得九分之八，故云并之得三九分之八。米一、菽二當粟三九分之八，此其粟率也。於今有術，米一、菽二皆爲所求率，當粟三九分之八爲所有率，粟二斛爲所有數。凡言率者當相與通之，則爲米九、菽十八當粟三十五也。亦有置米一、菽二，求其爲粟之率，以爲列衰，副并爲法，以粟乘列衰爲實，所得卽米一、菽二所求粟也。以米、菽本率而今有之，卽合所問。}\end{unitzh}
\begin{uniten}Method: Set up rice 1, beans 2; find the amounts of millet they make; sum them to get 3 8/9, and take this as the divisor. Also set up rice 1, beans 2, and multiply them by the millet 2 hu, each separately making the dividend; divide the dividend by the divisor to get [the result in] hu. [Li Chunfeng et al.: Your servants respectfully note: set up the millet rate 5 and multiply rice 1 by it, divide by the rice rate 3, getting 1 2/3: this is the millet of rice 1. Multiply beans 2 by the millet rate 10, divide by the bean rate 9, getting 2 2/9: this is the millet of beans 2. Summing the integers gives 3; equalizing the numerators and summing them gives 24; homogenizing the denominators gives 27; reducing gives 8/9; hence it says 'sum them to get 3 8/9'. Rice 1 and beans 2 are equivalent to millet 3 8/9: this is their millet rate. In the jinyou rule, rice 1 and beans 2 are both sought rates, the equivalent millet 3 8/9 is the given rate, and the millet 2 hu is the given number. Whenever one speaks of rates, they should be made commensurable (tong 通) with one another: then they become rice 9, beans 18, equivalent to millet 35. One may also set up rice 1, beans 2, find their rates as millet, and take these as the arrayed proportions (liecui); sum them on the side to make the divisor; multiply the arrayed proportions by the millet to make the dividend; what is obtained is the millet sought for rice 1 and beans 2. Applying jinyou to it with the original rates of rice and beans then matches what is asked.]\end{uniten}

\begin{unitzh}\unitlabel{jiuzhang-025.1}今有取傭負鹽二斛，行一百里，與錢四十。今負鹽一斛七斗三升少半升，行八十里，問與錢幾何？荅曰：二十七錢十五分錢之十一。\end{unitzh}
\begin{uniten}Suppose there is [a man] who takes hired work carrying 2 hu (斛) of salt; he travels 100 li (里) and is given 40 qian (錢). Now [he] carries 1 hu 7 dou (斗) 3 1/3 sheng (升) of salt and travels 80 li. Question: how many qian is he given? Answer: 27 11/15 qian.\end{uniten}

\begin{unitzh}\unitlabel{jiuzhang-025.2}術曰：置鹽二斛升數，以一百里乘之，爲法。{\zhsmall 按：此術以負鹽二斛升數乘所行一百里，得二萬里，是爲負鹽一升行二萬里，得錢四十，於今有術爲所有率。升數乘所行里爲法，於今有術爲所有數也。}以四十錢乘今負鹽升數，又以八十里乘之，爲實。實如法得一錢。{\zhsmall 以今負鹽升數乘所行里，今負鹽一升凡所行里也。於今有術，以所有四十錢爲所求率也。《衰分》章「貸人千錢」與此同。}\end{unitzh}
\begin{uniten}Method: Set out the number of sheng in 2 hu of salt, multiply it by 100 li, and make this the divisor (fa). [Liu Hui: Note: In this method one multiplies the 100 li traveled by the number of sheng in the 2 hu of salt carried, obtaining 20000 li; this means carrying 1 sheng of salt for 20000 li and being given 40 qian, which in the "suppose there is" method (jinyou shu) is the given rate (suoyou lü). The li traveled multiplied by the number of sheng makes the divisor; in the "suppose there is" method it is the given number (suoyou shu).] Multiply the number of sheng of salt now carried by 40 qian, again multiply it by 80 li, and make this the dividend (shi). For each time the dividend contains the divisor, one obtains 1 qian. [Liu Hui: The li traveled multiplied by the number of sheng of salt now carried is the total li traveled by 1 sheng of the salt now carried. In the "suppose there is" method, the given 40 qian is taken as the sought rate (suoqiu lü). The [problem of] "lending someone 1000 qian" in the chapter Cuifen ("Distribution by Proportion") is the same as this.]\end{uniten}

\begin{unitzh}\unitlabel{jiuzhang-025.3}今有負籠重一石一十七斤，行七十六步，五十返。今負籠重一石，行百步，問返幾何？荅曰：五十七返二千六百三分返之一千六百二十九。\end{unitzh}
\begin{uniten}Suppose there is [a man] carrying a basket weighing 1 dan (石) 17 jin (斤), who walks 76 bu (步) and [makes] 50 return trips. Now [he] carries a basket weighing 1 dan and walks 100 bu. Question: how many return trips? Answer: 57 1629/2603 return trips.\end{uniten}

\begin{unitzh}\unitlabel{jiuzhang-025.4}術曰：以故所行步數乘故籠重斤數，爲法。{\zhsmall 此法謂負一斤一返所行之積步也。}今籠重斤數乘今步，又以返數乘之，爲實。實如法得一返。{\zhsmall 按：此法，負一斤一返所行之積步；此實者，一斤一日所行之積步。故以一返之課除終日之程，卽是返數也。}\end{unitzh}
\begin{uniten}Method: Multiply the number of jin in the weight of the original basket by the number of bu originally walked, and make this the divisor. [Liu Hui: This divisor means the accumulated bu (jibu) walked in 1 return trip carrying 1 jin.] Multiply the present bu by the number of jin in the weight of the present basket, again multiply it by the number of return trips, and make this the dividend. For each time the dividend contains the divisor, one obtains 1 return trip. [Liu Hui: Note: this divisor is the accumulated bu walked in 1 return trip carrying 1 jin; this dividend is the accumulated bu walked with 1 jin in 1 day. Hence dividing the quota (cheng) of the whole day by the assessment (ke) of one return trip gives the number of return trips.]\end{uniten}

\begin{unitzh}\unitlabel{jiuzhang-025.5}{\zhsmall 臣淳風等謹按：此術所行步多者得返少，所行步少者得返多。然則所行者，今返率也。故令所得返乘今返之率爲實，而以故返之率爲法，今有術也。按：此負籠又有輕重，於是爲術者因令重者得返少，輕者得返多，故又因其率以乘法實者，重今有之義也。然此意非也。按：此籠雖輕而行有限，籠過重則人力遺。力有遺而術無窮，人行有限而籠輕重不等，使其有限之力隨彼無窮之變，故知此術率乖理也。若故所行有空行返數，設以問者，當因其所負以爲返率，則今返之數可得而知也。假令空行一日六十里，負重一斛行四十里，減重一斗進二里半，負重三斗以下與空行同。今負籠重六斗，往還行一百步，問返幾何？荅曰：一百五十返。術曰：置重行率，加十里，以里法通之，爲實；以一返之步爲法。實如法而一，卽得也。}\end{unitzh}
\begin{uniten}[Li Chunfeng et al.: Your servants respectfully note: In this method, where the bu walked are many, the return trips obtained are few; where the bu walked are few, the return trips obtained are many. That being so, what is walked is the rate of the present return trips. Hence one lets the return trips obtained multiply the rate of the present return trips as dividend, and takes the rate of the original return trips as divisor: this is the "suppose there is" method. Note: in this basket-carrying there is moreover lightness and heaviness; thereupon the maker of the method accordingly made the heavy obtain few return trips and the light obtain many; hence that one again uses their rates to multiply the divisor and the dividend is the sense of the double "suppose there is" (chong jinyou). But this idea is wrong. Note: though this basket be light, the walking has its limit; if the basket is too heavy, a person's strength fails. Strength fails, but the method is without limit; a person's walking is limited, but baskets are unequal in lightness and heaviness; making that limited strength follow those limitless variations — hence one knows that the rates of this method are contrary to reason. If, for what was originally walked, there were a number of return trips walking empty, and one set up a question with it, one should base the rate of return trips on what is carried; then the number of present return trips can be obtained and known. Suppose that walking empty one goes 60 li in 1 day; carrying a load of 1 hu one goes 40 li; for each dou by which the load is reduced one advances 2 1/2 li; carrying a load of 3 dou or less is the same as walking empty. Now one carries a basket weighing 6 dou and walks 100 bu going and returning. Question: how many return trips? Answer: 150 return trips. Method: Set out the rate of walking loaded, add 10 li, convert it by the li standard (lifa), and make this the dividend; take the bu of one return trip as divisor. Divide the dividend by the divisor, and one obtains [the answer].]\end{uniten}

\begin{unitzh}\unitlabel{jiuzhang-025.6}今有程傳委輸，空車日行七十里，重車日行五十里。今載太倉粟輸上林，五日三返。問太倉去上林幾何？荅曰：四十八里十八分里之十一。\end{unitzh}
\begin{uniten}Suppose there is a quota (cheng) for relay transport (zhuan weishu): an empty cart travels 70 li a day, a loaded cart travels 50 li a day. Now [one] loads millet (su) of the Great Granary (Taicang) and transports it to Shanglin, [making] 3 return trips in 5 days. Question: how far is the Great Granary from Shanglin? Answer: 48 11/18 li.\end{uniten}

\begin{unitzh}\unitlabel{jiuzhang-025.7}術曰：并空重里數，以三返乘之，爲法。令空重相乘，又以五日乘之，爲實。實如法得一里。{\zhsmall 此亦如上術。率：一百七十五里之路，往返用六日也。於今有術，則五日爲所有數，一百七十五里爲所求率，六日爲所有率。以此所得，則三返之路。今求一返，當以三約之，因令乘法而并除也。爲術亦可各置空重行一里用日之率，以爲列衰，副并爲法。以五日乘列衰爲實。實如法，所得卽各空重行日數也。各以一日所行以乘，爲凡日所行。三返約之，爲上林去太倉之數。}\end{unitzh}
\begin{uniten}Method: Sum the numbers of li of the empty and the loaded [cart], multiply it by the 3 return trips, and make this the divisor. Let the empty and the loaded [li] multiply each other, again multiply it by the 5 days, and make this the dividend. For each time the dividend contains the divisor, one obtains 1 li. [Liu Hui: This too is like the method above. As rates: a road of 175 li takes 6 days going and returning. In the "suppose there is" method, then, the 5 days are the given number, the 175 li the sought rate, and the 6 days the given rate. What is obtained by this is the road of 3 return trips. Now to find one return trip, one should reduce it by 3; accordingly one lets [it] multiply the divisor and divides in one go (bing chu). To make the method, one may also set out separately the rates of days used by the empty and by the loaded [cart] to travel 1 li, take them as ranked proportions (liecui), and their sum in an auxiliary position (fu) as divisor. Multiply the ranked proportions by the 5 days as dividend. Dividing the dividend by the divisor, what is obtained is the number of days traveled by the empty and by the loaded [cart] respectively. Multiply each by what is traveled in 1 day, making what is traveled in all the days. Reducing by the 3 return trips gives the distance from Shanglin to the Great Granary.]\end{uniten}

\begin{unitzh}\unitlabel{jiuzhang-025.8}{\zhsmall 臣淳風等謹按：此術重往空還，一輸再還道。置空行一里用七十分日之一，重行一里用五十分日之一。齊而同之，空重行一里之路，往返用一百七十五分日之六。定言之者，一百七十五里之路，往返用六日。故并空重者，并齊也；空重相乘者，同其母也。於今有術，五日爲所有數，一百七十五爲所求率，六爲所有率。以此所得，則三返之路。今求一返者，當以三約之，故令乘法而并除，亦當約之也。}\end{unitzh}
\begin{uniten}[Li Chunfeng et al.: Your servants respectfully note: In this method [the cart] goes loaded and comes back empty; one delivery goes over the road twice, [there and] back. Set out: traveling 1 li empty uses 1/70 day; traveling 1 li loaded uses 1/50 day. Homogenizing (qi) and equalizing (tong) them, a road of 1 li traveled empty and loaded uses 6/175 day going and returning. Stated in definite terms: a road of 175 li uses 6 days going and returning. Hence summing the empty and the loaded is summing the homogenized [terms]; multiplying the empty and the loaded together is equalizing their denominators. In the "suppose there is" method, the 5 days are the given number, 175 the sought rate, 6 the given rate. What is obtained by this is the road of 3 return trips. Now to find one return trip, one should reduce it by 3; hence letting [it] multiply the divisor and dividing in one go also amounts to reducing it.]\end{uniten}

\begin{unitzh}\unitlabel{jiuzhang-025.9}今有絡絲一斤爲練絲十二兩，練絲一斤爲靑絲一斤十二銖。今有靑絲一斤，問本絡絲幾何？荅曰：一斤四兩十六銖三十三分銖之十六。\end{unitzh}
\begin{uniten}Suppose there is 1 jin of raw silk (luosi) that makes 12 liang (兩) of refined silk (liansi), and 1 jin of refined silk makes 1 jin 12 zhu (銖) of blue silk (qingsi). Now there is 1 jin of blue silk. Question: how much was the original raw silk? Answer: 1 jin 4 liang 16 16/33 zhu.\end{uniten}

\begin{unitzh}\unitlabel{jiuzhang-025.10}術曰：以練絲十二兩乘靑絲一斤十二銖，爲法。以靑絲一斤銖數乘練絲一斤兩數，又以絡絲一斤乘之，爲實。實如法得一斤。\end{unitzh}
\begin{uniten}Method: Multiply 1 jin 12 zhu of blue silk by 12 liang of refined silk, and make this the divisor. Multiply the number of liang in 1 jin of refined silk by the number of zhu in 1 jin of blue silk, again multiply it by 1 jin of raw silk, and make this the dividend. For each time the dividend contains the divisor, one obtains 1 jin.\end{uniten}

\begin{unitzh}\unitlabel{jiuzhang-025.11}{\zhsmall 按：練絲一斤爲靑絲一斤十二銖，此練率三百八十四，靑率三百九十六也。又絡絲一斤爲練絲十二兩，此絡率十六，練率十二也。置今有靑絲一斤，以練率三百八十四乘之，爲實，實如靑絲率三百九十六而一，所得，靑絲一斤練絲之數也。又以絡率十六乘之，所得爲實，以練率十二爲法，所得卽練絲用絡絲之數也。是謂重今有也。雖各有率，不用中間，故令後實乘前實，後法乘前法，而并除也。故以練絲兩數爲實，靑絲銖數爲法。}\end{unitzh}
\begin{uniten}[Liu Hui: Note: 1 jin of refined silk makes 1 jin 12 zhu of blue silk: this [gives] the refined rate 384 and the blue rate 396. Again, 1 jin of raw silk makes 12 liang of refined silk: this [gives] the raw rate 16 and the refined rate 12. Set out the 1 jin of blue silk one now has, multiply it by the refined rate 384, and make this the dividend; divide the dividend by the blue-silk rate 396; what is obtained is the amount of refined silk for 1 jin of blue silk. Again multiply it by the raw rate 16, take what is obtained as dividend and the refined rate 12 as divisor; what is obtained is the amount of raw silk used for the refined silk. This is called the double "suppose there is". Although each has its rates, the intermediate is not used; hence one lets the later dividend multiply the former dividend and the later divisor multiply the former divisor, and divides in one go. Hence the number of liang of refined silk is made the dividend and the number of zhu of blue silk the divisor.]\end{uniten}

\begin{unitzh}\unitlabel{jiuzhang-025.12}{\zhsmall 一曰：又置絡絲一斤兩數與練絲十二兩，約之，絡得四，練得三，此其相與之率。又置練絲一斤銖數與靑絲一斤十二銖，約之，練得三十二，靑得三十三，亦其相與之率。齊其靑絲、絡絲，同其二練，絡得一百二十八，靑得九十九，練得九十六，卽三率悉通矣。今有練絲一斤爲所有數，絡絲一百二十八爲所求率，靑絲九十九爲所有率。爲率之意猶此，但不先約諸率耳。凡率錯互不通者，皆積齊同用之。放此，雖四五轉不異也。言同其二練者，以明三率之相與通耳，於術無以異也。}\end{unitzh}
\begin{uniten}[Liu Hui: Another [way]: Again set out the number of liang in 1 jin of raw silk and the 12 liang of refined silk, and reduce them: raw gets 4, refined gets 3; these are their mutual rates. Again set out the number of zhu in 1 jin of refined silk and the 1 jin 12 zhu of blue silk, and reduce them: refined gets 32, blue gets 33; these too are their mutual rates. Homogenize the blue silk and the raw silk and equalize their two [rates of] refined silk: raw gets 128, blue gets 99, refined gets 96; then the three rates all communicate. The 1 jin of refined silk one now has is the given number, the raw silk 128 the sought rate, the blue silk 99 the given rate. [Translator's note: the problem gives 1 jin of blue silk.] The idea of making the rates is like this, only the rates are not reduced first. Whenever rates are interlaced and do not communicate, one always uses them by accumulated homogenizing and equalizing. Following this pattern, even with four or five transfers it is no different. Saying "equalize their two [rates of] refined silk" is only to make clear the mutual communication of the three rates; in the method there is no difference.]\end{uniten}

\begin{unitzh}\unitlabel{jiuzhang-025.13}{\zhsmall 又一術：今有靑絲一斤銖數乘練絲一斤兩數，爲實，以靑絲一斤十二銖爲法，所得卽用練絲兩數。以絡絲一斤乘所得，爲實，以練絲十二兩爲法，所得卽用絡絲斤數也。}\end{unitzh}
\begin{uniten}[Liu Hui: Another method: Multiply the number of liang in 1 jin of refined silk by the number of zhu in the 1 jin of blue silk one now has, and make this the dividend; take 1 jin 12 zhu of blue silk as divisor; what is obtained is the number of liang of refined silk used. Multiply what is obtained by 1 jin of raw silk, and make this the dividend; take 12 liang of refined silk as divisor; what is obtained is the number of jin of raw silk used.]\end{uniten}

\begin{unitzh}\unitlabel{jiuzhang-026.1}今有惡粟二十斗，春之，得糲米九斗。今欲求粺米十斗，問惡粟幾何？荅曰：二十四斗六升八十一分升之七十四。\end{unitzh}
\begin{uniten}Suppose there is 20 dou (斗) of bad millet (esu); hulling it gives 9 dou of coarse rice (limi). Now one wishes to obtain 10 dou of polished rice (baimi). Question: how much is the bad millet? Answer: 24 dou 6 sheng (升) 74/81 sheng. [Translator's note: the text writes 春 'spring' for 舂 'to hull'.]\end{uniten}

\begin{unitzh}\unitlabel{jiuzhang-026.2}術曰：置糲米九斗，以九乘之，爲法。亦置粺米十斗，以十乘之，又以惡粟二十斗乘之，爲實。實如法得一斗。{\zhsmall 按：此術置今有求粺米十斗，以糲米率十乘之，如糲率九而一，卽糲亦化爲惡粟矣。此亦重今有之義。爲術之意猶絡絲也。雖各有率，不問中間，故令後實乘前實，後法乘前法，而并除之也。}\end{unitzh}
\begin{uniten}Method: Set up the 9 dou of coarse rice and multiply it by 9 to make the divisor (fa). Also set up the 10 dou of polished rice, multiply it by 10, and further multiply it by the 20 dou of bad millet to make the dividend (shi). Each time the dividend contains the divisor, one obtains 1 dou. [Liu Hui: Examining this method: set up the 10 dou of polished rice now sought, multiply it by the coarse-rice rate (lü) 10, and divide by the coarse-rice rate 9; then the coarse rice too is transformed into bad millet. This too is the sense of the repeated 'Suppose there is' (chong jinyou). The idea of the method is like that for the raw silk (luosi). Although each has its rate, one does not ask about what lies in between; therefore one makes the later dividend multiply the earlier dividend and the later divisor multiply the earlier divisor, and divides by them together.] [Translator's note: in the commentary, words appear to be missing after 'the coarse-rice rate 10'.]\end{uniten}

\begin{unitzh}\unitlabel{jiuzhang-026.3}今有善行者行一百步，不善行者行六十步。今不善行者先行一百步，善行者追之，問幾何步及之？荅曰：二百五十步。\end{unitzh}
\begin{uniten}Suppose there is a good walker who walks 100 bu (步) while a poor walker walks 60 bu. Now the poor walker walks 100 bu ahead first, and the good walker pursues him. Question: in how many bu does he catch up with him? Answer: 250 bu.\end{uniten}

\begin{unitzh}\unitlabel{jiuzhang-026.4}術曰：置善行者一百步，減不善行者六十步，餘四十步，以爲法。以善行者之一百步乘不善行者先行一百步，爲實。實如法得一步。{\zhsmall 按：此術以六十步減一百步，餘四十步，卽不善行者先行率也。善行者行一百步，追及率。約之，追及率得五，先行率得二。於今有術，不善行者先行一百步爲所有數，五爲所求率，二爲所有率，而今有之，得追及步也。}\end{unitzh}
\begin{uniten}Method: Set up the good walker's 100 bu and subtract the poor walker's 60 bu; the remainder, 40 bu, is taken as the divisor. With the good walker's 100 bu multiply the 100 bu the poor walker walked ahead, to make the dividend. Each time the dividend contains the divisor, one obtains 1 bu. [Liu Hui: Examining this method: subtracting 60 bu from 100 bu, the remainder, 40 bu, is the rate of the poor walker's walking ahead (xianxing lü). The 100 bu the good walker walks: the rate of catching up (zhuiji lü). Reduce them (yue): the rate of catching up gives 5, the rate of walking ahead gives 2. In the method of 'Suppose there is' (jinyou shu), the 100 bu the poor walker walked ahead is the number one has (suoyou shu), 5 is the sought rate (suoqiu lü), 2 is the rate one has (suoyou lü); apply 'Suppose there is' to them, and one obtains the bu of catching up.]\end{uniten}

\begin{unitzh}\unitlabel{jiuzhang-026.5}今有不善行者先行十里，善行者追之一百里，先至不善行者二十里。問善行者幾何里及之？荅曰：三十三里少半里。\end{unitzh}
\begin{uniten}Suppose there is a poor walker who walks 10 li (里) ahead first; a good walker pursues him for 100 li and arrives first, 20 li ahead of the poor walker. Question: in how many li does the good walker catch up with him? Answer: 33 1/3 li.\end{uniten}

\begin{unitzh}\unitlabel{jiuzhang-026.6}術曰：置不善行者先行十里，以善行者先至二十里增之，以爲法。以不善行者先行十里乘善行者一百里，爲實。實如法得一里。{\zhsmall 按：此術不善行者旣先行十里，後不及二十里，并之，得三十里也，謂之先行率。善行者一百里爲追及率。約之，先行率得三，三爲所有率，而今有之，卽得也。其意如上術。}\end{unitzh}
\begin{uniten}Method: Set up the poor walker's 10 li walked ahead, add to it the 20 li by which the good walker arrives first, and take this as the divisor. With the poor walker's 10 li walked ahead multiply the good walker's 100 li, to make the dividend. Each time the dividend contains the divisor, one obtains 1 li. [Liu Hui: Examining this method: the poor walker has already walked 10 li ahead, and afterwards falls short by 20 li; adding them gives 30 li, which is called the rate of walking ahead. The good walker's 100 li is the rate of catching up. Reduce them: the rate of walking ahead gives 3; 3 is the rate one has; apply 'Suppose there is' to it, and one obtains it. Its idea is like that of the method above.] [Translator's note: words appear to be missing after 'the rate of walking ahead gives 3'.]\end{uniten}

\begin{unitzh}\unitlabel{jiuzhang-026.7}今有兔先走一百步，犬追之二百五十步，不及三十步而止。問犬不止，復行幾何步及之？荅曰：一百七步七分步之一。\end{unitzh}
\begin{uniten}Suppose there is a hare that runs 100 bu ahead first; a dog pursues it for 250 bu, falls 30 bu short of it and stops. Question: if the dog does not stop, how many more bu does it go to catch up with it? Answer: 107 1/7 bu.\end{uniten}

\begin{unitzh}\unitlabel{jiuzhang-026.8}術曰：置兔先走一百步，以犬走不及三十步減之，餘爲法。以不及三十步乘犬追步數，爲實。實如法得一步。{\zhsmall 按：此術以不及三十步減先走一百步，餘七十步，爲兔先走率。犬行二百五十步爲追及率。約之，先走率得七，追及率得二十五。於今有術，不及三十步爲所有數，二十五爲所求率，七爲所有率，而今有之，卽得也。}\end{unitzh}
\begin{uniten}Method: Set up the hare's 100 bu run ahead and subtract from it the 30 bu by which the dog's running falls short; the remainder is the divisor. With the 30 bu shortfall multiply the number of bu of the dog's pursuit, to make the dividend. Each time the dividend contains the divisor, one obtains 1 bu. [Liu Hui: Examining this method: subtracting the 30 bu shortfall from the 100 bu run ahead, the remainder, 70 bu, is the rate of the hare's running ahead. The 250 bu the dog went is the rate of catching up. Reduce them: the rate of running ahead gives 7, the rate of catching up gives 25. In the method of 'Suppose there is', the 30 bu shortfall is the number one has, 25 is the sought rate, 7 is the rate one has; apply 'Suppose there is' to them, and one obtains it.]\end{uniten}

\begin{unitzh}\unitlabel{jiuzhang-026.9}今有人持金十二斤出關，關稅之，十分而取一。今關取金二斤，償錢五千。問金一斤値錢幾何？荅曰：六千二百五十。\end{unitzh}
\begin{uniten}Suppose there is a man who carries 12 jin (斤) of gold out through a pass; the pass taxes it, taking 1 out of 10 parts. Now the pass takes 2 jin of gold and pays back 5000 qian (錢, cash). Question: how many qian is 1 jin of gold worth? Answer: 6250.\end{uniten}

\begin{unitzh}\unitlabel{jiuzhang-026.10}術曰：以十乘二斤，以十二斤減之，餘爲法。以十乘五千，爲實。實如法得一。{\zhsmall 按：此術置十二斤，以一乘之，十而一，得一斤五分斤之一，卽所當稅者也。減二斤，餘卽關取盈金。以盈除所償錢，卽金値也。今術旣以十二斤爲所稅，則是以十爲母，故以十乘二斤及所償錢，通其率。於今有術，五千錢爲所有數，十爲所求率，八爲所有率，而今有之，卽得也。}\end{unitzh}
\begin{uniten}Method: Multiply 2 jin by 10 and subtract 12 jin from it; the remainder is the divisor. Multiply 5000 by 10 to make the dividend. Each time the dividend contains the divisor, one obtains 1. [Liu Hui: Examining this method: set up 12 jin, multiply it by 1 and divide by 10, obtaining 1 1/5 jin; this is what should be taxed. Subtract it from 2 jin: the remainder is the excess gold taken by the pass. Dividing the cash paid back by the excess gives the value of the gold. Now since the method takes 12 jin as what is taxed, this means taking 10 as the denominator (mu); therefore one multiplies the 2 jin and the cash paid back by 10, so as to make their rates commensurate (tong). In the method of 'Suppose there is', 5000 qian is the number one has, 10 is the sought rate, 8 is the rate one has; apply 'Suppose there is' to them, and one obtains it.]\end{uniten}

\begin{unitzh}\unitlabel{jiuzhang-026.11}今有客馬日行三百里。客去忘持衣，日已三分之一，主人乃覺，持衣追，及與之而還，至家視日四分之三。問主人馬不休，日行幾何？荅曰：七百八十里。\end{unitzh}
\begin{uniten}Suppose there is a guest's horse that travels 300 li a day. The guest leaves and forgets to take his clothes. When 1/3 of the day has already passed, the host then notices; he takes the clothes and pursues him, catches up, gives them to him and returns; reaching home, he looks at the sun: 3/4 of the day. Question: if the host's horse does not rest, how far does it travel in a day? Answer: 780 li.\end{uniten}

\begin{unitzh}\unitlabel{jiuzhang-026.12}術曰：置四分日之三，除三分日之一，{\zhsmall 按：此術「置四分日之三，除三分日之一」者，除卽減也。減之，餘有十二分之五，卽是主人追客還用日率也。}半其餘以爲法。{\zhsmall 去其還，存其往。率之者，子不可半，故倍母，二十四分之五，是爲主人與客均行用日之率也。}副置法，增三分日之一，{\zhsmall 法二十四分之五者，主人往追用日之分也。三分之一者，客去主人未覺之前獨行用日之分也。并連此數，得二十四分日之十三，則主人追及前用日之分也，是爲客用日率也。然則主人用日率者，客馬行率也；客用日率者，主人馬行率也。母同則子齊，是爲客馬行率五，主人馬行率十三。於今有術，三百里爲所有數，十三爲所求率，五爲所有率，而今有之，卽得也。}以三百里乘之，爲實。實如法得主人馬一日行。{\zhsmall 欲知主人追客所行里者，以三百里乘主人均行日分子十三，以母二十四而一，得一百六十二里半。以此乘主人均行日分母二十四，如客馬與主人均行用日分子五而一，亦得主人馬一日行七百八十里也。}\end{unitzh}
\begin{uniten}Method: Set up 3/4 day and remove (chu) 1/3 day from it; [Liu Hui: Examining this method: in 'set up 3/4 day and remove 1/3 day', removing is subtracting. Subtracting, the remainder is 5/12; this is the rate of the days used by the host in pursuing the guest and returning.] halve the remainder and take it as the divisor. [Liu Hui: Take away his return and keep his going. In making a rate of it, the numerator (zi) cannot be halved, so the denominator is doubled: 5/24. This is the rate of the days used by the host and the guest for equal travel.] Set up the divisor in a secondary position (fuzhi) and add 1/3 day to it; [Liu Hui: The divisor 5/24 is the fraction of a day used by the host in going in pursuit. The 1/3 is the fraction of a day used in travelling alone after the guest left and before the host noticed. Joining these numbers together gives 13/24 day, which is the fraction of a day used before the host caught up; this is the rate of the days used by the guest. Thus the rate of the days used by the host is the rate of travel of the guest's horse; the rate of the days used by the guest is the rate of travel of the host's horse. When the denominators are the same, the numerators are homogeneous (qi): this gives the rate of travel of the guest's horse, 5, and the rate of travel of the host's horse, 13. In the method of 'Suppose there is', 300 li is the number one has, 13 is the sought rate, 5 is the rate one has; apply 'Suppose there is' to them, and one obtains it.] multiply it by 300 li to make the dividend. Dividing the dividend by the divisor, one obtains what the host's horse travels in 1 day. [Liu Hui: If one wishes to know the li the host travelled in pursuing the guest: with 300 li multiply 13, the numerator of the host's equal-travel day fraction, and divide by the denominator 24, obtaining 162 1/2 li. Multiply this by 24, the denominator of the host's equal-travel day fraction, and divide by 5, the numerator of the day fraction used by the guest's horse and the host for equal travel; one also obtains that the host's horse travels 780 li in 1 day.]\end{uniten}

\begin{unitzh}\unitlabel{jiuzhang-026.13}今有金箠，長五尺，斬本一尺，重四斤；斬末一尺，重二斤。問次一尺各重幾何？荅曰：末一尺重二斤；次一尺重二斤八兩；次一尺重三斤；次一尺重三斤八兩；次一尺重四斤。\end{unitzh}
\begin{uniten}Suppose there is a golden rod (chui), 5 chi (尺) long. Cutting off 1 chi at the base, it weighs 4 jin; cutting off 1 chi at the tip, it weighs 2 jin. Question: how much does each successive 1 chi weigh? Answer: the 1 chi at the tip weighs 2 jin; the next 1 chi weighs 2 jin 8 liang (兩); the next 1 chi weighs 3 jin; the next 1 chi weighs 3 jin 8 liang; the next 1 chi weighs 4 jin.\end{uniten}

\begin{unitzh}\unitlabel{jiuzhang-026.14}術曰：令末重減本重，餘卽差率也。又置本重，以四間乘之，爲下第一衰。副置，以差率減之，每尺各自爲衰。{\zhsmall 按：此術五尺有四間者，有四差也。令本末相減，餘卽四差之凡數也。以四約之，卽得每尺之差。以差數減本重，餘卽次尺之重也。爲術所置，如是而已。今此率以四爲母，故令母乘本爲衰，通其率也。亦可置末重，以四問乘之，爲上第一衰，以差重率加之，爲次下衰也。}副置下第一衰，以爲法。以本重四斤徧乘列衰，各自爲實。實如法得一斤。{\zhsmall 以下第一衰爲法，以本重乘其分母之數，而又反此率乘本重爲實。一乘一除，勢無損益，故惟本存焉。衆衰相推爲率，則其餘可知也。亦可副置末衰爲法，而以末重二斤乘列衰爲實。此雖迂廽，然是其舊，故就新而言之也。}\end{unitzh}
\begin{uniten}Method: Subtract the weight of the tip from the weight of the base; the remainder is the difference rate (cha lü). Further set up the weight of the base and multiply it by the 4 intervals (jian) to make the first share (cui) from the bottom. Set it up in a secondary position and subtract the difference rate from it; each chi separately makes a share. [Liu Hui: Examining this method: that 5 chi have 4 intervals means that there are 4 differences. Subtract the base and the tip from each other: the remainder is the total of the 4 differences. Reducing it by 4, one obtains the difference for each chi. Subtracting the difference from the weight of the base, the remainder is the weight of the next chi. What the method sets up is just this. Now these rates take 4 as the denominator; therefore one makes the denominator multiply the base to make the share, so as to make their rates commensurate. One may also set up the weight of the tip, multiply it by the 4 intervals to make the first share from the top, and add the difference-of-weight rate to it to make the next share below.] [Translator's note: in the commentary, 'the 4 intervals' renders 四問; 問 'ask' evidently stands for 間 'interval'.] Set up the first share from the bottom in a secondary position and take it as the divisor. With the weight of the base, 4 jin, multiply all the listed shares (lie cui), each separately making a dividend. Each time the dividend contains the divisor, one obtains 1 jin. [Liu Hui: Taking the first share from the bottom as the divisor is multiplying the weight of the base by the number of its denominator; and further, conversely, these rates multiply the weight of the base to make the dividends. Multiplying once and dividing once, in effect there is neither loss nor gain, so only the base remains. When the shares are deduced from one another as rates, the rest can be known. One may also set up the share of the tip in a secondary position as the divisor, and multiply the listed shares by the weight of the tip, 2 jin, to make the dividends. Although this is roundabout, it is nevertheless the old one; therefore one adopts the new and states it.]\end{uniten}

\begin{unitzh}\unitlabel{jiuzhang-027.1}今有五人分五錢，令上二人所得與下三人等。問各得幾何？荅曰：甲得一錢六分錢之二，乙得一錢六分錢之一，丙得一錢，丁得六分錢之五，戊得六分錢之四。\end{unitzh}
\begin{uniten}Suppose there are 5 persons who divide 5 qian (錢, cash), so that what the upper 2 persons obtain is equal to [what] the lower 3 persons [obtain]. Question: how much does each obtain? Answer: A (jia) obtains 1 2/6 qian; B (yi) obtains 1 1/6 qian; C (bing) obtains 1 qian; D (ding) obtains 5/6 qian; E (wu) obtains 4/6 qian.\end{uniten}

\begin{unitzh}\unitlabel{jiuzhang-027.2}術曰：置錢錐行衰，{\zhsmall 按此術錐行者，謂如立錐：初一、次二、次三、次四、次五，各均爲一例者也。}并上二人爲九，并下三人爲六。六少於九三，{\zhsmall 數不得等，但以五、四、三、二、一爲率也。}以三均加焉，副并爲法。以所分錢乘未并者，各自爲實。實如法得一錢。\end{unitzh}
\begin{uniten}Method: Set out the cash [according to] awl-shaped (zhuixing) proportional parts (shuai). [Liu Hui: According to this method, "awl-shaped" means like an upright awl: first 1, next 2, next 3, next 4, next 5, each evenly forming a single series.] Sum the upper 2 persons, making 9; sum the lower 3 persons, making 6. 6 is less than 9 by 3; [Liu Hui: The numbers cannot be equal; one merely takes 5, 4, 3, 2, 1 as the rates (lü).] add 3 evenly to them, and sum them in an auxiliary position (fu) to make the divisor (fa). Multiply the unsummed ones by the cash to be divided, each separately making a dividend (shi). Divide the dividend by the divisor, obtaining [the result in] qian.\end{uniten}

\begin{unitzh}\unitlabel{jiuzhang-027.3}{\zhsmall 此問者，令上二人與下三人等，上下部差一人，其差三。均加上部，則得二三；均加下部，則得三三。下部猶差一人，差得一，以通於本率，卽上下部等也。於今有術，副并爲所有率，未并者各爲所求率，五錢爲所有數，而今有之，卽得等耳。假令七人分七錢，欲令上二人與下五人等，則上下部差三人。并上部爲十三，下部爲十五，下多上少，下不足減上，當似上下部列差而後均減，乃合所問耳。此可倣下術：令上二人分二錢半爲上率，令下三人分二錢半爲下率，上下二率以少減多，餘爲實；置二人、三人，各半之，減五人，餘爲法。實如法得一錢，卽衰相去也。下衰率六分之五者，丁所得錢數也。}\end{unitzh}
\begin{uniten}[Liu Hui: In this question, the upper 2 persons are made equal to the lower 3 persons; the upper and lower parts differ by 1 person, and their difference is 3. Adding [3] evenly to the upper part, one obtains 2 [times] 3; adding [3] evenly to the lower part, one obtains 3 [times] 3. The lower part still differs by 1 person; the difference obtains 1 [times 3], and by this it is made to accord with the original rates: then the upper and lower parts are equal. In terms of the "Suppose there is" method (jinyou shu), the sum in the auxiliary position is the given rate (suoyou lü), the unsummed ones are each a sought rate (suoqiu lü), and the 5 qian is the given number (suoyou shu); applying the "Suppose there is" [method] to them, one obtains the equality. Suppose 7 persons divide 7 qian, and one wishes the upper 2 persons to be equal to the lower 5 persons; then the upper and lower parts differ by 3 persons. Summing the upper part makes 13, and the lower part makes 15; the lower is more and the upper less, and the lower is insufficient for subtracting the upper; it should be as if one set out the difference of the upper and lower parts and then subtracted evenly; only then does it agree with what is asked. This can be modeled on the method below: let the upper 2 persons divide 2 1/2 qian to make the upper rate; let the lower 3 persons divide 2 1/2 qian to make the lower rate; of the 2 rates, upper and lower, subtract the smaller from the larger, and the remainder is the dividend; set out 2 persons and 3 persons, halve each, and subtract [them] from 5 persons; the remainder is the divisor. Divide the dividend by the divisor, obtaining [the result in] qian: this is the distance between [successive] proportional parts. The lower proportional rate, 5/6, is the number of qian that D obtains.]\end{uniten}

\begin{unitzh}\unitlabel{jiuzhang-027.4}今有竹九節，下三節容四升，上四節容三升。問中間二節欲均容，各多少？荅曰：下初一升六十六分升之二十九，次一升六十六分升之二十二，次一升六十六分升之一十五，次一升六十六分升之八，次一升六十六分升之一，次六十六分升之六十，次六十六分升之五十三，次六十六分升之四十六，次六十六分升之三十九。\end{unitzh}
\begin{uniten}Suppose there is a bamboo of 9 sections; the lower 3 sections hold 4 sheng (升), and the upper 4 sections hold 3 sheng. Question: the 2 sections in the middle are to hold [amounts] evenly [graded]; how much does each hold? Answer: lowest, first: 1 29/66 sheng; next: 1 22/66 sheng; next: 1 15/66 sheng; next: 1 8/66 sheng; next: 1 1/66 sheng; next: 60/66 sheng; next: 53/66 sheng; next: 46/66 sheng; next: 39/66 sheng.\end{uniten}

\begin{unitzh}\unitlabel{jiuzhang-027.5}術曰：以下三節分四升爲下率，以上四節分三升爲上率，{\zhsmall 此二率者，各其平率也。}上下率以少減多，餘爲實，{\zhsmall 按此上下節各分所容爲率者，各其平率。上下以少減多者，餘爲中間五節半之凡差，故以爲實也。}置四節、三節，各半之，以減九節，餘爲法。實如法得一升，卽衰相去也。{\zhsmall 按此術上下節所容已定之節，中間相去節數也。實者，中間五節半之凡差也。故實如法而一，則每節之差也。}下率一升少半升者，下第二節容也。{\zhsmall 一升少半升者，下三節通分四升之平率。平率卽爲中分節之容也。}\end{unitzh}
\begin{uniten}Method: With the lower 3 sections divide 4 sheng to make the lower rate; with the upper 4 sections divide 3 sheng to make the upper rate. [Liu Hui: These 2 rates are each their average rate (ping lü).] Of the upper and lower rates, subtract the smaller from the larger; the remainder is the dividend. [Liu Hui: According to this, [the rates made by] the upper and lower sections each dividing what they hold are each their average rate. Subtracting the smaller from the larger of the upper and lower, the remainder is the total difference (fan cha) over the 5 1/2 sections in the middle; therefore it is taken as the dividend.] Set out 4 sections and 3 sections, halve each, and subtract them from 9 sections; the remainder is the divisor. Divide the dividend by the divisor, obtaining [the result in] sheng: this is the distance between [successive] proportional parts. [Liu Hui: According to this method, [the divisor is] the number of sections by which the middles of the upper and lower sections, the sections whose capacities are already determined, are apart. The dividend is the total difference over the 5 1/2 sections in the middle. Therefore dividing the dividend by the divisor gives the difference for each section.] The lower rate, 1 1/3 sheng, is the capacity of the second section from the bottom. [Liu Hui: 1 1/3 sheng is the average rate of the lower 3 sections dividing the 4 sheng in common (tongfen). The average rate is the capacity of the section at the middle division.]\end{uniten}

\begin{unitzh}\unitlabel{jiuzhang-027.6}今有鳬起南海，七日至北海；鴈起北海，九日至南海。今鳬、鴈俱起，問何日相逢？荅曰：三日十六分日之十五。\end{unitzh}
\begin{uniten}Suppose a wild duck (fu) sets out from the South Sea and arrives at the North Sea in 7 days; a wild goose (yan) sets out from the North Sea and arrives at the South Sea in 9 days. Now the duck and the goose set out together. Question: on what day do they meet? Answer: 3 15/16 days.\end{uniten}

\begin{unitzh}\unitlabel{jiuzhang-027.7}術曰：并日數爲法，日數相乘爲實，實如法得一日。{\zhsmall 按此術置鳬七日一至，鴈九日一至。齊其至，同其日，定六十三日鳬九至，鴈七至。令鳬鴈俱起而問相逢者，是爲共至。并齊以除同，卽得相逢日。故并日數爲法者，并齊之意；日數相乘爲實者，猶以同爲實也。一日鳬飛七分至之一，鴈飛九分至之一。齊而同之，鳬飛定日行六十三分至之九，鴈飛定日行六十三分至之七。是爲南北海相去六十三分，鳬日行九分，鴈日行七分也。并鳬鴈一日所行，以除南北相去，而得相逢日也。}\end{unitzh}
\begin{uniten}Method: Sum the numbers of days to make the divisor; multiply the numbers of days together to make the dividend. Divide the dividend by the divisor, obtaining [the result in] days. [Liu Hui: According to this method, set out: the duck arrives once in 7 days, the goose arrives once in 9 days. Homogenize (qi) their arrivals and equalize (tong) their days: fixing 63 days, the duck arrives 9 times and the goose arrives 7 times. Letting the duck and goose set out together and asking when they meet: this is a joint arrival. Sum the homogenized and with it divide the equalized; then one obtains the day of meeting. Therefore summing the numbers of days to make the divisor is the idea of summing the homogenized; multiplying the numbers of days together to make the dividend is like taking the equalized as the dividend. In 1 day the duck flies 1/7 of an arrival, and the goose flies 1/9 of an arrival. Homogenizing and equalizing them, the duck's flight is fixed at a daily travel of 9/63 of an arrival, and the goose's flight is fixed at a daily travel of 7/63 of an arrival. This is [to say] that the South and North Seas are 63 parts apart, the duck travels 9 parts a day, and the goose travels 7 parts a day. Sum what the duck and goose travel in 1 day, divide the distance between South and North by it, and one obtains the day of meeting.]\end{uniten}

\begin{unitzh}\unitlabel{jiuzhang-027.8}今有甲發長安，五日至齊；乙發齊，七日至長安。今乙發已先二日，甲乃發長安。問幾何日相逢？荅曰：二日十二分日之一。\end{unitzh}
\begin{uniten}Suppose A sets out from Chang'an and arrives at Qi in 5 days; B sets out from Qi and arrives at Chang'an in 7 days. Now B has already set out 2 days earlier, and then A sets out from Chang'an. Question: in how many days do they meet? Answer: 2 1/12 days.\end{uniten}

\begin{unitzh}\unitlabel{jiuzhang-027.9}術曰：并五日、七日以爲法。{\zhsmall 按此術并五日、七日爲法者，猶并齊爲法。置甲五日一至，乙七日一至，齊而同之，定三十五日甲七至，乙五至。并之爲十二至者，用三十五日也。謂甲乙與發之率耳。然則日化爲至，當除日，故以爲法也。}以乙先發二日減七日，{\zhsmall 減七日者，言甲乙俱發，今以發爲始發之端，於本道里則餘分也。}餘以乘甲日數爲實。{\zhsmall 七者，長安去齊之率也；五者，後發相去之率也。今問後發，故舍七用五，以乘甲五日爲二十五日。言甲七至、乙五至，更相去用此二十五日也。}實如法得一日。\end{unitzh}
\begin{uniten}Method: Sum 5 days and 7 days to make the divisor. [Liu Hui: According to this method, summing 5 days and 7 days to make the divisor is like summing the homogenized to make the divisor. Set out: A arrives once in 5 days, B arrives once in 7 days; homogenizing and equalizing them, fixing 35 days, A arrives 7 times and B arrives 5 times. Summing them makes 12 arrivals, which use 35 days. This refers to the rate of A and B setting out together. Thus days are transformed into arrivals, [which] should divide the days; therefore [their sum] is taken as the divisor.] Subtract the 2 days by which B set out first from 7 days; [Liu Hui: "Subtracting from 7 days" speaks of A and B setting out together: now, taking the setting out [of A] as the start of setting out, with respect to the original distance (daoli) it is the remaining part.] with the remainder multiply A's number of days to make the dividend. [Liu Hui: 7 is the rate of the distance from Chang'an to Qi; 5 is the rate of the distance apart at the later setting out. Now the question concerns the later setting out, so discard the 7 and use the 5, multiplying A's 5 days by it to make 25 days. This means that, with A arriving 7 times and B arriving 5 times, the renewed distance apart uses these 25 days.] Divide the dividend by the divisor, obtaining [the result in] days.\end{uniten}

\begin{unitzh}\unitlabel{jiuzhang-027.10}{\zhsmall 一日甲行五分至之一，乙行七分至之一。齊而同之，甲定日行三十五分至之七，乙定日行三十五分至之五。是爲齊去長安三十五分，甲日行七分，乙日行五分也。今乙先行發二日，巳行十分，餘相去二十五分，故減乙二日，餘令相乘，爲二十五分。}\end{unitzh}
\begin{uniten}[Liu Hui: In 1 day A travels 1/5 of an arrival, and B travels 1/7 of an arrival. Homogenizing and equalizing them, A is fixed at a daily travel of 7/35 of an arrival, and B is fixed at a daily travel of 5/35 of an arrival. This is [to say] that Qi is 35 parts distant from Chang'an, A travels 7 parts a day, and B travels 5 parts a day. Now B set out and traveled 2 days in advance and has already traveled 10 parts; the remaining distance apart is 25 parts. Therefore subtract B's 2 days, and let the remainder be multiplied [with A's number of days], making 25 parts.]\end{uniten}

\begin{unitzh}\unitlabel{jiuzhang-027.11}今有一人三日爲牡瓦三十八枚，一人二日爲牝瓦七十六枚。今令一人一日作瓦，牝牡相半，問成瓦幾何？荅曰：二十五枚少半枚。\end{unitzh}
\begin{uniten}Suppose 1 person makes 38 male tiles (mu wa) in 3 days, and 1 person makes 76 female tiles (pin wa) in 2 days. Now let 1 person make tiles in 1 day, female and male half and half. Question: how many tiles are made? Answer: 25 1/3 pieces (枚). [Translator's note: this answer does not agree with the 3 days and 2 days stated in the question.]\end{uniten}

\begin{unitzh}\unitlabel{jiuzhang-027.12}術曰：并牝牡爲法，牝牡相乘爲實，實如法得一枚。{\zhsmall 此術亦與鳬鴈術同，牝牡瓦相并，猶如鳬鴈日飛相并也。}{\zhsmall 臣淳風等謹按：此術并牝牡爲法者，并齊之意；牝牡相乘爲實者，猶以同爲實也。故實如法卽得也。}\end{unitzh}
\begin{uniten}Method: Sum the female and male [tiles] to make the divisor; multiply the female and male [tiles] together to make the dividend. Divide the dividend by the divisor, obtaining [the result in] pieces. [Liu Hui: This method is also the same as the duck-and-goose method: summing the female and male tiles is just like summing the daily flights of the duck and goose.] [Li Chunfeng et al.: We respectfully note: in this method, summing the female and male to make the divisor is the idea of summing the homogenized; multiplying the female and male together to make the dividend is like taking the equalized as the dividend. Therefore dividing the dividend by the divisor gives [the answer] at once.]\end{uniten}

\begin{unitzh}\unitlabel{jiuzhang-027.13}今有一人一日矯矢五十，一人一日羽矢三十，一人一日筈矢十五。今令一人一日自矯、羽、筈，\end{unitzh}
\begin{uniten}Suppose 1 person in 1 day straightens 50 arrows (jiao shi); 1 person in 1 day fletches 30 arrows (yu shi); 1 person in 1 day nocks 15 arrows (kuo shi). Now let 1 person in 1 day himself straighten, fletch and nock [continues]\end{uniten}

\begin{unitzh}\unitlabel{jiuzhang-028.1}問成矢幾何？荅曰：八矢少半矢。\end{unitzh}
\begin{uniten}[continued] Question: how many arrows are completed? Answer: 8 1/3 arrows.\end{uniten}

\begin{unitzh}\unitlabel{jiuzhang-028.2}術曰：矯矢五十，用徒一人；羽矢五十，用徒一人太半人；筈矢五十，用徒三人少半人。并之，得六人，以爲法；以五十矢爲實。實如法得一矢。{\zhsmall 按：此術言成矢五十，用徒六人一日工也。此同功共作，猶鳬鴈共至之類，亦以同爲實，并齊爲法。可令矢互乘一人爲齊，矢相乘爲同。今先令同於五十矢，矢同則徒齊，其歸一也。以此術爲鳬鴈者，當鴈飛九日而一至，鳬飛七日而一至七分至之二，并之，得二至七分至之二，以爲法；以九日爲實，實如法而一，得一人日矯矢之數也。}\end{unitzh}
\begin{uniten}Method: Straightening 50 arrows uses 1 laborer; feathering 50 arrows uses 1 2/3 laborers; notching 50 arrows uses 3 1/3 laborers. Sum them, obtaining 6 persons, and take this as the divisor (fa); take 50 arrows as the dividend (shi). Divide the dividend by the divisor to obtain [the number of] arrows. [Liu Hui: Considering this method: it says that completing 50 arrows uses one day's work of 6 laborers. This is the same task worked at jointly, like the type in which the wild duck and the goose arrive together; it likewise takes the common (tong) [number] as the dividend and the sum of the homogenized (qi) [numbers] as the divisor. One may let the arrows cross-multiply (hucheng) the 1 person to make the homogenized [numbers], and multiply the arrows together to make the common [number]. Here one first makes [them] common at 50 arrows; when the arrows are common, the laborers are homogenized; it comes to the same thing. If this method is applied to the wild duck and goose, then: the goose flies 9 days and arrives once; the wild duck flies 7 days and arrives 1 2/7 times [Translator's note: the computation requires 9 days here]; summing them gives 2 2/7 arrivals, which is taken as the divisor; take 9 days as the dividend; divide the dividend by the divisor, and one obtains the number of arrows straightened by one person in one day.]\end{uniten}

\begin{unitzh}\unitlabel{jiuzhang-028.3}今有假田，初假之歲三畝一錢，明年四畝一錢，後年五畝一錢，凡三歲得一百。問田幾何？荅曰：一頃二十七畝四十七分畝之三十一。\end{unitzh}
\begin{uniten}Suppose there is a field leased out (jia tian). In the first year of the lease, 3 mu (畝) for 1 qian (錢, cash); the next year, 4 mu for 1 qian; the year after, 5 mu for 1 qian. In all, over 3 years, 100 is obtained. Question: how much field is there? Answer: 1 qing (頃) 27 31/47 mu.\end{uniten}

\begin{unitzh}\unitlabel{jiuzhang-028.4}術曰：置畝數及錢數，令畝數互乘錢數，并以爲法；畝數相乘，又以百錢乘之，爲實。實如法得一畝。{\zhsmall 按：此術令畝互乘錢者，齊其錢；畝數相乘者，同其畝。同於六十，則初假之歲得錢二十，明年得錢十五，後年得錢十二也。凡三歲得錢一百，爲所有數；同畝爲所求率；四十七錢爲所有率。今有之，卽得也。齊其錢，同其畝，亦如鳬鴈術。於今有術，百錢爲所有數，同畝爲所求率，并齊爲所有率。臣淳風等謹按：假田六十畝，初歲得錢二十，明年得錢十五，後年得錢十二，并之得錢四十七，是爲得田六十畝三歲所治。於今有術，百錢爲所有數，六十畝爲所求率，四十七爲所有率，而今有之，卽合問也。}\end{unitzh}
\begin{uniten}Method: Set out the numbers of mu and the numbers of qian; let the numbers of mu cross-multiply the numbers of qian, and sum them to make the divisor. Multiply the numbers of mu together, and further multiply by 100 qian, to make the dividend. Divide the dividend by the divisor to obtain [the number of] mu. [Liu Hui: Considering this method: letting the mu cross-multiply the qian homogenizes the qian; multiplying the numbers of mu together makes the mu common. When they are made common at 60, then in the first year of the lease 20 qian are obtained, the next year 15 qian, and the year after 12 qian. The 100 qian obtained in all over 3 years is the given number (suoyou shu); the common mu is the sought rate (suoqiu lü); the 47 qian are the given rate (suoyou lü). Apply 'suppose there is' (jinyou) to them, and one obtains [the answer]. Homogenizing the qian and making the mu common is also like the method of the wild duck and goose. In the 'suppose there is' method (jinyou shu), the 100 qian are the given number, the common mu is the sought rate, and the sum of the homogenized [numbers] is the given rate.] [Li Chunfeng et al.: For a leased field of 60 mu, in the first year 20 qian are obtained, the next year 15 qian, the year after 12 qian; summing them gives 47 qian; this is what is obtained for 60 mu of field managed over 3 years. In the 'suppose there is' method, 100 qian are the given number, 60 mu the sought rate, 47 the given rate; and applying 'suppose there is' to them matches the question.]\end{uniten}

\begin{unitzh}\unitlabel{jiuzhang-028.5}今有程耕，一人一日發七畝，一人一日耕三畝，一人一日耰種五畝。今令一人一日自發、耕、耰種之，問治田幾何？荅曰：一畝一百一十四步七十一分步之六十六。\end{unitzh}
\begin{uniten}Suppose there is a standard (cheng) for tilling: one person in one day breaks up 7 mu; one person in one day ploughs 3 mu; one person in one day harrows and sows 5 mu. Now let one person in one day himself break up, plough, harrow and sow it. Question: how much field is worked? Answer: 1 mu 114 66/71 bu (步).\end{uniten}

\begin{unitzh}\unitlabel{jiuzhang-028.6}術曰：置發、耕、耰畝數，令互乘人數，并以爲法；畝數相乘爲實。實如法得一畝。{\zhsmall 此猶鳬鴈術也。}{\zhsmall 臣淳風等謹按：此術亦發、耕、耰種畝數互乘人者，齊其人；畝數相乘者，同其畝。故并齊爲法，以同爲實。計田一百五畝，發用十五人，耕用三十五人，種用二十一人，并之得七十一工，治得一百五畝，故以爲實；而一人一日所治，故以人數爲法除之，卽得也。}\end{unitzh}
\begin{uniten}Method: Set out the numbers of mu for breaking up, ploughing and harrowing; let them cross-multiply the numbers of persons, and sum them to make the divisor. Multiply the numbers of mu together to make the dividend. Divide the dividend by the divisor to obtain [the number of] mu. [Liu Hui: This is like the method of the wild duck and goose.] [Li Chunfeng et al.: In this method, likewise, letting the numbers of mu for breaking up, ploughing, and harrowing and sowing cross-multiply the persons homogenizes the persons; multiplying the numbers of mu together makes the mu common. Therefore the sum of the homogenized [numbers] is the divisor and the common [number] is the dividend. Reckoning a field of 105 mu: breaking up uses 15 persons, ploughing uses 35 persons, sowing uses 21 persons; summing them gives 71 [units of] labor (gong), which work 105 mu; therefore this is taken as the dividend. And [what is sought is] what one person works in one day; therefore the number of persons is taken as the divisor to divide it, and one obtains [the answer].]\end{uniten}

\begin{unitzh}\unitlabel{jiuzhang-028.7}今有池，五渠注之。其一渠開之，少半日一滿；次，一日一满；次，二日半一滿；次，三日一满；次，五日一滿。今皆決之，問幾何日滿池？荅曰：七十四分日之十五。\end{unitzh}
\begin{uniten}Suppose there is a pool into which 5 channels flow. When the first channel is opened, [the pool] fills once in 1/3 day; the next, once in 1 day; the next, once in 2 1/2 days; the next, once in 3 days; the next, once in 5 days. Now if they are all released, in how many days will the pool be filled? Answer: 15/74 day.\end{uniten}

\begin{unitzh}\unitlabel{jiuzhang-028.8}術曰：各置渠一日滿池之數，并以爲法。{\zhsmall 按：此術其一渠少半日滿者，是一日三滿也；次一日一滿；次二日半滿者，是一日五分滿之二也；次三日滿者，是一日三分滿之一也；次五日滿者，是一日五分滿之一也。并之，得四滿十五分滿之十四也。}以一日爲實，實如法得一日。{\zhsmall 此猶矯矢之術也。先令同於一日，日同則滿齊。自鳬鴈至此，其爲同齊有二術焉，可隨率宜也。}\end{unitzh}
\begin{uniten}Method: Set out, for each channel, the number of times it fills the pool in 1 day; sum them to make the divisor. [Liu Hui: Considering this method: the first channel, which fills in 1/3 day, fills 3 times in 1 day; the next fills once in 1 day; the next, which fills in 2 1/2 days, fills 2/5 in 1 day; the next, which fills in 3 days, fills 1/3 in 1 day; the next, which fills in 5 days, fills 1/5 in 1 day. Summing them gives 4 14/15 fillings.] Take 1 day as the dividend; divide the dividend by the divisor to obtain [the number of] days. [Liu Hui: This is like the method for straightening arrows. First make [them] common at 1 day; when the days are common, the fillings are homogenized. From the wild duck and goose up to here, there are two methods for making common and homogenizing; one may follow whichever suits the rates (lü).]\end{uniten}

\begin{unitzh}\unitlabel{jiuzhang-028.9}其一術：列置日數及滿數，{\zhsmall 其一渠少半日滿者，是一日三滿也；次一日一滿；次二日半滿者，是五日二滿；次三日一滿；次五日一滿。此謂之列置日數及滿數也。}令日互相乘滿，并以爲法；日數相乘爲實。實如法得一日。{\zhsmall 亦如鳬鴈術也。}{\zhsmall 臣淳風等謹按：此其一渠少半日滿池者，是一日三滿池也；次一日一滿；次二日半滿者，是五日再滿；次三日一滿；次五日一滿。此謂列置日數於右行，及滿數於左行。以日互乘滿者，齊其滿；日數相乘者，同其日。滿齊而日同，故并齊以除同，卽得也。}\end{unitzh}
\begin{uniten}Another method: Lay out the numbers of days and the numbers of fillings. [Liu Hui: The first channel, which fills in 1/3 day, fills 3 times in 1 day; the next fills once in 1 day; the next, which fills in 2 1/2 days, fills 2 times in 5 days; the next fills once in 3 days; the next fills once in 5 days. This is what is called laying out the numbers of days and the numbers of fillings.] Let the days cross-multiply the fillings, and sum them to make the divisor; multiply the numbers of days together to make the dividend. Divide the dividend by the divisor to obtain [the number of] days. [Liu Hui: This too is like the method of the wild duck and goose.] [Li Chunfeng et al.: Here the first channel, which fills the pool in 1/3 day, fills the pool 3 times in 1 day; the next fills once in 1 day; the next, which fills in 2 1/2 days, fills twice in 5 days; the next fills once in 3 days; the next fills once in 5 days. This means laying out the numbers of days in the right column and the numbers of fillings in the left column. Cross-multiplying the fillings by the days homogenizes the fillings; multiplying the numbers of days together makes the days common. The fillings being homogenized and the days common, therefore the sum of the homogenized [numbers] divides the common [number], and one obtains [the answer].]\end{uniten}

\begin{unitzh}\unitlabel{jiuzhang-028.10}今有人持米出三關，外關三而取一，中關五而取一，内關七而取一，餘米五斗。問本持米幾何？荅曰：十斗九升八分升之三。\end{unitzh}
\begin{uniten}Suppose a person carries rice out through 3 passes. At the outer pass, 1 out of 3 is taken; at the middle pass, 1 out of 5 is taken; at the inner pass, 1 out of 7 is taken. 5 dou (斗) of rice remain. Question: how much rice was originally carried? Answer: 10 dou 9 3/8 sheng (升).\end{uniten}

\begin{unitzh}\unitlabel{jiuzhang-028.11}術曰：置米五斗，以所稅者三之、五之、七之，爲實；以餘不稅者二、四、六互相乘，爲法。實如法得一斗。{\zhsmall 此亦重今有術也。所稅者，謂今所當稅之定。三、五、七皆爲所求率，二、四、六皆爲所有率。置今有餘米五斗，以七乘之，六而一，卽内關未稅之本米也。又以三乘之，二而一，卽外關未稅之本米也。今從末求本，不問中間，故令中率轉相乘而同之，亦如絡絲術。}\end{unitzh}
\begin{uniten}Method: Set out the rice, 5 dou; multiply it by those taxed, by 3, by 5 and by 7, to make the dividend. Multiply together the remainders not taxed, 2, 4 and 6, to make the divisor. Divide the dividend by the divisor to obtain [the number of] dou. [Liu Hui: This too is the repeated 'suppose there is' method (chong jinyou shu). 'Those taxed' refers to the fixed [figures] at which tax is now to be levied. 3, 5 and 7 are all sought rates; 2, 4 and 6 are all given rates. Set out the 5 dou of rice now remaining, multiply it by 7 and divide by 6: this is the original rice at the inner pass, before it was taxed. Again multiply by 3 and divide by 2: this is the original rice at the outer pass, before it was taxed. Now one seeks the beginning from the end and does not ask about what is in between; therefore one lets the middle rates be multiplied successively (zhuan xiang cheng) and makes them common. This is also like the method for raw silk (luo si).]\end{uniten}

\begin{unitzh}\unitlabel{jiuzhang-028.12}{\zhsmall 又一術：外關三而取一，則其餘本米三分之二也。求外關所稅之餘，則當置三分乘之二而一。欲知中關，以四乘之，五而一；欲知内關，以六乘之，七而一。凡餘分者，乘其母，而以三、五、七相乘，得一百五，爲分母；二、四、六相乘，得四十八，爲分子。約而言之，則是餘米於本所持三十五分之十六也。於今有術，餘米五斗爲所有數，分母三十五爲所求率，分子十六爲所有率也。}\end{unitzh}
\begin{uniten}[Liu Hui: Another method: at the outer pass 1 out of 3 is taken; then what remains is 2/3 of the original rice. To find what remains after the tax at the outer pass, one should set out [it], multiply it by 3 parts, and divide by 2 [Translator's note: this clause appears garbled]. To know [the remainder at] the middle pass, multiply by 4 and divide by 5; to know [the remainder at] the inner pass, multiply by 6 and divide by 7. In general, for the remaining fractions, multiply their denominators: multiplying 3, 5 and 7 together gives 105, as the denominator; multiplying 2, 4 and 6 together gives 48, as the numerator. Reduced and expressed, the remaining rice is 16/35 of what was originally carried. In the 'suppose there is' method, the remaining rice, 5 dou, is the given number; the denominator 35 is the sought rate; the numerator 16 is the given rate.]\end{uniten}

\begin{unitzh}\unitlabel{jiuzhang-028.13}今有人持金出五關，前關二而稅一，次關三而稅一，次關四而稅一，次關五而稅一，次關六而稅一。并五關所稅，適重一斤。問本持金幾何？荅曰：一斤三兩四銖五分銖之四。\end{unitzh}
\begin{uniten}Suppose a person carries gold out through 5 passes. At the first pass, 1 out of 2 is taxed; at the next pass, 1 out of 3 is taxed; at the next pass, 1 out of 4 is taxed; at the next pass, 1 out of 5 is taxed; at the next pass, 1 out of 6 is taxed. The taxes of the 5 passes together weigh exactly 1 jin (斤). Question: how much gold was originally carried? Answer: 1 jin 3 liang (兩) 4 4/5 zhu (銖).\end{uniten}

\begin{unitzh}\unitlabel{jiuzhang-028.14}術曰：置一斤，通所稅者以乘之，爲實；亦通其不稅者，以減所通，餘爲法。實如法得一斤。{\zhsmall 此意猶上術也。置一斤通所稅者，謂令二、三、四、五、六相乘，爲分母，七百二十也。通其所不稅者，謂令所稅之餘一、二、三、四、五相乘，爲分子，一百二十也。約而言之，是爲餘金於本所持六分之一也。以子減母，凡五關所稅六分之五也。於今有術，所稅一斤爲所有數，分母六爲所求率，分子五爲所有率。此亦重今有之義。又雖各有率，不問中間，故令中率轉相乘而連除之，卽得也。置一以爲持金之本率，以稅率乘之、除之，則其率亦成積分也。}\end{unitzh}
\begin{uniten}Method: Set out 1 jin; combine (tong) those taxed and multiply it by them, to make the dividend. Also combine those not taxed, and subtract them from what was combined; the remainder is the divisor. Divide the dividend by the divisor to obtain [the number of] jin. [Liu Hui: The idea of this is like the preceding method. 'Set out 1 jin and combine those taxed' means letting 2, 3, 4, 5 and 6 be multiplied together to make the denominator, 720. 'Combine those not taxed' means letting the remainders after the tax, 1, 2, 3, 4 and 5, be multiplied together to make the numerator, 120. Reduced and expressed, this makes the remaining gold 1/6 of what was originally carried. Subtracting the numerator from the denominator, what the 5 passes tax in all is 5/6. In the 'suppose there is' method, the tax of 1 jin is the given number, the denominator 6 is the sought rate, the numerator 5 is the given rate. This too is the sense of the repeated 'suppose there is'. Moreover, although each has its rate, one does not ask about what is in between; therefore one lets the middle rates be multiplied successively and divides continuously (lian chu), and one obtains [the answer]. Set out 1 as the original rate of the gold carried, and multiply and divide it by the tax rates; then its rate too becomes an accumulated fraction (jifen).]\end{uniten}

\begin{unitzh}\unitlabel{jiuzhang-028.15}《九章算術》卷第六終\end{unitzh}
\begin{uniten}End of Juan 6 of the Nine Chapters on the Mathematical Art.\end{uniten}

\begin{unitzh}\unitlabel{jiuzhang-028.16}《九章算術》卷六訂訛{\zhsmall 《算經十書》之二} 休寧戴震東原\end{unitzh}
\begin{uniten}Correction of Errors (Ding'e) for Juan 6 of the Nine Chapters on the Mathematical Art [Note: The second of the Ten Mathematical Classics (Suanjing shishu).] Dai Zhen, Dongyuan, of Xiuning\end{uniten}

\begin{unitzh}\unitlabel{jiuzhang-028.17}臣淳風等謹按：爲衰。於今有術，副并爲所有率，未并者各爲所求率，以賦卒人數爲所有數。此術以别考，則意同，以廣異聞，故存之也。{\zhsmall 「爲衰」二字上有脫文，當云「各令居所及行道曰數約縣卒爲衰」。}\end{unitzh}
\begin{uniten}Your servants [Li] Chunfeng et al. respectfully note: to make the graded shares (shuai). In the 'suppose there is' method, the auxiliary sum (fu bing) is the given rate, those not summed are each sought rates, and the number of soldiers levied is the given number. This method, examined by a separate [route], has the same meaning; to broaden what is heard of differing [approaches], it is therefore preserved. [Note: Above the two characters 爲衰 ('to make the graded shares') there is missing text; it should read: 'for each, let the numbers of days at the station and of travel on the road divide the soldiers of the county, to make the graded shares' (各令居所及行道曰數約縣卒爲衰).] [Translator's note: 曰 here evidently stands for 日 'day'.]\end{uniten}

\begin{unitzh}\unitlabel{jiuzhang-028.18}計經賦之率，旣有戸算之率，亦有遠近、貴賤之率。此二率者，各自相與通。通則甲二十，乙十二，丙七，丁十三，戊五。一斛之費謂之錢率。錢率約戸率者，則錢爲母，戸爲子。子不齊，令母互乘爲齊，卽衰也。若其不然，以一斛之費約戸數取衰，並有分，當通分内子約之，於算甚繁。此一章皆相與通功共率，畧相依似。以上二率、下一率，亦可放此，從其簡易而巳。又以分言之，使甲一戸出二十分斛之一，乙一戸出十八分斛之一，各以戸數乘之，亦可得一縣凡所當輸，俱未衰也。乘之者，乘其子，母報除之。以此觀之，則以一斛之費約戸數者，其意不異矣。然則可置一斛之費而反衰之，約戸以乘戸數爲衰也。合分注曰：「母除爲率，率乘子爲齊。」反衰法曰：「先同其母，各以分母約其子爲反衰。」施其率爲算，旣約且不妨處下也。{\zhsmall 計經賦之率至此多舛誤不可通。}\end{unitzh}
\begin{uniten}In reckoning the rates of the regular levy (jingfu), there are the rates of the household assessment (husuan), and there are also the rates of distance and nearness and of dearness and cheapness. These two [kinds of] rates are each combined with one another. Combined, they are: Jia 20, Yi 12, Bing 7, Ding 13, Wu 5. The cost of 1 hu (斛) is called the cash rate (qianlü). When the cash rate divides the household rate, the cash is the denominator and the households are the numerator. The numerators are not homogenized; let the denominators cross-multiply [them] to homogenize them: these are the graded shares. If it is not done so, but the cost of 1 hu divides the numbers of households to take the graded shares, then all have fractions, and one must convert to parts and add in the numerators (tongfen nazi) and reduce them, which in calculation is very cumbersome. In this one chapter, all [the problems] pool work and share rates (tonggong gonglü) with one another, and roughly resemble one another. With two rates above and one rate below, one may also imitate this, simply following what is simple and easy. Again, speaking in terms of fractions: let one household of Jia contribute 1/20 hu and one household of Yi contribute 1/18 hu, and multiply each by the numbers of households; one can also obtain all that one county should deliver, but none of it is yet graded shares. 'Multiplying them' is multiplying their numerators, the denominators dividing in return (baochu). Seen in this light, dividing the numbers of households by the cost of 1 hu does not differ in meaning. Thus one may set out the costs of 1 hu and grade them inversely, divide the households, and multiply the numbers of households to make the graded shares. The commentary on combining fractions (hefen) says: 'The denominators divide to make the rates; the rates multiply the numerators to make the homogenized [numerators].' The rule for inverse grading (fanshuai) says: 'First make their denominators common; with each denominator divide its numerator to make the inverse graded shares.' Applying these rates in the calculation is both concise and no hindrance to placing [them] below. [Note: From 計經賦之率 ('In reckoning the rates of the regular levy') to here there are many errors, and it cannot be made sense of.]\end{uniten}

\begin{unitzh}\unitlabel{jiuzhang-028.19}術曰：以下三節分四升爲下率，以上四節分三升爲上率。{\zhsmall 謂之率者，以四節、三節爲分母，三升、四升爲分子，母互乘子，得上率九，下率十六，母相乘得十二。}上下率以少減多，餘爲實。{\zhsmall 十六減九，餘七。}置四節、三節，各半之，以減九節，餘爲法。實如法得一升，卽衰相去也。{\zhsmall 以十二通五節半，得六十六，爲一升之率。}\end{unitzh}
\begin{uniten}Method: Divide 4 sheng among the lower 3 sections to make the lower rate; divide 3 sheng among the upper 4 sections to make the upper rate. [Note: As for calling them rates: take 4 sections and 3 sections as the denominators and 3 sheng and 4 sheng as the numerators; cross-multiplying the numerators by the denominators gives the upper rate 9 and the lower rate 16; multiplying the denominators together gives 12.] Of the upper and lower rates, subtract the smaller from the larger; the remainder is the dividend. [Note: 16 minus 9 leaves 7.] Set out the 4 sections and the 3 sections, halve each, and subtract them from the 9 sections; the remainder is the divisor. Divide the dividend by the divisor to obtain [the number of] sheng: this is the amount by which the graded shares differ from one another. [Note: Combining 5 1/2 sections with 12 gives 66, which is the rate for 1 sheng.]\end{uniten}

\section*{卷第七}\addcontentsline{toc}{section}{卷第七}

\begin{unitzh}\unitlabel{jiuzhang-029.1}九章算術卷七{\zhsmall 算經十書之二} 魏劉徽注 唐朝議大夫行太史令上輕車都尉臣李淳風等奉勅注釋\end{unitzh}
\begin{uniten}The Nine Chapters on the Mathematical Art, Juan 7 [Note: the second of the Ten Mathematical Classics]. Commentary by Liu Hui of the Wei. Annotated by imperial command by Your servant Li Chunfeng et al., Grand Master for Court Discussion, acting Grand Astrologer, Senior Commandant of Light Chariots, of the Tang.\end{uniten}

\begin{unitzh}\unitlabel{jiuzhang-029.2}盈不足{\zhsmall 以御隱雜互見}\end{unitzh}
\begin{uniten}Excess and Deficit (ying buzu) [Liu Hui: for dealing with the hidden and the mixed, which appear by turns]\end{uniten}

\begin{unitzh}\unitlabel{jiuzhang-029.3}今有共買物，人出八，盈三；人出七，不足四。問人數、物價各幾何？荅曰：七人，物價五十三。\end{unitzh}
\begin{uniten}Suppose there is a joint purchase of a thing: each person pays 8, and there is an excess of 3; each person pays 7, and there is a deficit of 4. What are the number of people and the price of the thing, respectively? Answer: 7 people; price of the thing, 53.\end{uniten}

\begin{unitzh}\unitlabel{jiuzhang-029.4}今有共買鷄，人出九，盈十一；人出六，不足十六。問人數、鷄價各幾何？苔曰：九人，鷄價七十。\end{unitzh}
\begin{uniten}Suppose there is a joint purchase of chickens: each person pays 9, and there is an excess of 11; each person pays 6, and there is a deficit of 16. What are the number of people and the price of the chickens, respectively? Answer: 9 people; price of the chickens, 70.\end{uniten}

\begin{unitzh}\unitlabel{jiuzhang-029.5}今有共買璡，人出半，盈四；人出少半，不足三。問人數、璡價各幾何？荅曰：四十二人，璡價十七。\end{unitzh}
\begin{uniten}Suppose there is a joint purchase of a jade-stone (jin 璡): each person pays 1/2, and there is an excess of 4; each person pays 1/3, and there is a deficit of 3. What are the number of people and the price of the jade-stone, respectively? Answer: 42 people; price of the jade-stone, 17.\end{uniten}

\begin{unitzh}\unitlabel{jiuzhang-029.6}今有共買牛，七家共出一百九十，不足三百三十；九家共出二百七十，盈三十。問家數、牛價各幾何？荅曰：一百二十六家，牛價三千七百五十。{\zhsmall 按此術并盈不足者爲衆家之差，故以爲實。置所出率，各以家數除之，各得一家所出率，以少減多者，得一家之差。以除，卽家數。以出率乘之，減盈，故得牛價也。}\end{unitzh}
\begin{uniten}Suppose there is a joint purchase of an ox: 7 families together pay 190, and there is a deficit of 330; 9 families together pay 270, and there is an excess of 30. What are the number of families and the price of the ox, respectively? Answer: 126 families; price of the ox, 3750. [Liu Hui: Note that in this method the sum of the excess and the deficit is the difference for all the families; therefore it is taken as the dividend (shi). Set out the rates paid (suochu lü) and divide each by its number of families: each gives the rate paid by one family; subtracting the smaller from the larger gives the difference for one family. Dividing by it gives the number of families. Multiplying by the rate paid and subtracting the excess thus gives the price of the ox.]\end{uniten}

\begin{unitzh}\unitlabel{jiuzhang-029.7}盈不足{\zhsmall 按盈者謂之朓，不足者謂之朒，所出率謂之假令。盈朒維乘兩設者，欲爲齊同之意。}\end{unitzh}
\begin{uniten}Excess and deficit [Liu Hui: Note that the excess is called the "surplus" (tiao 朓), the deficit is called the "shortfall" (nü 朒), and the rates paid are called the "suppositions" (jialing). That the excess and shortfall cross-multiply (weicheng) the two posits (liang she) is intended for the purpose of homogenizing and equalizing (qi tong).]\end{uniten}

\begin{unitzh}\unitlabel{jiuzhang-029.8}術曰：盈不足相與同共買物者，置所出率，盈、不足各居其下。令維乘所出率，并以爲實，并盈、不足爲法。{\zhsmall 據共買物，人出八，盈三；人出七，不足四。齊其假令，同其盈朒，盈朒俱十二。通計齊則不盈不朒之正數，故可并以爲實，并盈不足爲法。齊之三十二者，是四假令，有盈十二；齊之二十一者，是三假令，亦朒十二。并七假令合爲一實，故并三四爲法。}有分者通之。{\zhsmall 若兩設有分者，齊其子，同其母。此問兩設俱見零分，故齊其子，同其母。}副置所出率，以少減多，餘以約法、實。實爲物價，法爲人數。{\zhsmall 令下維乘上訖，以同約之。不可約，故以乘同之。所出率以少減多者，餘謂之設差，以爲少設，則并盈朒是爲定實。故以少設約法則爲人數，約實則爲物價。盈朒當與少設相通，不可徧約，亦當分母乘設差爲約法、實。}\end{unitzh}
\begin{uniten}Method: For excess and deficit together, in jointly buying a thing, set out the rates paid, with the excess and the deficit each placed beneath [its rate]. Have [them] cross-multiply the rates paid, and sum to make the dividend; sum the excess and the deficit to make the divisor (fa). [Liu Hui: Take the joint purchase of a thing in which each person pays 8, with an excess of 3, and each person pays 7, with a deficit of 4. Homogenize the suppositions and equalize the excess and shortfall: the excess and shortfall are both 12. Reckoned together, the homogenized [amounts] are then the correct number (zhengshu), neither in excess nor short; therefore they may be summed to make the dividend, and the excess and deficit summed to make the divisor. The homogenized 32 is 4 suppositions, with an excess of 12; the homogenized 21 is 3 suppositions, likewise short by 12. Summing, the 7 suppositions are combined into one dividend; therefore sum 3 and 4 to make the divisor.] Where there are fractions, make them commensurable (tong 通). [Liu Hui: If the two posits have fractions, homogenize their numerators and equalize their denominators. In this problem both posits show fractional parts; therefore homogenize their numerators and equalize their denominators.] Set out the rates paid again in a separate place; subtract the smaller from the larger, and with the remainder reduce (yue) the divisor and the dividend. The dividend becomes the price of the thing; the divisor becomes the number of people. [Liu Hui: Have the lower [numbers] cross-multiply the upper ones; when this is done, reduce them by the common [denominator] (tong 同). If they cannot be reduced, then equalize them by multiplication. Subtract the smaller of the rates paid from the larger: the remainder is called the difference of the posits (shecha). Taking it as a small posit (shaoshe), the sum of the excess and shortfall is then the fixed dividend (dingshi). Therefore reducing the divisor by the small posit gives the number of people, and reducing the dividend gives the price of the thing. The excess and shortfall should be made commensurable with the small posit; where they cannot all be reduced, one should likewise multiply the difference of the posits by the denominator, to reduce the divisor and the dividend.]\end{uniten}

\begin{unitzh}\unitlabel{jiuzhang-029.9}其一術曰：并盈、不足爲實，以所出率以少減多，餘爲法。實如法得一人。以所出率乘之，減盈、增不足，卽物價。{\zhsmall 此術意謂盈不足爲衆人之差，以所出率以少減多餘爲一人之差，以一人之差約衆人之差，故得人數也。}\end{unitzh}
\begin{uniten}Another method: Sum the excess and the deficit to make the dividend; with the rates paid, subtract the smaller from the larger, and the remainder is the divisor. The dividend divided by the divisor gives 1 person [per divisor]. Multiply [this] by the rate paid, subtract the excess or add the deficit, and this is the price of the thing. [Liu Hui: The idea of this method is that the excess and deficit are the difference for all the people; with the rates paid, subtract the smaller from the larger, and the remainder is the difference for one person; reducing the difference for all the people by the difference for one person therefore gives the number of people.]\end{uniten}

\begin{unitzh}\unitlabel{jiuzhang-029.10}今有共買金，人出四百，盈三千四百；人出三百，盈一百。問人數、金價各幾何？荅曰：三十三人，金價九千八百。\end{unitzh}
\begin{uniten}Suppose there is a joint purchase of gold: each person pays 400, and there is an excess of 3400; each person pays 300, and there is an excess of 100. What are the number of people and the price of the gold, respectively? Answer: 33 people; price of the gold, 9800.\end{uniten}

\begin{unitzh}\unitlabel{jiuzhang-029.11}今有共買羊，人出五，不足四十五；人出七，不足三。問人數、羊價各幾何？荅曰：二十一人，羊價一百五十。\end{unitzh}
\begin{uniten}Suppose there is a joint purchase of sheep: each person pays 5, and there is a deficit of 45; each person pays 7, and there is a deficit of 3. What are the number of people and the price of the sheep, respectively? Answer: 21 people; price of the sheep, 150.\end{uniten}

\begin{unitzh}\unitlabel{jiuzhang-029.12}兩盈兩不足\end{unitzh}
\begin{uniten}Two excesses, two deficits\end{uniten}

\begin{unitzh}\unitlabel{jiuzhang-029.13}術曰：兩盈、兩不足相與同共買物者，置所出率，盈、不足各居其下。令維乘所出率，以少減多，餘爲實。兩盈、兩不足以少減多，餘爲法。有分者通之。副置所出率，以少減多，餘以約法、實。實爲物價，法爲人數。\end{unitzh}
\begin{uniten}Method: For two excesses or two deficits together, in jointly buying a thing, set out the rates paid, with the excess or deficit each placed beneath [its rate]. Have [them] cross-multiply the rates paid; subtract the smaller from the larger, and the remainder is the dividend. Of the two excesses or two deficits, subtract the smaller from the larger; the remainder is the divisor. Where there are fractions, make them commensurable. Set out the rates paid again in a separate place; subtract the smaller from the larger, and with the remainder reduce the divisor and the dividend. The dividend becomes the price of the thing; the divisor becomes the number of people.\end{uniten}

\begin{unitzh}\unitlabel{jiuzhang-029.14}{\zhsmall 按此術兩盈者，兩設皆逾於正數；兩不足者，兩設皆不足於正數。其所以變化，則兩不足亦猶兩盈。而或有勢同而情違者，當其爲實，俱令不足，其相乘減，則遺其所不足焉。故其餘所以爲實者，無朒數以損焉。葢出而有餘，兩盈兩設皆逾於正數。假令與共買物，人出八，盈三；人出九，盈十。齊其假令，同其兩盈，兩盈俱三十。舉齊則兼去，其餘所以爲實者，無盈數。兩盈以少減多，餘爲法。齊之八十者，是十假令，而凡盈三十，是三以十齊之；齊之二十七者，是三假令，而凡盈三十，是十以三齊之。兩盈共十三，以三減十，餘七假令爲一實，故令以三減十，餘七爲法。所出率以少減多，餘謂之設差，因設差爲少設，則兩盈之差是爲定實。故以少設約法則爲人數，約實則得物價。}\end{unitzh}
\begin{uniten}[Liu Hui: Note that in this method, with two excesses both posits exceed the correct number; with two deficits both posits fall short of the correct number. As for the way they are transformed, two deficits are just like two excesses. Yet there are cases where the form is the same but the substance differs: when they are made into the dividend, both are made deficient; when they are multiplied and subtracted from each other, what they lack is left aside. Therefore in the remainder that is taken as the dividend there is no shortfall to diminish it. For when what is paid leaves something over, with two excesses both posits exceed the correct number. Suppose, in joining a joint purchase of a thing, each person pays 8, with an excess of 3, and each person pays 9, with an excess of 10. Homogenize the suppositions and equalize the two excesses: the two excesses are both 30. When all are homogenized, [the excesses] are removed together, and in the remainder that is taken as the dividend there is no excess. Of the two excesses, subtract the smaller from the larger; the remainder is the divisor. The homogenized 80 is 10 suppositions, and the total excess of 30 is 3 homogenized by 10; the homogenized 27 is 3 suppositions, and the total excess of 30 is 10 homogenized by 3. The two excesses together are 13; subtracting 3 from 10, the remaining 7 suppositions make one dividend; therefore subtract 3 from 10 and take the remaining 7 as the divisor. Subtract the smaller of the rates paid from the larger; the remainder is called the difference of the posits. Taking the difference of the posits as a small posit, the difference of the two excesses is then the fixed dividend. Therefore reducing the divisor by the small posit gives the number of people, and reducing the dividend gives the price of the thing.]\end{uniten}

\begin{unitzh}\unitlabel{jiuzhang-029.15}其一術曰：置所出率，以少減多，餘爲法。兩盈、兩不足以少減多，餘爲實。實如法而一，得人數。以所出率乘之，減盈、増不足，卽物價。{\zhsmall 置所出率，以少減多，得一人之差。兩盈、兩不足相減，餘爲衆人之差。故以一人之差除之，得人數。以所出率乘之，減盈、增不足，卽物價。}\end{unitzh}
\begin{uniten}Another method: Set out the rates paid; subtract the smaller from the larger, and the remainder is the divisor. Of the two excesses or two deficits, subtract the smaller from the larger; the remainder is the dividend. Divide the dividend by the divisor, obtaining the number of people. Multiply [this] by the rate paid, subtract the excess or add the deficit, and this is the price of the thing. [Liu Hui: Set out the rates paid; subtracting the smaller from the larger gives the difference for one person. Subtract the two excesses or two deficits from each other; the remainder is the difference for all the people. Therefore dividing by the difference for one person gives the number of people. Multiply [this] by the rate paid, subtract the excess or add the deficit, and this is the price of the thing.]\end{uniten}

\begin{unitzh}\unitlabel{jiuzhang-029.16}今有共買豕，人出一百，盈一百；人出九十，適足。問人數、豕價各幾何？荅曰：一十人，豕價九百。\end{unitzh}
\begin{uniten}Suppose there is a joint purchase of a pig: each person pays 100, and there is an excess of 100; each person pays 90, and it is exactly sufficient (shizu). What are the number of people and the price of the pig, respectively? Answer: 10 people; price of the pig, 900.\end{uniten}

\begin{unitzh}\unitlabel{jiuzhang-029.17}今有共買犬，人出五，不足九十；人出五十，適足。問人數、犬價各幾何？荅曰：二人，犬價一百。\end{unitzh}
\begin{uniten}Suppose there is a joint purchase of a dog: each person pays 5, and there is a deficit of 90; each person pays 50, and it is exactly sufficient. What are the number of people and the price of the dog, respectively? Answer: 2 people; price of the dog, 100.\end{uniten}

\begin{unitzh}\unitlabel{jiuzhang-029.18}盈適足、不足適足\end{unitzh}
\begin{uniten}Excess and exactly sufficient; deficit and exactly sufficient\end{uniten}

\begin{unitzh}\unitlabel{jiuzhang-029.19}術曰：以盈及不足之數爲實。置所出率，以少減多，餘爲法。實如法得一人。其求物價者，以適足乘人數，得物價。{\zhsmall 盈不足數爲實者，數單見，卽衆人差，故以爲實。所出率以少減多者，餘卽一人差，故以爲法。以除衆人差，得人數。以適足乘人數，卽得物價也。此術意謂所出率以少減多者，餘是一人不足之差，不足數爲衆人之差，以一人差約之，故得人之數也。}\end{unitzh}
\begin{uniten}Method: Take the amount of the excess or of the deficit as the dividend. Set out the rates paid; subtract the smaller from the larger, and the remainder is the divisor. The dividend divided by the divisor gives 1 person [per divisor]. To find the price of the thing, multiply the number of people by the exactly sufficient [amount], obtaining the price of the thing. [Liu Hui: As for the amount of the excess or deficit being the dividend: the amount appears singly and is the difference for all the people; therefore it is taken as the dividend. Subtracting the smaller of the rates paid from the larger, the remainder is the difference for one person; therefore it is taken as the divisor. Dividing the difference for all the people by it gives the number of people. Multiplying the number of people by the exactly sufficient [amount] gives the price of the thing. The idea of this method is that, subtracting the smaller of the rates paid from the larger, the remainder is the difference of one person's deficit, and the amount of the deficit is the difference for all the people; reducing it by the difference for one person therefore gives the number of the people.]\end{uniten}

\begin{unitzh}\unitlabel{jiuzhang-029.20}今有米在十斗桶中，不知其數。滿中添粟而舂之，得米七斗。問故米幾何？荅曰：二斗五升。\end{unitzh}
\begin{uniten}Suppose there is rice in a 10-dou (斗) bucket, its amount unknown. It is filled up by adding millet (su 粟), and this is pounded, yielding 7 dou of rice. How much was the original rice? Answer: 2 dou 5 sheng (升).\end{uniten}

\begin{unitzh}\unitlabel{jiuzhang-029.21}術曰：以盈不足術求之。假令故米二斗，不足二升；令之三斗，有餘二升。{\zhsmall 按桶受一斛。若使故米二斗，須添粟八斗以滿之。八斗得糲米四斗八升，課於七斗，是爲不足二升。若使故米三斗，須添粟七斗以滿之。七斗得糲米四斗二升，課於七斗，是爲有餘二升。以盈不足維乘假令之數者，欲爲齊同之意。實如法，卽得故米斗數，乃不盈不朒之正數也。}\end{unitzh}
\begin{uniten}Method: Seek it by the method of excess and deficit. Suppose the original rice is 2 dou: there is a deficit of 2 sheng; let it be 3 dou: there are 2 sheng left over. [Liu Hui: Note that the bucket holds 1 hu (斛). If the original rice were 2 dou, 8 dou of millet would have to be added to fill it. 8 dou gives 4 dou 8 sheng of coarse rice (limi 糲米); measured against 7 dou, this makes a deficit of 2 sheng. If the original rice were 3 dou, 7 dou of millet would have to be added to fill it. 7 dou gives 4 dou 2 sheng of coarse rice; measured against 7 dou, this makes 2 sheng left over. Cross-multiplying the numbers of the suppositions by the excess and deficit is intended for the purpose of homogenizing and equalizing. Dividing the dividend by the divisor then gives the number of dou of the original rice, which is the correct number, neither in excess nor short.]\end{uniten}

\begin{unitzh}\unitlabel{jiuzhang-030.1}今有垣高九尺，瓜生其上，蔓日長七寸；瓠生其下，蔓日長一尺。問：幾何日相逢？瓜、瓠各長幾何？荅曰：五日十七分日之五。瓜長三尺七寸十七分寸之一。瓠長五尺二寸十七分寸之十六。\end{unitzh}
\begin{uniten}Suppose there is a wall 9 chi (尺) high. A melon grows on top of it, and its vine grows 7 cun (寸) a day; a bottle gourd grows at its foot, and its vine grows 1 chi a day. Question: in how many days do they meet? How much does each of the melon and the bottle gourd grow? Answer: 5 5/17 days. The melon grows 3 chi 7 1/17 cun. The bottle gourd grows 5 chi 2 16/17 cun.\end{uniten}

\begin{unitzh}\unitlabel{jiuzhang-030.2}術曰：假令五日，不足五寸；令之六日，有餘一尺二寸。\end{unitzh}
\begin{uniten}Method: Assume (jialing) 5 days: a deficit of 5 cun. Let it be 6 days: an excess of 1 chi 2 cun.\end{uniten}

\begin{unitzh}\unitlabel{jiuzhang-030.3}{\zhsmall 按：假令五日不足五寸者，瓜生五日，下垂蔓三尺五寸；⿰夸𤓰 -- 瓠生五日，上延蔓五尺；課於九尺之垣，是爲不足五寸。令之六日有餘一尺二寸者，若使瓜生六日，下垂蔓四尺二寸；⿰夸𤓰 -- 瓠生六日，上延蔓六尺；課於九尺之垣，是爲有餘一尺二寸。以盈不足維乘假令之數者，欲爲齊同之意。實如法而一，卽設差不盈不朒之正數，卽得日數。以瓜、瓠一日之長乘之，故各得其長之數也。}\end{unitzh}
\begin{uniten}[Liu Hui: Note: as for "assume 5 days: a deficit of 5 cun": the melon, growing for 5 days, has a vine hanging down 3 chi 5 cun; the bottle gourd [Translator's note: here and below the character is transcribed as the variant ⿰夸𤓰 followed by "-- 瓠"], growing for 5 days, has a vine extending upward 5 chi; compared with the 9-chi wall, this is a deficit of 5 cun. As for "let it be 6 days: an excess of 1 chi 2 cun": if the melon grows for 6 days, its vine hangs down 4 chi 2 cun; the bottle gourd, growing for 6 days, has a vine extending upward 6 chi; compared with the 9-chi wall, this is an excess of 1 chi 2 cun. Cross-multiplying (weicheng) the assumed numbers by the excess and deficit is intended as the idea of homogenizing and equalizing (qitong). Divide the dividend (shi) by the divisor (fa): this is the correct number at which the posited differences are neither in excess nor in shortfall (bu ying bu nü), and one thus obtains the number of days. Multiply it by the growth of the melon and of the bottle gourd in 1 day; therefore one obtains the number of each one's length.]\end{uniten}

\begin{unitzh}\unitlabel{jiuzhang-030.4}今有蒲生一日長三尺，莞生一日長一尺。蒲生日自半，莞生日自倍。問：幾何日而長等？荅曰：二日十三分日之六。各長四尺八寸十三分寸之六。\end{unitzh}
\begin{uniten}Suppose there is a cattail (pu) that grows 3 chi in 1 day, and a rush (guan) that grows 1 chi in 1 day. The cattail's growth halves itself each day; the rush's growth doubles itself each day. Question: after how many days are their lengths equal? Answer: 2 6/13 days. Each is 4 chi 8 6/13 cun long.\end{uniten}

\begin{unitzh}\unitlabel{jiuzhang-030.5}術曰：假令二日，不足一尺五寸；令之三日，有餘一尺七寸半。\end{unitzh}
\begin{uniten}Method: Assume 2 days: a deficit of 1 chi 5 cun. Let it be 3 days: an excess of 1 chi 7 1/2 cun.\end{uniten}

\begin{unitzh}\unitlabel{jiuzhang-030.6}{\zhsmall 按：假令二日不足一尺五寸者，蒲生二日，長四尺五寸；莞生二日，長三尺；是爲未相及一尺五寸，故曰不足。令之三日有餘一尺七寸半者，蒲增前七寸半，莞增前四尺，是爲過一尺七寸半，故日有餘。以盈不足乘除之。又以後一日所長各乘日分子，如日分母而一者，各得日分子之長也。故各增二長，卽得其數。}\end{unitzh}
\begin{uniten}[Liu Hui: Note: as for "assume 2 days: a deficit of 1 chi 5 cun": the cattail, growing for 2 days, is 4 chi 5 cun long; the rush, growing for 2 days, is 3 chi long; this means they have not yet caught up with each other, by 1 chi 5 cun; therefore it says "deficit". As for "let it be 3 days: an excess of 1 chi 7 1/2 cun": the cattail increases on the previous by 7 1/2 cun, the rush increases on the previous by 4 chi; this means going beyond by 1 chi 7 1/2 cun; therefore it says "excess" [Translator's note: the text writes 日 "day" where 曰 "it says" is evidently meant]. Multiply and divide it by excess and deficit. Further, as for multiplying the numerator of the day-fraction (ri fenzi) by what each grows in the later 1 day and dividing by the denominator of the day-fraction (ri fenmu): this gives for each the growth in the numerator of the day-fraction. Therefore, increasing each by the 2 lengths, one obtains their numbers.]\end{uniten}

\begin{unitzh}\unitlabel{jiuzhang-030.7}今有垣厚五尺，兩鼠對穿。大鼠日一尺，小鼠亦日一尺。大鼠日自倍，小鼠日自半。問：幾何日相逢？各穿幾何？荅曰：二日十七分日之二。大鼠穿三尺四寸十七分寸之十二。小鼠穿一尺五寸十七分寸之五。\end{unitzh}
\begin{uniten}Suppose there is a wall 5 chi thick, and 2 rats bore through it from opposite sides. The large rat: 1 chi a day; the small rat: also 1 chi a day. The large rat doubles itself each day; the small rat halves itself each day. Question: in how many days do they meet? How much does each bore? Answer: 2 2/17 days. The large rat bores 3 chi 4 12/17 cun. The small rat bores 1 chi 5 5/17 cun.\end{uniten}

\begin{unitzh}\unitlabel{jiuzhang-030.8}術曰：假令二日，不足五寸；令之三日，有餘三尺七寸半。大鼠日倍，二日合穿三尺；小鼠日自半，合穿一尺五寸；并大鼠所穿，合尺五寸；課於垣厚五尺，是爲不足五寸。令之三日，大鼠穿得七尺，小鼠穿得一尺七寸半，并之以減垣厚五尺，有餘三尺七寸半。以盈不足術求之，卽得。以後一日所穿乘日分子，如日分母而一，卽各得日分子之中所穿。故各増二日定穿，卽合所問也。\end{unitzh}
\begin{uniten}Method: Assume 2 days: a deficit of 5 cun. Let it be 3 days: an excess of 3 chi 7 1/2 cun. The large rat doubles each day: in 2 days it bores altogether 3 chi; the small rat halves itself each day: altogether it bores 1 chi 5 cun; adding what the large rat bores, together chi 5 cun [Translator's note: a numeral is evidently missing before chi]; compared with the wall thickness of 5 chi, this is a deficit of 5 cun. Let it be 3 days: the large rat bores 7 chi, the small rat bores 1 chi 7 1/2 cun; add them and subtract the wall thickness of 5 chi: there is an excess of 3 chi 7 1/2 cun. Seek it by the method of excess and deficit, and one obtains it. Multiply the numerator of the day-fraction by what is bored in the later 1 day and divide by the denominator of the day-fraction: then one obtains for each what is bored within the numerator of the day-fraction. Therefore, increasing each by the fixed boring of 2 days, it then agrees with what is asked.\end{uniten}

\begin{unitzh}\unitlabel{jiuzhang-030.9}今有醇酒一斗，直錢五十；行酒一斗，直錢一十。今將錢三十，得酒二斗。問：醇、行酒各得幾何？荅曰：醇酒二升半，行酒一斗七升半。\end{unitzh}
\begin{uniten}Suppose there is pure wine (chunjiu): 1 dou (斗) is worth 50 qian (錢); ordinary wine (xingjiu): 1 dou is worth 10 qian. Now one takes 30 qian and obtains 2 dou of wine. Question: how much of the pure and of the ordinary wine does one obtain? Answer: pure wine 2 1/2 sheng (升); ordinary wine 1 dou 7 1/2 sheng.\end{uniten}

\begin{unitzh}\unitlabel{jiuzhang-030.10}術曰：假令醇酒五升，行酒一斗五升，有餘一十；令之醇酒二升，行酒一斗八升，不足二。\end{unitzh}
\begin{uniten}Method: Assume pure wine 5 sheng and ordinary wine 1 dou 5 sheng: an excess of 10. Let it be pure wine 2 sheng and ordinary wine 1 dou 8 sheng: a deficit of 2.\end{uniten}

\begin{unitzh}\unitlabel{jiuzhang-030.11}{\zhsmall 據醇酒五升直錢二十五，行酒一斗五升直錢十五，課於三十，是爲有餘十。據醇酒二升直錢一十，行酒一斗八升直錢一十八，課於三十，是爲不足二。以盈不足術求之。此問已有重說及其齊同之意也。}\end{unitzh}
\begin{uniten}[Liu Hui: Taking pure wine 5 sheng, worth 25 qian, and ordinary wine 1 dou 5 sheng, worth 15 qian: compared with 30, this is an excess of 10. Taking pure wine 2 sheng, worth 10 qian, and ordinary wine 1 dou 8 sheng, worth 18 qian: compared with 30, this is a deficit of 2. Seek it by the method of excess and deficit. This problem has already had a repeated explanation, together with its idea of homogenizing and equalizing.]\end{uniten}

\begin{unitzh}\unitlabel{jiuzhang-030.12}今有大器五、小器一，容三斛；大器一、小器五，容二斛。問：大、小器各容幾何？荅曰：大器容二十四分斛之十三，小器二十四分斛之七。\end{unitzh}
\begin{uniten}Suppose there are 5 large vessels and 1 small vessel, which hold 3 hu (斛); 1 large vessel and 5 small vessels, which hold 2 hu. Question: how much does each of the large and the small vessels hold? Answer: a large vessel holds 13/24 hu; a small vessel 7/24 hu.\end{uniten}

\begin{unitzh}\unitlabel{jiuzhang-030.13}術曰：假令大器五斗，小器亦五斗，盈十斗；令之大器五斗五升，小器二斗五升，不足二斗。\end{unitzh}
\begin{uniten}Method: Assume the large vessel 5 dou and the small vessel also 5 dou: an excess of 10 dou. Let it be the large vessel 5 dou 5 sheng and the small vessel 2 dou 5 sheng: a deficit of 2 dou.\end{uniten}

\begin{unitzh}\unitlabel{jiuzhang-030.14}{\zhsmall 按：大器容五斗，大器五容二斛五斗，以減三斛，餘五斗，卽小器一所容，故曰小器亦五斗。小器五容二斛五斗，大器一，合爲三斛，課於兩斛，乃多十斗。令之大器五斗五升，大器五合容二斛七斗五升，以減三斛，餘二斗五升，卽小器一所容，故曰小器二斗五升。大器一容五斗五升，小器五合容一斛二斗五升，合爲一斛八斗，課於二斛，少二斗，故曰不足二斗。以盈不足維乘之，各并爲實，并盈不足爲法，除之。}\end{unitzh}
\begin{uniten}[Liu Hui: Note: a large vessel holds 5 dou; 5 large vessels hold 2 hu 5 dou; subtracting this from 3 hu leaves 5 dou, which is what 1 small vessel holds; therefore it says "the small vessel also 5 dou". 5 small vessels hold 2 hu 5 dou; with 1 large vessel, together this makes 3 hu; compared with 2 hu, it is then 10 dou too much. Let the large vessel be 5 dou 5 sheng: 5 large vessels together hold 2 hu 7 dou 5 sheng; subtracting this from 3 hu leaves 2 dou 5 sheng, which is what 1 small vessel holds; therefore it says "the small vessel 2 dou 5 sheng". 1 large vessel holds 5 dou 5 sheng, 5 small vessels together hold 1 hu 2 dou 5 sheng; together this makes 1 hu 8 dou; compared with 2 hu, it is 2 dou short; therefore it says "a deficit of 2 dou". Cross-multiply them by the excess and deficit, add each to make a dividend (shi), add the excess and deficit to make the divisor (fa), and divide.]\end{uniten}

\begin{unitzh}\unitlabel{jiuzhang-030.15}今有漆三得油四，油四和漆五。今有漆三斗，欲令分以易油，還自和餘漆。問：出漆、得油、和漆各幾何？荅曰：出漆一斗一升四分升之一，得油一斗五升，和漆一斗八升四分升之三。\end{unitzh}
\begin{uniten}Suppose there is lacquer: 3 of it obtain 4 of oil, and 4 of oil mix with 5 of lacquer. Suppose there are 3 dou of lacquer; one wishes to divide it so as to exchange for oil, which in turn is mixed with the remaining lacquer. Question: how much is each of the lacquer paid out, the oil obtained and the lacquer mixed? Answer: lacquer paid out 1 dou 1 1/4 sheng; oil obtained 1 dou 5 sheng; lacquer mixed 1 dou 8 3/4 sheng.\end{uniten}

\begin{unitzh}\unitlabel{jiuzhang-030.16}術曰：假令出漆九升，不足六升；令之出漆一斗二升，有餘二升。\end{unitzh}
\begin{uniten}Method: Assume the lacquer paid out is 9 sheng: a deficit of 6 sheng. Let the lacquer paid out be 1 dou 2 sheng: an excess of 2 sheng.\end{uniten}

\begin{unitzh}\unitlabel{jiuzhang-030.17}{\zhsmall 按：此術三斗之漆，出九升，得油一斗二升，可和漆一斗五升，餘有二斗一升，則六升無油可和，故曰不足六升。令之出漆一斗二升，則易得油一斗六升，可和漆二斗。於三斗之中已出一斗二升，餘有一斗八升，見在油合和得漆二斗，則是有餘二升。以盈不足維乘之爲實，并盈不足爲法，實如法而一，得出漆升數。求油及和漆者，四、五各爲所求率，四、三各爲所有率，而今有之，卽得也。}\end{unitzh}
\begin{uniten}[Liu Hui: Note: in this method, of the 3 dou of lacquer, 9 sheng are paid out, obtaining 1 dou 2 sheng of oil, which can mix 1 dou 5 sheng of lacquer; there remain 2 dou 1 sheng, so 6 sheng have no oil with which they can be mixed; therefore it says "a deficit of 6 sheng". Let the lacquer paid out be 1 dou 2 sheng: then one obtains in exchange 1 dou 6 sheng of oil, which can mix 2 dou of lacquer. Of the 3 dou, 1 dou 2 sheng have already been paid out, and there remain 1 dou 8 sheng; the oil now at hand, combined in mixing, obtains 2 dou of lacquer; so this is an excess of 2 sheng. Cross-multiply them by the excess and deficit to make the dividend (shi); add the excess and deficit to make the divisor (fa); divide the dividend by the divisor, and one obtains the number of sheng of lacquer paid out. As for finding the oil and the mixed lacquer: 4 and 5 each serve as the sought rate (suoqiu lü), 4 and 3 each as the given rate (suoyou lü), and applying "suppose there is" (jinyou) to them, one obtains them.]\end{uniten}

\begin{unitzh}\unitlabel{jiuzhang-030.18}今有玉方一寸，重七兩；石方一寸，重六兩。今有石立方三寸，中有玉，并重十一斤。問：玉、石重各幾何？荅曰：玉十四寸，重六斤二兩；石十三寸，重四斤十四兩。\end{unitzh}
\begin{uniten}Suppose there is jade 1 cun square (fang), weighing 7 liang (兩), and stone 1 cun square, weighing 6 liang. Suppose there is a stone, a cube (lifang) of 3 cun, with jade inside it, weighing together 11 jin (斤). Question: what is the weight of each of the jade and the stone? Answer: jade 14 cun, weighing 6 jin 2 liang; stone 13 cun, weighing 4 jin 14 liang.\end{uniten}

\begin{unitzh}\unitlabel{jiuzhang-030.19}術曰：假令皆玉，多十三兩；令之皆石，不足十四兩。不足爲玉，多爲石。各以一寸之重乘之，得玉、石之積重。\end{unitzh}
\begin{uniten}Method: Assume it is all jade: 13 liang too much. Let it be all stone: a deficit of 14 liang. The deficit makes the jade; the "too much" makes the stone. Multiply each by the weight of 1 cun, and one obtains the total weights (jizhong) of the jade and the stone.\end{uniten}

\begin{unitzh}\unitlabel{jiuzhang-030.20}{\zhsmall 立方三寸，是一靣之方，計積二十七寸。玉方一寸重七兩，石方一寸重六兩，是爲玉、石重差一兩。假令皆玉，合有一百八十九兩，課於十一斤，有餘十三兩。玉重而石輕，故有此多。卽二十七寸之中有十三寸，寸損一兩，則以爲石重，故言多爲石。言多之數，出於石以爲玉。假令皆石，合有一百六十二兩，課於十一斤，少十四兩，故曰不足。此不足卽以重爲輕，故令減少數於石重，卽二十七寸之中有十四寸，寸增一兩，則以爲玉重也。}\end{unitzh}
\begin{uniten}[Liu Hui: "A cube of 3 cun" is the square of 1 face; reckoning the volume, it is 27 cun. Jade 1 cun square weighs 7 liang, stone 1 cun square weighs 6 liang: this makes the difference in weight between jade and stone 1 liang. Assume it is all jade: altogether there are 189 liang; compared with 11 jin, there is an excess of 13 liang. Jade is heavy and stone light, therefore there is this "too much". That is, among the 27 cun there are 13 cun which, each cun reduced by 1 liang, are then taken as the weight of stone; therefore it says "the 'too much' makes the stone". The number called "too much" comes from stone having been taken as jade. Assume it is all stone: altogether there are 162 liang; compared with 11 jin, it is 14 liang short; therefore it says "deficit". This deficit is from taking the heavy as light; therefore it makes the weight of the stone reduced by the number short; that is, among the 27 cun there are 14 cun which, each cun increased by 1 liang, are then taken as the weight of jade.]\end{uniten}

\begin{unitzh}\unitlabel{jiuzhang-031.1}今有善田一畝價三百，惡田七畝價五百。今并買一頃，價錢一萬。問善、惡田各幾何？荅曰：善田十二畝半，惡田八十七畝半。\end{unitzh}
\begin{uniten}Suppose there is good field, 1 mu (畝) priced at 300, and bad field, 7 mu priced at 500. Now 1 qing (頃) is bought altogether, the price being 10000 qian (錢). Question: how much is the good field and the bad field each? Answer: good field 12 1/2 mu; bad field 87 1/2 mu.\end{uniten}

\begin{unitzh}\unitlabel{jiuzhang-031.2}術曰：假令善田二十畝，惡田八十畝，多一千七百一十四錢七分錢之二；令之善田十畝，惡田九十畝，不足五百七十一錢七分錢之三。{\zhsmall 按：善田二十畝直錢六千，惡田八十畝直錢五千七百一十四七分錢之二，謂於一萬，是多一千七百一十四錢七分錢之二。令之善田十畝直錢三千，惡田九十畝直錢六千四百二十八七分錢之四，課於一萬，是爲不足五百七十一錢七分錢之三。以盈不足術求之也。}\end{unitzh}
\begin{uniten}Method: Suppose good field 20 mu and bad field 80 mu: excess 1714 2/7 qian; let it instead be good field 10 mu and bad field 90 mu: deficit 571 3/7 qian. [Liu Hui: Examining: 20 mu of good field are worth 6000 qian, and 80 mu of bad field are worth 5714 2/7 qian; that is, with respect to 10000, this is an excess of 1714 2/7 qian. Let it instead be 10 mu of good field, worth 3000 qian, and 90 mu of bad field, worth 6428 4/7 qian; compared with 10000, this makes a deficit of 571 3/7 qian. Seek it by the method of excess and deficit.]\end{uniten}

\begin{unitzh}\unitlabel{jiuzhang-031.3}今有黃金九枚，白銀十一枚，稱之重適等。交易其一，金輕十三兩。問金、銀一枚各重幾何？荅曰：金重二斤三兩十八銖，銀重一斤十三兩六銖。\end{unitzh}
\begin{uniten}Suppose there are 9 pieces (mei 枚) of yellow gold and 11 pieces of white silver; weighed, their weights are exactly equal. When one of them is exchanged, the gold is lighter by 13 liang (兩). Question: what does 1 piece of gold and of silver each weigh? Answer: the gold weighs 2 jin (斤) 3 liang 18 zhu (銖); the silver weighs 1 jin 13 liang 6 zhu.\end{uniten}

\begin{unitzh}\unitlabel{jiuzhang-031.4}術曰：假令黃金三斤，白銀二斤十一分斤之五，不足四十九，於右行；令之黃金二斤，白銀一斤十一分斤之七，多十五，於左行。以分母各乘其行内之數，以盈不足維乘所出率，并以爲實；并盈不足爲法。實如法得黃金重；分母乘法以除，得銀重。約之得分也。\end{unitzh}
\begin{uniten}Method: Suppose the yellow gold is 3 jin and the white silver 2 5/11 jin: deficit 49, in the right column (hang 行); let it instead be yellow gold 2 jin and white silver 1 7/11 jin: excess 15, in the left column. Multiply each of the numbers within its column by the denominator; with the excess and deficit cross-multiply (weicheng 維乘) the rates (lü) set out, and add them to make the dividend (shi); add the excess and deficit to make the divisor (fa). Dividing the dividend by the divisor gives the weight of the yellow gold; multiply the divisor by the denominator and divide, which gives the weight of the silver. Reduce (yue 約) them to obtain the fractions.\end{uniten}

\begin{unitzh}\unitlabel{jiuzhang-031.5}{\zhsmall 按：此術假令黃金九、白銀十一俱重二十七斤，金九約之得三斤，銀十一約之得二斤十一分斤之五，各爲金、銀一枚重數。就金重二十七斤之中減一金之重以益銀，銀重二十七斤之中減一銀之重以益金，則金重二十六斤十一分斤之五，銀重二十七斤十一分斤之六。以少減多，則金輕十七兩十一分兩之五，課於十三兩，多四兩十一分兩之五，通分内子言之，是爲不足四十九。又令之黃金九，一枚重二斤，九枚重十八斤，白銀十一亦合重十八斤也，乃以十一除之，得一枚一斤十一分斤之七，爲銀一枚之重數。今就金重十八斤之中減一枚金以益銀，復減一枚銀以益金，則金重十七斤十一分斤之七，銀重十八斤十一分斤之四。以少減多，卽金輕十一分斤之八，課於十三兩，少一兩十一分兩之四，通分内子言之，是爲多十五。以盈不足術求之，實如法得金重。「分母乘法以除」者，謂銀兩分母同，須通法而後乃除，得銀重。餘皆約之者，術省故也。}\end{unitzh}
\begin{uniten}[Liu Hui: Examining: in this method, suppose that the yellow gold, 9, and the white silver, 11, both weigh 27 jin; the gold reduced by 9 gives 3 jin, and the silver reduced by 11 gives 2 5/11 jin; these are respectively the weights of 1 piece of gold and of silver. Taking the gold's weight of 27 jin, subtract from it the weight of one gold to augment the silver; from the silver's weight of 27 jin subtract the weight of one silver to augment the gold; then the gold weighs 26 5/11 jin and the silver weighs 27 6/11 jin. Subtracting the smaller from the larger, the gold is lighter by 17 5/11 liang; compared with 13 liang, it is more by 4 5/11 liang; expressed by converting to parts and incorporating the numerator (tongfen nazi 通分内子), this makes the deficit 49. Again, let the yellow gold be 9, one piece weighing 2 jin, 9 pieces weighing 18 jin; the white silver, 11, likewise ought to weigh 18 jin; then divide by 11, which gives for one piece 1 7/11 jin, as the weight of one piece of silver. Now, taking the gold's weight of 18 jin, subtract from it one piece of gold to augment the silver, and again subtract one piece of silver to augment the gold; then the gold weighs 17 7/11 jin and the silver weighs 18 4/11 jin. Subtracting the smaller from the larger, the gold is then lighter by 8/11 jin; compared with 13 liang, it is less by 1 4/11 liang; expressed by converting to parts and incorporating the numerator, this makes the excess 15. Seek it by the method of excess and deficit: dividing the dividend by the divisor gives the weight of the gold. As for "multiply the divisor by the denominator and divide": this means that the two denominators of the silver are the same, so one must convert the divisor (tongfa 通法) and only then divide, which gives the weight of the silver. As for reducing all the remainders: this is because the method is abbreviated.]\end{uniten}

\begin{unitzh}\unitlabel{jiuzhang-031.6}今有良馬與駑馬發長安至齊，齊去長安三千里。良馬初日行一百九十三里，日增十三里；駑馬初日行九十七里，日減半里。良馬先至齊，復還迎駑馬。問幾何日相逢及各行幾何？荅曰：十五日一百九十一分日之一百三十五而相逢。良馬行四千五百三十四里一百九十一分里之四十六。駑馬行一千四百六十五里一百九十一分里之一百四十五。\end{unitzh}
\begin{uniten}Suppose there are a good horse and a nag that set out from Chang'an for Qi; Qi is 3000 li (里) from Chang'an. The good horse travels 193 li on the first day and increases by 13 li each day; the nag travels 97 li on the first day and decreases by 1/2 li each day. The good horse reaches Qi first and turns back to meet the nag. Question: after how many days do they meet, and how far does each travel? Answer: they meet after 15 135/191 days. The good horse travels 4534 46/191 li. The nag travels 1465 145/191 li.\end{uniten}

\begin{unitzh}\unitlabel{jiuzhang-031.7}術曰：假令十五日，不足三百三十七里半；令之十六日，多一百四十里。以盈不足維乘假令之數，并而爲實；并盈不足爲法。實如法而一，得日數。不盡者，以等數除之而命分。{\zhsmall 求良馬行者：十四乘益疾里數而半之，加良馬初日之行里數，以乘日分子，如日分母而一，所得及其不盡而命分，加於前良馬凡行里數，卽得。求駑馬行者：以十四乘半里，以減駑馬初日之行里數，餘以并初日之行，又半之，乘十五日，得駑馬十五日之凡行。又以十五日乘半里，以減駑馬初日之行，餘以乘日分子，如日分母而一，所得加前里，卽駑馬定行里數。其奇半里者，爲半法，以半法增殘分，卽得。其不盡者而命分。}\end{unitzh}
\begin{uniten}Method: Suppose 15 days: deficit 337 1/2 li; let it instead be 16 days: excess 140 li. With the excess and deficit cross-multiply the supposed numbers, and add them to make the dividend (shi); add the excess and deficit to make the divisor (fa). Divide the dividend by the divisor, which gives the number of days. As for what is not exhausted, divide by the equal number (dengshu 等數) and name the fraction (mingfen 命分). [Liu Hui: To find the travel of the good horse: multiply the li of speeding up (yiji 益疾) by 14 and halve it; add the li travelled by the good horse on the first day; with this multiply the numerator of the day fraction, and divide by the denominator of the day fraction; what is obtained, together with what is not exhausted named as a fraction, add to the previous li of the good horse's total travel, and you have it. To find the travel of the nag: multiply 1/2 li by 14, and subtract it from the li travelled by the nag on the first day; add the remainder to the first day's travel, and then halve it; multiply by 15 days, which gives the nag's total travel in 15 days. Further, multiply 1/2 li by 15 days, and subtract it from the nag's first day's travel; with the remainder multiply the numerator of the day fraction, and divide by the denominator of the day fraction; add what is obtained to the previous li, and this is the nag's definite travel (dingxing 定行) in li. As for its odd half li, make it the half-divisor (banfa 半法); augment the remaining parts (canfen 殘分) by the half-divisor, and you have it. What is not exhausted is named as a fraction.] [Translator's note: no earlier total for the good horse has been computed in this passage; the text appears defective here.]\end{uniten}

\begin{unitzh}\unitlabel{jiuzhang-031.8}{\zhsmall 按：「令十五日不足三百三十七里半」者，據良馬十五日凡行四千二百六十里，除先去齊三千里，定還迎駑馬一千二百六十里；駑馬凡十五日凡行一千四百二里半。并良、駑二馬所行，得二千六百六十二里半，課於三千里，少三百三十七里半，故曰不足。「令之十六日多一百四十里」者，據良馬十六日凡行四千六百四十八里，先除去齊三千里，定還迎駑馬一千六百四十八里；駑馬十六日凡行一千四百九十二里。并良、駑二馬所行，得三千一百四十里，課於三千里，餘有一百四十里，故謂之多也。}\end{unitzh}
\begin{uniten}[Liu Hui: Examining: as for "let it be 15 days, deficit 337 1/2 li": it is based on the good horse travelling in all 4260 li in 15 days; removing the 3000 li of first going to Qi, the definite return to meet the nag is 1260 li; the nag in all 15 days travels in all 1402 1/2 li. Adding what the two horses, the good horse and the nag, travelled gives 2662 1/2 li; compared with 3000 li, it is short by 337 1/2 li; hence it is called "deficit". As for "let it instead be 16 days, excess 140 li": it is based on the good horse travelling in all 4648 li in 16 days; first removing the 3000 li of going to Qi, the definite return to meet the nag is 1648 li; the nag in 16 days travels in all 1492 li. Adding what the two horses, the good horse and the nag, travelled gives 3140 li; compared with 3000 li, there is a remainder of 140 li; hence it is called "excess".]\end{uniten}

\begin{unitzh}\unitlabel{jiuzhang-031.9}{\zhsmall 「以盈不足維乘假令之數，并而爲實；并盈不足爲法；實如法而一，得日數」者，卽設差不盈不朒之正數。以二馬初日所行里乘十五日，爲十五日平行數。求初末益疾減遲之數者，并一與十四，以十四乘而半之，爲中平之積；又令益疾減遲里數乘之，各爲減益之中平里，故各減益平行數，得十五日定行里。若求後一日，以十六日之定行里數乘日分子，如日分母而一，各得日分子之定行里數，故各并十五日定行里，卽得。其駑馬奇半里者，法爲全里之分，故破半里爲半法，以增殘分，卽合所問也。}\end{unitzh}
\begin{uniten}[Liu Hui: As for "with the excess and deficit cross-multiply the supposed numbers, and add them to make the dividend; add the excess and deficit to make the divisor; divide the dividend by the divisor, which gives the number of days": this is the correct number (zhengshu 正數) at which the supposed difference is neither in excess nor in shortfall (buying bunü 不盈不朒). Multiply 15 days by the li travelled by the two horses on the first day, to make the level-travel (pingxing 平行) numbers for 15 days. To find the amounts of speeding up and slowing down (jianchi 減遲) from beginning to end: add 1 and 14, multiply by 14 and halve it, which makes the area (ji 積) of the mean level (zhongping 中平); further let the li of speeding up and slowing down multiply it, which makes in each case the mean-level li of decrease and increase; therefore decrease and increase the respective level-travel numbers, which gives the definite li travelled in 15 days. If one seeks the last day: multiply the numerator of the day fraction by the definite li travelled on the 16th day, and divide by the denominator of the day fraction; each obtains the definite li travelled for the numerator of the day fraction; therefore add each to the definite li travelled in 15 days, and you have it. As for the nag's odd half li: the divisor constitutes the parts of a whole li; therefore break the half li into the half-divisor, and augment the remaining parts with it; then it accords with what is asked.]\end{uniten}

\begin{unitzh}\unitlabel{jiuzhang-032.1}今有人持錢之蜀賈，利十三。初返歸一萬四千，次返歸一萬三千，次返歸一萬二千，次返歸一萬一千，後返歸一萬。凡五返歸錢，本利俱盡。問本持錢及利各幾何？荅曰：本三萬四百六十八錢三十七萬一千二百九十三分錢之八萬四千八百七十六；利二萬九千五百三十一錢三十七萬一千二百九十三分錢之二十八萬六千四百一十七。\end{unitzh}
\begin{uniten}Suppose there is a person who carries cash (qian 錢) to Shu to trade, making a profit of 3 for every 10. On the first return he brings home 14000; on the next return he brings home 13000; on the next return he brings home 12000; on the next return he brings home 11000; on the last return he brings home 10000. In all, he brings cash home on 5 returns, and principal and profit are both exhausted. Question: how much are the cash originally carried and the profit, each? Answer: principal 30468 84876/371293 qian; profit 29531 286417/371293 qian.\end{uniten}

\begin{unitzh}\unitlabel{jiuzhang-032.2}術曰：假令本錢三萬，不足一千七百三十八錢半；令之四萬，多三萬五千三百九十錢八分。\end{unitzh}
\begin{uniten}Method: Suppose the principal cash is 30000: deficit 1738 1/2 qian; let it be 40000: over by 35390 qian 8 fen (分).\end{uniten}

\begin{unitzh}\unitlabel{jiuzhang-032.3}{\zhsmall 按：假令本錢三萬，并利爲三萬九千，除初返歸留，餘加利爲三萬二千五百；除二返歸留，餘又加利爲二萬五千三百五十；除第三返歸留，餘又加利爲一萬七千三百五十五；除第四返歸留，餘又加利爲八千二百六十一錢半；除第五返歸留，合一萬錢，不足一千七百三十八錢半。若使本錢四萬，并利爲五萬二千，除初返歸留，餘加利爲四萬九千四百；除第二返歸留，餘又加利爲四萬七千三百二十；除第三返歸留，餘又加利爲四萬五千九百一十六；除第四返歸留，餘又加利爲四萬五千三百九十錢八分；除第五返歸留，合一萬，餘三萬五千三百九十錢八分，故曰多。}\end{unitzh}
\begin{uniten}[Liu Hui: Consider: suppose the principal cash is 30000; with the profit added it becomes 39000. Subtract what is kept at the first return home; the remainder, with profit added, becomes 32500. Subtract what is kept at the 2nd return home; the remainder, with profit added again, becomes 25350. Subtract what is kept at the 3rd return home; the remainder, with profit added again, becomes 17355. Subtract what is kept at the 4th return home; the remainder, with profit added again, becomes 8261 1/2 qian. Subtracting what is kept at the 5th return home, which should come to 10000 qian, there is a deficit of 1738 1/2 qian. If the principal cash is made 40000, with the profit added it becomes 52000. Subtract what is kept at the first return home; the remainder, with profit added, becomes 49400. Subtract what is kept at the 2nd return home; the remainder, with profit added again, becomes 47320. Subtract what is kept at the 3rd return home; the remainder, with profit added again, becomes 45916. Subtract what is kept at the 4th return home; the remainder, with profit added again, becomes 45390 qian 8 fen. Subtracting what is kept at the 5th return home, which comes to 10000, there remain 35390 qian 8 fen; hence it is called "over".]\end{uniten}

\begin{unitzh}\unitlabel{jiuzhang-032.4}{\zhsmall 又術：置後返歸一萬，以十乘之，十三而一，卽後所持之本；加一萬一千，又以十乘之，十三而一，卽第四返之本；加一萬二千，又以十乘之，十三而一，卽第三返之本；加一萬三千，又以十乘之，十三而一，卽第二返之本；加一萬四千，又以十乘之，十三而一，卽初持之本。并五返之錢，以減之，卽利也。}\end{unitzh}
\begin{uniten}[Liu Hui: Another method: set out the last return home, 10000; multiply it by 10 and divide by 13: this is the principal carried last. Add 11000, again multiply by 10 and divide by 13: this is the principal of the 4th return. Add 12000, again multiply by 10 and divide by 13: this is the principal of the 3rd return. Add 13000, again multiply by 10 and divide by 13: this is the principal of the 2nd return. Add 14000, again multiply by 10 and divide by 13: this is the principal carried at first. Sum the cash of the 5 returns and subtract the principal from it: this is the profit.]\end{uniten}

\begin{unitzh}\unitlabel{jiuzhang-032.5}九章算術卷七終\end{unitzh}
\begin{uniten}End of Chapter (juan) 7 of the Nine Chapters on the Mathematical Art.\end{uniten}

\begin{unitzh}\unitlabel{jiuzhang-032.6}九章算術卷七訂訛{\zhsmall 算經十書之二}休寧戴震東原\end{unitzh}
\begin{uniten}Corrections of Errors (dinge) in Chapter (juan) 7 of the Nine Chapters on the Mathematical Art [Note: The second of the Ten Mathematical Classics.] Dai Zhen (Dongyuan) of Xiuning\end{uniten}

\begin{unitzh}\unitlabel{jiuzhang-032.7}副置所出率，以少減多，餘以約法、實，實爲物價，法爲人數。{\zhsmall 此術前列四問，本通爲一，徽注辭不達意。據第二問：人出九，盈十一；人出六，不足十六。以十六互乘九，得一百四十四；以十一互乘六，得六十六；并之，二百一十，爲實，所謂「相與同共買物者，置所出率，盈、不足各居其下，令維乘所出率，并以爲實」也。盈十一及不足十六，并之，二十七，爲法，所謂「幷盈、不足爲法」也。於出九減出六，餘三，以三約實二百一十，得價七十；以三約法二十七，得九人，所謂「副置所出率，以少減多，餘以約法、實，實爲物價，法爲人數」也。其第一問所出率相減餘一，則法、實無庸約，故維乘并之，徑得物價；并盈、不足，徑得人數，偶省一算耳。}\end{unitzh}
\begin{uniten}Set out the rates of payment (suochu lü) separately (fuzhi); subtract the smaller from the larger; with the remainder reduce (yue) the divisor (fa) and the dividend (shi): the dividend is the price of the thing; the divisor is the number of people. [Dai Zhen: For this method, the four problems listed before it are fundamentally unified as one; Hui's commentary is worded so that it does not convey the meaning. According to the second problem: each person pays 9: excess 11; each person pays 6: deficit 16. Multiplying 9 crosswise (hucheng) by 16 gives 144; multiplying 6 crosswise by 11 gives 66; summed, 210, as the dividend: this is what is meant by "for those who jointly buy a thing together, set out the rates of payment, with the excess and the deficit each placed below them; let them cross-multiply (weicheng) the rates of payment, and sum to make the dividend". The excess 11 and the deficit 16, summed, 27, make the divisor: this is what is meant by "sum the excess and deficit to make the divisor". From the payment of 9 subtract the payment of 6, remainder 3; reducing the dividend 210 by 3 gives the price 70; reducing the divisor 27 by 3 gives 9 people: this is what is meant by "set out the rates of payment separately; subtract the smaller from the larger; with the remainder reduce the divisor and the dividend: the dividend is the price of the thing; the divisor is the number of people". In the first problem, the rates of payment subtracted from each other leave 1, so there is no need to reduce the divisor and dividend; therefore cross-multiplying and summing directly gives the price of the thing, and summing the excess and deficit directly gives the number of people; it just happens to save one computation.]\end{uniten}

\begin{unitzh}\unitlabel{jiuzhang-032.8}{\zhsmall 第三問：人出半，盈四；人出少半，不足三。半卽二分之一，少半卽三分之一也。又卽爲二人共出一，盈四；三人共出一，不足三。以二分通三之一，得二；以三分通二之一，得三，是齊其子也。二分、三分相乘，得六，是同其母也。以盈四互乘二，得八；以不足三互乘三，得九；并之，十七，爲實，是盈、不足維乘所出率，并以爲實也。盈四及不足三并而爲七，以分母相乘所得六通而齊之，得四十二，爲法，是并盈、不足爲法也。二、三相減，餘一，以約法、實，亦無庸約，故徑得物價、人數。}\end{unitzh}
\begin{uniten}[Dai Zhen: The third problem: each person pays a half: excess 4; each person pays a lesser half (shaoban): deficit 3. A half is 1/2; a lesser half is 1/3. That is, again, 2 people together pay 1: excess 4; 3 people together pay 1: deficit 3. Converting (tong) 1/3 by the 2 parts gives 2; converting 1/2 by the 3 parts gives 3: this is homogenizing (qi) their numerators (zi). The 2 parts and the 3 parts multiplied together give 6: this is equalizing (tong) their denominators (mu). Multiplying 2 crosswise by the excess 4 gives 8; multiplying 3 crosswise by the deficit 3 gives 9; summed, 17, as the dividend: this is the excess and deficit cross-multiplying the rates of payment and being summed to make the dividend. The excess 4 and the deficit 3 summed make 7; converting and homogenizing it by the 6 obtained from multiplying the denominators (fenmu) together gives 42, as the divisor: this is summing the excess and deficit to make the divisor. 2 and 3 subtracted from each other leave 1; reducing the divisor and dividend by it, again there is no need to reduce; hence the price of the thing and the number of people are obtained directly.]\end{uniten}

\begin{unitzh}\unitlabel{jiuzhang-032.9}{\zhsmall 第四問：七家共出一百九十，不足三百三十；九家共出二百七十，盈三十。以九家互乘一百九十，得一千七百一十；以七家互乘二百七十，得一千八百九十，猶之齊其子也。七家、九家相乘，得六十三，猶之同其母也。以盈三十互乘一千七百一十，得五萬一千三百；以不足三百三十互乘一千八百九十，得六十二萬三千七百；并之，六十七萬五千，爲實，是盈、不足維乘所出率，并以爲實也。盈三十及不足三百三十并而爲三百六十，以六十三通而齊之，得二萬二千六百八十，爲法，是并盈、不足爲法也。於一千八百九十減一千七百一十，餘一百八十，以一百八十約實六十七萬五千，得價三千七百五十；以一百八十約法二萬二千六百八十，得一百二十六家，是副置所出率，以少減多，餘以約法、實也。三、四兩問皆所謂「有分者通之」也。}\end{unitzh}
\begin{uniten}[Dai Zhen: The fourth problem: 7 households together pay 190: deficit 330; 9 households together pay 270: excess 30. Multiplying 190 crosswise by the 9 households gives 1710; multiplying 270 crosswise by the 7 households gives 1890: this is like homogenizing their numerators. The 7 households and the 9 households multiplied together give 63: this is like equalizing their denominators. Multiplying 1710 crosswise by the excess 30 gives 51300; multiplying 1890 crosswise by the deficit 330 gives 623700; summed, 675000, as the dividend: this is the excess and deficit cross-multiplying the rates of payment and being summed to make the dividend. The excess 30 and the deficit 330 summed make 360; converting and homogenizing it by 63 gives 22680, as the divisor: this is summing the excess and deficit to make the divisor. From 1890 subtract 1710, remainder 180; reducing the dividend 675000 by 180 gives the price 3750; reducing the divisor 22680 by 180 gives 126 households: this is setting out the rates of payment separately, subtracting the smaller from the larger, and with the remainder reducing the divisor and dividend. The third and fourth problems are both what is meant by "where there are fractions, convert them".]\end{uniten}

\begin{unitzh}\unitlabel{jiuzhang-032.10}按此術：兩盈者，兩設皆逾於正數；兩不足者，兩設皆不足於正數。其所以變化，則兩不足亦猶兩盈。而或有勢同而情違者，當其爲實，俱令不足，其相乘減，則遺其所不足焉。故其餘所以爲實者，無朒數以損焉。葢出而有餘兩盈，兩設皆逾於正數。{\zhsmall 以上訛舛不可通}\end{unitzh}
\begin{uniten}Consider this method: with two excesses, both suppositions (she) exceed the true number (zhengshu); with two deficits, both suppositions fall short of the true number. As for how it is transformed, two deficits are just like two excesses. But there may be cases where the configuration is the same but the circumstances differ: when they make the dividend, both are made deficits; when they are multiplied and subtracted, what they are short of is left out. Therefore what remains, which makes the dividend, has no deficient (nü) number to diminish it. For when the payment leaves a surplus, with two excesses, both suppositions exceed the true number. [Dai Zhen: The above is corrupt and cannot be made sense of.]\end{uniten}

\begin{unitzh}\unitlabel{jiuzhang-032.11}術曰：以盈不足術求之。假令故米二斗，不足二升；令之三斗，有餘二升。{\zhsmall 此下有脫文，當云：以盈、不足維乘假令之數，并爲實；并盈、不足爲法；實如法而一，得故米斗數。}\end{unitzh}
\begin{uniten}Method: Find it by the excess and deficit method. Suppose the original rice (gumi) is 2 dou (斗): deficit 2 sheng (升); let it be 3 dou: surplus 2 sheng. [Dai Zhen: Below this there is missing text; it should read: 'Cross-multiply the supposed numbers by the excess and deficit, and sum to make the dividend; sum the excess and deficit to make the divisor; divide the dividend by the divisor, giving the number of dou of original rice.']\end{uniten}

\begin{unitzh}\unitlabel{jiuzhang-032.12}術曰：假令五日，不足五寸；令之六日，有餘一尺二寸。{\zhsmall 此下有脫文，當云：以盈、不足維乘假令之數，并爲實；并盈、不足爲法；實如法而一，得日數。}\end{unitzh}
\begin{uniten}Method: Suppose 5 days: deficit 5 cun (寸); let it be 6 days: surplus 1 chi (尺) 2 cun. [Dai Zhen: Below this there is missing text; it should read: 'Cross-multiply the supposed numbers by the excess and deficit, and sum to make the dividend; sum the excess and deficit to make the divisor; divide the dividend by the divisor, giving the number of days.']\end{uniten}

\begin{unitzh}\unitlabel{jiuzhang-032.13}術曰：假令二日，不足一尺五寸；令之三日，有餘一尺七寸半。{\zhsmall 此下有脫文，當云：以盈、不足維乘假令之數，并爲實；并盈、不足爲法；實如法而一，得日數。不盡者，等數約之而命分。以後一日所長乘日分子，如日分母而一，各增二長，爲二物共長齊等之數。}\end{unitzh}
\begin{uniten}Method: Suppose 2 days: deficit 1 chi 5 cun; let it be 3 days: surplus 1 chi 7 1/2 cun. [Dai Zhen: Below this there is missing text; it should read: 'Cross-multiply the supposed numbers by the excess and deficit, and sum to make the dividend; sum the excess and deficit to make the divisor; divide the dividend by the divisor, giving the number of days. If it does not divide exactly, reduce it by the equal number (dengshu) and name the fraction (mingfen). Multiply the numerator of the day-fraction (rifen) by what is grown on the last day, divide by the denominator of the day-fraction, and add to each of the two lengths: this gives the number of the equal length common to the two things.']\end{uniten}

\begin{unitzh}\unitlabel{jiuzhang-032.14}術曰：假令醇酒五升，行酒一斗五升，有餘一十；令之醇酒二升，行酒一斗八升，不足二。{\zhsmall 此下有脫文，當云：各以盈、不足維乘之，并爲實；并盈、不足爲法；實如法而一，得二酒之數。}\end{unitzh}
\begin{uniten}Method: Suppose pure wine (chunjiu) 5 sheng and common wine (xingjiu) 1 dou 5 sheng: surplus 10; let it be pure wine 2 sheng and common wine 1 dou 8 sheng: deficit 2. [Dai Zhen: Below this there is missing text; it should read: 'Cross-multiply each by the excess and deficit, and sum to make the dividend; sum the excess and deficit to make the divisor; divide the dividend by the divisor, giving the amounts of the two wines.']\end{uniten}

\begin{unitzh}\unitlabel{jiuzhang-032.15}術曰：假令一十五日，不足三百三十七里半；令之一十六日，多一百四十里。以盈、不足維乘假令之數，并而爲實；并盈、不足爲法；實如法而一，得日數。不盡者，以等數除之而命分。{\zhsmall 此下有脫文，當云：求二馬各行里數，以一十五日初、末之行并而半之，一十五乘之；又以後一日之定行乘日分子，日分母爲法而一，各加於前；破半里爲半法，以入殘分。}\end{unitzh}
\begin{uniten}Method: Suppose 15 days: deficit 337 1/2 li (里); let it be 16 days: over by 140 li. Cross-multiply the supposed numbers by the excess and deficit, and sum to make the dividend; sum the excess and deficit to make the divisor; divide the dividend by the divisor, giving the number of days. If it does not divide exactly, divide it by the equal number and name the fraction. [Dai Zhen: Below this there is missing text; it should read: 'To find the li traveled by each of the two horses: sum the travel of the first and the last of the 15 days and halve it, and multiply by 15; also multiply the numerator of the day-fraction by the definite travel (dingxing) of the last day, and divide by the denominator of the day-fraction as divisor; add each to the preceding; break the half li into half the divisor (banfa), to enter it into the remaining fraction (canfen).']\end{uniten}

\begin{unitzh}\unitlabel{jiuzhang-032.16}求良馬行者，一十四乘益疾里數而半之，加良馬初日之行里數。{\zhsmall 此下有脫文，當云：以乘一十五日，得一十五日之凡行；又以一十五日乘益疾里數，加良馬初日之行。}\end{unitzh}
\begin{uniten}To find the travel of the good horse: multiply the li of increase in speed (yiji) by 14 and halve it, and add the li of the good horse's travel on the first day. [Dai Zhen: Below this there is missing text; it should read: 'Multiply by 15 days, giving the total travel (fanxing) of 15 days; also multiply the li of increase in speed by 15 days and add the good horse's travel on the first day.']\end{uniten}

\section*{卷第八}\addcontentsline{toc}{section}{卷第八}

\begin{unitzh}\unitlabel{jiuzhang-033.1}九章算術卷八{\zhsmall 算經十書之二} 魏劉徽注 唐朝議大夫、行太史令、上輕車都尉臣李淳風等奉勅注釋\end{unitzh}
\begin{uniten}The Nine Chapters on the Mathematical Art, juan 8 [Note: The second of the Ten Mathematical Classics.] Commentary by Liu Hui of the Wei. Of the Tang, Grand Master for Court Discussion, acting Grand Astrologer, Superior Commandant of Light Chariots, Your servant Li Chunfeng et al., annotated by imperial order.\end{uniten}

\begin{unitzh}\unitlabel{jiuzhang-033.2}方程{\zhsmall 以御錯糅正負}\end{unitzh}
\begin{uniten}Rectangular arrays (fangcheng) [Liu Hui: For dealing with intermingled (cuorou) positives and negatives.]\end{uniten}

\begin{unitzh}\unitlabel{jiuzhang-033.3}今有上禾三秉，中禾二秉，下禾一秉，實三十九斗；上禾二秉，中禾三秉，下禾一秉，實三十四斗；上禾一秉，中禾二秉，下禾三秉，實二十六斗。問上、中、下禾實一秉各幾何？荅曰：上禾一秉九斗四分斗之一，中禾一秉四斗四分斗之一，下禾一秉二斗四分斗之三。\end{unitzh}
\begin{uniten}Suppose there are 3 bundles (bing 秉) of top-grade grain (shang he 上禾), 2 bundles of middle-grade grain and 1 bundle of low-grade grain, with a yield (shi) of 39 dou (斗); 2 bundles of top-grade grain, 3 bundles of middle-grade grain and 1 bundle of low-grade grain, with a yield of 34 dou; 1 bundle of top-grade grain, 2 bundles of middle-grade grain and 3 bundles of low-grade grain, with a yield of 26 dou. Question: what is the yield of 1 bundle each of top-, middle- and low-grade grain? Answer: 1 bundle of top-grade grain, 9 1/4 dou; 1 bundle of middle-grade grain, 4 1/4 dou; 1 bundle of low-grade grain, 2 3/4 dou.\end{uniten}

\begin{unitzh}\unitlabel{jiuzhang-033.4}方程{\zhsmall 程，課程也。羣物總雜，各列有數，總言其實。令每行爲率，二物者再程，三物者三程，皆如物數程之，並列爲行，故謂之方程。行之左右無所同存，且爲有所據而言耳。此都術也，以空言難曉，故特繫之禾以決之。又列中行如右行也。}\end{unitzh}
\begin{uniten}Rectangular arrays (fangcheng) [Liu Hui: Cheng (程) is kecheng (課程), "assessing". When many things are all mixed together, each is laid out with its number, and their yield is stated in total. Make each column (xing) a rate (lü): for two things assess twice, for three things assess three times; in every case assess according to the number of things, and lay them out side by side as columns; therefore it is called fangcheng. That nothing is identically present in the columns on the left and right is said just so that there is something to rely on. This is a general method (dushu); since it is hard to understand in empty words, it is specially attached to grain in order to settle it. Also, lay out the middle column like the right column.]\end{uniten}

\begin{unitzh}\unitlabel{jiuzhang-033.5}術曰：置上禾三秉，中禾二秉，下禾一秉，實三十九斗，於右方。中、左禾列如右方。以右行上禾徧乘中行，而以直除。{\zhsmall 爲術之意，令少行減多行，反覆相減，則頭位必先盡。上無一位，則此行亦闕一物矣。然而舉率以相減，不害餘數之課也。若消去頭位，則下去一物之實。如是叠令左右行相減，審其正負，則可得而知。先令右行上禾乘中行，爲齊同之意。爲齊同者，謂中行上禾亦乘右行也。從簡易，雖不言齊同，以齊同之意觀之，其義然矣。}又乘其次，亦以直除。{\zhsmall 復去左行首。}然以中行中禾不盡者{\zhsmall 徧乘左行}而以直除。{\zhsmall 亦令兩行相乘，去行之中禾也。}左方下禾不盡者，上爲法，下爲實。實卽下禾之實。{\zhsmall 上、中禾皆去，故餘數是下禾實，非但一秉。欲約衆秉之實，當以禾秉數爲法。列此下禾之秉實，乘兩行以直除，則下禾之位自決矣。各以其餘一位之秉除其下實，卽斗數矣。用算繁而不省，所以别爲法約也。然猶不如自用其舊，廣異法也。}\end{unitzh}
\begin{uniten}Method: Set up 3 bundles of top-grade grain, 2 bundles of middle-grade grain, 1 bundle of low-grade grain and the yield of 39 dou on the right. Lay out the grain of the middle and the left like the right. With the top-grade grain of the right column multiply the middle column throughout, and eliminate directly (zhichu 直除). [Liu Hui: The intention of the method is to make the smaller column subtract from the larger column; subtracting from each other back and forth, the head position (touwei) is necessarily exhausted first. When the top is short of one position, this column also lacks one thing. Nevertheless, taking up the rates to subtract from each other does no harm to the assessment of the remaining numbers. If the head position is cancelled and removed, then below, the yield of one thing is removed. If in this way the left and right columns are made to subtract from each other in turn, and their positives and negatives are examined, then it can be obtained and known. First making the top-grade grain of the right column multiply the middle column has the sense of homogenizing and equalizing (qitong 齊同). What makes it homogenizing and equalizing is that the top-grade grain of the middle column also multiplies the right column. For simplicity, although it does not speak of homogenizing and equalizing, when it is viewed in the sense of homogenizing and equalizing, its meaning is so.] Again multiply the next one, and also eliminate directly. [Liu Hui: Again remove the head of the left column.] Then with the middle-grade grain of the middle column that is not exhausted [Liu Hui: multiply the left column throughout] and eliminate directly. [Liu Hui: Also make the two columns multiply each other, to remove the middle-grade grain of the column.] For the low-grade grain on the left that is not exhausted, the upper is the divisor (fa) and the lower is the dividend (shi). The dividend is the yield of the low-grade grain. [Liu Hui: The top- and middle-grade grain have both been removed, so the remaining number is the yield of the low-grade grain, but not of only 1 bundle. If one wishes to reduce the yield of the many bundles, one should take the number of bundles of grain as the divisor. Lay out the bundles and the yield of this low-grade grain, multiply the two columns and eliminate directly; then the positions of the low-grade grain are settled by themselves. In each case, dividing its yield at the bottom (xiashi) by the bundles of its single remaining position gives the number of dou. The use of counting rods (suan) is complicated and not economical; that is why one separately makes a divisor to reduce. Yet it is still not as good as using its old way by itself; it broadens the different methods.]\end{uniten}

\begin{unitzh}\unitlabel{jiuzhang-033.6}求中禾，以法乘中行下實，而除下禾之實。{\zhsmall 此謂中、下兩禾實。下禾一秉實數先見，將中秉求中禾，其列實以減下實。而左方下禾不唯一秉，下禾實旣以法爲母，則中行下實不以法爲母，於率不通。故先以法乘其實而同之，俱令法爲母，而除下禾實。以下禾先見之實令乘下禾秉數，卽得下禾一位之列實。減於下實，則其數是中禾之實也。}餘如中禾秉數而一，卽中禾之實。{\zhsmall 餘中禾一位之實也。故以一位秉數約之，乃得一秉之實也。}求上禾，亦以法乘右行下實，而除下禾、中禾之實。{\zhsmall 此右行三禾共實，合三位之實，故以二位秉數約之，乃得上禾一秉之實。此右行三禾共實，合中、下禾之實，其數並見，以中、下禾先見之實令乘右行中、下禾秉數以減之，故亦如前各求列實以減下實也。}餘如上禾秉數而一，卽上禾之實。實皆如法，各得一斗。{\zhsmall 三實同用。不滿法者，以法命之。母實皆當除之。}\end{unitzh}
\begin{uniten}To find the middle-grade grain, multiply the yield at the bottom of the middle column by the divisor, and eliminate the yield of the low-grade grain. [Liu Hui: This refers to the yield of both the middle- and low-grade grain. The number for the yield of 1 bundle of low-grade grain has been found first; in taking the middle bundles to find the middle-grade grain, its arrayed yield (lieshi) is subtracted from the yield at the bottom. But the low-grade grain on the left is not just 1 bundle; since the yield of the low-grade grain takes the divisor as denominator (mu), while the yield at the bottom of the middle column does not take the divisor as denominator, they do not communicate in terms of rates. Therefore first multiply its yield by the divisor to make them the same, letting both take the divisor as denominator, and eliminate the yield of the low-grade grain. Making the yield of the low-grade grain found first multiply the number of bundles of low-grade grain gives the arrayed yield of the single position of low-grade grain. Subtract it from the yield at the bottom, and the number is the yield of the middle-grade grain.] Divide the remainder by the number of bundles of middle-grade grain; this is the yield of the middle-grade grain. [Liu Hui: The remainder is the yield of the single position of middle-grade grain. Therefore, reducing it by the number of bundles of the single position, one obtains the yield of 1 bundle.] To find the top-grade grain, also multiply the yield at the bottom of the right column by the divisor, and eliminate the yields of the low-grade and middle-grade grain. [Liu Hui: This is the joint yield of the three grains of the right column, joining the yields of three positions; therefore, reducing it by the numbers of bundles of two positions, one obtains the yield of 1 bundle of top-grade grain. This joint yield of the three grains of the right column joins the yields of the middle- and low-grade grain, whose numbers have both been found; making the yields of the middle- and low-grade grain found first multiply the numbers of bundles of middle- and low-grade grain of the right column, one subtracts them; therefore, also as before, in each case one finds the arrayed yield to subtract from the yield at the bottom.] Divide the remainder by the number of bundles of top-grade grain; this is the yield of the top-grade grain. For all the dividends, each amount equal to the divisor gives 1 dou. [Liu Hui: The three dividends are used alike. What does not fill the divisor is named by the divisor. The denominators and dividends should all be divided.]\end{uniten}

\begin{unitzh}\unitlabel{jiuzhang-033.7}今有上禾七秉，損實一斗，益之下禾二秉，而實十斗；下禾八秉，益實一斗，與上禾二秉，而實十斗。問上、下禾實一秉各幾何？荅曰：上禾一秉實一斗五十二分斗之十八，下禾一秉實五十二分斗之四十一。\end{unitzh}
\begin{uniten}Suppose there are 7 bundles of top-grade grain: diminishing (sun) the yield by 1 dou and augmenting (yi) it with 2 bundles of low-grade grain, the yield is 10 dou; 8 bundles of low-grade grain: augmenting the yield by 1 dou, together with 2 bundles of top-grade grain, the yield is 10 dou. Question: what is the yield of 1 bundle each of top- and low-grade grain? Answer: 1 bundle of top-grade grain yields 1 18/52 dou; 1 bundle of low-grade grain yields 41/52 dou.\end{uniten}

\begin{unitzh}\unitlabel{jiuzhang-033.8}術曰：如方程。損之曰益，益之曰損。{\zhsmall 問者之辭雖以損益爲說，今按實云上禾七秉、下禾二秉，實十一斗；上禾二秉、下禾八秉，實九斗也。「損之日益」，言損一斗，餘當十斗，今欲全其實，當加所損也。「益之曰損」，言益實一斗，乃滿十斗，今欲加本實，當減所加，卽得也。}損實一斗者，其實過十斗也；益實一斗者，其實不滿十斗也。{\zhsmall 重諭損益數者，各以損益之數損益之也。}\end{unitzh}
\begin{uniten}Method: As in rectangular arrays. Diminishing, call it augmenting; augmenting, call it diminishing. [Liu Hui: Although the questioner's wording is couched in terms of diminishing and augmenting, now, according to the facts, it says: 7 bundles of top-grade grain and 2 bundles of low-grade grain yield 11 dou; 2 bundles of top-grade grain and 8 bundles of low-grade grain yield 9 dou. "Diminishing, call it augmenting" [Translator's note: here the text has 日 "day" for 曰 "call".] says that, diminishing by 1 dou, the remainder amounts to 10 dou; now, wishing to make its yield whole, one should add what was diminished. "Augmenting, call it diminishing" says that, augmenting the yield by 1 dou, it then fills 10 dou; now, wishing to add the original yield, one should subtract what was added, and then it is obtained.] As for "diminishing the yield by 1 dou", the yield exceeds 10 dou; as for "augmenting the yield by 1 dou", the yield does not fill 10 dou. [Liu Hui: Explaining once more the numbers diminished and augmented: in each case one diminishes or augments by the number diminished or augmented.]\end{uniten}

\begin{unitzh}\unitlabel{jiuzhang-033.9}今有上禾二秉，中禾三秉，下禾四秉，實皆不滿斗。上取中，中取下，下取上，各一秉，而實滿斗。問上、中、下禾實一秉各幾何？荅曰：上禾一秉實二十五分斗之九，中禾一秉實二十五分斗之七，下禾一秉實二十五分斗之四。\end{unitzh}
\begin{uniten}Suppose there are 2 bundles of top-grade grain, 3 bundles of middle-grade grain and 4 bundles of low-grade grain, whose yields each do not fill 1 dou. If the top-grade takes from the middle-grade, the middle-grade from the low-grade, and the low-grade from the top-grade, 1 bundle each, then the yields fill 1 dou. Question: what is the yield of 1 bundle each of top-, middle- and low-grade grain? Answer: 1 bundle of top-grade grain yields 9/25 dou; 1 bundle of middle-grade grain yields 7/25 dou; 1 bundle of low-grade grain yields 4/25 dou.\end{uniten}

\begin{unitzh}\unitlabel{jiuzhang-033.10}術曰：如方程，各置所取。{\zhsmall 置上禾二秉爲右行之上，中禾三秉爲中行之中，下禾四秉爲左行之下，所取一秉及實一斗各從其位。諸行相借取之物，皆依此例。}以正負術入之。\end{unitzh}
\begin{uniten}Method: As in rectangular arrays, set up in each case what is taken. [Liu Hui: Set 2 bundles of top-grade grain as the top of the right column, 3 bundles of middle-grade grain as the middle of the middle column, and 4 bundles of low-grade grain as the bottom of the left column; the 1 bundle taken and the yield of 1 dou each follow their positions. Things that the columns borrow and take from each other all follow this example.] Apply the method of positive and negative (zhengfu shu) to it.\end{uniten}

\begin{unitzh}\unitlabel{jiuzhang-033.11}正負術曰：{\zhsmall 今兩算得失相反，要令正負以名之。正算赤，負算黑，否則以邪正爲異。方程自有赤黑相取，左右數相推求之術，而其并減之勢不得交通，故使赤黑相消奪之。於算或減或益，同行異位殊爲二品，各有并減之差，見於下焉。著此二條，特繫之禾，以成此二條之意。故赤黑相雜，足以定上下之程；減益雖殊，足以通左右之數；差實雖分，足以應同異之率。然則其正無入負之，負無入正之，其率不妄也。}\end{unitzh}
\begin{uniten}The method of positive and negative says: [Liu Hui: Now, when two counting-rod numbers (suan) are opposite as gain and loss, one must name them by means of positive and negative. Positive counting rods are red, negative counting rods are black; otherwise, they are distinguished as slanting and upright. Rectangular arrays have of themselves a method in which red and black take from each other and the numbers on the left and right are deduced from each other, but its configurations (shi 勢) of combining and subtracting cannot interconnect; therefore one makes red and black cancel and take away (xiaoduo) from each other. On the counting rods one sometimes subtracts and sometimes augments; different positions in the same column are distinguished as two kinds (pin), each having its differences in combining and subtracting, which are shown below. These two clauses (tiao) are set down and specially attached to grain, so as to complete the meaning of these two clauses. Therefore the intermixing of red and black suffices to settle the assessments (cheng) of top and bottom; although subtracting and augmenting differ, they suffice to make the numbers on the left and right communicate; although the differences and the yields are separate, they suffice to answer to the rates of same and different. So then, with its "if a positive has nothing to enter (wuru), make it negative; if a negative has nothing to enter, make it positive", the rates are not wrong.]\end{uniten}

\begin{unitzh}\unitlabel{jiuzhang-033.12}同名相除，{\zhsmall 此爲以赤除赤，以黑除黑。行求相減者，爲法頭位也。然則頭位同名者，當用此條；頭位異名者，當用下條。}異名相益，{\zhsmall 益行減行，當各以其類矣。其異名者，非其類也。非其類者，猶無對也，非所得減也。故赤用黑對則除黑，無對則除赤，赤黑并於本數。此爲相益之，皆所以爲消奪。消奪之與減益成一實也。術本取要，必除行首，至於他位，不嫌多少，故或令相減，或令相并，理無同異，一也。}正無入負之，負無入正之。{\zhsmall 無入爲無對也。無所得減，則使消奪者居位也。其當以列實或減下實，而中行正負雜者，亦用此條。此條者，同名減實，異名益實，正無入負之，負無入正之也。}\end{unitzh}
\begin{uniten}Those with the same name (tongming) eliminate each other; [Liu Hui: This is eliminating red with red and black with black. When columns seek to subtract from each other, it is the head position that serves as divisor. Thus, when the head positions have the same name, one should use this clause; when the head positions have different names (yiming), one should use the following clause.] those with different names augment each other; [Liu Hui: In augmenting a column or subtracting from a column, one should in each case go by its category (lei). Those with different names are not of its category. What is not of its category is like having no counterpart (wudui), and is not something that can be subtracted. Therefore, when red uses black as counterpart, the black is eliminated; when there is no counterpart, the red is eliminated, and red and black are combined in the original number. This is making them augment each other; all of this is a means of cancelling and taking away. Cancelling and taking away and subtracting and augmenting form one and the same reality. The method fundamentally takes what is essential: it must eliminate the head of the column; as for the other positions, it does not mind whether they are more or less; therefore it sometimes makes them subtract from each other and sometimes makes them combine with each other; in principle there is no same or different: it is one.] if a positive has nothing to enter, make it negative; if a negative has nothing to enter, make it positive. [Liu Hui: "Nothing to enter" means "no counterpart". When there is nothing from which it can be subtracted, then make what cancels and takes away occupy the position. When one should, with the arrayed yield, subtract as the case may be from the yield at the bottom, and positives and negatives are mixed in the middle column, one also uses this clause. This clause means: with the same name, subtract from the yield; with different names, augment the yield; if a positive has nothing to enter, make it negative; if a negative has nothing to enter, make it positive.]\end{uniten}

\begin{unitzh}\unitlabel{jiuzhang-033.13}其異名相除，同名相益，正無入正之，負無入負之。{\zhsmall 此條異名相除爲例，故亦與上條互取。凡正負所以記其同異，使二品互相取而已矣。言負者未必負於少，言正者未必正於多。故毎一行之中，雖復赤黑異算，無傷。然則可得使頭位常相與異名。此條之實兼通矣，遂以二條反覆一率。觀其毎與上下互相取位，則隨算而言耳，猶一術也。又本設諸行，欲因減數以相去耳，故其多少無限，令上下相命而已。若以正負相減，如數有舊増法者，每行可均之，不但數物左右之也。}\end{unitzh}
\begin{uniten}In the other case: those with different names eliminate each other, those with the same name augment each other; if a positive has nothing to enter, make it positive; if a negative has nothing to enter, make it negative. [Liu Hui: This clause takes "those with different names eliminate each other" as its model (li); therefore it and the clause above also take from each other. In general, positive and negative serve to record their sameness and difference, making the two kinds take from each other, and that is all. What is said to be negative is not necessarily negative by being less; what is said to be positive is not necessarily positive by being more. Therefore, within each single column, even if the red and black counting rods are swapped, there is no harm. Thus one can make the head positions always have different names from each other. The substance of this clause thus applies to both, and so the two clauses are one rate (lü) turned over. Observing how each time it takes positions mutually with the top and bottom, it is merely stated according to the counting rods; it is still one method. Moreover, the columns were originally set up only in order to remove one another by means of subtracting numbers; therefore there is no limit to how much or how little they are, and one only makes the top and bottom name each other (xiangming). If one subtracts by means of positive and negative, as when numbers have the old method of increasing, each column can be equalized (jun); it is not only for moving the numbers of things left and right.]\end{uniten}

\begin{unitzh}\unitlabel{jiuzhang-034.1}今有上禾五秉，損實一斗一升，當下禾七秉；上禾七秉，損實二斗五升，當下禾五秉。問上、下禾實一秉各幾何？荅曰：上禾一秉五升，下禾一秉二升。\end{unitzh}
\begin{uniten}Suppose there are 5 bundles (bing 秉) of top-grade grain which, with their yield (shi 實) diminished by 1 dou (斗) 1 sheng (升), equal 7 bundles of low-grade grain; and 7 bundles of top-grade grain which, with their yield diminished by 2 dou 5 sheng, equal 5 bundles of low-grade grain. Question: how much is the yield of 1 bundle each of top-grade and low-grade grain? Answer: 1 bundle of top-grade grain, 5 sheng; 1 bundle of low-grade grain, 2 sheng.\end{uniten}

\begin{unitzh}\unitlabel{jiuzhang-034.2}術曰：如方程。置上禾五秉正，下禾七秉負，損實一斗一升正。{\zhsmall 言上禾五秉之實多，減其一斗一升，餘是與下禾七秉相當數也。故互其算，令相折除，以一斗一升爲差。爲差者，上禾之餘實也。}次置上禾七秉正，下禾五秉負，損實二斗五升正。以正負術八之。{\zhsmall 按：正負之術，本設列行，物程之數不限多少，必令與實上下相次，而以每行各自爲率。然而或減或益，同行異位，殊爲二品，并減之差見於下也。}\end{unitzh}
\begin{uniten}Method: as in rectangular arrays (fangcheng). Set up 5 bundles of top-grade grain, positive; 7 bundles of low-grade grain, negative; the yield diminution, 1 dou 1 sheng, positive. [Liu Hui: This says that the yield of the 5 bundles of top-grade grain is more; when 1 dou 1 sheng of it is subtracted, the remainder is the amount equal to 7 bundles of low-grade grain. Therefore one sets their counts (suan) against each other (hu), making them offset and eliminate each other (zhe chu), and takes 1 dou 1 sheng as the difference (cha). What is taken as the difference is the surplus yield of the top-grade grain.] Next set up 7 bundles of top-grade grain, positive; 5 bundles of low-grade grain, negative; the yield diminution, 2 dou 5 sheng, positive. Ba (八) it by the positive and negative method. [Translator's note: 八 "eight" is evidently an error for 入 "enter", the reading of this phrase in the following problems.] [Liu Hui: Note: the positive and negative method fundamentally sets up arrayed columns (lie hang); the numbers of the measures of things (wu cheng) are not limited as to more or fewer, but they must be made to follow one another from top to bottom together with the shi, and each column is taken by itself as a rate (lü). Yet, whether there is subtraction or increase, in the same column but in different positions they are distinguished as 2 kinds (pin), and the difference of combining and subtracting appears below.]\end{uniten}

\begin{unitzh}\unitlabel{jiuzhang-034.3}今有上禾六秉，損實一斗八升，當下禾十秉；下禾十五秉，損實五升，當上禾五秉。問上、下禾實一秉各幾何？荅曰：上禾一秉實八升，下禾一秉實三升。\end{unitzh}
\begin{uniten}Suppose there are 6 bundles of top-grade grain which, with their yield diminished by 1 dou 8 sheng, equal 10 bundles of low-grade grain; and 15 bundles of low-grade grain which, with their yield diminished by 5 sheng, equal 5 bundles of top-grade grain. Question: how much is the yield of 1 bundle each of top-grade and low-grade grain? Answer: 1 bundle of top-grade grain, yield 8 sheng; 1 bundle of low-grade grain, yield 3 sheng.\end{uniten}

\begin{unitzh}\unitlabel{jiuzhang-034.4}術曰：如方程。置上禾六秉正，下禾十秉負，損實一斗八升正；次置上禾五秉負，下禾十五秉正，損實五升正。以正負術入之。{\zhsmall 言上禾六秉之實多，減損其一斗八升，餘是與下禾十秉相當之數，故亦互其算，而以一斗八升爲差實。差實者，下禾之餘實。}\end{unitzh}
\begin{uniten}Method: as in rectangular arrays. Set up 6 bundles of top-grade grain, positive; 10 bundles of low-grade grain, negative; the yield diminution, 1 dou 8 sheng, positive; next set up 5 bundles of top-grade grain, negative; 15 bundles of low-grade grain, positive; the yield diminution, 5 sheng, positive. Enter it (ru) by the positive and negative method. [Liu Hui: This says that the yield of the 6 bundles of top-grade grain is more; when it is reduced by 1 dou 8 sheng, the remainder is the amount equal to 10 bundles of low-grade grain; therefore again one sets their counts against each other, and takes 1 dou 8 sheng as the difference shi (cha shi). The difference shi is the surplus yield of the low-grade grain.]\end{uniten}

\begin{unitzh}\unitlabel{jiuzhang-034.5}今有上禾三秉，益實六斗，當下禾十秉；下禾五秉，益實一斗，當上禾二秉。問上、下禾實一秉各幾何？荅曰：上禾一秉實八斗，下禾一秉實三斗。\end{unitzh}
\begin{uniten}Suppose there are 3 bundles of top-grade grain which, with their yield increased by 6 dou, equal 10 bundles of low-grade grain; and 5 bundles of low-grade grain which, with their yield increased by 1 dou, equal 2 bundles of top-grade grain. Question: how much is the yield of 1 bundle each of top-grade and low-grade grain? Answer: 1 bundle of top-grade grain, yield 8 dou; 1 bundle of low-grade grain, yield 3 dou.\end{uniten}

\begin{unitzh}\unitlabel{jiuzhang-034.6}術曰：如方程。置上禾三秉正，下禾十秉負，益實六斗正；次置上禾二秉負，下禾五秉正，益實一斗正。以正負術入之。{\zhsmall 言上禾三秉之實少，益其六斗，然後於下禾十秉相當也。故亦互其算，而以六斗爲差實。差實者，下禾之餘實。}\end{unitzh}
\begin{uniten}Method: as in rectangular arrays. Set up 3 bundles of top-grade grain, positive; 10 bundles of low-grade grain, negative; the yield increase, 6 dou, positive; next set up 2 bundles of top-grade grain, negative; 5 bundles of low-grade grain, positive; the yield increase, 1 dou, positive. Enter it by the positive and negative method. [Liu Hui: This says that the yield of the 3 bundles of top-grade grain is less; only after it has been increased by 6 dou does it equal the 10 bundles of low-grade grain. Therefore again one sets their counts against each other, and takes 6 dou as the difference shi. The difference shi is the surplus yield of the low-grade grain.]\end{uniten}

\begin{unitzh}\unitlabel{jiuzhang-034.7}今有牛五、羊二，直金十兩；牛二、羊五，直金八兩。問牛、羊各直金幾何？荅曰：牛一直金一兩二十一分兩之十三，羊一直金二十一分兩之二十。\end{unitzh}
\begin{uniten}Suppose there are 5 oxen and 2 sheep, worth 10 liang (兩) of gold; and 2 oxen and 5 sheep, worth 8 liang of gold. Question: how much gold are an ox and a sheep each worth? Answer: 1 ox is worth 1 13/21 liang of gold; 1 sheep is worth 20/21 liang of gold.\end{uniten}

\begin{unitzh}\unitlabel{jiuzhang-034.8}術曰：如方程。{\zhsmall 假令爲同齊，頭位爲牛，左右行相乘定，更置牛十、羊四，直金二十兩；左行牛十、羊二十五，直金四十兩。牛數等同，金多二十兩者，羊差二十一使之然也。以少行減多行，則牛數盡，惟羊與直金之數見，可得而知也。以小推大，雖四五行不異也。}\end{unitzh}
\begin{uniten}Method: as in rectangular arrays. [Liu Hui: Suppose one homogenizes and equalizes them (tong qi), the head position being the oxen: when the left and right columns have been multiplied by each other and fixed, one sets up anew oxen 10, sheep 4, worth 20 liang of gold; the left column: oxen 10, sheep 25, worth 40 liang of gold. The numbers of oxen being equal, that the gold is more by 20 liang is caused by the difference of 21 sheep. Subtract the smaller column from the larger column: then the number of oxen is exhausted, and only the numbers of the sheep and of the value in gold appear, which can then be known. Extending from the small to the large, even with 4 or 5 columns it is no different.]\end{uniten}

\begin{unitzh}\unitlabel{jiuzhang-034.9}今有賣牛二、羊五，以買十三豕，有餘錢一千；賣牛三、豕三，以買九羊，錢適足；賣羊六、豕八，以買五牛，錢不足六百。問牛、羊、豕價各幾何？荅曰：牛價一千二百，羊價五百，豕價三百。\end{unitzh}
\begin{uniten}Suppose there are 2 oxen and 5 sheep sold to buy 13 pigs, with 1000 qian (錢) left over; 3 oxen and 3 pigs sold to buy 9 sheep, the money being exactly sufficient; and 6 sheep and 8 pigs sold to buy 5 oxen, the money being 600 short. Question: what is the price of each of an ox, a sheep and a pig? Answer: price of an ox, 1200; price of a sheep, 500; price of a pig, 300.\end{uniten}

\begin{unitzh}\unitlabel{jiuzhang-034.10}術曰：如方程。置牛二、羊五正，豕十三負，餘錢數正；次置牛三正，羊九負，豕三正；次置牛五負，羊六正，豕八正，不足錢負。以正負術入之。{\zhsmall 此中行買賣相折，錢適足，但互買賣算而已，故下無錢直也。設欲以此行如方程法，先令牛二徧乘左行，而以右行直除之，是終於下實虚缺矣。故注日：「正無實負，負無實正」，方爲類也。方將以别實加不足之數與實物作實。《盈不足》章「黃金白銀」與此相當：「假令黃金九，白銀十一，稱之重適等。交易其一，金輕十三兩。問金、銀一枚各重幾何？」與此同。}\end{unitzh}
\begin{uniten}Method: as in rectangular arrays. Set up 2 oxen and 5 sheep, positive; 13 pigs, negative; the amount of money left over, positive; next set up 3 oxen, positive; 9 sheep, negative; 3 pigs, positive; next set up 5 oxen, negative; 6 sheep, positive; 8 pigs, positive; the money short, negative. Enter it by the positive and negative method. [Liu Hui: In this middle column buying and selling offset each other (xiang zhe) and the money is exactly sufficient; one merely sets the counts of buying and selling against each other, and that is all; therefore below there is no money value. Suppose one wished to treat this column by the method of rectangular arrays: first let the 2 oxen multiply the left column throughout, and subtract the right column from it directly (zhi chu); this ends with the shi below empty and lacking. Therefore the note says: "positive without shi, negative; negative without shi, positive"; this is exactly of that kind (lei). Only then will one take a separate shi, add the amount of the deficit, and with the actual things (shi wu) make the shi. The "yellow gold and white silver" of the chapter Excess and Deficit (ying buzu) corresponds to this: "Suppose there are 9 of yellow gold and 11 of white silver; weighed, their weights are exactly equal. When 1 of them is exchanged, the gold is lighter by 13 liang. Question: how much does 1 piece each of gold and silver weigh?" It is the same as this.] [Translator's note: in 故注日, 日 is evidently an error for 曰 "says".]\end{uniten}

\begin{unitzh}\unitlabel{jiuzhang-034.11}今有五雀、六燕，集稱之衡，雀俱重，燕俱輕。一雀一燕交而處，衡適平。并燕、雀重一斤。問燕、雀一枚各重幾何？荅曰：雀重一兩十九分兩之十三，燕重一兩十九分兩之五。\end{unitzh}
\begin{uniten}Suppose there are 5 sparrows and 6 swallows, gathered and weighed on a balance (heng): the sparrows together are heavy, the swallows together are light. When 1 sparrow and 1 swallow exchange places, the balance is exactly level. Together, the swallows and sparrows weigh 1 jin (斤). Question: how much do 1 swallow and 1 sparrow each weigh? Answer: a sparrow weighs 1 13/19 liang; a swallow weighs 1 5/19 liang.\end{uniten}

\begin{unitzh}\unitlabel{jiuzhang-034.12}術曰：如方程，交易質之，各重八兩。{\zhsmall 此四雀一燕與一雀五燕衡適平，并重一斤，故各八兩。列兩行程數，左行頭位其數是一，可省乘。令右行徧乘左行，而取其法實於左。左行數多，以右行取其數。左頭位減盡，中法下實，卽每枚當重宜可知也。按：此四雀一燕與一雀五燕其重等，是三雀四燕重相當，雀率重四，燕率重三也。諸再程之率，皆可異術求之，卽其數也。}\end{unitzh}
\begin{uniten}Method: as in rectangular arrays; exchange them and weigh them (zhi): each weighs 8 liang. [Liu Hui: Here, with 4 sparrows and 1 swallow against 1 sparrow and 5 swallows, the balance is exactly level, and together they weigh 1 jin; therefore each is 8 liang. Lay out the measure-numbers (cheng shu) in 2 columns: in the left column the number in the head position is 1, so the multiplication can be saved. Let the right column multiply the left column throughout, and obtain its divisor (fa) and dividend (shi) in the left. The numbers of the left column being larger, take away its numbers with the right column. When the head position on the left has been subtracted to nothing, the divisor is in the middle and the dividend below, and then the weight each piece should have can properly be known. Note: here 4 sparrows and 1 swallow are equal in weight to 1 sparrow and 5 swallows; this means that 3 sparrows and 4 swallows are equivalent in weight: the rate of the sparrow is weight 4, the rate of the swallow is weight 3. The rates of all twofold measures (zai cheng) can be found by a different method, and these are then their numbers.]\end{uniten}

\begin{unitzh}\unitlabel{jiuzhang-034.13}今有甲、乙二人持錢，不知其數。甲得乙半而錢五十，乙得甲太半而亦錢五十。問甲、乙持錢各幾何？荅曰：甲持三十七錢半，乙持二十五錢。\end{unitzh}
\begin{uniten}Suppose there are 2 persons, A (jia 甲) and B (yi 乙), holding money, the amounts of which are not known. If A obtains half of B's, he has 50 qian; if B obtains the large half (taiban) of A's, he also has 50 qian. [Translator's note: the large half is 2/3.] Question: how much money does each of A and B hold? Answer: A holds 37 1/2 qian; B holds 25 qian.\end{uniten}

\begin{unitzh}\unitlabel{jiuzhang-034.14}術曰：如方程，損妓之。{\zhsmall 此問者言一甲半乙而五十，太半甲一乙亦五十也。各以分母乘其全，内子，行定：二甲一乙而錢一百，二甲三乙而錢一百五十。於是乃如方程。諸物有分者倣此。}\end{unitzh}
\begin{uniten}Method: as in rectangular arrays: diminish and ji (妓) it. [Translator's note: 妓 "courtesan" is evidently an error for 益 "increase"; cf. 損益之 "diminish and increase it" in the next problem.] [Liu Hui: This question says: 1 A and half a B make 50; the large half of A and 1 B also make 50. In each case multiply the integer part (quan) by the denominator (fen mu) and add in the numerator (zi); the columns are fixed: 2 A and 1 B make 100 qian; 2 A and 3 B make 150 qian. Thereupon proceed as in rectangular arrays. All things that have fractions follow this model.]\end{uniten}

\begin{unitzh}\unitlabel{jiuzhang-034.15}今有二馬、一牛價過一萬，如半馬之價；一馬、二牛價不滿一萬，如半牛之價。問牛、馬價各幾何？荅曰：馬價五千四百五十四錢十一分錢之六，牛價一千八百一十八錢十一分錢之二。\end{unitzh}
\begin{uniten}Suppose there are 2 horses and 1 ox whose price exceeds 10000 by as much as the price of half a horse; and 1 horse and 2 oxen whose price falls short of 10000 by as much as the price of half an ox. Question: what is the price of each of an ox and a horse? Answer: price of a horse, 5454 6/11 qian; price of an ox, 1818 2/11 qian.\end{uniten}

\begin{unitzh}\unitlabel{jiuzhang-034.16}術曰：如方程，損益之。{\zhsmall 此一馬半與一牛價直一萬也，二牛半與一馬亦直一萬也。一馬半與一牛，通分内子，右行爲三馬二牛，直錢二萬；二牛半於一馬，通分内子，左行爲二馬五牛，直錢二萬也。}\end{unitzh}
\begin{uniten}Method: as in rectangular arrays: diminish and increase it (sun yi). [Liu Hui: Here the price of 1 1/2 horses and 1 ox is worth 10000; 2 1/2 oxen and 1 horse are also worth 10000. For 1 1/2 horses and 1 ox, convert into parts (tong fen) and add in the numerator: the right column becomes 3 horses and 2 oxen, worth 20000 qian; for 2 1/2 oxen with 1 horse, convert into parts and add in the numerator: the left column becomes 2 horses and 5 oxen, worth 20000 qian.]\end{uniten}

\begin{unitzh}\unitlabel{jiuzhang-034.17}今有武馬一匹，中馬二匹，下馬三匹，皆載四十石，至阪，皆不能上。武馬借中馬一匹，中馬借下馬一匹，下馬借武馬一匹，乃皆上。問武、中、下馬一匹各力引幾何？荅曰：武馬一匹力引二十二石七分石之六，中馬一匹力引十七石七分石之一，下馬一匹力引五石七分石之五。\end{unitzh}
\begin{uniten}Suppose there are 1 war horse (wu ma), 2 middle-grade horses and 3 low-grade horses, all carrying 40 dan (石); arriving at a slope, none of them can go up. The war horse borrows 1 middle-grade horse, the middle-grade horses borrow 1 low-grade horse, the low-grade horses borrow 1 war horse; then they all go up. Question: how much can 1 each of the war, middle-grade and low-grade horses pull with its strength? Answer: 1 war horse pulls with its strength 22 6/7 dan; 1 middle-grade horse pulls with its strength 17 1/7 dan; 1 low-grade horse pulls with its strength 5 5/7 dan.\end{uniten}

\begin{unitzh}\unitlabel{jiuzhang-034.18}術曰：如方程，各置所借，以正負術入之。\end{unitzh}
\begin{uniten}Method: as in rectangular arrays: set up in each case what is borrowed, and enter it by the positive and negative method.\end{uniten}

\begin{unitzh}\unitlabel{jiuzhang-035.1}今有五家共井，甲二綆不足如乙一綆，乙三綆不足如丙一綆，丙四綆不足如丁一綆，丁五綆不足如戊一綆，戊六綆不足如甲一綆。各得所不足一綆，皆逮。問井深、綆長各幾何？荅曰：井深七丈二尺一寸，甲綆長二丈六尺五寸，乙綆長一丈九尺一寸，丙綆長一丈四尺八寸，丁綆長一丈二尺九寸，戊綆長七尺六寸。\end{unitzh}
\begin{uniten}Suppose there are 5 families sharing a well. A's (甲) 2 ropes (geng 綆) fall short by as much as B's (乙) 1 rope; B's 3 ropes fall short by as much as C's (丙) 1 rope; C's 4 ropes fall short by as much as D's (丁) 1 rope; D's 5 ropes fall short by as much as E's (戊) 1 rope; E's 6 ropes fall short by as much as A's 1 rope. If each obtains the 1 rope by which it falls short, all reach. Question: what are the depth of the well and the length of each rope? Answer: The depth of the well is 7 zhang (丈) 2 chi (尺) 1 cun (寸). A's rope is 2 zhang 6 chi 5 cun long. B's rope is 1 zhang 9 chi 1 cun long. C's rope is 1 zhang 4 chi 8 cun long. D's rope is 1 zhang 2 chi 9 cun long. E's rope is 7 chi 6 cun long.\end{uniten}

\begin{unitzh}\unitlabel{jiuzhang-035.2}術曰：如方程，以正負術入之。{\zhsmall 此率初如方程爲之，名各一逮井。其後法得七百二十一，實七十六，是爲七百二十一綆而七十六逮井，而戊一綆逮之數定，逮七百二十一分之七十六。是故七百二十一爲井深，七十六爲戊綆之長，舉率以言之。}\end{unitzh}
\begin{uniten}Method: As in rectangular arrays (fangcheng), enter it with the positive and negative method. [Liu Hui: This rate (lü) is at first worked as in rectangular arrays, naming each [column] as 1 reaching of the well. Afterwards the divisor (fa) obtained is 721 and the dividend (shi) 76: this means 721 ropes and 76 reachings of the well, and the amount by which E's 1 rope reaches it is fixed: it reaches 76/721. Therefore 721 is the depth of the well and 76 the length of E's rope; it is stated by citing the rates.]\end{uniten}

\begin{unitzh}\unitlabel{jiuzhang-035.3}今有白禾二步，靑禾三步，黃禾四步，黑禾五步，實各不滿斗。白取靑、黃，靑取黃、黑，黃取黑、白，黑取白、靑，各一步，而實滿斗。問白、靑、黃、黑禾實一步各幾何？荅曰：白禾一步實一百一十一分斗之三十三，靑禾一步實一百一十一分斗之二十八，黃禾一步實一百一十一分斗之十七，黑禾一步實一百一十一分斗之十。\end{unitzh}
\begin{uniten}Suppose there are 2 bu (步) of white grain, 3 bu of green grain, 4 bu of yellow grain and 5 bu of black grain, and the yield (shi 實) of each does not fill a dou (斗). If the white takes the green and yellow, the green takes the yellow and black, the yellow takes the black and white, and the black takes the white and green, 1 bu of each, then the yield fills a dou. Question: what is the yield of 1 bu of each of the white, green, yellow and black grain? Answer: 1 bu of white grain yields 33/111 dou. 1 bu of green grain yields 28/111 dou. 1 bu of yellow grain yields 17/111 dou. 1 bu of black grain yields 10/111 dou.\end{uniten}

\begin{unitzh}\unitlabel{jiuzhang-035.4}術曰：如方程，各置所取，以正負術入之。\end{unitzh}
\begin{uniten}Method: As in rectangular arrays, set out in each [column] what is taken, and enter it with the positive and negative method.\end{uniten}

\begin{unitzh}\unitlabel{jiuzhang-035.5}今有甲禾二秉，乙禾三秉，丙禾四秉，重皆過於石。甲二重如乙一，乙三重如丙一，丙四重如甲一。問甲、乙、丙禾一秉各重幾何？荅曰：甲禾一秉重二十三分石之十七，乙禾一秉重二十三分石之十一，丙禾一秉重二十三分石之十。\end{unitzh}
\begin{uniten}Suppose there are 2 sheaves (bing 秉) of grain A, 3 sheaves of grain B and 4 sheaves of grain C, and their weights all exceed a dan (石). [The excess] weight of A's 2 is as B's 1; of B's 3, as C's 1; of C's 4, as A's 1. Question: what does 1 sheaf of each of grains A, B and C weigh? Answer: 1 sheaf of grain A weighs 17/23 dan. 1 sheaf of grain B weighs 11/23 dan. 1 sheaf of grain C weighs 10/23 dan.\end{uniten}

\begin{unitzh}\unitlabel{jiuzhang-035.6}術曰：如方程，置重過於石之物爲負。{\zhsmall 此問者言甲禾二秉之重過於一石也。其過者幾何？如乙一秉重矣。互其算，令相折除，以石爲之差實。差實者，如甲禾餘實，故置算相與同也。}以正負術入之。{\zhsmall 此入頭位異名相除者，正無入正之，負無入負之也。}\end{unitzh}
\begin{uniten}Method: As in rectangular arrays, set the thing that is the weight in excess of the dan as negative. [Liu Hui: What this question says is that the weight of 2 sheaves of grain A exceeds 1 dan. How much is the excess? It is as the weight of 1 sheaf of B. Interchange their counting rods (suan 算) so that they offset (zhechu 折除) each other, and take the dan as their difference dividend (chashi 差實). The difference dividend is like the remaining dividend of grain A; therefore the counting rods are set out in the same way.] Enter it with the positive and negative method. [Liu Hui: In this entering, where the head positions (touwei 頭位) are of different names and subtract from each other (xiangchu 相除), a positive [entering against] nothing is made positive, and a negative [entering against] nothing is made negative.]\end{uniten}

\begin{unitzh}\unitlabel{jiuzhang-035.7}今有令一人，吏五人，從者十人，食雞十；令十人，吏一人，從者五人，食雞八；令五人，吏十人，從者一人，食雞六。問令、吏、從者食雞各幾何？荅曰：令一人食一百二十二分雞之四十五，吏一人食一百二十二分雞之四十一，從者一人食一百二十二分雞之九十七。\end{unitzh}
\begin{uniten}Suppose 1 magistrate (ling 令), 5 clerks (li 吏) and 10 attendants (congzhe 從者) eat 10 chickens; 10 magistrates, 1 clerk and 5 attendants eat 8 chickens; 5 magistrates, 10 clerks and 1 attendant eat 6 chickens. Question: how much chicken does each of magistrate, clerk and attendant eat? Answer: 1 magistrate eats 45/122 of a chicken. 1 clerk eats 41/122 of a chicken. 1 attendant eats 97/122 of a chicken.\end{uniten}

\begin{unitzh}\unitlabel{jiuzhang-035.8}術曰：如方程，以正負術入之。\end{unitzh}
\begin{uniten}Method: As in rectangular arrays, enter it with the positive and negative method.\end{uniten}

\begin{unitzh}\unitlabel{jiuzhang-035.9}今有五羊、四犬、三雞、二兔，直錢一千四百九十六；四羊、二犬、六雞、三兔，直錢一千一百七十五；三羊、一犬、七雞、五兔，直錢九百五十八；二羊、三犬、五雞、一兔，直錢八百六十一。問羊、犬、雞、兔價各幾何？荅曰：羊價一百七十七，犬價一百二十一，雞價二十三，兔價二十九。\end{unitzh}
\begin{uniten}Suppose 5 sheep, 4 dogs, 3 chickens and 2 rabbits are worth 1496 qian (錢); 4 sheep, 2 dogs, 6 chickens and 3 rabbits are worth 1175 qian; 3 sheep, 1 dog, 7 chickens and 5 rabbits are worth 958 qian; 2 sheep, 3 dogs, 5 chickens and 1 rabbit are worth 861 qian. Question: what is the price of each of sheep, dog, chicken and rabbit? Answer: Sheep price 177. Dog price 121. Chicken price 23. Rabbit price 29.\end{uniten}

\begin{unitzh}\unitlabel{jiuzhang-035.10}術曰：如方程，以正負術入之。\end{unitzh}
\begin{uniten}Method: As in rectangular arrays, enter it with the positive and negative method.\end{uniten}

\begin{unitzh}\unitlabel{jiuzhang-035.11}今有麻九斗，麥七斗，菽三斗，荅二斗，𮮐五斗，直錢一百四十；麻七斗，麥六斗，菽四斗，荅五斗，𮮐三斗，直錢一百二十八；麻三斗，麥五斗，菽七斗，荅六斗，𮮐四斗，直錢一百一十六；麻二斗，麥五斗，菽三斗，答九斗，𮮐四斗，直錢一百一十二；麻一斗，麥三斗，菽二斗，荅八斗，𮮐五斗，直錢九十五。問一斗直幾何？荅曰：麻一斗七錢，麥一斗四錢，菽一斗三錢，荅一斗五錢，𮮐一斗六錢。\end{unitzh}
\begin{uniten}Suppose 9 dou of hemp (ma 麻), 7 dou of wheat (mai 麥), 3 dou of beans (shu 菽), 2 dou of small beans (da 荅) and 5 dou of millet (shu 𮮐) are worth 140 qian; 7 dou of hemp, 6 dou of wheat, 4 dou of beans, 5 dou of small beans and 3 dou of millet are worth 128 qian; 3 dou of hemp, 5 dou of wheat, 7 dou of beans, 6 dou of small beans and 4 dou of millet are worth 116 qian; 2 dou of hemp, 5 dou of wheat, 3 dou of beans, 9 dou of small beans and 4 dou of millet are worth 112 qian; 1 dou of hemp, 3 dou of wheat, 2 dou of beans, 8 dou of small beans and 5 dou of millet are worth 95 qian. Question: what is 1 dou worth? Answer: Hemp, 1 dou: 7 qian. Wheat, 1 dou: 4 qian. Beans, 1 dou: 3 qian. Small beans, 1 dou: 5 qian. Millet, 1 dou: 6 qian.\end{uniten}

\begin{unitzh}\unitlabel{jiuzhang-035.12}術曰：如方程，以正負術入之。{\zhsmall 此麻麥與均輸、少廣章之重衰、積分，皆爲大事。其拙於精理徒按本術者，或用算而布氊，方好煩而喜誤，曾不知其非，反欲以多爲貴。故其算也，莫不闇於設通而專於一端。至於此類，苟務其成，然或失之，不可謂要約。更有異術者，庖丁解牛，游刅理間，故能歷久其刅如新。夫數，猶刅也，易簡用之則動中庖丁之理，故能和神愛刅，速而寡尤。凡《九章》爲大事，按法皆不盡一百算也。雖布算不多，然足以算多。世人多以方程爲難，或盡布算之象在綴正負而已，未暇以論其設動無方，斯膠柱調瑟之類。聊復恢演，爲作新術，著之於此，將亦啓導疑意。網羅道精，豈傳之空言？記其施用之例，著䇿之數，每舉一隅焉。}\end{unitzh}
\begin{uniten}Method: As in rectangular arrays, enter it with the positive and negative method. [Liu Hui: This [problem of] hemp and wheat, together with the repeated proportions (chongcui 重衰) and the accumulated fractions (jifen 積分) of the chapters Junshu ('Fair Levies') and Shaoguang ('Lesser Breadth'), are all major tasks. Those who are clumsy in refined principles and merely follow the original method sometimes use counting rods spread out on a felt [mat]; they are fond of complication and prone to error, never realizing that this is wrong, and on the contrary wish to regard quantity as valuable. Therefore in their computations none fail to be in the dark about setting up general [procedures] (shetong 設通), and they concentrate on a single end. When it comes to this kind [of problem], if they strive for its completion, they may still miss it; this cannot be called concise. Moreover, as for having a different method: Cook Ding cut up an ox, roaming his blade among the interstices of the grain, and therefore after a long time his blade was as if new. Now numbers are like the blade: if one uses them with ease and simplicity, every movement accords with Cook Ding's principle. Therefore one can harmonize the spirit and cherish the blade, being quick and having few faults. In all the major tasks of the Nine Chapters, following the rules, none uses as many as 100 counting rods. Although the counting rods laid out are not many, they suffice to compute a great deal. People of the world mostly regard rectangular arrays as difficult, or [hold that] the whole image of laying out the counting rods lies merely in arranging positive and negative, and have no leisure to discuss how its setting-up and moving have no fixed pattern; this is of the kind of tuning a se zither with glued bridges. So I have once more expanded upon it and made a new method, recording it here; it may also open up and guide the doubting mind. Gathering in the essence of the Way as with a net — how could it be transmitted in empty words? I record examples of its application and write down the number of counting sticks (ce 䇿), in each case citing one corner.]\end{uniten}

\begin{unitzh}\unitlabel{jiuzhang-035.13}{\zhsmall 方程新術曰：以正負術入之。令左右相減，先去下實，又轉去物位，求其一行二物正負相借者，易其相當之率。又令二物與他行互相去取，轉其二物相借之數，卽皆相當之率也。各據二物相當之率，對易其數，卽各當之率也。更置減行及其下實，各以其物本率今有之，求其所同，并以爲法。其當相并而行中正負雜者，同名相從，異名相消，餘以爲法。以下實爲實。實如法，卽合所問也。一物各以本率今有之，卽皆合所問也。率不通者，齊之。}\end{unitzh}
\begin{uniten}[Liu Hui: The new method for rectangular arrays says: Enter it with the positive and negative method. Let the left and right [columns] subtract from each other: first eliminate the bottom dividend (xiashi 下實), then in turn eliminate the positions of the things (wuwei 物位), seeking a column (hang 行) in which two things, positive and negative, borrow from each other (xiangjie 相借); interchange [them] for their corresponding rates (xiangdang zhi lü 相當之率). Again let the two things and the other columns eliminate and take from each other; turning over the numbers in which the two things borrow from each other, these are all corresponding rates. According in each case to the corresponding rates of two things, exchange their numbers crosswise: these are the rates to which each corresponds. Again set up the subtracted column (jianhang 減行) and its bottom dividend; apply the rule of 'Suppose' (jinyou 今有) to each with the basic rate (benlü 本率) of its thing, seeking that in which they are the same, and add them together as the divisor. Where they should be added together but positive and negative are mixed in the column, those of the same name are combined (xiangcong 相從) and those of different names cancel each other (xiangxiao 相消); the remainder is taken as the divisor. Take the bottom dividend as the dividend. Divide the dividend by the divisor; then it agrees with what is asked. For the one thing, apply the rule of 'Suppose' to each with the basic rate; then all agree with what is asked. Where the rates are not common (tong 通), make them uniform (qi 齊).]\end{uniten}

\begin{unitzh}\unitlabel{jiuzhang-035.14}{\zhsmall 其一術曰：置羣物通率爲列衰。更置減行羣物之數，各以其率乘之，并以爲法。其當相并而行中正負雜者，同名相從，異名相消，餘爲法。以減行下實乘列衰，各自爲實。實如法而一，卽得。}\end{unitzh}
\begin{uniten}[Liu Hui: One method says: Set up the common rates (tonglü 通率) of the various things as the arrayed proportions (liecui 列衰). Again set up the numbers of the various things in the subtracted column, multiply each by its rate, and add them together as the divisor. Where they should be added together but positive and negative are mixed in the column, those of the same name are combined and those of different names cancel each other; the remainder is the divisor. Multiply the arrayed proportions by the bottom dividend of the subtracted column, each separately as a dividend. Divide the dividend by the divisor, and then it is obtained.]\end{uniten}

\begin{unitzh}\unitlabel{jiuzhang-035.15}{\zhsmall 以舊術爲之，凡應置五行。今欲要約，先置第三行，以減第四行及減第三行。次置第二行，以第二行減第三行，去其頭位。次置右行，去其頭位。次以第四行減左行頭位。次以左行去第四行及第二行頭位。次以第五行減第二行頭位，餘可半。次以第二行去第四行頭位，餘約之爲法、實，如法而一，得空，卽有𮮐價。以法減第二行，得荅價；左行得麥價；第三行麻價；右行得菽價。如此凡用七十七算。}\end{unitzh}
\begin{uniten}[Liu Hui: Doing it by the old method, altogether 5 columns should be set up. Now, wishing to be concise: first set up the third column, and with it subtract the fourth column and subtract the third column [Translator's note: thus in the text; evidently corrupt.]. Next set up the second column, and with the second column subtract the third column, eliminating its head position. Next set up the right column and eliminate its head position. Next, with the fourth column subtract the head position of the left column. Next, with the left column eliminate the head positions of the fourth and second columns. Next, with the fifth column subtract the head position of the second column; the remainder can be halved. Next, with the second column eliminate the head position of the fourth column; reduce the remainder to make the divisor and dividend; dividing by the divisor one obtains 'empty' (kong 空) [Translator's note: 空 'empty' is evidently corrupt here.], and then there is the price of millet. With the divisor subtract from the second column, obtaining the price of small beans; the left column gives the price of wheat; the third column, the price of hemp; the right column gives the price of beans. Done in this way, altogether 77 counting rods are used.]\end{uniten}

\begin{unitzh}\unitlabel{jiuzhang-035.16}{\zhsmall 以新術爲此：先以第四行減第三行。次以第三行去右行及第二行、第四行下位，又以減右行下位，不足減乃止。次以左行減第三行下位。次以第三行去左行下位。訖，廢去第三行。次以第四行去左行下位，右行當左行下位。次以右行去第二行及第四行下位。次以第二行減第四行及左行願位。次以第四行減右行菽位，不足減乃止。次以左行減第二行頭位，餘可再半。次以第四行去右行及第二行頭位。次以第二行去右行頭位，餘約之，上得五，下得三，是菽五當荅三。}\end{unitzh}
\begin{uniten}[Liu Hui: Doing this by the new method: first, with the fourth column subtract the third column. Next, with the third column eliminate the bottom positions (xiawei 下位) of the right column and of the second and fourth columns, and also with it subtract the bottom position of the right column, stopping when it is insufficient to subtract. Next, with the left column subtract the bottom position of the third column. Next, with the third column eliminate the bottom position of the left column. When this is done, discard the third column. Next, with the fourth column eliminate the bottom position of the left column; the right column corresponds to the bottom position of the left column. Next, with the right column eliminate the bottom positions of the second and fourth columns. Next, with the second column subtract the yuan (願, 'wish') positions of the fourth and left columns [Translator's note: 願位 is evidently garbled; elsewhere the text has 頭位 'head position'.]. Next, with the fourth column subtract the beans position of the right column, stopping when it is insufficient to subtract. Next, with the left column subtract the head position of the second column; the remainder can be halved twice. Next, with the fourth column eliminate the head positions of the right and second columns. Next, with the second column eliminate the head position of the right column; reduce the remainder: the top gives 5 and the bottom gives 3; that is, 5 of beans correspond to 3 of small beans.]\end{uniten}

\begin{unitzh}\unitlabel{jiuzhang-035.17}{\zhsmall 次以左行去第三行菽位，又以減第四行及右行菽位，不足減乃止。次以右行減第二行頭位，不足減乃止。次以第三行去左行頭位。次以左行去右行頭位，餘上得六，下得五，是爲荅六當𮮐五。次以右行去左行荅位，餘約之，上爲二，下爲三。次以左行去第二行下位。以第二行去第四行下位，又以減左行下位。次右行去第二行下位，餘上得三，下得四，是爲麥三當菽四。次以第二行減第四行下位。次以第四行去第二行下位，餘上得四，下得七，是爲麻四當麥七。是爲相當之率舉矣。據麻四當麥七，卽爲麻價率七而麥價率四；又麥三當菽四，卽爲麥價率四而菽價率三；又荅六當𮮐五，卽爲荅價率五而𮮐價率六；而率通矣。}\end{unitzh}
\begin{uniten}[Liu Hui: Next, with the left column eliminate the beans position of the third column, and also with it subtract the beans positions of the fourth and right columns, stopping when it is insufficient to subtract. Next, with the right column subtract the head position of the second column, stopping when it is insufficient to subtract. Next, with the third column eliminate the head position of the left column. Next, with the left column eliminate the head position of the right column; of the remainder, the top gives 6 and the bottom gives 5: this means 6 of small beans correspond to 5 of millet. Next, with the right column eliminate the small-beans position of the left column; reduce the remainder: the top is 2 and the bottom is 3. Next, with the left column eliminate the bottom position of the second column. With the second column eliminate the bottom position of the fourth column, and also with it subtract the bottom position of the left column. Next, the right column eliminates the bottom position of the second column; of the remainder, the top gives 3 and the bottom gives 4: this means 3 of wheat correspond to 4 of beans. Next, with the second column subtract the bottom position of the fourth column. Next, with the fourth column eliminate the bottom position of the second column; of the remainder, the top gives 4 and the bottom gives 7: this means 4 of hemp correspond to 7 of wheat. Thus the corresponding rates are all set forth. On the basis that 4 of hemp correspond to 7 of wheat, the hemp price rate is 7 and the wheat price rate 4; again, 3 of wheat correspond to 4 of beans, so the wheat price rate is 4 and the beans price rate 3; again, 6 of small beans correspond to 5 of millet, so the small-beans price rate is 5 and the millet price rate 6; and the rates are common.]\end{uniten}

\begin{unitzh}\unitlabel{jiuzhang-035.18}{\zhsmall 更置第三行，以第四行減之，餘有麻一斗、菽四斗、荅三斗負，下實四正。求其同爲麻之數，以菽率三、荅率五各乘菽、荅斗數，如麻率七而一，菽得一斗七分斗之五正，荅得二斗七分斗之一負，則菽、荅化爲麻。以并之，令同名相從，異名相消，餘爲定麻七分斗之四，以爲法。置下實四爲實，以分母乘之，實得二十八，而分子化爲法矣。以法除，得七，卽麻一斗之價。置麥率四、菽率三、荅率五、𮮐率六，皆以其斗數乘之，各自爲實，以麻率七爲法，所得卽同爲麻之數。}\end{unitzh}
\begin{uniten}[Liu Hui: Again set up the third column and subtract the fourth column from it: the remainder has hemp 1 dou, beans 4 dou, small beans 3 dou negative, and the bottom dividend 4 positive. To seek the amount of them that is the same as hemp, multiply the numbers of dou of beans and small beans by the beans rate 3 and the small-beans rate 5 respectively, and divide by the hemp rate 7: the beans give 1 5/7 dou positive and the small beans give 2 1/7 dou negative; thus the beans and small beans are transformed into hemp. Add them together, letting those of the same name be combined and those of different names cancel each other; the remainder is the settled hemp, 4/7 dou, which is taken as the divisor. Set up the bottom dividend 4 as the dividend and multiply it by the denominator: the dividend becomes 28, and the numerator is transformed into the divisor. Divide by the divisor, obtaining 7, which is the price of 1 dou of hemp. Set up the wheat rate 4, the beans rate 3, the small-beans rate 5 and the millet rate 6, multiply each by its number of dou, each separately as a dividend, and take the hemp rate 7 as the divisor; what is obtained is the amount that is the same as hemp.]\end{uniten}

\begin{unitzh}\unitlabel{jiuzhang-035.19}{\zhsmall 亦可使置本行實與物，同通之，各以本率今有之，求其本率所得，并以爲法。如此則無正負之異矣，擇異同而已。又可以一術爲之：置五行通率，爲麻七、麥四、菽三、荅五、𮮐六，以爲列衰。減行麻一斗，菽四斗正，荅三斗負，各以其率乘之。訖，令同名相從，異名相消，餘爲法。又置下實乘列衰，所得各爲實。此可以實約法，則不復乘列衰，各以列衰如所約，知其價。如此則凡用一百二十四算也。}\end{unitzh}
\begin{uniten}[Liu Hui: One may also set up the dividend and the things of the original column, make them all common, apply the rule of 'Suppose' to each with its basic rate, seek what is obtained with their basic rates, and add these together as the divisor. In this way there is no difference of positive and negative; one only distinguishes the different and the same. It can also be done by the one method: set up the common rates of the five columns, namely hemp 7, wheat 4, beans 3, small beans 5, millet 6, as the arrayed proportions. In the subtracted column, hemp 1 dou, beans 4 dou positive and small beans 3 dou negative: multiply each by its rate; when this is done, let those of the same name be combined and those of different names cancel each other; the remainder is the divisor. Again set up the bottom dividend and multiply the arrayed proportions by it; what is obtained is in each case a dividend. Here the divisor can be reduced by the dividend, and then one does not multiply the arrayed proportions again: by each arrayed proportion, as reduced, one knows its price. In this way altogether 124 counting rods are used.]\end{uniten}

\begin{unitzh}\unitlabel{jiuzhang-035.20}《九章算術》卷八訂訛 休寧 載震 東原\end{unitzh}
\begin{uniten}Nine Chapters on the Mathematical Art, Chapter 8: Corrections of Errors (Ding'e 訂訛). Xiuning, Zai Zhen, Dongyuan. [Translator's note: 載 is evidently an error for 戴; the author is Dai Zhen, courtesy name Dongyuan, of Xiuning.]\end{uniten}

\begin{unitzh}\unitlabel{jiuzhang-035.21}術曰：置上禾三秉，中禾二秉，下禾一秉，實三十九斗，於右方。中、左禾列如右方。以右行上禾徧乘中行，而以直除。{\zhsmall 古字直、値通用，直除猶言對減也。以右行上禾徧乘中行，復以中行上禾徧乘右行，然後可相對減。古人文省，故但舉一以該之。}\end{unitzh}
\begin{uniten}Method: Set out 3 sheaves of top-grade grain, 2 sheaves of middle-grade grain, 1 sheaf of low-grade grain, and the yield, 39 dou, in the right column (youfang 右方). The middle and left [columns of] grain are arrayed like the right column. Multiply the middle column throughout by the top-grade grain of the right column, and then eliminate directly (zhichu 直除). [Dai Zhen: In the ancient script 直 and 値 were used interchangeably; zhichu is like saying 'subtract against each other' (duijian 對減). After multiplying the middle column throughout by the top-grade grain of the right column, one again multiplies the right column throughout by the top-grade grain of the middle column; only then can they be subtracted against each other. The ancients' wording was sparing, so they cited only one to cover them.]\end{uniten}

\begin{unitzh}\unitlabel{jiuzhang-036.1}又乘其次，亦以直除。{\zhsmall 此謂右行上禾徧乘左行，復以左行上禾徧乘右行，亦相對減。古人文多省畧。}\end{unitzh}
\begin{uniten}Again multiply the next one, and likewise eliminate directly (zhichu). [Dai Zhen: This means that the top-grade grain of the right column multiplies throughout the left column, and again the top-grade grain of the left column multiplies throughout the right column, and these too are subtracted against each other. The writing of the ancients is often abbreviated.]\end{uniten}

\begin{unitzh}\unitlabel{jiuzhang-036.2}然以中行中禾不盡者{\zhsmall 古人單用然字，不曰然後者，然猶乃也。}徧乘左行。{\zhsmall 此以中行、左行所減之餘，如前徧乘，文亦從省。}\end{unitzh}
\begin{uniten}Then (ran), with the middle-grade grain of the middle column that is not exhausted, [Dai Zhen: The ancients used the character ran on its own; where it does not say ranhou ("after that"), ran is like nai ("thereupon").] multiply throughout the left column. [Dai Zhen: This means that with the remainders left after the subtractions in the middle column and the left column one multiplies throughout as before; the wording, too, is abbreviated.]\end{uniten}

\begin{unitzh}\unitlabel{jiuzhang-036.3}用算繁而不省，所以别爲法，約也。然猶不如自用其舊，廣異法也。{\zhsmall 「用法繁而不省」以下，訛舛衍文。}\end{unitzh}
\begin{uniten}The use of counting rods is cumbersome and not economical; that is why a separate method is made: for economy. Yet it is still not as good as simply using the old one; it serves to broaden the different methods. [Dai Zhen: From "the use of the method is cumbersome and not economical" onward is corrupt, garbled and superfluous text.]\end{uniten}

\begin{unitzh}\unitlabel{jiuzhang-036.4}益行減行，當各以其類矣。其異名者，非其類也。非其類者，猶無對也，非所得減也。故赤用黑對則除黑，無對則除赤。赤黑并於本數，此爲相益之。皆所以爲消奪，消奪之與減益成一實也。{\zhsmall 此注多訛舛。據方程術，無論物有幾品，遞減至一物乃止；又以赤黑别正負。首位赤減赤、黑減黑者，同名伯除也；次位以下遇赤用黑對則相益，此條是也。首位赤減黑、黑減赤者，異名相除也；次位以下遇皆赤皆黑則相益，後條是也。由此言之，注之謬顯然。葢傳寫失真，後人復妄加改竄，遂不可通。}\end{unitzh}
\begin{uniten}In augmenting columns and diminishing columns, each should go according to its kind. Those of unlike names (yiming) are not of its kind. Those not of its kind are as if without counterpart; they are not something that can be subtracted. Therefore, when red has black as its counterpart, eliminate the black; when there is no counterpart, eliminate the red. Red and black are combined with the original number: this is adding them to one another. All these are means of cancelling and taking away (xiaoduo); cancelling and taking away, together with subtracting and adding, form one single substance (shi). [Dai Zhen: This commentary has many corruptions and confusions. According to the method of rectangular arrays (fangcheng), no matter how many kinds of things there are, one subtracts successively until one thing is reached, and then stops; further, red and black distinguish positive and negative. When in the first position red subtracts red or black subtracts black, this is "like names (tongming) bo eliminate"; from the second position downward, when one meets red having black as counterpart, they are added to one another: that is the present clause. When in the first position red subtracts black or black subtracts red, this is "unlike names eliminate one another"; from the second position downward, when one meets both red or both black, they are added to one another: that is the following clause. Speaking from this, the error of the commentary is evident. Probably it lost its true form through copying, and later people further rashly altered it, so that it became unintelligible.] [Translator's note: 伯 (bo) in 同名伯除 is evidently an error for 相 "one another".]\end{uniten}

\begin{unitzh}\unitlabel{jiuzhang-036.5}正無入負之。{\zhsmall 入字原本訛作人，下文及注並同。據注云「無入爲無對也」，無對之說亦未分曉。釋方程者專爲遇空位起例，而左右兩行相對減，或正宜變爲負，或負宜變爲正，往往不得其義例。今考同名相除、異名相益者，如下實左右俱正，所減之餘屬左行則去右行，屬右行則去左行。其物品以正減正，餘在所去之行，爲正無入；以負減負，餘在所去之行，爲負無入；以正從負，爲正無入；以負從正，爲負無入。負對空位而負數在所去之行，與以負減負同例；正對空位而正數在所去之行，與以正從負同例。此皆所謂正無入負之、負無入正之也。異名相除、同名相益者，如下實左右俱正，并爲一數，則無分於左右。其物品以負減正，餘或左或右，爲正無入；以正減負，餘或左或右，爲負無入；以正從正，爲正無入；以負從負，爲負無入。正對空位，與以負減正同例；負對空位，與以負從負同例。此皆所謂正無入正之、負無入負之也。由是言之，在所去之行則其數無入，而或左或右以與無分於左右合爲一行，因亦謂之無入。人字乃傳寫之誤，明矣。今悉改正。}\end{unitzh}
\begin{uniten}Positive with no entry (wu ru): make it negative. [Dai Zhen: The character ru (入, "enter") was miswritten in the original edition as ren (人, "person"); likewise in the text below and in the commentary. According to the commentary, which says "No entry means having no counterpart", the explanation by "no counterpart" is also not clear. Those who explain rectangular arrays set up the rule specifically for meeting an empty position (kongwei), whereas when the two columns, left and right, are subtracted against each other, sometimes a positive ought to change into a negative and sometimes a negative ought to change into a positive, and they often fail to grasp its meaning and rule. Now, examining "like names eliminate one another, unlike names add to one another": for example, if the lower dividends (xia shi) on the left and right are both positive, then if the remainder of the subtraction belongs to the left column, the right column is removed, and if it belongs to the right column, the left column is removed. As for the kinds of things (wupin): subtracting a positive from a positive, with the remainder in the column that is removed, is "positive with no entry"; subtracting a negative from a negative, with the remainder in the column that is removed, is "negative with no entry"; joining a positive to a negative is "positive with no entry"; joining a negative to a positive is "negative with no entry". A negative facing an empty position, with the negative number in the column that is removed, is the same case as subtracting a negative from a negative; a positive facing an empty position, with the positive number in the column that is removed, is the same case as joining a positive to a negative. These are all what is meant by "positive with no entry, make it negative; negative with no entry, make it positive". As for "unlike names eliminate one another, like names add to one another": for example, if the lower dividends on the left and right are both positive and are combined into one number, then there is no distinction between left and right. As for the kinds of things: subtracting a negative from a positive, with the remainder either on the left or on the right, is "positive with no entry"; subtracting a positive from a negative, with the remainder either on the left or on the right, is "negative with no entry"; joining a positive to a positive is "positive with no entry"; joining a negative to a negative is "negative with no entry". A positive facing an empty position is the same case as subtracting a negative from a positive; a negative facing an empty position is the same case as joining a negative to a negative. These are all what is meant by "positive with no entry, make it positive; negative with no entry, make it negative". Speaking from this: when it is in the column that is removed, its number has no entry; and when it is either on the left or on the right, since, with no distinction between left and right, it is combined into a single column, it is therefore also called "no entry". That the character ren is a copying error is clear. Now all have been corrected.]\end{uniten}

\begin{unitzh}\unitlabel{jiuzhang-036.6}無入爲無對也，無所得減，則使消奪者居位也。\end{unitzh}
\begin{uniten}"No entry" means having no counterpart. When there is nothing one can subtract from, one lets that which cancels and takes away occupy the position.\end{uniten}

\begin{unitzh}\unitlabel{jiuzhang-036.7}其當以列實或減下實{\zhsmall 此句訛舛。據後注内方程新術，以列衰乘下實，又以列衰乘羣物之數，并爲法；其當相并而行中正負雜者，同名相從，異名相消。似卽此所舉。}\end{unitzh}
\begin{uniten}Where one should use the arrayed dividends (lie shi) or reduce the lower dividend [Dai Zhen: This sentence is corrupt. According to the new method for rectangular arrays (fangcheng xin shu) in the later commentary, one multiplies the lower dividend by the arrayed proportions (lie shuai), and also multiplies the numbers of the various things by the arrayed proportions and sums them to make the divisor (fa); where they should be summed but positive and negative are mixed within the column, like names join one another and unlike names cancel one another. This seems to be precisely what is cited here.]\end{uniten}

\begin{unitzh}\unitlabel{jiuzhang-036.8}此條異名相除爲例，故亦與上條互取。凡正負所以記其同異，使二品互相取而已矣。言負者未必負於少，言正者未必正於多。故每一行之中，雖復赤黑異算，無傷。然則可得使頭位常相與異名，此條之實兼通矣。遂以二條反覆一率。觀其每與上下互相取位，則隨算而言耳，猶一術也。又本設諸行，欲因減數以相去耳，故其多少無限，令上下相命而已。若以正負相減，如數有舊增法者，每行可均之，不但數物左右之也。{\zhsmall 此注皆不分曉，當由傳寫失真，後人又妄加改竄，遂不可通。}\end{unitzh}
\begin{uniten}This clause takes "unlike names eliminate one another" as its pattern; therefore it too is taken reciprocally with the clause above. In general, positive and negative serve to record their sameness and difference, so as to make the two kinds take from one another, and that is all. What is called negative is not necessarily negative by being less; what is called positive is not necessarily positive by being more. Therefore, within each single column, even if the red and black rods differ, there is no harm. That being so, one can make the head positions always be of unlike names to each other, and the substance (shi) of this clause then applies to both. Thus the two clauses, taken back and forth, are one rate (lü). Observe how each time it takes positions reciprocally with those above and below: it is then merely spoken of according to the counting rods; it is like a single method. Furthermore, the columns were originally set up wishing only to remove one another by subtracting numbers; therefore their amounts are without limit, and one merely lets the upper and lower name one another. If one subtracts with positive and negative, as when the numbers have an old method of augmenting, each column can be equalized; it is not merely placing the several things left and right. [Dai Zhen: This commentary is entirely unclear; it must be because it lost its true form through copying and later people further rashly altered it, so that it became unintelligible.]\end{uniten}

\begin{unitzh}\unitlabel{jiuzhang-036.9}此中行買賣相折，錢適足，但互買賣算而已，故下無錢直也。設欲以此行如方程法，先令牛二徧乘左行，而以右行直除之，是終於下實虚缺矣。故注曰「正無實負，負無實正」，方爲類也。方將以别實加不足之數與實物作實。{\zhsmall 此注訛脫不分曉。據術意，應列三行。先令右行牛二徧乘中行，復令中行牛三徧乘右行，而以直除，得羊三十三正，豕四十五負，餘錢三千正。此同名相除，異名相益，正無入負之，負無入正之也。次令右行牛二徧乘左行，復令左行牛五徧乘右行，而以直除，得羊三十七正，豕四十九負，餘錢三千八百正。此異名相除，同名相益，正無入正之，負無入負之也。重列爲左右兩行，先令右行羊三十三徧乘左行，復令左行羊三十七徧乘右行，而以直除，得豕四十八正，以爲法；錢一萬四千四百正，爲實。實如法而一，得豕價三百。轉減而上，得羊價五百，牛價一千二百。此亦同名相除，異名相益，正無入負之，負無入正之也。中行下實虛缺，本無庸論。葢注文傳寫失眞，後人又妄加改竄，遂不可通。}\end{unitzh}
\begin{uniten}In this middle column, buying and selling offset each other and the cash (qian, 錢) is just sufficient; there are only the counting rods of the mutual buying and selling, and therefore at the bottom there is no cash value. Suppose one wishes to treat this column by the method of rectangular arrays: first let the 2 oxen multiply throughout the left column, and eliminate it directly with the right column; this ends with the lower dividend being empty and lacking. Therefore the commentary says "positive with no dividend, negative; negative with no dividend, positive"; it is precisely of this kind. One is then to take a separate dividend, add the amount of the deficit, and together with the actual things make the dividend. [Dai Zhen: This commentary is corrupt, has lacunae and is unclear. According to the intent of the method, three columns should be set out. First let the 2 oxen of the right column multiply throughout the middle column, then let the 3 oxen of the middle column multiply throughout the right column, and eliminate directly: one obtains sheep 33 positive, pigs 45 negative, surplus cash 3000 positive. This is "like names eliminate one another, unlike names add to one another; positive with no entry, make it negative; negative with no entry, make it positive". Next let the 2 oxen of the right column multiply throughout the left column, then let the 5 oxen of the left column multiply throughout the right column, and eliminate directly: one obtains sheep 37 positive, pigs 49 negative, surplus cash 3800 positive. This is "unlike names eliminate one another, like names add to one another; positive with no entry, make it positive; negative with no entry, make it negative". Set them out again as two columns, left and right. First let the 33 sheep of the right column multiply throughout the left column, then let the 37 sheep of the left column multiply throughout the right column, and eliminate directly: one obtains pigs 48 positive, which is taken as the divisor (fa); cash 14400 positive is the dividend (shi). Dividing the dividend by the divisor, one obtains the price of a pig, 300. Subtracting in turn upward, one obtains the price of a sheep, 500, and the price of an ox, 1200. This too is "like names eliminate one another, unlike names add to one another; positive with no entry, make it negative; negative with no entry, make it positive". That the lower dividend of the middle column is empty and lacking was in itself not worth discussing. Probably the commentary text lost its true form through copying, and later people further rashly altered it, so that it became unintelligible.]\end{uniten}

\begin{unitzh}\unitlabel{jiuzhang-036.10}此四雀一燕與一雀五燕衡適平，并重一斤，故各八兩。列兩行程數，左行頭位其數是一，可省乘，令右行徧乘左行，而取其法實於左。左行數多，以右行取其數，左頭位減盡，中法下實，卽每枚當重，宜可知也。按此四雀一燕與一雀五燕其重等，是三雀四燕重相當，雀率重四，燕率重三也。諸再程之率，皆可異術求之，卽其數也。{\zhsmall 此卽末條所言新術。以下實俱重八兩徧乘上雀燕，以左雀八減右雀三十二，餘二十四；以右燕八減左燕四十，餘三十二。是爲二物正負相借，因而約之，雀得三，燕得四，乃三雀當四燕也。對易其數，卽雀率重四，而燕率重三。注於此突入異術，幾不解其所謂。}\end{unitzh}
\begin{uniten}These 4 sparrows with 1 swallow and 1 sparrow with 5 swallows: the balance is exactly level, and together they weigh 1 jin (斤); hence each is 8 liang (兩). Set out the measured numbers (cheng shu) in two columns. In the left column the number in the head position is 1, so the multiplication can be dispensed with; let the right column multiply throughout the left column, and take the divisor (fa) and dividend (shi) on the left. The numbers of the left column are larger; with the right column take from its numbers. The head position on the left is subtracted until exhausted; with the middle as divisor and the lower as dividend, the weight that each piece amounts to should then be knowable. Consider: these 4 sparrows with 1 swallow and 1 sparrow with 5 swallows are equal in weight; this means that 3 sparrows and 4 swallows correspond in weight: the sparrow's rate (lü) is a weight of 4, and the swallow's rate a weight of 3. The rates of all double measures (zai cheng) can all be found by a different method; these are precisely their numbers. [Dai Zhen: This is precisely the new method spoken of in the last item. With the lower dividends, both of weight 8 liang, multiply throughout the sparrows and swallows above: subtract the 8 sparrows of the left from the 32 sparrows of the right, remainder 24; subtract the 8 swallows of the right from the 40 swallows of the left, remainder 32. These are the two things borrowing from one another (xiangjie) as positive and negative. Reducing them accordingly, the sparrows give 3 and the swallows 4; that is, 3 sparrows correspond to 4 swallows. Exchanging their numbers crosswise (duiyi), the sparrow's rate is a weight of 4 and the swallow's rate a weight of 3. Here the commentary abruptly brings in a different method; one can hardly understand what it means.]\end{uniten}

\begin{unitzh}\unitlabel{jiuzhang-037.1}今有五家共井，甲二綆不足，如乙一綆；乙三綆不足，如丙一綆；丙四綆不足，如丁一綆；丁五綆不足，如戊一綆；戊六綆不足，如甲一綆。各得所不足一綆，皆逮。問井深、綆長各幾何？荅曰：井深七丈二尺一寸，甲綆長二丈六尺五寸，乙綆長一丈九尺一寸，丙綆長一丈四尺八寸，丁綆長一丈二尺九寸，戊綆長七尺六寸。\end{unitzh}
\begin{uniten}Suppose there are 5 families sharing a well. A's 2 well-ropes (geng 綆) fall short by as much as B's 1 rope; B's 3 ropes fall short by as much as C's 1 rope; C's 4 ropes fall short by as much as D's 1 rope; D's 5 ropes fall short by as much as E's 1 rope; E's 6 ropes fall short by as much as A's 1 rope. If each gets the 1 rope by which it falls short, all reach (dai 逮) [the bottom]. What are the depth of the well and the length of each rope? Answer: The depth of the well is 7 zhang (丈) 2 chi (尺) 1 cun (寸); A's rope is 2 zhang 6 chi 5 cun long; B's rope is 1 zhang 9 chi 1 cun long; C's rope is 1 zhang 4 chi 8 cun long; D's rope is 1 zhang 2 chi 9 cun long; E's rope is 7 chi 6 cun long.\end{uniten}

\begin{unitzh}\unitlabel{jiuzhang-037.2}{\zhsmall 此問不言丈尺，無由知井深、綆長於丈尺幾何。使井深半之，爲三丈六尺有半寸，則甲綆一丈三尺二寸半，乙綆九尺五寸半，丙綆七尺四寸，丁綆六尺四寸半，戊綆三尺八寸；使井深倍之，爲十四丈四尺二寸，則甲綆五丈三尺，乙綆三丈八尺二寸，丙綆二丈九尺六寸，丁綆二丈五尺八寸，戊綆一丈五尺二寸，皆合所問。由是言之，問旣不定以丈尺，依術推求，先得七百二十一，無以定百爲丈、十爲尺也。問井深、綆長之率各幾何，荅以井深之率七百二十一，甲綆長率二百六十五，乙綆長率一百九十一，丙綆長率一百四十八，丁綆長率一百二十九，戊綆長率七十六，於義乃通。}\end{unitzh}
\begin{uniten}[Note: This problem does not speak of zhang and chi, so there is no way of knowing how many zhang and chi the depth of the well and the lengths of the ropes are. Suppose the depth of the well is halved, making 3 zhang 6 chi and 1/2 cun; then A's rope is 1 zhang 3 chi 2 1/2 cun, B's rope 9 chi 5 1/2 cun, C's rope 7 chi 4 cun, D's rope 6 chi 4 1/2 cun, E's rope 3 chi 8 cun. Suppose the depth of the well is doubled, making 14 zhang 4 chi 2 cun; then A's rope is 5 zhang 3 chi, B's rope 3 zhang 8 chi 2 cun, C's rope 2 zhang 9 chi 6 cun, D's rope 2 zhang 5 chi 8 cun, E's rope 1 zhang 5 chi 2 cun. All agree with what is asked. Speaking from this: since the problem is not determined in zhang and chi, when one deduces according to the method one first obtains 721, and there is no way to fix that the hundreds are zhang and the tens are chi. If one asked "What are the rates (lü) of the depth of the well and of the lengths of the ropes?" and answered with: the rate of the depth of the well 721, the rate of the length of A's rope 265, the rate of the length of B's rope 191, the rate of the length of C's rope 148, the rate of the length of D's rope 129, the rate of the length of E's rope 76, only then would it make sense.]\end{uniten}

\begin{unitzh}\unitlabel{jiuzhang-037.3}此率，初如方程爲之，名各一逮井。其後，法得七百二十一，實七十六。{\zhsmall 此上訛舛不可通。據術，先得七百二十一，爲所列五行之通率，卽井深率也。以此率列各行下，爲各行之下實，重求之，法得七百二十一，實得五萬四千七百九十六，以法除實，得用逮之數七十六。}\end{unitzh}
\begin{uniten}These rates are at first worked as in rectangular arrays (fangcheng), each being named "1 reaching the well". After that, the divisor (fa) obtained is 721 and the dividend (shi) 76. [Note: The above is corrupt and confused and cannot be made sense of. According to the method, one first obtains 721, which is the common rate (tonglü 通率) of the 5 rows set out, that is, the rate of the depth of the well. Set this rate at the bottom of each row as the lower dividend (xiashi 下實) of each row and seek it again: the divisor obtained is 721 and the dividend obtained 54796; dividing the dividend by the divisor gives the number used for reaching, 76.]\end{uniten}

\begin{unitzh}\unitlabel{jiuzhang-037.4}此麻麥與均輸、少廣章之重衰、積分，皆爲大事。其拙於精理、徒按本術者，或用算而布氊，方好煩而喜誤，曾不知其非，反欲以多爲貴。故其算也，莫不闇於設通而專於一端。至於此類，苟務其成，然或失之，不可謂要約。更有異術者，庖丁解牛，游刃理閒，故能歷久其刃如新。夫數猶刃也，易簡用之，則動中庖丁之理，故能和神愛刃，速而寡尤。凡九章爲大事，按法皆不盡一百算也。雖布算不多，然足以算多。世人多以方程爲難，或盡布算之象在綴正負而已，未暇以論其設動無方，斯膠柱調瑟之類。聊復恢演，爲作新術，著之於此。將亦啓導疑意，網羅道精，豈傳之空言？記其施用之例，著䇿之數，每舉一隅焉。{\zhsmall 以上字句多誤，又皆屬虛辭，非有實義可考，無從訂正，姑仍之。}\end{unitzh}
\begin{uniten}This [problem of] hemp (ma 麻) and wheat (mai 麥), and the repeated proportions (chongcui 重衰) and accumulated fractions (jifen 積分) of the chapters Fair Transport (Junshu 均輸) and Short Width (Shaoguang 少廣), are all major tasks (dashi 大事). Those who are clumsy in refined principles (jingli 精理) and merely follow the basic method sometimes use counting rods (suan 算) and spread out a felt mat; they are just fond of complication and prone to error, and, never realizing that this is wrong, on the contrary wish to take the many as valuable. Therefore, in their calculating, all without exception are in the dark about general setting-up (shetong 設通) and fixed on a single aspect. When it comes to this kind [of problem], they merely strive to complete it, yet sometimes they fail at it; this cannot be called concise (yaoyue 要約). Moreover, as for having a different method (yishu 異術): Cook Ding, cutting up an ox, let his blade roam in the interstices of the grain (li 理); therefore he could go on for a long time with his blade as if new. Now, numbers are like the blade: if one uses them with ease and simplicity, then every movement hits Cook Ding's grain; therefore one can keep the spirit harmonious and spare the blade, being quick and with few faults. All the major tasks of the Nine Chapters, done according to the rules (fa 法), never use up 100 calculations (suan). Although the calculations laid out are not many, yet they suffice to calculate much. People of the world mostly take rectangular arrays (fangcheng) to be difficult; for some, the whole configuration (xiang 象) of the laid-out calculation lies merely in arranging (zhui 綴) positive and negative, and they have had no leisure to discuss that its setting-up and moves have no fixed pattern; this is of the kind of "tuning the se zither with glued bridges". For what it is worth, I have further expanded and developed [this], made a new method (xinshu 新術) for it, and set it down here. It will also open up and guide doubting minds and net in the essence of the Way; how could it be the handing down of empty words? I record examples of its application and set down the numbers of rods (ce 䇿), in each case raising one corner. [Note: The words and phrases above contain many errors and, moreover, all belong to empty rhetoric; there is no substantive meaning that can be examined and no way to correct them; we provisionally keep them as they are.]\end{uniten}

\begin{unitzh}\unitlabel{jiuzhang-037.5}其一術曰：置羣物通率爲列衰，更置減行羣物之數，各以其率乘之，并以爲法。其當相并而行中正負雜者，同名相從，異名相消，餘爲法。以減行下實乘列衰，各自爲實。實如法而一，卽得。\end{unitzh}
\begin{uniten}One of its methods says: Set the common rates of the various things as the arrayed proportions (liecui 列衰). Again set the numbers of the various things in the subtraction row (jianxing 減行), multiply each by its rate, and add them together to make the divisor. Where they are to be added together but positive and negative are mixed within the row, those of the same name are joined (cong 從) and those of different names cancel each other (xiao 消); the remainder is the divisor. Multiply the arrayed proportions by the lower dividend of the subtraction row; each separately makes a dividend. Divide the dividend by the divisor, and one obtains [the result].\end{uniten}

\begin{unitzh}\unitlabel{jiuzhang-037.6}以舊術爲之，凡應置五行。今欲要約，先置第三行，以減第四行及減第三行。次置第二行，以第二行減第三行，去其頭位。次置右行，去其頭位。次以第四行減左行頭位。次以左行去第四行及第二行頭位。次以第五行減第二行頭位，餘可半。次以第二行去第四行頭位，餘約之爲法、，實如法而一，得空，卽有𮮐價。以法減第二行得荅價，左行得麥價，第三行麻價，右行得菽價。如此凡用七十七算。\end{unitzh}
\begin{uniten}Doing it by the old method (jiushu 舊術), in all one should set out 5 rows. Now, wishing to be concise: first set out the 3rd row, and with it subtract from the 4th row and subtract from the 3rd row. Next set out the 2nd row, and subtract the 2nd row from the 3rd row, eliminating its head position (touwei 頭位). Next set out the right row, eliminating its head position. Next, with the 4th row subtract from the head position of the left row. Next, with the left row eliminate the head positions of the 4th row and the 2nd row. Next, with the 5th row subtract from the head position of the 2nd row; the remainder can be halved. Next, with the 2nd row eliminate the head position of the 4th row; simplify (yue 約) the remainder to make the divisor and the dividend. Divide the dividend by the divisor, obtaining "empty" (kong 空); then there is the price of millet (𮮐). Subtracting from the 2nd row by the rule (fa) gives the price of small beans (da 荅); the left row gives the price of wheat, the 3rd row the price of hemp, and the right row gives the price of beans (shu 菽). Done in this way, in all 77 calculations (suan) are used.\end{uniten}

\begin{unitzh}\unitlabel{jiuzhang-037.7}{\zhsmall 以上所言舊術訛舛不可通。據方程術，凡五物及總價求其各物之價者，應列五行，行五位及價直。以上一位互乘，因徧乘次位以下及價直，兩兩相對減去其頭位，所減之餘重列之，減至一物一價乃止。物爲法，價爲實，實如法而一，得一物之價。轉減而上，以知各價。先化五爲四，次化四爲三，次化三爲二，次化二爲一，凡用十算；兼乘減除言之，則一百四十五算。凡上一位互乘，其數必同，可省乘：若遇上一位數同，則省徧乘；或上一位遇一，則省其與對行徧乘。考問意，左行上一位是一，先以左行減右行，次減第二行，次減第三行，次減第四行，所減之餘重列爲四行。其左行上一位又是一，以左行減右行，次減第二行，次減第三行，所減之餘重列爲三行。其上一位數皆同，卽以本數減之，餘列爲兩行，依術得荅價，轉而上，求得菽價及麥價、麻價。凡用九十九算。}\end{unitzh}
\begin{uniten}[Note: The old method stated above is corrupt and confused and cannot be made sense of. According to the method of rectangular arrays, in all [cases] where, with 5 things and the total prices, one seeks the price of each thing, one should set out 5 rows, each row with 5 positions and the price value (jiazhi 價直). Cross-multiply (hucheng 互乘) by the top positions (shang yi wei 上一位), and thereupon multiply throughout (biancheng 徧乘) the next positions (ciwei 次位) downward and the price values; facing each other two by two, subtract to eliminate their head positions; set out anew (chonglie 重列) the remainders of the subtractions; subtract until there are 1 thing and 1 price, and then stop. The thing is the divisor, the price is the dividend; divide the dividend by the divisor to obtain the price of 1 thing. Subtract in turn going upward, so as to know each price. First transform 5 into 4, next transform 4 into 3, next transform 3 into 2, next transform 2 into 1: in all 10 calculations are used; speaking of multiplications, subtractions and divisions together, then 145 calculations. Whenever the top positions are cross-multiplied, their numbers are necessarily the same; multiplication can be saved: if one meets top positions whose numbers are the same, then the multiplying throughout is saved; or if a top position happens to be 1, then the multiplying throughout of the row paired with it is saved. Examining the intent of the problem: the top position of the left row is 1; first subtract the left row from the right row, next subtract from the 2nd row, next subtract from the 3rd row, next subtract from the 4th row; the remainders of the subtractions are set out anew as 4 rows. The top position of their left row is again 1; subtract the left row from the right row, next subtract from the 2nd row, next subtract from the 3rd row; the remainders of the subtractions are set out anew as 3 rows. The numbers in their top positions are all the same; then subtract with the numbers as they stand, and the remainders are set out as 2 rows; following the method one obtains the price of small beans; turning upward, one finds the price of beans and the prices of wheat and hemp. In all, 99 calculations are used.]\end{uniten}

\begin{unitzh}\unitlabel{jiuzhang-037.8}以新術爲此：先以第四行減第三行。次以第三行去右行及第二行、第四行下位，又以減右行下位，不足減乃止。次以左行減第三行下位。次以第三行去左行下位，訖，廢去第三行。次以第四行去左行下位，右行當左行下位。次以右行去第二行及第四行下位。次以第二行減第四行及左行頭位。次以第四行減右行菽位，不足減乃止。次以左行減第二行頭位，餘可再半。次以第四行去右行及第二行頭位。次以第二行去右行頭位，餘約之，上得五，下得三，是菽五當荅三。\end{unitzh}
\begin{uniten}Doing this by the new method: first subtract the 4th row from the 3rd row. Next, with the 3rd row eliminate the lower positions (xiawei 下位) of the right row and of the 2nd and 4th rows, and further subtract it from the lower position of the right row, stopping when there is not enough to subtract. Next, with the left row subtract from the lower position of the 3rd row. Next, with the 3rd row eliminate the lower position of the left row; when this is finished, discard the 3rd row. Next, with the 4th row eliminate the lower position of the left row; the right row is equivalent (dang 當) to the lower position of the left row. Next, with the right row eliminate the lower positions of the 2nd row and the 4th row. Next, with the 2nd row subtract from the head positions of the 4th row and the left row. Next, with the 4th row subtract from the bean position of the right row, stopping when there is not enough to subtract. Next, with the left row subtract from the head position of the 2nd row; the remainder can be halved twice. Next, with the 4th row eliminate the head positions of the right row and the 2nd row. Next, with the 2nd row eliminate the head position of the right row; simplify the remainder: the upper gives 5 and the lower gives 3. This is: 5 beans are equivalent to 3 small beans.\end{uniten}

\begin{unitzh}\unitlabel{jiuzhang-037.9}次以左行去第三行菽位，又以減第四行及右行菽位，不足減乃止。次以右行減第二行頭位，不足減乃止。次以第三行去左行頭位。次以左行去右行頭位，餘上得六，下得五，是爲荅六當𮮐五。次以右行去左行荅位，餘約之，上爲二，下爲三。次以左行去第二行下位，以第二行去第四行下位，又以減左行下位。次右行去第二行下位，餘上得三，下得四，是爲麥三當菽四。次以第二行減第四行下位。次以第四行去第二行下位，餘上得四，下得七，是爲麻四當麥七。是爲相當之率舉矣。\end{unitzh}
\begin{uniten}Next, with the left row eliminate the bean position of the 3rd row, and further subtract it from the bean positions of the 4th row and the right row, stopping when there is not enough to subtract. Next, with the right row subtract from the head position of the 2nd row, stopping when there is not enough to subtract. Next, with the 3rd row eliminate the head position of the left row. Next, with the left row eliminate the head position of the right row; of the remainder, the upper gives 6 and the lower gives 5: this is 6 small beans being equivalent to 5 millet. Next, with the right row eliminate the small-bean position of the left row; simplify the remainder: the upper makes 2 and the lower makes 3. Next, with the left row eliminate the lower position of the 2nd row; with the 2nd row eliminate the lower position of the 4th row, and further subtract it from the lower position of the left row. Next, the right row eliminates the lower position of the 2nd row; of the remainder, the upper gives 3 and the lower gives 4: this is 3 wheat being equivalent to 4 beans. Next, with the 2nd row subtract from the lower position of the 4th row. Next, with the 4th row eliminate the lower position of the 2nd row; of the remainder, the upper gives 4 and the lower gives 7: this is 4 hemp being equivalent to 7 wheat. With this the rates of mutual equivalence (xiangdang 相當) are all brought out.\end{uniten}

\begin{unitzh}\unitlabel{jiuzhang-037.10}{\zhsmall 以上所言新術亦訛舛不可通。據其術求之，先以左行減第三行，去其次位；次并右行、左行，亦并第三行、第四行，以減之，去其次位；次倍左行，以第二行減之，去其次位；所減之餘重列之爲三行。次以第四行減第二行，去次位及下位；次以重列之中行減右行，去其下位；次以重列之左行減右行，去其下位；所減之餘又重列之爲三行。次以此右行減中行，去其頭位；次以此右行減左行，去其頭位；所減之餘兩行兩物，減去下實，餘約之，上得五，下得三，是菽五當荅三。前云令左右相減，先去下實，又轉去物位，求其一行二物正負相借者，易其相當之率，謂菽五當荅三，卽菽價率三、荅價率五也。或先減下實乃減物位，或先減物位乃減下實，各從省便，本無一定之先後。其先求菽與荅相當之率，次求荅與𮮐相當之率，次求麥與菽相當之率，次求麻與麥相當之率，亦無一定之先後。然非遇正負相借者，則二物相當之率不可得，往往窮而復推，輾轉滋繁，遠不若舊術之究歸易簡也。}\end{unitzh}
\begin{uniten}[Note: The new method stated above is likewise corrupt and confused and cannot be made sense of. Seeking it according to its method: first subtract the left row from the 3rd row, eliminating its next position; next add the right row and the left row, and also add the 3rd row and the 4th row, and subtract [the one] from [the other], eliminating their next positions; next double the left row and subtract the 2nd row from it, eliminating its next position; set out the remainders of the subtractions anew as 3 rows. Next subtract the 4th row from the 2nd row, eliminating the next position and the lower position; next subtract the middle row of those set out anew from the right row, eliminating its lower position; next subtract the left row of those set out anew from the right row, eliminating its lower position; set out the remainders of the subtractions anew once more as 3 rows. Next subtract this right row from the middle row, eliminating its head position; next subtract this right row from the left row, eliminating its head position; the remainders of the subtractions are 2 rows with 2 things; subtract away the lower dividends and simplify the remainder: the upper gives 5 and the lower gives 3; this is 5 beans being equivalent to 3 small beans. What was said earlier — to make the left and right [rows] subtract from each other, first eliminating the lower dividends and then in turn eliminating the positions of the things (wuwei 物位), seeking the case where a row has 2 things whose positive and negative borrow from each other (xiangjie 相借), and to exchange their rates of equivalence — means: 5 beans are equivalent to 3 small beans, that is, the price rate of beans is 3 and the price rate of small beans is 5. One may first subtract the lower dividends and then subtract the positions of the things, or first subtract the positions of the things and then subtract the lower dividends; each follows what is economical and convenient, and there is fundamentally no fixed order. Its seeking first the rate of equivalence of beans and small beans, next the rate of equivalence of small beans and millet, next the rate of equivalence of wheat and beans, next the rate of equivalence of hemp and wheat, likewise has no fixed order. However, unless one meets [a row whose] positive and negative borrow from each other, the rate of equivalence of 2 things cannot be obtained; often one comes to a dead end and must push on again, and going round and round it grows ever more complicated; it is far from as good as the old method, which in the end comes down to ease and simplicity.]\end{uniten}

\begin{unitzh}\unitlabel{jiuzhang-037.11}九章算術卷八終\end{unitzh}
\begin{uniten}End of juan (卷) 8 of the Nine Chapters on the Mathematical Art.\end{uniten}

\section*{卷第九}\addcontentsline{toc}{section}{卷第九}

\begin{unitzh}\unitlabel{jiuzhang-038.1}《九章算術》卷九{\zhsmall 《算經十書》之二} 魏劉徽注 唐朝議大夫行太史令上輕車都尉臣李淳風等奉勅注釋\end{unitzh}
\begin{uniten}The Nine Chapters on the Mathematical Art, Juan 9 [Note: The second of the Ten Mathematical Classics] — Commentary by Liu Hui of the Wei — Of the Tang: Grand Master for Court Discussion, acting Grand Astrologer, Supreme Commandant of Light Chariots, Your servant Li Chunfeng and others, annotation and explanation by imperial command\end{uniten}

\begin{unitzh}\unitlabel{jiuzhang-038.2}句股{\zhsmall 以御高深廣遠}\end{unitzh}
\begin{uniten}Base (gou) and altitude (gu) [Liu Hui: By which to deal with height and depth, breadth and distance.]\end{uniten}

\begin{unitzh}\unitlabel{jiuzhang-038.3}今有句三尺，股四尺，問爲弦幾何？荅曰：五尺。\end{unitzh}
\begin{uniten}Suppose there is a base of 3 chi (尺) and an altitude of 4 chi. What is the hypotenuse (xian)? Answer: 5 chi.\end{uniten}

\begin{unitzh}\unitlabel{jiuzhang-038.4}今有弦五尺，句三尺，問爲股幾何？荅曰：四尺。\end{unitzh}
\begin{uniten}Suppose there is a hypotenuse of 5 chi and a base of 3 chi. What is the altitude? Answer: 4 chi.\end{uniten}

\begin{unitzh}\unitlabel{jiuzhang-038.5}今有股四尺，弦五尺，問爲句幾何？荅曰：三尺。\end{unitzh}
\begin{uniten}Suppose there is an altitude of 4 chi and a hypotenuse of 5 chi. What is the base? Answer: 3 chi.\end{uniten}

\begin{unitzh}\unitlabel{jiuzhang-038.6}句股{\zhsmall 短面日句，長面日股，相與結角日弦。句短其股，股短其弦。將以施於諸率，故先具此術以見其源也。}\end{unitzh}
\begin{uniten}Base and altitude [Liu Hui: The short side is called the base, the long side is called the altitude, and what joins with them at the corners is called the hypotenuse. [Translator's note: in all three places the text writes 日 "sun", evidently for 曰 "is called".] The base is shorter than its altitude, and the altitude is shorter than its hypotenuse. Since this is going to be applied to the various rates (lü), this method is therefore set out first, so as to show its source.]\end{uniten}

\begin{unitzh}\unitlabel{jiuzhang-038.7}術曰：句股各自乘，并而開方除之，卽弦。{\zhsmall 句自乘爲朱方，股自乘爲靑方，令出入相補，各從其類，因就其餘不移動也，合成弦方之幂，開方除之，卽弦也。}\end{unitzh}
\begin{uniten}Method: Multiply the base and the altitude each by itself, add them, and extract the square root: this is the hypotenuse. [Liu Hui: The base multiplied by itself makes the vermilion square (zhu fang); the altitude multiplied by itself makes the blue-green square (qing fang). Let what goes out and what comes in make up for each other (churu xiangbu), each following its own kind, relying on the rest, which is not moved; together they form the area (mi) of the hypotenuse square. Extract its square root: this is the hypotenuse.]\end{uniten}

\begin{unitzh}\unitlabel{jiuzhang-038.8}又股自乘，以減弦自乘，其餘開方除之，卽句。{\zhsmall 臣淳風等謹按：此術以句股幂合成弦幂，句方於内，則句短於股。令股自乘，以減弦自乘，餘者卽句幂也，故開方除之，卽句也。}\end{unitzh}
\begin{uniten}Again, multiply the altitude by itself and subtract it from the hypotenuse multiplied by itself; extract the square root of the remainder: this is the base. [Li Chunfeng et al.: We respectfully note: in this method the areas of the base and the altitude combine to form the area of the hypotenuse. The base square lies inside, so the base is shorter than the altitude. Let the altitude be multiplied by itself and subtracted from the hypotenuse multiplied by itself; what remains is the area of the base. Therefore extracting its square root gives the base.]\end{uniten}

\begin{unitzh}\unitlabel{jiuzhang-038.9}又句自乘，以減弦自乘，其餘開方除之，卽股。{\zhsmall 句股幂合以成弦幂，令去其一，則餘在者皆可得而知之。}\end{unitzh}
\begin{uniten}Again, multiply the base by itself and subtract it from the hypotenuse multiplied by itself; extract the square root of the remainder: this is the altitude. [Liu Hui: The areas of the base and the altitude combine to form the area of the hypotenuse. If one of them is removed, then what remains can in every case be obtained and known.]\end{uniten}

\begin{unitzh}\unitlabel{jiuzhang-038.10}今有圓材徑二尺五寸，欲爲方版，令厚七寸，問廣幾何？荅曰：二尺四寸五分。\end{unitzh}
\begin{uniten}Suppose there is a round timber with a diameter of 2 chi 5 cun (寸). One wishes to make a squared board from it, letting it be 7 cun thick. What is the width (guang)? Answer: 2 chi 4 cun 5 fen (分).\end{uniten}

\begin{unitzh}\unitlabel{jiuzhang-038.11}術曰：令徑二尺五寸自乘，以七寸自乘減之，其餘開方除之，卽廣。{\zhsmall 此以圓徑二尺五寸爲弦，版厚七寸爲句，所求廣爲股也。}\end{unitzh}
\begin{uniten}Method: Let the diameter, 2 chi 5 cun, be multiplied by itself, and subtract from it 7 cun multiplied by itself; extract the square root of the remainder: this is the width. [Liu Hui: Here the diameter of the circle, 2 chi 5 cun, is taken as the hypotenuse, the thickness of the board, 7 cun, as the base, and the width sought is the altitude.]\end{uniten}

\begin{unitzh}\unitlabel{jiuzhang-038.12}今有木長二丈，圍之三尺，葛生其下，纒木七周，上與木齊，問葛長幾何？荅曰：二丈九尺。\end{unitzh}
\begin{uniten}Suppose there is a tree 2 zhang (丈) long and 3 chi in girth. A kudzu vine (ge) grows at its bottom, winds around the tree 7 times, and at the top is level with the tree. What is the length of the vine? Answer: 2 zhang 9 chi.\end{uniten}

\begin{unitzh}\unitlabel{jiuzhang-038.13}術曰：以七周乘三圍爲股，木長爲句，爲之求弦。弦者，葛之長。{\zhsmall 據圍廣、木長求葛之長。其形葛卷裏袤。以筆管靑線宛轉，有似葛之纒木。解而觀之，則每周之間自有相間成句股弦。則其間木長爲股，圍之爲句，葛長爲弦。弦七周乘三圍，是并合衆句以爲一句，則句長而股短。故術以木長謂之句，圍之謂之股，言之倒互。句與股求弦，亦如前圖。句三自乘爲朱幂，股四自乘爲靑幂，合朱靑二十五，爲弦五自乘幂，出上第一圖。句股幂合爲弦幂，明矣。然二幂之數謂倒在於弦幂之中而已。可更相裏者則成方幂，其居表者則成矩幂。二表裏形訛而數均。又按此圖，句幂之矩靑卷白表，是其幂以股弦差爲廣，股弦并爲袤，而股幂方其裏；股幂之矩靑卷白表，是其幂以句弦差爲廣，句弦并爲袤，而句幂方其裏。是故差之與并用除之，短長互相乘也。}\end{unitzh}
\begin{uniten}Method: Multiply the 3 of girth by the 7 turns to make the altitude, take the length of the tree as the base, and for these find the hypotenuse. The hypotenuse is the length of the vine. [Liu Hui: On the basis of the width of the girth and the length of the tree, the length of the vine is found. Its shape: the vine coils round within the length (mao). If a blue-green thread is wound around a writing-brush tube, it resembles the vine winding round the tree. Unwind it and look at it: then between each turn there are naturally intervals forming base, altitude and hypotenuse. So in each interval the length of the tree is the altitude, its girth is the base, and the length of the vine is the hypotenuse. As for the hypotenuse, multiplying the 3 of girth by the 7 turns is to combine all the bases to make one base; then the base is long and the altitude short. Therefore the method calls the length of the tree the base and the girth the altitude, speaking of them the other way round. Finding the hypotenuse from the base and the altitude is likewise as in the preceding figure. The base 3 multiplied by itself makes the vermilion area; the altitude 4 multiplied by itself makes the blue-green area; vermilion and blue-green combined, 25, make the area of the hypotenuse 5 multiplied by itself; this comes from the first figure above. That the areas of the base and the altitude combine to make the area of the hypotenuse is clear. However, the numbers of the two areas are only said to be turned over within the area of the hypotenuse. Those that can alternately lie inside form a square area; those lying on the outside form a carpenter's-square area (ju mi). The two, outer and inner, are altered in shape but equal in number. Further, according to this figure: the carpenter's square of the base area, with blue-green rolled up and white outside, is an area that takes the difference of the altitude and the hypotenuse as width and the sum of the altitude and the hypotenuse as length, while the altitude area is a square on its inside; the carpenter's square of the altitude area, with blue-green rolled up and white outside, is an area that takes the difference of the base and the hypotenuse as width and the sum of the base and the hypotenuse as length, while the base area is a square on its inside. Therefore the difference and the sum are used to divide it: the short and the long multiply each other.]\end{uniten}

\begin{unitzh}\unitlabel{jiuzhang-038.14}今有池方一丈，葭生其中央，出水一尺。引葭赴岸，適與岸齊。問水深、葭長各幾何？荅曰：水深一丈二尺，葭長一丈三尺。\end{unitzh}
\begin{uniten}Suppose there is a pond 1 zhang square, with a reed (jia) growing at its centre and rising 1 chi out of the water. When the reed is pulled toward the bank, it comes exactly level with the bank. What are the depth of the water and the length of the reed, each? Answer: The depth of the water is 1 zhang 2 chi; the length of the reed is 1 zhang 3 chi.\end{uniten}

\begin{unitzh}\unitlabel{jiuzhang-038.15}術曰：半池方自乘，{\zhsmall 此以池方半之，得五尺爲句，水深爲股，葭長爲弦。以句及股弦差求股弦，故令句自乘，先見矩幂也。}以出水一尺自乘減之，{\zhsmall 出水者，股弦差。減此差幂於矩幂，餘爲倍股弦差乘股長之矩幂。}餘，倍出水除之，卽得水深。{\zhsmall 倍差爲矩幂之廣，水深是股。欲先見葭長者，出水一尺自乘，以加於半池方自乘，倍出水除之，卽得。令此幂得出水一尺爲袤，故爲矩而得葭長也。}加出水數，得葭長。{\zhsmall 臣淳風等謹按：此葭本出水一尺，旣見水深，故加出水尺數而得葭長也。}\end{unitzh}
\begin{uniten}Method: Halve the side of the pond and multiply it by itself; [Liu Hui: Here the side of the pond is halved, giving 5 chi, as the base; the depth of the water is the altitude; the length of the reed is the hypotenuse. From the base and the difference of the altitude and the hypotenuse one finds the altitude and the hypotenuse; therefore the base is multiplied by itself, first to exhibit the carpenter's-square area.] subtract from it the 1 chi out of the water multiplied by itself; [Liu Hui: The part out of the water is the difference of the altitude and the hypotenuse. When the area of this difference is subtracted from the carpenter's-square area, the remainder is the carpenter's-square area of twice the difference of the altitude and the hypotenuse multiplying the length of the altitude.] divide the remainder by twice the part out of the water: this gives the depth of the water. [Liu Hui: Twice the difference is the width of the carpenter's-square area; the depth of the water is the altitude. If one wishes to find the length of the reed first: multiply the 1 chi out of the water by itself, add it to the half side of the pond multiplied by itself, and divide by twice the part out of the water; this gives it. Let this area obtain the 1 chi out of the water as its length; therefore it forms a carpenter's square and the length of the reed is obtained.] Add the amount out of the water: this gives the length of the reed. [Li Chunfeng et al.: We respectfully note: this reed originally rises 1 chi out of the water; once the depth of the water has been found, the number of chi out of the water is therefore added, and the length of the reed is obtained.]\end{uniten}

\begin{unitzh}\unitlabel{jiuzhang-038.16}今有立木，繫索其末，委地三尺。引索郤行，去本八尺而索盡。問索長幾何？荅曰：一丈二尺二十一分尺之一。\end{unitzh}
\begin{uniten}Suppose there is an upright post with a rope tied to its top, 3 chi of which lies on the ground. Pulling the rope, one walks back, and at 8 chi from the foot (ben) the rope is used up. What is the length of the rope? Answer: 1 zhang 2 1/21 chi.\end{uniten}

\begin{unitzh}\unitlabel{jiuzhang-038.17}術曰：以去本自乘，{\zhsmall 此以去本八尺爲句，所求索者弦也。引而索盡，與開門去閫者，句及股弦差求股弦，同一術。去本自乘者，先張矩幂。}令如委數而一，{\zhsmall 委地者，股弦差也。以除矩幂，卽是股弦并也。}所得加委地數而半之，卽索長。{\zhsmall 子不可半者，倍其母。加差於并，則成兩索長，故又半之。其減差於并而半之，得木長也。}\end{unitzh}
\begin{uniten}Method: Multiply the distance from the foot by itself; [Liu Hui: Here the distance from the foot, 8 chi, is taken as the base; the rope sought is the hypotenuse. Pulling until the rope is used up is, together with the door opened away from the threshold, one and the same method: from the base and the difference of the altitude and the hypotenuse, find the altitude and the hypotenuse. Multiplying the distance from the foot by itself first spreads out the carpenter's-square area.] divide it by the amount lying down; [Liu Hui: What lies on the ground is the difference of the altitude and the hypotenuse. Dividing the carpenter's-square area by it gives the sum of the altitude and the hypotenuse.] add the amount lying on the ground to the result and halve it: this is the length of the rope. [Liu Hui: If the numerator (zi) cannot be halved, double its denominator (mu). Adding the difference to the sum makes two lengths of the rope; therefore it is halved again. Subtracting the difference from the sum and halving it gives the length of the post.]\end{uniten}

\begin{unitzh}\unitlabel{jiuzhang-038.18}今有垣高一丈，倚木於垣，高與垣齊。引木郤行一尺，其木至地。問木幾何？荅曰：五丈五寸。\end{unitzh}
\begin{uniten}Suppose there is a wall 1 zhang high, with a pole leaning against the wall, its height level with the wall. When the pole is pulled back 1 chi, the pole comes down to the ground. How long is the pole? Answer: 5 zhang 5 cun.\end{uniten}

\begin{unitzh}\unitlabel{jiuzhang-038.19}術曰：以垣高十尺自乘，如郤行尺數而一，所得以加郤行尺數而半之，卽木長數。{\zhsmall 此以垣高一丈爲句，所求倚木者爲弦，引郤行一尺爲股弦差。其爲術之意，與繫索問同也。}\end{unitzh}
\begin{uniten}Method: Multiply the height of the wall, 10 chi, by itself and divide by the number of chi moved back; add the number of chi moved back to the result and halve it: this is the length of the pole. [Liu Hui: Here the height of the wall, 1 zhang, is taken as the base; the leaning pole sought is the hypotenuse; the 1 chi it is pulled back is the difference of the altitude and the hypotenuse. The idea of this method is the same as in the problem of the tied rope.]\end{uniten}

\begin{unitzh}\unitlabel{jiuzhang-038.20}今有圓材埋在壁中，不知大小。以鐻鐻之，深一寸，鐻道長一尺。問徑幾何？荅曰：材徑二尺六寸。\end{unitzh}
\begin{uniten}Suppose there is a round timber buried in a wall, its size unknown. One saws it with a saw (ju): the depth is 1 cun and the saw path (ju dao) is 1 chi long. What is the diameter? Answer: The diameter of the timber is 2 chi 6 cun.\end{uniten}

\begin{unitzh}\unitlabel{jiuzhang-038.21}術曰：半鐻道自乘，{\zhsmall 此術以鐻道一尺爲句，材徑爲弦，鐻深一寸爲股弦差之一半，故鐻長亦半之也。}{\zhsmall 臣淳風等謹按：下鐻深得一寸爲半股弦差，注云爲股弦差者，鐻道也。}如深寸而一，以深寸增之，卽材徑。{\zhsmall 亦以半增之。如上術，去本當半之，今此皆同半差，不復半也。}\end{unitzh}
\begin{uniten}Method: Halve the saw path and multiply it by itself; [Liu Hui: In this method the saw path, 1 chi, is taken as the base, the diameter of the timber as the hypotenuse, and the saw depth, 1 cun, as one half of the difference of the altitude and the hypotenuse; therefore the saw length is also halved.] [Li Chunfeng et al.: We respectfully note: the depth sawn downward, obtained as 1 cun, is half the difference of the altitude and the hypotenuse; where the commentary says "is the difference of the altitude and the hypotenuse", it is the saw path.] divide by the cun of depth and increase it by the cun of depth: this is the diameter of the timber. [Liu Hui: It too is increased by the half. As in the method above with the distance from the foot, one ought to halve it; here, now, all alike are the same half-difference, so it is not halved again.]\end{uniten}

\begin{unitzh}\unitlabel{jiuzhang-038.22}今有開門去閫一尺，不合二寸，問門廣幾何？荅曰：一丈一寸。\end{unitzh}
\begin{uniten}Suppose there is a door opened 1 chi away from the threshold (kun), with 2 cun not closed. What is the width of the door? Answer: 1 zhang 1 cun.\end{uniten}

\begin{unitzh}\unitlabel{jiuzhang-038.23}術曰：以去閫一尺自乘，所得以不合二寸半之而一，所得增不合之半，卽得門廣。{\zhsmall 此去閫一尺爲句，門廣爲股，不合二寸以半之，得一寸爲股弦差，求弦，故當半之。今卽以兩弦爲廣數，故不復半之也。}\end{unitzh}
\begin{uniten}Method: Multiply the 1 chi away from the threshold by itself; divide the result by the 2 cun not closed, halved; increase the result by half of the amount not closed: this gives the width of the door. [Liu Hui: Here the 1 chi away from the threshold is the base and the width of the door is the altitude; the 2 cun not closed is halved, giving 1 cun as the difference of the altitude and the hypotenuse, to find the hypotenuse; therefore it ought to be halved. Now the two hypotenuses are directly taken as the number for the width, therefore it is not halved again.]\end{uniten}

\begin{unitzh}\unitlabel{jiuzhang-039.1}今有戸，高多於廣六尺八寸，兩隅相去適一丈。問戸高、廣各幾何？荅曰：廣二尺八寸，高九尺六寸。\end{unitzh}
\begin{uniten}Suppose there is a door whose height exceeds its width by 6 chi (尺) 8 cun (寸), and whose two corners are exactly 1 zhang (丈) apart. Question: what are the height and the width of the door each? Answer: width 2 chi 8 cun; height 9 chi 6 cun.\end{uniten}

\begin{unitzh}\unitlabel{jiuzhang-039.2}術曰：令一丈自乘爲實。半相多，令自乘，倍之，減實，半其餘，以開方除之。所得減相多之半，卽戸廣；加相多之半，卽戸高。{\zhsmall 令戸廣爲句，高爲股，兩隅相去一丈爲弦，高多於廣六尺八寸爲句股差。按圖爲位，弦幂適滿萬寸。倍之，減句股差幂，開方除之，其所得卽高廣并數。以差減并而半之，卽戸廣；加相多之數，卽戸高也。今此術先求其半。一丈自乘爲朱幂四、黃幂一。半差自乘，又倍之，爲黃幂四分之二。減實，半其餘，有朱幂二、黃幂四分之一。其於大方棄四分之三，適得四分之一。故開方除之，得高廣并數之半。減差半，得廣；加，得戸高。}\end{unitzh}
\begin{uniten}Method: Let 1 zhang be multiplied by itself as the dividend (shi). Halve the excess, let it be multiplied by itself, double it, subtract it from the dividend, halve the remainder, and extract the square root. Subtract half the excess from what is obtained: this is the width of the door; add half the excess: this is the height of the door. [Liu Hui: Let the width of the door be the base (gou), the height the altitude (gu), the distance between the two corners, 1 zhang, the hypotenuse (xian), and the excess of the height over the width, 6 chi 8 cun, the difference of base and altitude. Setting out the positions according to the figure, the area (mi) of the hypotenuse is exactly a full 10000 cun. Double it, subtract the area of the difference of base and altitude, and extract the square root: what is obtained is the sum of height and width. Subtract the difference from the sum and halve it: this is the width of the door; add the amount of the excess: this is the height of the door. Now this method first seeks the half. 1 zhang multiplied by itself makes 4 red areas and 1 yellow area. Half the difference multiplied by itself and then doubled makes 2/4 of the yellow area. Subtract it from the dividend and halve the remainder: there are 2 red areas and 1/4 of the yellow area. Relative to the large square, 3/4 is discarded, and exactly 1/4 is obtained. Therefore extracting the square root gives half the sum of height and width. Subtract half the difference: this gives the width; add it: this gives the height of the door.]\end{uniten}

\begin{unitzh}\unitlabel{jiuzhang-039.3}{\zhsmall 又按此圖幂：句股并自乘，加差幂，爲兩弦幂，半之，開方得弦。今倍弦幂，減差幂，求句股并。葢先見其弦，然後知其句與股也。句股適等者，并而自乘，卽爲兩弦幂，皆各爲方。先見其弦，然後知其句與股者，倍弦幂，卽爲句股適等者并而自乘之幂。半相多，自乘，倍之，又半股并，自乘，亦倍之，合爲弦幂。其無差數者，句股各自乘，并之爲實，與句股相乘倍之爲實，皆開方得弦。弦幂半之爲實，開方卽得句股。及股長句短，同源而分流焉。假令句股各五，弦幂五十，開方除之，得七尺，有餘一，不盡。假令弦十，其幂有百，半之爲句股二幂，各得五十，當亦不可開。故曰「圓三徑一，方五斜七」，雖不正得盡理，亦可言相近耳。}\end{unitzh}
\begin{uniten}[Liu Hui: Further, examining this figure of areas: the sum of base and altitude multiplied by itself, with the area of the difference added, makes two areas of the hypotenuse; halve it and extract the square root to obtain the hypotenuse. Now double the area of the hypotenuse and subtract the area of the difference to seek the sum of base and altitude. For one first sees the hypotenuse, and only afterwards knows its base and altitude. When base and altitude are exactly equal, their sum multiplied by itself is then two areas of the hypotenuse, all of them being squares. Where one first sees the hypotenuse and afterwards knows its base and altitude, double the area of the hypotenuse: this is the area of the sum of an exactly equal base and altitude multiplied by itself. Halve the excess, multiply it by itself and double it; further halve the altitude sum, multiply it by itself and also double it: together they make the area of the hypotenuse. [Translator's note: "altitude sum" (股并) evidently stands for the sum of base and altitude.] Where there is no difference, base and altitude each multiplied by itself and added together make the dividend, and base and altitude multiplied together and doubled make the dividend; in both cases extracting the square root gives the hypotenuse. Halve the area of the hypotenuse as the dividend; extracting the square root then gives base and altitude. And [the case where] the altitude is long and the base short has the same source, but the streams divide. Suppose base and altitude are each 5: the area of the hypotenuse is 50; extracting its square root gives 7 chi, with 1 left over, not exhausted. Suppose the hypotenuse is 10: its area is 100; halving it gives the two areas of base and altitude, each 50, which likewise cannot be extracted. Hence it is said, "circle 3, diameter 1; square 5, diagonal 7": although these do not exactly attain the full principle, they may still be said to be close.]\end{uniten}

\begin{unitzh}\unitlabel{jiuzhang-039.4}{\zhsmall 其句股合而自乘之幂，令弦自乘倍之爲兩弦幂，以減之，其餘開方除之，爲句股差。加差於合而半之，爲股；減差於合而半之，爲。句、股、弦卽高、廣、袤。其出此圖也，其倍弦爲廣袤合，矩句卽爲幂，得廣卽句股差。其矩句之幂，倍爲從法，開之，亦句股差。其餘以句股幂減，半其餘，差爲從法，開方除之，卽句也。}\end{unitzh}
\begin{uniten}[Liu Hui: As for the area of base and altitude combined and multiplied by itself: let the hypotenuse be multiplied by itself and doubled to make two areas of the hypotenuse, and subtract it from this; extract the square root of the remainder: this is the difference of base and altitude. Add the difference to the combined [amount] and halve it: this is the altitude; subtract the difference from the combined [amount] and halve it: this is the base. Base, altitude and hypotenuse are the height, width and length (mao). As for what comes out of this figure: the doubled hypotenuse makes the combined width and length; the gnomon-base (ju gou) then makes the area; the width obtained is the difference of base and altitude. The area of the gnomon-base: double to make the joined divisor (congfa), and extract it: this is again the difference of base and altitude. From the remainder subtract the area of base and altitude, halve the remainder, take the difference as the joined divisor, and extract the square root: this is the base. [Translator's note: the text from "As for what comes out of this figure" onward is evidently garbled.]]\end{uniten}

\begin{unitzh}\unitlabel{jiuzhang-039.5}今有戸，不知高、廣；竿，不知長短。橫之不出四尺，從之不出二尺，邪之適出。問戸高、廣、袤各幾何？荅曰：廣六尺，高八尺，衺一丈。\end{unitzh}
\begin{uniten}Suppose there is a door whose height and width are not known, and a pole whose length is not known. Held crosswise, it fails to pass through by 4 chi; held upright, it fails to pass through by 2 chi; held slantwise, it just passes through. Question: what are the height, width and length (mao) of the door each? Answer: width 6 chi; height 8 chi; diagonal (xie) 1 zhang.\end{uniten}

\begin{unitzh}\unitlabel{jiuzhang-039.6}術曰：從、橫不出相乘，倍，而開方除之。所得加從不出，卽戸廣；{\zhsmall 此以戸廣爲句，戸高爲股，戸袤爲弦。凡并句股之幂，卽爲弦幂。或矩於表，或方於裏。連之者，舉表矩而方之。又從句方裏令爲靑矩之表，未滿黃方。满此方，則兩端之亷重於隅中，各以股弦差爲廣，句弦差爲袤。故兩端差相乘，又倍之，則成黄方之幂。開方除之，得黃方之面。其外之靑，知亦以股弦差爲廣，故以股弦差加之，則爲句也。}加橫不出，卽戸高；兩不出加之，得戸袤。\end{unitzh}
\begin{uniten}Method: Multiply together the amounts by which it fails to pass upright and crosswise, double, and extract the square root. Add to what is obtained the amount by which it fails to pass upright: this is the width of the door; [Liu Hui: Here the width of the door is taken as the base, the height of the door as the altitude, and the length of the door as the hypotenuse. In general, the areas of base and altitude combined make the area of the hypotenuse. One is a gnomon (ju) on the outside, the other a square on the inside. In joining them, one takes up the outer gnomon and makes it square. Further, from inside the square on the base, let it be made the outside of the blue gnomon; it does not fill the yellow square. To fill this square, the side-strips (lian) at the two ends overlap at the corner, each having the difference of altitude and hypotenuse as width and the difference of base and hypotenuse as length. Therefore multiply together the differences at the two ends and double it: this forms the area of the yellow square. Extract its square root to obtain the side of the yellow square. As for the blue outside it, one knows that it also has the difference of altitude and hypotenuse as width; therefore adding the difference of altitude and hypotenuse to it gives the base.] add the amount by which it fails to pass crosswise: this is the height of the door; add to it both amounts by which it fails to pass: this gives the length of the door.\end{uniten}

\begin{unitzh}\unitlabel{jiuzhang-039.7}今有竹高一丈，末折抵地，去本三尺。問折者高幾何？荅曰：四尺二十分尺之十一。\end{unitzh}
\begin{uniten}Suppose there is a bamboo 1 zhang high; its top is broken and touches the ground 3 chi from the foot. Question: how high is the break? Answer: 4 11/20 chi.\end{uniten}

\begin{unitzh}\unitlabel{jiuzhang-039.8}術曰：以去本自乘，{\zhsmall 此去本三尺爲句，折之餘高爲股，末折抵地爲弦。以句及股弦并求股，故先令句自乘，見矩幂。}令如高而一。{\zhsmall 竹高一丈爲股弦并，以除此幂，得差。}所得以減竹高，而半其餘，卽折者之高也。{\zhsmall 此術與繫索者之類更相反覆也。亦可如上術，令高自乘爲股弦并幂，去本自乘爲矩幂，減之，餘爲實；倍高爲法，則得折之高數也。}\end{unitzh}
\begin{uniten}Method: Multiply the distance from the foot by itself, [Liu Hui: Here the distance from the foot, 3 chi, is the base, the remaining height at the break is the altitude, and the broken top touching the ground is the hypotenuse. One seeks the altitude from the base and the sum of altitude and hypotenuse; therefore one first lets the base be multiplied by itself, to show the area of the gnomon.] and let it be divided by the height. [Liu Hui: The height of the bamboo, 1 zhang, is the sum of altitude and hypotenuse; dividing this area by it gives the difference.] Subtract what is obtained from the height of the bamboo and halve the remainder: this is the height of the break. [Liu Hui: This method and those of the type with the tied rope are reverses of each other. One may also, as in the method above, let the height multiplied by itself be the area of the sum of altitude and hypotenuse, and the distance from the foot multiplied by itself the area of the gnomon; subtract it, and the remainder is the dividend; double the height as the divisor (fa): then one obtains the height of the break.]\end{uniten}

\begin{unitzh}\unitlabel{jiuzhang-039.9}今有二人同所立，甲行率七，乙行率三。乙東行，甲南行十步而邪東北與乙會。問甲、乙行各幾行？荅曰：乙東行十步半，甲邪行十四步半及之。\end{unitzh}
\begin{uniten}Suppose there are two people standing at the same place. A's walking rate (lü) is 7; B's walking rate is 3. B walks east; A walks south 10 bu (步) and then slantwise to the northeast, meeting B. Question: how much do A and B each walk? Answer: B walks east 10 1/2 bu; A walks slantwise 14 1/2 bu and reaches him.\end{uniten}

\begin{unitzh}\unitlabel{jiuzhang-039.10}術曰：令七自乘，三亦自乘，并而半之，以爲甲邪行率。邪行率減於七自乘，餘爲南行率。以三乘七，爲乙東行率。{\zhsmall 此以南行爲句，東行爲股，邪行爲弦。句弦并七。欲知弦者，當以股自乘爲幂，如并而一，所得爲句弦差。加差於并而半之，爲弦。以弦減差，餘爲句。如是或有分，當通而約之乃定。術以句弦并爲分母，差爲分子，故令句弦并自乘爲朱黄相連之方，股自乘爲靑幂之矩。令其矩引之直，加損同之。以句弦并爲袤，差爲廣。其圖大體以兩弦爲袤，句弦并爲廣。引黃斷其半爲弦率。七自乘者，句弦并之率，故弦減之，餘爲句率。同立處是中停也。列用率，皆句弦并爲袤，弦與句各爲之廣，故亦以股率同其袤也。}置南行十步，以甲邪行率乘之；副置十步，以乙東行率乘之；各自爲實。實如南行率而一，各得行數。{\zhsmall 南行十步者，所有見句，求見弦、股，故以弦、股率如句率而一。}\end{unitzh}
\begin{uniten}Method: Let 7 be multiplied by itself, and 3 also be multiplied by itself; add them and halve, taking this as A's slantwise rate. Subtract the slantwise rate from 7 multiplied by itself: the remainder is the southward rate. Multiply 7 by 3 as B's eastward rate. [Liu Hui: Here the southward walk is the base, the eastward walk the altitude, and the slantwise walk the hypotenuse. The sum of base and hypotenuse is 7. To know the hypotenuse, one should multiply the altitude by itself to make an area and divide it by the sum; what is obtained is the difference of base and hypotenuse. Add the difference to the sum and halve it: this is the hypotenuse. Subtract the hypotenuse from the difference: the remainder is the base. Done in this way there may be fractions; one should bring them to common terms (tong) and reduce them (yue), and then it is settled. The method takes the sum of base and hypotenuse as denominator and the difference as numerator; therefore let the sum of base and hypotenuse multiplied by itself make the square in which red and yellow are joined, and the altitude multiplied by itself make the gnomon of the blue area. Let the gnomon be drawn out straight, equalizing it by adding and taking away. Take the sum of base and hypotenuse as length and the difference as width. The figure, in its overall form, takes two hypotenuses as length and the sum of base and hypotenuse as width. Draw out the yellow and cut off its half as the rate of the hypotenuse. 7 multiplied by itself is the rate of the sum of base and hypotenuse; therefore, subtracting the hypotenuse from it, the remainder is the rate of the base. The place where they stand together is the midpoint. Of the rates set out for use, all have the sum of base and hypotenuse as length, and the hypotenuse and the base each as width; therefore the rate of the altitude too is made the same in its length.] Set out the southward walk of 10 bu and multiply it by A's slantwise rate; set out 10 bu again in an auxiliary position and multiply it by B's eastward rate; each by itself is a dividend. Divide the dividends by the southward rate: each gives the distance walked. [Liu Hui: The southward walk of 10 bu is the actual base that is given; one seeks the actual hypotenuse and altitude; therefore one takes the rates of hypotenuse and altitude and divides by the rate of the base.]\end{uniten}

\begin{unitzh}\unitlabel{jiuzhang-039.11}今有句五步，股十二步。問句中容方幾何？荅曰：方三步十七分步之九。\end{unitzh}
\begin{uniten}Suppose there is a base of 5 bu and an altitude of 12 bu. Question: how large is the square contained within the base (gou)? Answer: the square is 3 9/17 bu.\end{uniten}

\begin{unitzh}\unitlabel{jiuzhang-039.12}術曰：并句股爲法，句股相乘爲實，實如法而一，得方一步。{\zhsmall 句股相乘爲朱、靑、黃幂各二。令黃幂連於下隅，朱、靑各以類合，共成修幂。中方黃爲廣，并句股爲袤，故并句股爲法。幂圖方在句中，則方之兩亷各自成小句股，而其相與之勢不失本率也。句面之小股，股面之小句，從横相連合而成中方。令股爲中方率，并句股爲廣率，據見句五步而今有之，得中方也。復令句爲中方率，以并句股爲袤率，據股十二步而今有之，則中方又可知。此則雖不效，而法、實有法由生矣。不容圓率以今有、衰分言之，可以見之也。}\end{unitzh}
\begin{uniten}Method: Add base and altitude as the divisor; multiply base and altitude together as the dividend; divide the dividend by the divisor, obtaining 1 bu of the square. [Liu Hui: Base and altitude multiplied together make red, blue and yellow areas, 2 of each. Let the yellow areas be joined at the lower corner, and the red and the blue each combine with their own kind, together forming an oblong area (xiu mi). The yellow of the central square is the width and the sum of base and altitude the length; therefore the sum of base and altitude is the divisor. In the figure of areas the square is within the base; then the two side-strips of the square each form a small base and altitude by themselves, and their mutual configuration does not lose the original rates. The small altitude on the face of the base and the small base on the face of the altitude, joined upright and crosswise, combine to form the central square. Let the altitude be the rate of the central square and the sum of base and altitude the rate of the width; on the basis of the actual base of 5 bu, apply jinyou ("suppose there is") to it, and one obtains the central square. Again let the base be the rate of the central square, with the sum of base and altitude as the rate of the length; on the basis of the altitude of 12 bu, apply jinyou to it, and the central square can again be known. Thus, although this is not verified, the divisor and the dividend have the source from which the method arises. The rates of not containing a circle, expressed by jinyou and proportional distribution (cuifen), can be seen. [Translator's note: "not containing a circle" (不容圓) is evidently corrupt.]]\end{uniten}

\begin{unitzh}\unitlabel{jiuzhang-040.1}今有句八步，股十五步。問句中容圓，徑幾何？荅曰：六步。\end{unitzh}
\begin{uniten}Suppose there is a base (gou) of 8 bu (步) and an altitude (gu) of 15 bu. Question: if a circle is contained within the base [and altitude], what is its diameter? Answer: 6 bu.\end{uniten}

\begin{unitzh}\unitlabel{jiuzhang-040.2}術曰：八步爲句，十五步爲股，爲之求弦。三位并之爲法，以句乘股，倍之爲實。實如法得徑一步。\end{unitzh}
\begin{uniten}Method: Take 8 bu as the base and 15 bu as the altitude, and find the hypotenuse (xian) for them. Add the three positions together as the divisor (fa); multiply the altitude by the base and double it as the dividend (shi). The dividend divided by the divisor gives the diameter [in units of] 1 bu.\end{uniten}

\begin{unitzh}\unitlabel{jiuzhang-040.3}{\zhsmall 句股相乘爲圖之本體，朱、靑、黃幂各二，則倍之爲各四。可用畫於小𥿄，分裁邪正之會，令顛倒相補，各以𩔖合，成脩幂：圓徑爲廣，并句、股、弦爲袤。故并句、股、弦以爲法。又以圖之大體言之：股中靑必令立規於橫廣，句、股又邪三徑均，而復連規，從橫量度句股，必合而成小方矣。又畫中弦以觀其會，則句股之中成小句股弦者四，句面之小股、股面之小句皆小方之面，皆圓徑之半。其數故可衰。以句、股、弦爲列衰，副并爲法。以小句乘未并者，各自爲實。實如法而一，得句面之小股可知也。以股乘列衰爲實，則得股面之小句可知。言雖異矣，及其所以成法之實，則同歸矣。則圓徑又可以句乘之差并句弦差減股爲圓徑。又弦減句股并，餘爲圓徑。以句弦差乘股弦差而倍之，開方除之，亦圓徑也。}\end{unitzh}
\begin{uniten}[Liu Hui: The base and the altitude multiplied together make the basic body of the figure, [in which] the red, blue and yellow areas (mi) are two each; doubling it then makes four of each. One can draw it on a small [piece of] paper, cut it apart at the junctions of the slanting and the straight [lines], and let [the pieces] be turned upside down to complement one another, each joined with its own kind, forming an oblong area (xiumi): the diameter of the circle is the width, and the sum of base, altitude and hypotenuse is the length. Therefore the sum of base, altitude and hypotenuse is taken as the divisor. Further, speaking of it in terms of the general form of the figure: for the blue in the altitude, one must set up the compass (gui) on the horizontal width; [toward] the base, the altitude and also the slant, the three diameters are equal; and when the compass [curve] is joined up again and the base and altitude are measured vertically and horizontally, they must combine to form a small square. Further, draw the middle hypotenuse to observe their junctions; then within the base and altitude there are formed four small [figures of] base, altitude and hypotenuse; the small altitude on the side of the base and the small base on the side of the altitude are both sides of the small square, and both are half the diameter of the circle. Hence their numbers can be [obtained by] proportional distribution (cui). Take base, altitude and hypotenuse as the arranged proportions (liecui), and add them in an auxiliary position as the divisor. Multiply the ones not added by the small base, each separately as a dividend. Dividing the dividend by the divisor, the small altitude on the side of the base is obtained, as can be known. Take the arranged proportions multiplied by the altitude as the dividend; then the small base on the side of the altitude is obtained, as can be known. Although the statements differ, when it comes to the substance by which the method is formed, they return to the same. Then the diameter of the circle can moreover be [had as follows]: the difference from multiplying it by the base, added, [and] the difference of base and hypotenuse subtracted from the altitude, as the diameter of the circle. [Translator's note: this sentence appears garbled.] Also, subtract the hypotenuse from the sum of base and altitude; the remainder is the diameter of the circle. Multiply the difference of altitude and hypotenuse by the difference of base and hypotenuse and double it; extract the square root; this too is the diameter of the circle.]\end{uniten}

\begin{unitzh}\unitlabel{jiuzhang-040.4}今有邑方二百步，各中開門。出東門十五步有木，問出南門幾何步而見木？荅曰：六百六十六步太半步。\end{unitzh}
\begin{uniten}Suppose there is a town (yi) 200 bu square, with a gate opened in the middle of each [side]. 15 bu out of the east gate there is a tree. Question: going out of the south gate, after how many bu does one see the tree? Answer: 666 2/3 bu.\end{uniten}

\begin{unitzh}\unitlabel{jiuzhang-040.5}術曰：出東門步數爲法，{\zhsmall 以句率爲法也。}半邑方自乘爲實，實如法得一步。{\zhsmall 此以出東門十五步爲句率，東門南至隅一百步爲股率，南門東至隅一百步爲見句步。欲以見句求股，以爲出南門數。正合半邑方自乘者，股率當乘見句，此二者數同也。}\end{unitzh}
\begin{uniten}Method: Take the number of bu out of the east gate as the divisor, [Liu Hui: This takes the rate (lü) of the base as the divisor.] and half the side of the town multiplied by itself as the dividend; the dividend divided by the divisor gives [the result in units of] 1 bu. [Liu Hui: Here the 15 bu out of the east gate is taken as the rate of the base, the 100 bu from the east gate south to the corner as the rate of the altitude, and the 100 bu from the south gate east to the corner as the bu of the given base (jian gou). One wishes to find the altitude from the given base, to make it the number [of bu] out of the south gate. As for its agreeing exactly with half the side of the town multiplied by itself: the rate of the altitude should multiply the given base, and these two are the same in number.]\end{uniten}

\begin{unitzh}\unitlabel{jiuzhang-040.6}今有邑東西七里，南北九里，各中開門。出東門十五里有木，問出南門幾何步而見木？荅曰：三百一十五步。\end{unitzh}
\begin{uniten}Suppose there is a town 7 li (里) from east to west and 9 li from south to north, with a gate opened in the middle of each [side]. 15 li out of the east gate there is a tree. Question: going out of the south gate, after how many bu does one see the tree? Answer: 315 bu.\end{uniten}

\begin{unitzh}\unitlabel{jiuzhang-040.7}術曰：東門南至隅步數，以乘南門東至隅步數爲實，以木去門步數爲法，實如法而一。{\zhsmall 此以東門南至隅四里半爲句率，出東門十五里爲股率，南門東至隅三里半爲見股。所問出南門，卽見股之句。爲術之意與上同也。}\end{unitzh}
\begin{uniten}Method: Take the number of bu from the east gate south to the corner, and with it multiply the number of bu from the south gate east to the corner, as the dividend; take the number of bu of the tree's distance from the gate as the divisor; divide the dividend by the divisor. [Liu Hui: Here the 4 1/2 li from the east gate south to the corner is taken as the rate of the base, the 15 li out of the east gate as the rate of the altitude, and the 3 1/2 li from the south gate east to the corner as the given altitude (jian gu). What is asked, [the distance] out of the south gate, is the base of the given altitude. The idea of making the method is the same as above.]\end{uniten}

\begin{unitzh}\unitlabel{jiuzhang-040.8}今有邑方不知大小，各中開門。出北門三十步有木，出西門七百五十步見木。問邑方幾何？荅曰：一里。\end{unitzh}
\begin{uniten}Suppose there is a square town whose size is not known, with a gate opened in the middle of each [side]. 30 bu out of the north gate there is a tree; going 750 bu out of the west gate one sees the tree. Question: what is the side of the town? Answer: 1 li.\end{uniten}

\begin{unitzh}\unitlabel{jiuzhang-040.9}術曰：令兩出門步數相乘，因而四之，爲實。開方除之，卽得邑方。{\zhsmall 按前術：半邑方自乘，出東門步數除之，卽出南門步數。今兩出門相乘，爲半方邑自乘居一隅之積分。因而四之，卽得四隅之積分。故以爲實，開方除之，卽邑方也。}\end{unitzh}
\begin{uniten}Method: Let the numbers of bu out of the two gates be multiplied together, and then multiply by 4, as the dividend. Extract the square root; this gives the side of the town. [Liu Hui: According to the preceding method, half the side of the town multiplied by itself, divided by the number of bu out of the east gate, is the number of bu out of the south gate. Now the [distances] out of the two gates multiplied together make the area (jifen) of half the [side of the] square town multiplied by itself, occupying one corner. Multiplying it by 4 then gives the area of the four corners. Therefore it is taken as the dividend; extracting the square root then gives the side of the town.]\end{uniten}

\begin{unitzh}\unitlabel{jiuzhang-040.10}今有邑方不知大小，各中開門。出北門二十步有木，出南門十四步，折而西行一千七百七十五步見木。問邑方幾何？荅曰：二百五十步。\end{unitzh}
\begin{uniten}Suppose there is a square town whose size is not known, with a gate opened in the middle of each [side]. 20 bu out of the north gate there is a tree. Going 14 bu out of the south gate and turning to walk west 1775 bu, one sees the tree. Question: what is the side of the town? Answer: 250 bu.\end{uniten}

\begin{unitzh}\unitlabel{jiuzhang-040.11}術曰：以出北門步數乘西行步數，倍之，爲實。{\zhsmall 此以折而西行爲股，自木至邑南十四步爲句，以出北門二十步爲句率，北門至西隅爲股率，卽半廣數。故以出北門句率乘西行股，得半廣股率乘句之幂。然此幂居半以西，故又倍之，合半以東也。}并出南門步數爲從法，開方除之，卽邑方。{\zhsmall 此術之幂，東西廣如邑方，南北自木盡邑南十四步爲袤，合南北步數爲廣袤差，故連并兩步數爲從法，以爲隅外之幂也。}\end{unitzh}
\begin{uniten}Method: Multiply the number of bu of the westward walk by the number of bu out of the north gate, and double it, as the dividend. [Liu Hui: Here the westward walk after turning is taken as the altitude, and [the distance] from the tree to 14 bu south of the town as the base; the 20 bu out of the north gate is taken as the rate of the base, and [the distance] from the north gate to the west corner as the rate of the altitude, which is the number of half the width. Therefore, multiplying the altitude of the westward walk by the rate of the base out of the north gate gives the area of the rate of the altitude, half the width, multiplied by the base. But this area lies in the half to the west; therefore double it again, to combine [it with] the half to the east.] Add the number of bu out of the south gate as the following divisor (congfa); extract the square root; this is the side of the town. [Liu Hui: In the area of this method, the east–west width is equal to the side of the town; south to north, from the tree to the end of the 14 bu south of the town, is the length; the south and north numbers of bu combined are the difference of the width and the length. Therefore the two numbers of bu are joined and added as the following divisor, to make the area outside the corner.]\end{uniten}

\begin{unitzh}\unitlabel{jiuzhang-040.12}今有邑方十里，各中開門。甲、乙俱從邑中央而出：乙東出；甲南出，出門不知步數，邪向東門，磨邑，適與乙會。率：甲行五，乙行三。問甲、乙行各幾何？荅曰：甲出南門八百步，邪東北行四千八百八十七步半及乙。乙東行四千三百一十二步半。\end{unitzh}
\begin{uniten}Suppose there is a town 10 li square, with a gate opened in the middle of each [side]. A and B both go out from the centre of the town: B goes out eastward; A goes out southward, [goes] out of the gate an unknown number of bu, heads slantwise toward the east gate, grazing the town, and just meets B. The rates: A walks 5, B walks 3. Question: how far does each of A and B walk? Answer: A goes 800 bu out of the south gate and walks slantwise northeast 4887 1/2 bu to reach B. B walks east 4312 1/2 bu.\end{uniten}

\begin{unitzh}\unitlabel{jiuzhang-040.13}術曰：令五自乘，三亦自乘，并而半之，爲邪行率。邪行率減於五自乘者，餘爲南行率。以三乘五爲乙東行率。{\zhsmall 求三率之意與上甲乙同。}置邑方，半之，以南行率乘之，如東行率而一，卽得出南門步數。{\zhsmall 邑半方自南門至東隅五里，以爲小股。求出南門步數，爲小股之句。以東行爲股率，南行爲句率。故置邑方半之，以南行句率乘之，如股率而一。}以增邑方半，卽南行。{\zhsmall 半邑者，謂從邑心中停也。}置南行步，求弦者，以邪行率乘之；求東者，以東行率乘之，各自爲實。實如南行率得一步。{\zhsmall 此術與上甲乙同。}\end{unitzh}
\begin{uniten}Method: Let 5 be multiplied by itself, and 3 also be multiplied by itself; add them and halve, as the rate of the slant walk. Subtract the rate of the slant walk from 5 multiplied by itself; the remainder is the rate of the southward walk. Multiply 5 by 3 as the rate of B's eastward walk. [Liu Hui: The idea of finding the three rates is the same as for A and B above.] Set up the side of the town and halve it; multiply it by the rate of the southward walk and divide by the rate of the eastward walk; this gives the number of bu out of the south gate. [Liu Hui: Half the side of the town, from the south gate to the east corner, 5 li, is taken as the small altitude. The number of bu out of the south gate that is sought is the base of the small altitude. The eastward walk is taken as the rate of the altitude, and the southward walk as the rate of the base. Therefore set up the side of the town, halve it, multiply it by the base rate of the southward walk, and divide by the rate of the altitude.] Increase it by half the side of the town; this is the southward walk. [Liu Hui: 'Half the town' means [the distance] from the centre of the town, divided in the middle.] Set up the bu of the southward walk; to find the hypotenuse, multiply it by the rate of the slant walk; to find the eastward [walk], multiply it by the rate of the eastward walk; each separately as the dividend. The dividend divided by the rate of the southward walk gives [the result in units of] 1 bu. [Liu Hui: This method is the same as that for A and B above.]\end{uniten}

\begin{unitzh}\unitlabel{jiuzhang-040.14}有木去人，不知遠近。立四表，相去各一丈，令左兩表與所望參相直。從後右表望之，入前右表三寸。問木去人幾何？荅曰：三十三丈三尺三寸少半寸。\end{unitzh}
\begin{uniten}There is a tree at a distance from a person; it is not known whether it is far or near. Set up four poles (biao), each 1 zhang (丈) apart, and let the two left poles and what is sighted be aligned, the three in a straight line. Sighting it from the rear right pole, [the line of sight] enters 3 cun (寸) inside the front right pole. Question: how far is the tree from the person? Answer: 33 zhang 3 chi (尺) 3 1/3 cun.\end{uniten}

\begin{unitzh}\unitlabel{jiuzhang-040.15}術曰：令一丈自乘爲實，以三寸爲法，實如法而一。{\zhsmall 此以入前右表三寸爲句率，右兩表相去一丈爲股率，左右兩表相去一丈爲見句。所問木去人者，見句之股。股率當乘見句，此二率俱一丈，故曰自乘。以三寸爲法，實如法得一寸。}\end{unitzh}
\begin{uniten}Method: Let 1 zhang multiplied by itself be the dividend, take 3 cun as the divisor, and divide the dividend by the divisor. [Liu Hui: Here the 3 cun entering [inside] the front right pole is taken as the rate of the base, the 1 zhang by which the two right poles are apart as the rate of the altitude, and the 1 zhang by which the left and right poles are apart as the given base. What is asked, the distance of the tree from the person, is the altitude of the given base. The rate of the altitude should multiply the given base; these two rates are both 1 zhang, therefore it says 'multiplied by itself'. Take 3 cun as the divisor; the dividend divided by the divisor gives [the result in units of] 1 cun.]\end{uniten}

\begin{unitzh}\unitlabel{jiuzhang-040.16}有山居木西，不知其高。山去木五十三里，木高九丈五尺。人立木東三里，望木末，適與山峯斜平。人目高七尺。問山高幾何？荅曰：一百六十四丈九尺六寸太半寸。\end{unitzh}
\begin{uniten}There is a mountain lying west of a tree; its height is not known. The mountain is 53 li from the tree, and the tree is 9 zhang 5 chi high. A person stands 3 li east of the tree and sights the tip of the tree, which is just in line, on the slant, with the peak of the mountain. The person's eye is 7 chi high. Question: how high is the mountain? Answer: 164 zhang 9 chi 6 2/3 cun.\end{uniten}

\begin{unitzh}\unitlabel{jiuzhang-040.17}術曰：置木高，減人目高七尺，餘以乘五十三里爲實，以人去木三里爲法，實如法而一，所得加木高，卽山高。{\zhsmall 此術句股之義。以木高減人目高七尺，餘有八丈八尺，爲句率；去人目三里爲股率；山去木五十三里爲見股。以句率乘見股，如股率而一，得句，加木之高，故爲山高也。}\end{unitzh}
\begin{uniten}Method: Set up the height of the tree and subtract the height of the person's eye, 7 chi; with the remainder multiply 53 li, as the dividend; take the person's distance from the tree, 3 li, as the divisor; divide the dividend by the divisor; add what is obtained to the height of the tree; this is the height of the mountain. [Liu Hui: This method [follows] the meaning of base and altitude. Take the height of the tree and subtract the height of the person's eye, 7 chi; the remainder, 8 zhang 8 chi, is the rate of the base; the distance from the person's eye, 3 li, is the rate of the altitude; the mountain's distance from the tree, 53 li, is the given altitude. Multiply the given altitude by the rate of the base and divide by the rate of the altitude, obtaining the base; add the height of the tree; therefore it is the height of the mountain.]\end{uniten}

\begin{unitzh}\unitlabel{jiuzhang-041.1}今有井徑五尺，不知其深。立五尺木於井上，從木末望水岸，入徑四寸。問井深幾何？荅曰：五丈七尺五寸。\end{unitzh}
\begin{uniten}Suppose there is a well 5 chi (尺) in diameter, whose depth is not known. A post of 5 chi is set up on the well; looking from the tip of the post at the water's edge, [the line of sight] enters the diameter by 4 cun (寸). Question: How deep is the well? Answer: 5 zhang (丈) 7 chi 5 cun.\end{uniten}

\begin{unitzh}\unitlabel{jiuzhang-041.2}術曰：置井徑五尺，以入徑四寸減之，餘以乘立木五尺爲實。以入徑四寸爲法。實如法得一寸。{\zhsmall 此以入徑四寸爲句率，立木五尺爲股率，井徑四尺六寸爲見句，問井深者，見句之股也。}\end{unitzh}
\begin{uniten}Method: Set up the well's diameter, 5 chi, and subtract from it the 4 cun entered into the diameter; with the remainder multiply the standing post, 5 chi, making the dividend (shi). Take the 4 cun entered into the diameter as the divisor (fa). Divide the dividend by the divisor; [each unit] obtained is 1 cun. [Liu Hui: Here the 4 cun entered into the diameter is taken as the rate (lü) of the base (gou), the standing post 5 chi as the rate of the altitude (gu), and the well's diameter 4 chi 6 cun as the given base (jian gou); the depth of the well that is asked for is the altitude of the given base.]\end{uniten}

\begin{unitzh}\unitlabel{jiuzhang-041.3}《九章算術》卷之九終。錢唐莫濰校字。\end{unitzh}
\begin{uniten}The Nine Chapters on the Mathematical Art, Chapter 9, end. Characters collated by Mo Wei of Qiantang.\end{uniten}

\begin{unitzh}\unitlabel{jiuzhang-041.4}大淸乾隆三十八年癸巳秋，闕里孔氏依汲古閣影宋刻本重雕。\end{unitzh}
\begin{uniten}In the autumn of the 38th year of Qianlong of the Great Qing, a guisi year, the Kong family of Queli recut [the blocks] following the Jiguge facsimile of a Song printed edition.\end{uniten}

\begin{unitzh}\unitlabel{jiuzhang-041.5}秘書省。《九章算術》一部，共九冊。元豐七年九月 日。校定：降授宣德郎秘書省校書郞{\zhsmall 臣}葉祖洽上進。校定：承議郎行秘書省校書郎{\zhsmall 臣}王仲脩。校定：朝奉郎行秘書省校書郎{\zhsmall 臣}錢長卿。奉議郎守秘書郞丞{\zhsmall 臣}韓宗古。朝請郎試秘書少監{\zhsmall 臣}趙彦若。\end{unitzh}
\begin{uniten}Palace Library (Mishu sheng). The Nine Chapters on the Mathematical Art, 1 set, 9 fascicles (ce) in all. 7th year of Yuanfeng, 9th month, [blank] day. Collated and settled: demoted to Gentleman for Promoting Virtue (Xuande lang), Editor in the Palace Library, [Note: your servant] Ye Zuqia, presents it to the throne. Collated and settled: Gentleman for Comprehensive Counsel (Chengyi lang), serving as Editor in the Palace Library, [Note: your servant] Wang Zhongxiu. Collated and settled: Gentleman for Court Service (Chaofeng lang), serving as Editor in the Palace Library, [Note: your servant] Qian Changqing. Gentleman for Rendering Counsel (Fengyi lang), acting Assistant (lang cheng) of the Palace Library, [Note: your servant] Han Zonggu. Gentleman for Court Audiences (Chaoqing lang), probationary Vice Director of the Palace Library, [Note: your servant] Zhao Yanruo.\end{uniten}

\begin{unitzh}\unitlabel{jiuzhang-041.6}元豐七年九月二十八日進呈，奉御寶批：「宐依已校定鏤板。」朝奉郎祕書丞上騎都尉賜緋魚袋{\zhsmall 臣}韓治。朝散郎試祕書少監上騎都尉賜緋魚袋{\zhsmall 臣}顧臨。朝議大夫守祕書少監上䕶軍賜紫金魚袋{\zhsmall 臣}劉攽。\end{unitzh}
\begin{uniten}7th year of Yuanfeng, 9th month, 28th day: presented to the throne, and the endorsement under the imperial seal was received: "It is fitting to cut the blocks according to what has already been collated and settled." Gentleman for Court Service (Chaofeng lang), Assistant Director of the Palace Library, Superior Commandant of Cavalry, granted the crimson fish pouch, [Note: your servant] Han Zhi. Gentleman for Court Discussion (Chaosan lang), probationary Vice Director of the Palace Library, Superior Commandant of Cavalry, granted the crimson fish pouch, [Note: your servant] Gu Lin. Grand Master for Court Counsel (Chaoyi dafu), acting Vice Director of the Palace Library, Superior Protector-General, granted the purple and gold fish pouch, [Note: your servant] Liu Ban.\end{uniten}

\begin{unitzh}\unitlabel{jiuzhang-041.7}中大夫守尙書右丞䕶軍東平郡開國侯，食邑貳千叁百戸，賜紫金魚袋{\zhsmall 臣}呂大防。通議大夫守尙書左丞上柱國平原郡開國公，食邑貳千捌百戸，食實封伍伯戸{\zhsmall 臣}李淸臣。正議大夫守中書侍郞上柱國馮翊郡開國公，食邑貳千叁百戸，食實封伍伯戸{\zhsmall 臣}張璪。正議大夫守門下侍郎上柱國南陽郡開國公，食邑貳千壹百戸，食實封壹阡戸{\zhsmall 臣}韓維。金紫光祿大夫守尙書右僕射兼中書侍郎上柱國東平郡開國公，食邑六千二百戸，食實封壹阡玖伯戸{\zhsmall 臣}呂公著。正議大夫守尙書左僕射兼門下侍郎上柱國河内郡開國公，食邑四千一百戸，食實封壹阡伍伯戸{\zhsmall 臣}司馬光。\end{unitzh}
\begin{uniten}Grand Master of the Palace (Zhong dafu), acting Right Assistant Director of the Department of State Affairs, Protector-General, Dynasty-Founding Marquis of Dongping Commandery, with a fief of 2300 households, granted the purple and gold fish pouch, [Note: your servant] Lü Dafang. Grand Master for Thorough Counsel (Tongyi dafu), acting Left Assistant Director of the Department of State Affairs, Supreme Pillar of State, Dynasty-Founding Duke of Pingyuan Commandery, with a fief of 2800 households and an actual fief of 500 households, [Note: your servant] Li Qingchen. Grand Master for Upright Counsel (Zhengyi dafu), acting Vice Director of the Secretariat, Supreme Pillar of State, Dynasty-Founding Duke of Fengyi Commandery, with a fief of 2300 households and an actual fief of 500 households, [Note: your servant] Zhang Zao. Grand Master for Upright Counsel, acting Vice Director of the Chancellery, Supreme Pillar of State, Dynasty-Founding Duke of Nanyang Commandery, with a fief of 2100 households and an actual fief of 1000 households, [Note: your servant] Han Wei. Grand Master for Splendid Happiness with Gold and Purple (Jinzi guanglu dafu), acting Right Vice Director of the Department of State Affairs and concurrently Vice Director of the Secretariat, Supreme Pillar of State, Dynasty-Founding Duke of Dongping Commandery, with a fief of 6200 households and an actual fief of 1900 households, [Note: your servant] Lü Gongzhu. Grand Master for Upright Counsel, acting Left Vice Director of the Department of State Affairs and concurrently Vice Director of the Chancellery, Supreme Pillar of State, Dynasty-Founding Duke of Henei Commandery, with a fief of 4100 households and an actual fief of 1500 households, [Note: your servant] Sima Guang.\end{uniten}

\begin{unitzh}\unitlabel{jiuzhang-041.8}《九章算術》卷九訂訛補圖{\zhsmall 《算經十書》之二}。休寧戴震東原。\end{unitzh}
\begin{uniten}The Nine Chapters on the Mathematical Art, Chapter 9: Corrections of Errors and Supplementary Diagrams [Note: The 2nd of the Ten Mathematical Classics]. Dai Zhen (Dongyuan) of Xiuning.\end{uniten}

\begin{unitzh}\unitlabel{jiuzhang-041.9}句股弦互求圖 □□□□□□□□□ □□□□□□□□□ □□□□□□□□□ □□□□□□□□□ □□□□□□□□□ □□□□□□□□□{\zhsmall 注云：「句自乘爲朱方，股自乘爲靑方，令出入相補，合成弦方之幂。」又李淳風等釋云：「句方於内，則句短於股。」據此，則劉徽注此章舊有圖而缺，今並按注補圖。}\end{unitzh}
\begin{uniten}Diagram for finding base, altitude and hypotenuse from one another. [illegible: diagram, 6 rows × 9 characters] [Dai Zhen: The commentary says: "The base (gou) multiplied by itself makes the vermilion (zhu) square, the altitude (gu) multiplied by itself makes the azure (qing) square; let what goes out and what comes in make up for each other, together forming the area (mi) of the square on the hypotenuse (xian)." Further, Li Chunfeng et al. explain: "If the square on the base is inside, then the base is shorter than the altitude." According to this, Liu Hui's commentary on this chapter originally had diagrams, which are lost; now diagrams are supplied for all of them in accordance with the commentary.]\end{uniten}

\begin{unitzh}\unitlabel{jiuzhang-041.10}據圍廣木長求葛之長。其形：葛卷裏袤，以筆管靑線宛轉，有似葛之纒木。解而觀之，則每周之閒自有相閒成句股弦。則其閒木長爲股，圍之爲句，葛長爲弦。弦七周乘三圍，是并合衆句以爲一句，則句長而股短。故術以木長謂之句，圍之謂之股，言之倒互。句與股求弦，亦如前圖。\end{unitzh}
\begin{uniten}From the girth as width (guang) and the length of the tree, find the length of the vine. Its form: the vine coils around the length (mao); with a brush tube and a green thread wound round it, it resembles the vine wrapping the tree. Unwind it and look at it: then between each turn there are naturally [segments], alternating with one another, that form base, altitude and hypotenuse. Then the length of the tree between them is the altitude, the girth is the base, and the length of the vine is the hypotenuse. The hypotenuse: 7 turns multiplying the girth 3 — this is joining the many bases together to make 1 base; then the base is long and the altitude short. Therefore the method calls the length of the tree the base and the girth the altitude: they are spoken of inverted and interchanged. [Finding] the hypotenuse from base and altitude is also as in the preceding diagram.\end{uniten}

\begin{unitzh}\unitlabel{jiuzhang-041.11}句三自乘爲朱幂，股四自乘爲靑幂，合朱靑二十五，爲弦五自乘幂，出上第一圖。句股幂合爲弦幂，明矣。然二幂之數謂倒在於弦幂之中而已，可更相裏者則成方幂，其居表者則成矩幂，二表裏形訛而數均。又按此圖，句幂之矩靑卷白表，是其幂以句股弦差爲廣，股弦并爲袤，而股幂方其裏。股幂之矩靑卷白表，是其幂以句弦差爲廣，句弦并爲袤，而句幂方其裏。是故差之與并用除之，短長互相乘也。{\zhsmall 據趙君卿注《周髀算經》云：「凡并句股之實，卽成弦實。或矩於外，或方於内，形詭而量均，體殊而數齊。句實之矩，以股弦差爲廣，股弦并爲袤，而股實方其裏。股實之矩，以句弦差爲廣，句弦并爲袤，而句實方其裏。」君卿漢人，此注葢用其說，而傳寫失真，加以後人竄改，遂不可通。}\end{unitzh}
\begin{uniten}The base 3 multiplied by itself makes the vermilion area; the altitude 4 multiplied by itself makes the azure area; vermilion and azure combined, 25, make the area of the hypotenuse 5 multiplied by itself, which comes out of the 1st diagram above. That the areas of base and altitude combine to make the area of the hypotenuse is clear. However, the amounts of the 2 areas are said to be inverted within the area of the hypotenuse, that is all: what can in turn be put inside forms a square area, and what lies on the outside forms a carpenter's-square (ju) area; the 2, outer and inner, are altered in shape but equal in amount. Further, according to this diagram, the carpenter's square of the base area is azure, wrapping the white on the outside; its area takes the difference (cha) of base, altitude and hypotenuse as width and the sum (bing) of altitude and hypotenuse as length, and the altitude area is square within it. The carpenter's square of the altitude area is azure, wrapping the white on the outside; its area takes the difference of base and hypotenuse as width and the sum of base and hypotenuse as length, and the base area is square within it. Therefore the difference and the sum are used to divide it: short and long are multiplied by each other in turn. [Dai Zhen: According to Zhao Junqing's commentary on the Zhoubi suanjing, which says: "In general, combining the areas (shi) of base and altitude forms the area of the hypotenuse. Whether as a carpenter's square on the outside or as a square on the inside, the shapes are different but the quantities equal, the bodies distinct but the numbers the same. The carpenter's square of the base area takes the difference of altitude and hypotenuse as width and the sum of altitude and hypotenuse as length, and the altitude area is square within it. The carpenter's square of the altitude area takes the difference of base and hypotenuse as width and the sum of base and hypotenuse as length, and the base area is square within it." Junqing was a man of the Han; this commentary presumably used his explanation, but in copying it lost its true form, and in addition later people altered it, so that it became unintelligible.]\end{uniten}

\begin{unitzh}\unitlabel{jiuzhang-041.12}句股差句股并與弦互求之圖 □□□□□□□□□ □□□□□□□□□ □□□□□□□□□ □□□□□□□□□ □□□□□□□□□ □□□□□□□□□{\zhsmall 據注意，滿大方卽句股并自乘之幂，有朱幂八、黃幂一；内小方卽弦實，有朱幂四、黃幂一。倍弦實，滿外大方而多一黃幂。於弦實内減黃幂四分之二，半其餘，適得外大方四分之一。朱幂者，句股積也；黃幂之面，卽句股差。}\end{unitzh}
\begin{uniten}Diagram for finding the difference of base and altitude, the sum of base and altitude, and the hypotenuse from one another. [illegible: diagram, 6 rows × 9 characters] [Dai Zhen: According to the meaning of the commentary, the full large square is the area of the sum of base and altitude multiplied by itself, containing 8 vermilion areas and 1 yellow (huang) area; the inner small square is the area of the hypotenuse, containing 4 vermilion areas and 1 yellow area. Doubling the area of the hypotenuse fills the outer large square with 1 yellow area in excess. From the area of the hypotenuse subtract 2/4 of the yellow area and halve the remainder: this gives exactly 1/4 of the outer large square. The vermilion area is the area (ji) of base and altitude; the side (mian) of the yellow area is the difference of base and altitude.]\end{uniten}

\begin{unitzh}\unitlabel{jiuzhang-041.13}股實之矩圖 □□□□□□□ □□□□□□□ □□□□□□□ □□□□□□□{\zhsmall 趙君卿云：「股實之矩，以句弦差爲廣，句弦并爲袤，而句實方其裏。」劉注用句弦差、句弦并與句股弦互求，本此。}\end{unitzh}
\begin{uniten}Diagram of the carpenter's square of the altitude area. [illegible: diagram, 4 rows × 7 characters] [Dai Zhen: Zhao Junqing says: "The carpenter's square of the altitude area takes the difference of base and hypotenuse as width and the sum of base and hypotenuse as length, and the base area is square within it." Liu's commentary, in which the difference of base and hypotenuse and the sum of base and hypotenuse are used in finding base, altitude and hypotenuse from one another, is based on this.]\end{uniten}

\begin{unitzh}\unitlabel{jiuzhang-041.14}句實之矩圖 □□□□□□ □□□□□□ □□□□□□ □□□□□□{\zhsmall 趙君卿云：「句實之矩，以股弦差爲廣，股弦并爲袤，而股實方其裏。」劉注用股弦差、股弦并與句股弦互求，本此。}\end{unitzh}
\begin{uniten}Diagram of the carpenter's square of the base area. [illegible: diagram, 4 rows × 6 characters] [Dai Zhen: Zhao Junqing says: "The carpenter's square of the base area takes the difference of altitude and hypotenuse as width and the sum of altitude and hypotenuse as length, and the altitude area is square within it." Liu's commentary, in which the difference of altitude and hypotenuse and the sum of altitude and hypotenuse are used in finding base, altitude and hypotenuse from one another, is based on this.]\end{uniten}

\begin{unitzh}\unitlabel{jiuzhang-041.15}臣淳風等謹按：下鐻深得一寸爲半股弦差。注云：「爲股弦差者，鐻道也。」{\zhsmall 案：此言「下鐻深得一寸爲半股弦差」，卽注所謂「鐻深一寸爲股弦差之一半」也。更綴「注云爲股弦差者鐻道也」十字，舛誤不可通。據割圓術，鐻深一寸卽可爲股弦差，半鐻道五寸爲句，材半徑爲弦，若以此言之，尤合術意。}\end{unitzh}
\begin{uniten}Li Chunfeng et al. respectfully note: the depth of 1 cun obtained by sawing downward is half the difference of altitude and hypotenuse. The commentary says: "What is taken as the difference of altitude and hypotenuse is the saw cut (ju dao)." [Dai Zhen: Note: what this says, "the depth of 1 cun obtained by sawing downward is half the difference of altitude and hypotenuse", is what the commentary calls "the saw depth of 1 cun is 1/2 of the difference of altitude and hypotenuse". Further appending the 10 characters "The commentary says: what is taken as the difference of altitude and hypotenuse is the saw cut" is erroneous and unintelligible. According to the method of cutting the circle (geyuan shu), the saw depth of 1 cun can directly be the difference of altitude and hypotenuse, half the saw cut, 5 cun, the base, and the half-diameter of the timber the hypotenuse; if it were stated in this way, it would accord especially with the meaning of the method.]\end{uniten}

\begin{unitzh}\unitlabel{jiuzhang-041.16}股弦卽高廣袤，其出此圖也。其倍弦爲廣袤合，矩句卽爲幂，得廣卽句股差。{\zhsmall 股弦卽高廣袤以下訛舛。據趙君卿注《周髀算經》云：「其倍弦爲廣袤合，而令句股見者自乘，其爲實。四實以減之，開其餘，所得爲差。以差減合，半其餘爲廣。減廣於弦，卽所求也。」此注似用其說，而傳寫舛訛，後人又妄加改竄，遂不可通。就君卿說考之：倍弦自乘得弦實四，内有句實、股實各四。減四句實，餘卽四股實，開之得倍股；減四股實，餘卽四句實，開之得倍句：所謂「開其餘，所得爲差」也。減倍股於倍弦，半其餘爲股弦差；減倍句於倍弦，半其餘爲句弦差：所謂「以差減合，半其餘爲廣」也。減股弦差於弦卽股，減句弦差於弦卽句：所謂「減廣於弦，卽所求」也。凡股弦差爲廣，股弦并爲袤，其幂卽句幂；句弦差爲廣，句弦并爲袤，其幂卽股幂。合廣袤皆成倍弦，故曰「倍弦爲廣袤合」，而倍句、倍股卽廣袤差。}\end{unitzh}
\begin{uniten}The altitude and hypotenuse are then the height, width and length; this comes out of this diagram. Twice the hypotenuse makes the sum of width and length; the carpenter's-square base (ju gou) then makes the area; the width obtained is then the difference of base and altitude. [Dai Zhen: From "the altitude and hypotenuse are then the height, width and length" onward [the text] is erroneous. According to Zhao Junqing's commentary on the Zhoubi suanjing, which says: "Twice the hypotenuse is the sum of width and length; let whichever of base and altitude is given be multiplied by itself; this is the area. Subtract 4 areas from it, extract the root of the remainder: what is obtained is the difference. Subtract the difference from the sum and halve the remainder: this is the width. Subtract the width from the hypotenuse: this is what is sought." This commentary seems to use his explanation, but in copying it became garbled, and later people moreover recklessly altered it, so that it became unintelligible. Examining it according to Junqing's explanation: twice the hypotenuse multiplied by itself gives 4 hypotenuse areas, within which there are 4 each of base areas and altitude areas. Subtract 4 base areas: the remainder is 4 altitude areas, and extracting its root gives twice the altitude; subtract 4 altitude areas: the remainder is 4 base areas, and extracting its root gives twice the base: this is what is meant by "extract the root of the remainder: what is obtained is the difference". Subtract twice the altitude from twice the hypotenuse and halve the remainder: this is the difference of altitude and hypotenuse; subtract twice the base from twice the hypotenuse and halve the remainder: this is the difference of base and hypotenuse: this is what is meant by "subtract the difference from the sum and halve the remainder: this is the width". Subtract the difference of altitude and hypotenuse from the hypotenuse: this is the altitude; subtract the difference of base and hypotenuse from the hypotenuse: this is the base: this is what is meant by "subtract the width from the hypotenuse: this is what is sought". In general, with the difference of altitude and hypotenuse as width and the sum of altitude and hypotenuse as length, its area is the base area; with the difference of base and hypotenuse as width and the sum of base and hypotenuse as length, its area is the altitude area. Width and length combined always make twice the hypotenuse; therefore it says "twice the hypotenuse is the sum of width and length"; and twice the base and twice the altitude are the differences of width and length.]\end{uniten}

\begin{unitzh}\unitlabel{jiuzhang-041.17}其矩句之幂，倍爲從法，開之亦句股差。其餘以句股幂減，半其餘，差爲從法，開方除之，卽句也。{\zhsmall 據趙君卿云：「減矩句之實於弦實，開其餘卽股。倍股在兩邊爲從法，開矩句之角卽股弦差。減矩股之實於弦實，開其餘卽句。倍句在兩邊爲從法，開矩股之角卽句弦差。」此注亦用其說，而殘缺失次，遂不可通。}\end{unitzh}
\begin{uniten}The area of its carpenter's-square base: doubling makes the following divisor (congfa); extracting the root also gives the difference of base and altitude. From the remainder subtract the area of base and altitude, halve the remainder; the difference makes the following divisor; extract the square root: this is the base. [Dai Zhen: According to Zhao Junqing, who says: "Subtract the area of the carpenter's-square base from the area of the hypotenuse, extract the root of the remainder: this is the altitude. Twice the altitude, lying on the 2 sides, makes the following divisor; extracting the root of the corner of the carpenter's-square base gives the difference of altitude and hypotenuse. Subtract the area of the carpenter's-square altitude from the area of the hypotenuse, extract the root of the remainder: this is the base. Twice the base, lying on the 2 sides, makes the following divisor; extracting the root of the corner of the carpenter's-square altitude gives the difference of base and hypotenuse." This commentary also uses his explanation, but it is fragmentary and out of order, so that it became unintelligible.]\end{uniten}

\begin{unitzh}\unitlabel{jiuzhang-041.18}句實廣袤合圖 □□□□□□□□□□□□□ □□□□□□□□□□□□□ □□□□□□□□□□□□□ □□□□□□□□□□□□□ □□□□□□□□□□□□□ □□□□□□□□□□□□□ □□□□□□□□□□□□□ □□□□□□□□□□□□□{\zhsmall 據注意，句實以股弦差爲廣，股弦并爲袤，廣袤合亦卽倍弦。減四句實於倍弦自乘之幂，餘爲四股實，開方得倍股，卽句實之廣袤差也。}\end{unitzh}
\begin{uniten}Diagram of the combined width and length of the base area. [illegible: diagram, 8 rows × 13 characters] [Dai Zhen: According to the meaning of the commentary, the base area takes the difference of altitude and hypotenuse as width and the sum of altitude and hypotenuse as length; the sum of width and length is again just twice the hypotenuse. Subtract 4 base areas from the area of twice the hypotenuse multiplied by itself; the remainder is 4 altitude areas; extracting the square root gives twice the altitude, which is the difference of width and length of the base area.]\end{uniten}

\begin{unitzh}\unitlabel{jiuzhang-041.19}股□□□□□□□□□□□□□ □□□□□□□□□□□□□ □□□□□□□□□□□□□ □□□□□□□□□□□□□ □□□□□□□□□□□□□ □□□□□□□□□□□□□ □□□□□□□□□□□□□ □□□□□□□□□□□□□{\zhsmall 據注意，倍弦自乘之幂容句實、股實各四。股實以句弦差爲廣，句弦并爲袤，廣袤合卽倍弦。減四股實於倍弦自乘之幂，餘爲四句實，開方得倍句，卽股實之廣袤差也。}\end{unitzh}
\begin{uniten}Altitude (gu) [illegible ×13] [illegible: diagram, 7 further rows × 13 characters] [Dai Zhen: According to the meaning of the commentary, the area of twice the hypotenuse multiplied by itself contains 4 each of base areas and altitude areas. The altitude area takes the difference of base and hypotenuse as width and the sum of base and hypotenuse as length; the sum of width and length is twice the hypotenuse. Subtract 4 altitude areas from the area of twice the hypotenuse multiplied by itself; the remainder is 4 base areas; extracting the square root gives twice the base, which is the difference of width and length of the altitude area.]\end{uniten}

\begin{unitzh}\unitlabel{jiuzhang-041.20}句弦差股弦差求句股弦之圖 □□□□□□□□□□□□□ □□□□□□□□□□□□□ □□□□□□□□□□□□□ □□□□□□□□□□□□□ □□□□□□□□□□□□□ □□□□□□□□□□□□□ □□□□□□□□□□□□□{\zhsmall 句弦差乘股弦差，倍之，爲兩亷，其幂與黃方相等，故開方除之，得黃方之面。加股弦差卽句，加句弦差卽股，加兩差卽弦。}\end{unitzh}
\begin{uniten}Diagram for finding base, altitude and hypotenuse from the difference of base and hypotenuse and the difference of altitude and hypotenuse. [illegible: diagram, 7 rows × 13 characters] [Dai Zhen: Multiply the difference of base and hypotenuse by the difference of altitude and hypotenuse and double it: this makes the 2 edges (lian); their area is equal to the yellow square; therefore extract the square root, obtaining the side of the yellow square. Adding the difference of altitude and hypotenuse gives the base; adding the difference of base and hypotenuse gives the altitude; adding the 2 differences gives the hypotenuse.]\end{uniten}

\begin{unitzh}\unitlabel{jiuzhang-041.21}術曰：令七自乘，三亦自乘，并而半之，以爲甲邪行率。邪行率減於七自乘，餘爲南行率。以三乘七爲乙東行率。{\zhsmall 據問意，甲行率七者，設句弦并七也；乙行率三者，設股三也。術「令七自乘，三亦自乘，并而半之，爲甲邪行率」者，句弦并自乘加股自乘，半之，卽弦乘句弦并所得數，以爲弦率。「邪行率減於七自乘，餘爲南行率」者，句弦并自乘減弦乘句弦并，餘卽句乘句弦并所得數，以爲句率。弦率、句率皆通之以句弦并，則股三亦以乘句弦并七，乃爲股率。句股弦三率，皆句弦并通句股弦所得數。然則七自乘卽句弦并之率，三自乘卽句弦差之率，合差於并而半之爲弦，減弦於并爲句，其義一而已矣。}\end{unitzh}
\begin{uniten}Method: Let 7 be multiplied by itself, and 3 also be multiplied by itself; add them and halve, taking this as the rate of A's oblique travel (xie xing). Subtract the rate of oblique travel from 7 multiplied by itself; the remainder is the rate of southward travel. Multiply 7 by 3 to make the rate of B's eastward travel. [Dai Zhen: According to the meaning of the problem, A's rate of travel 7 means assuming the sum of base and hypotenuse to be 7; B's rate of travel 3 means assuming the altitude to be 3. As for the method's "let 7 be multiplied by itself, and 3 also be multiplied by itself, add them and halve, making A's rate of oblique travel": the sum of base and hypotenuse multiplied by itself plus the altitude multiplied by itself, halved, is the number obtained by the hypotenuse multiplying the sum of base and hypotenuse, which is taken as the rate of the hypotenuse. As for "subtract the rate of oblique travel from 7 multiplied by itself; the remainder is the rate of southward travel": the sum of base and hypotenuse multiplied by itself, minus the hypotenuse multiplied by the sum of base and hypotenuse, leaves the number obtained by the base multiplying the sum of base and hypotenuse, which is taken as the rate of the base. The rate of the hypotenuse and the rate of the base are both multiplied through (tong) by the sum of base and hypotenuse; then the altitude 3 is also used to multiply the sum of base and hypotenuse 7, which makes the rate of the altitude. The 3 rates of base, altitude and hypotenuse are all numbers obtained by multiplying base, altitude and hypotenuse through by the sum of base and hypotenuse. That being so, 7 multiplied by itself is the rate of the sum of base and hypotenuse, and 3 multiplied by itself is the rate of the difference of base and hypotenuse; adding the difference to the sum and halving gives the hypotenuse, and subtracting the hypotenuse from the sum gives the base: the meaning is one and the same.]\end{uniten}

\begin{unitzh}\unitlabel{jiuzhang-041.22}股□□□□□□□□□□□□□ □□□□□□□□□□□□□ □□□□□□□□□□□□□ □□□□□□□□□□□□□ □□□□□□□□□□□□□ □□□□□□□□□□□□□ □□□□□□□□□□□□□ □□□□□□□□□□□□□{\zhsmall 據注意，句弦并自乘加股自乘，成長方，以兩弦爲袤，句弦并爲廣，故半之爲弦率。減股自乘於弦率，餘爲句率。股自乘之矩，損上橫幂益於下，以句弦并爲袤，差爲廣。}\end{unitzh}
\begin{uniten}Altitude (gu) [illegible ×13] [illegible: diagram, 7 further rows × 13 characters] [Dai Zhen: According to the meaning of the commentary, the sum of base and hypotenuse multiplied by itself plus the altitude multiplied by itself forms a long rectangle (changfang), with 2 hypotenuses as length and the sum of base and hypotenuse as width; therefore halving it gives the rate of the hypotenuse. Subtract the altitude multiplied by itself from the rate of the hypotenuse; the remainder is the rate of the base. The carpenter's square of the altitude multiplied by itself: take away the upper horizontal area and add it below; it then takes the sum of base and hypotenuse as length and the difference as width.]\end{uniten}

\begin{unitzh}\unitlabel{jiuzhang-041.23}句股相乘爲朱靑黃幂各二，令黃幂連於下隅，朱靑各以類合，共成修幂：中方黃爲廣，并句股爲袤，故并句股爲法。幂圖：方在句中，則方之兩亷各自成小句股，而其相與之勢不失本率也。\end{unitzh}
\begin{uniten}Base and altitude multiplied together make vermilion, azure and yellow areas, 2 of each. Let the yellow areas be joined at the lower corner, and the vermilion and azure each be combined according to their kind; together they form an elongated area (xiu mi): the yellow of the middle square is the width, and the sum of base and altitude is the length; therefore the sum of base and altitude is taken as the divisor. In the area diagram the square lies within the base [triangle]; then the 2 edges of the square each form a small base-and-altitude [triangle] of their own, and their mutual configuration does not depart from the original rates.\end{uniten}

\begin{unitzh}\unitlabel{jiuzhang-041.24}句面之小股，股面之小句，從橫相連，合而成中方。令股爲中方率，并句股爲廣率，據見句五步而今有之，得中方也。復令句爲中方率，以并句股爲袤率，據股一十二步而今有之，則中方又可知。此則雖不效而法實有法由生矣。{\zhsmall 此十三字舛誤。據上以粟米章今有術及衰分章列衰之意解此術，大小句股互求，并句股卽所有率，中方率}\end{unitzh}
\begin{uniten}The small altitude on the side of the base and the small base on the side of the altitude, joined vertically and horizontally, combine to form the middle square. Let the altitude be the rate of the middle square and the sum of base and altitude the rate of the width; based on the given base of 5 bu (步), treat it by "suppose there is" (jinyou): this gives the middle square. Again let the base be the rate of the middle square, and take the sum of base and altitude as the rate of the length; based on the altitude of 12 bu, treat it by "suppose there is": then the middle square can also be known. This then, although not imitated, and the method — in fact there is that from which the method arises. [Dai Zhen: These 13 characters are erroneous. According to the above, this method is explained by the idea of the "suppose there is" method (jinyou shu) of the Millet and Grain chapter and of the arranged proportions (lie cui) of the Distribution by Proportion chapter. The large and small base-and-altitude [triangles] are found from one another; the sum of base and altitude is the rate of what one has (suoyou lü); the rate of the middle square] [continues]\end{uniten}

\begin{unitzh}\unitlabel{jiuzhang-042.1}{\zhsmall 卽所求率，見句、見股卽所有數。於事雖不同，而意相倣效，實術所由生也。注意當是如此。}\end{unitzh}
\begin{uniten}[continued] [Note: …is the rate sought (suoqiu lü); the known base (gou) and the known altitude (gu) are the given number (suoyou shu). Although they are not the same in matter, in idea they imitate one another: this is truly where the method arises. The intention of the commentary must have been thus.]\end{uniten}

\begin{unitzh}\unitlabel{jiuzhang-042.2}句股容方圖\end{unitzh}
\begin{uniten}Diagram of the square contained (rong fang) in base and altitude\end{uniten}

\begin{unitzh}\unitlabel{jiuzhang-042.3}□□□□□□□□□□□□□□□□□□□{\zhsmall 據注意，句股相乘得兩黄幂、兩朱幂、兩靑}□□□□□□□□□□□□□□□□□□□{\zhsmall 幂，令黃幂從黄幂，朱幂從朱幂，靑幂從靑}□□□□□□□□□□□□□□□□□□□{\zhsmall 幂，兩朱在上，兩黄在中，兩靑合而横之在}□□□□□□□□□□□□□□□□□□□{\zhsmall 下，則句股并爲幂，容方之面爲廣。其朱幂、}□□□□□□□□□□□□□□□□□□□{\zhsmall 靑幂各成小股，皆如大句股本率互求。}\end{unitzh}
\begin{uniten}[illegible] [Note: According to the intention of the commentary, base and altitude multiplied together give 2 yellow areas (mi), 2 vermilion areas and 2 blue-green] [illegible] [areas; let the yellow areas follow the yellow areas, the vermilion areas follow the vermilion areas, and the blue-green areas follow the blue-green] [illegible] [areas: the 2 vermilion are on top, the 2 yellow in the middle, and the 2 blue-green are joined and laid horizontally at the] [illegible] [bottom. Then the sum of base and altitude makes the area, and the side of the contained square is the width (guang). The vermilion areas and] [illegible] [blue-green areas each form a small altitude; all, as with the large base and altitude, are found from one another by the original rates (ben lü).]\end{uniten}

\begin{unitzh}\unitlabel{jiuzhang-042.4}句股相乘爲圖之本體，朱、靑、黄幂各二，則倍之爲各四。可用畫於小𥿄，分裁邪正之會，令顛倒相補，各以類合，成修幂：圓徑爲廣，並句股弦爲袤。故并句股弦以爲法。又以圖之大體言之，股中靑必令立規於橫廣，句股又邪三徑均而復連規{\zhsmall 此二十一字舛誤。據容圓之徑，卽減弦於句股并之餘也。取半徑規之，又以半徑減句股，其餘并之，適爲弦。如是截句股弦各爲二，三半徑均而復連於規之中央。注意葢以此爲言，而殘缺失次，遂不可通。}，從橫量度句股，必合而成小方矣。\end{unitzh}
\begin{uniten}Base and altitude multiplied together make the basic body of the diagram, with vermilion, blue-green and yellow areas, 2 of each; then doubling it makes 4 of each. It can be drawn on a small piece of paper: cut it apart at the meeting points of the oblique and the upright, let them be turned upside down to complement one another, and join each with its own kind, forming an elongated area (xiu mi): the diameter of the circle is the width, and base, altitude and hypotenuse (xian) summed are the length (mao). Therefore base, altitude and hypotenuse are summed to make the divisor (fa). Again, speaking in terms of the general form of the diagram, the blue-green in the altitude must be made to set up the compass (gui) on the horizontal width; base and altitude and also the oblique, the 3 diameters equal, and again joining the compass [Note: These 21 characters are erroneous. On examination, the diameter of the contained circle is precisely the remainder when the hypotenuse is subtracted from the sum of base and altitude. Take the half-diameter (ban jing) and draw a circle with it with the compass; also subtract the half-diameter from base and altitude, and sum the remainders: this is exactly the hypotenuse. In this way base, altitude and hypotenuse are each cut into 2, and the 3 half-diameters are equal and moreover join at the centre of the compass-drawn circle. The intention of the commentary was presumably to say this, but, being fragmentary and out of order, it cannot be understood.]; measuring base and altitude lengthwise and crosswise, they must combine to form a small square.\end{uniten}

\begin{unitzh}\unitlabel{jiuzhang-042.5}又畫中弦以觀其會，則句股之中成小句股弦四者。句面之小股、股面之小句，皆小方之面，皆圓徑之半。其數故可衰。以句股弦爲列衰，副并爲法。以小句乘未并者，各自爲實。實如法而一，得句面之小股可知也。以股乘列衰爲實，則得股面之小句可知。{\zhsmall 以小句乘未并者至此，訛舛不可通。或後人妄加改竄，又援衰分章之文入於此，遂漫無辨别。當是言：令股爲列衰，以見句乘之爲實，實如法而一，則句面之小股可知也；令句爲列衰，以見股乘之爲實，實如法而一，則股面之小句可知。此在粟米章，卽今有術，以所求率乘所有數，所有率除之，古算家謂之異乘同除。注以解大小句股互求，句率、股率謂所有率及所求率，見句或見股爲所有數，不可以見句乘句率、見股乘股率也。}\end{unitzh}
\begin{uniten}Further, draw the middle hypotenuses so as to observe where they meet; then within the base and altitude are formed 4 small base-altitude-hypotenuse figures. The small altitude on the base side and the small base on the altitude side are both the side of the small square, and both half the diameter of the circle. Therefore their numbers can be proportioned (shuai). Take base, altitude and hypotenuse as the arrayed proportions (lie shuai), and sum them on the side (fu bing) to make the divisor. Multiply those not summed by the small base, each separately making a dividend (shi). Divide the dividend by the divisor: that the small altitude on the base side is obtained can be known. Multiply the arrayed proportions by the altitude to make the dividend: then that the small base on the altitude side is obtained can be known. [Note: From "multiply those not summed by the small base" to here, it is corrupt and erroneous and cannot be understood. Perhaps later people rashly altered it, and also brought the text of the chapter Distribution by Proportion (Cuifen) into this place, so that it has become a muddle beyond distinguishing. It should be saying: let the altitude be the arrayed proportion, multiply it by the known base to make the dividend, divide the dividend by the divisor, and then the small altitude on the base side can be known; let the base be the arrayed proportion, multiply it by the known altitude to make the dividend, divide the dividend by the divisor, and then the small base on the altitude side can be known. This is in the chapter Millet and Grain (Sumi), namely the "Suppose there is" method (jinyou shu): multiply the given number by the rate sought, and divide by the given rate (suoyou lü); the ancient mathematicians called it "multiplying by the different, dividing by the same" (yi cheng tong chu). The commentary uses it to explain the finding of the large and small base and altitude from one another: the base rate and the altitude rate are the given rate and the rate sought, and the known base or the known altitude is the given number; one cannot multiply the base rate by the known base, or the altitude rate by the known altitude.]\end{uniten}

\begin{unitzh}\unitlabel{jiuzhang-042.6}言雖異矣，及其所以成法之實，則同歸矣。則圓徑又可以句乘之差并{\zhsmall 此句亦訛舛，當云：則又可以股弦差減句爲圓徑。}句弦差減股爲圓徑。{\zhsmall 此下有脫文，當補云：句弦差、股弦差并之，以減弦，餘爲圓徑。}\end{unitzh}
\begin{uniten}Although the wording differs, when it comes to the substance of how the method is formed, they come to the same end. Then the diameter of the circle again can, with the base multiplying it, the differences summed [Note: This phrase is also corrupt; it should say: "Then again one can subtract the difference of altitude and hypotenuse from the base to make the diameter of the circle."] subtract the difference of base and hypotenuse from the altitude to make the diameter of the circle. [Note: Below this there is missing text; it should be supplemented thus: "Sum the difference of base and hypotenuse and the difference of altitude and hypotenuse, and subtract this from the hypotenuse; the remainder is the diameter of the circle."]\end{uniten}

\begin{unitzh}\unitlabel{jiuzhang-042.7}□□□□□□□□{\zhsmall 據術意，句股相乘半之爲句股積，有朱、靑、}□□□□□□□□{\zhsmall 黃幂各一，則倍之有朱、靑、黃幂各四。截朱、}□□□□□□□□{\zhsmall 靑幂各成小句股者二，令倒順相補，各成}□□□□□□□□{\zhsmall 小長方。合四朱、四靑、四黃而成大長方，以}□□□□□□□□{\zhsmall 容圓半徑爲廣，并句股弦爲袤。}\end{unitzh}
\begin{uniten}[illegible] [Note: According to the intention of the method, base and altitude multiplied together and halved make the area of the base and altitude, in which there are vermilion, blue-green] [illegible] [and yellow areas, 1 of each; then doubling it, there are vermilion, blue-green and yellow areas, 4 of each. Cut the vermilion] [illegible] [and blue-green areas so that each forms 2 small base-and-altitudes; let them complement one another, inverted and upright, each forming] [illegible] [a small oblong (chang fang). Combine the 4 vermilion, 4 blue-green and 4 yellow to form a large oblong, with] [illegible] [the half-diameter of the contained circle as the width, and the sum of base, altitude and hypotenuse as the length.]\end{uniten}

\begin{unitzh}\unitlabel{jiuzhang-042.8}《九章》之術，乃算術之鼻祖，囊括後賢，胥不能度越範圍焉，猶六經之臨百氏也。《周官》保氏九數，鄭君以《九章》之方田、粟米、衰分、少廣、商功、均輸、方程、嬴不足、旁要釋之，綴曰：「今有重差、夕桀、句股也。」錢曉徵學士以爲夕桀乃互椉之脫誤，良然。蓋《九章》句股篇末有望遠、度高、測深七術，或析之，名曰《九章重差》。互椉卽方程術所謂維椉是也。句股卽旁要，疏所云「今《九章》以句股替旁要」。旁要云者，不必實有是形，可自旁假設以要取之，祖冲之謂之綴術。疏又引馬氏融注「今有重差、夕桀」，馬氏不連及句股者，以句股替旁要，故不重舉。\end{unitzh}
\begin{uniten}The methods of the Nine Chapters are the first ancestor of mathematics; they embrace later worthies, none of whom can step beyond their bounds, just as the Six Classics preside over the Hundred Schools. As to the nine numbers (jiu shu) of the Guardian (Baoshi) in the Zhou guan, Lord Zheng explained them with the Nine Chapters' Rectangular Fields (Fangtian), Millet and Grain (Sumi), Distribution by Proportion (Cuifen), Short Width (Shaoguang), Consultations on Works (Shanggong), Fair Transport (Junshu), Rectangular Arrays (fangcheng), Excess and Deficit (yingbuzu) and Pangyao, and added: "Now there are Double Differences (chongcha), Xijie and Base and Altitude (gougu)." The Academician Qian Xiaozheng held that xijie is a corrupt and defective form of hucheng (mutual multiplication); this is indeed so. For at the end of the Base and Altitude chapter of the Nine Chapters there are 7 methods of sighting the distant, measuring heights and gauging depths; someone separated them and named them Double Differences of the Nine Chapters (Jiuzhang chongcha). Mutual multiplication is precisely what the method of rectangular arrays calls cross multiplication (weicheng). Base and Altitude is Pangyao, as the subcommentary (shu) says: "The present Nine Chapters replaces Pangyao with Base and Altitude." As for what is called pangyao: it is not necessary that this figure really exist; one may posit it hypothetically from the side (pang) and thereby seek (yao) and obtain it. Zu Chongzhi called it zhuishu (the art of connecting). The subcommentary also cites the commentary of Mr. Ma Rong: "Now there are Double Differences and Xijie." That Mr. Ma does not go on to mention Base and Altitude is because Base and Altitude had replaced Pangyao; therefore he does not list it again.\end{uniten}

\begin{unitzh}\unitlabel{jiuzhang-042.9}劉徽序云：「漢張蒼、耿壽昌因舊文遺殘，各稱刪補，故校其目與古或異，而所論多近語。」所謂目與古異者，則句股替旁要是也。至唐王孝通云「校其條目，頗與古術不合」，則妄而敢矣。夫古今豈有異術哉？劉徽因其有望遠諸術，遂造重差，綴於句股之下，卽今《海島算》，引而申之，觸類而長之，事之宜也。舊有圖，今缺，余友休寧東原戴先生補之，今分附諸篇之末，亦猶劉徽之綴重差於句股焉。乾隆癸巳，闕里孔繼涵識於京師壽雲簃之敏事齋。\end{unitzh}
\begin{uniten}Liu Hui's preface says: "Zhang Cang and Geng Shouchang of the Han, on the basis of the remnants of the old text, each claimed to have made deletions and supplements; therefore, when its headings are examined, they sometimes differ from the ancient ones, and what is discussed is mostly in recent language." What is meant by the headings differing from the ancient ones is precisely the replacement of Pangyao by Base and Altitude. As for Wang Xiaotong of the Tang saying, "Examining its items, they considerably disagree with the ancient methods," this is reckless and presumptuous. Now, how could there be different methods in antiquity and the present? Liu Hui, because it had the various methods of sighting the distant, thereupon made the Double Differences and attached it below Base and Altitude; this is the present Haidao suan (Sea Island Computations). To draw it out and extend it, to touch on related kinds and enlarge it, was fitting to the matter. Formerly there were diagrams; now they are missing. My friend Mr. Dai Dongyuan of Xiuning supplied them, and they are now divided and appended at the ends of the various chapters, just as Liu Hui attached the Double Differences to Base and Altitude. Noted by Kong Jihan of Queli in the guisi year of Qianlong, in the Minshi Studio of the Shouyun Annex in the capital.\end{uniten}

\begin{unitzh}\unitlabel{jiuzhang-042.10}《九章算術》卷九終\end{unitzh}
\begin{uniten}End of juan 9 of the Nine Chapters on the Mathematical Art\end{uniten}
\clearpage\section*{Sources}\addcontentsline{toc}{section}{Sources}
The Chinese text is Sibu congkan photo-reprint of the Weiboxie edition (四部叢刊, 景上海涵芬樓藏微波榭刊本), as transcribed on Chinese Wikisource. The text of the work is in the public domain; the Wikisource transcription is credited here with the pages and revisions used. The punctuation, the division into units and the English translation are AI-produced and dedicated to the public domain (CC0 1.0).

\begin{itemize}
\item \href{https://zh.wikisource.org/wiki/\%E4\%B9\%9D\%E7\%AB\%A0\%E7\%AE\%97\%E8\%A1\%93_\%28\%E5\%9B\%9B\%E9\%83\%A8\%E5\%8F\%A2\%E5\%88\%8A\%E6\%9C\%AC\%29/\%E5\%BA\%8F}{序}, revision 2658132
\item \href{https://zh.wikisource.org/wiki/\%E4\%B9\%9D\%E7\%AB\%A0\%E7\%AE\%97\%E8\%A1\%93_\%28\%E5\%9B\%9B\%E9\%83\%A8\%E5\%8F\%A2\%E5\%88\%8A\%E6\%9C\%AC\%29/\%E7\%9B\%AE\%E9\%8C\%B2}{目録}, revision 2658133
\item \href{https://zh.wikisource.org/wiki/\%E4\%B9\%9D\%E7\%AB\%A0\%E7\%AE\%97\%E8\%A1\%93_\%28\%E5\%9B\%9B\%E9\%83\%A8\%E5\%8F\%A2\%E5\%88\%8A\%E6\%9C\%AC\%29/\%E5\%8D\%B7\%E7\%AC\%AC\%E4\%B8\%80}{卷第一}, revision 2658134
\item \href{https://zh.wikisource.org/wiki/\%E4\%B9\%9D\%E7\%AB\%A0\%E7\%AE\%97\%E8\%A1\%93_\%28\%E5\%9B\%9B\%E9\%83\%A8\%E5\%8F\%A2\%E5\%88\%8A\%E6\%9C\%AC\%29/\%E5\%8D\%B7\%E7\%AC\%AC\%E4\%BA\%8C}{卷第二}, revision 2658135
\item \href{https://zh.wikisource.org/wiki/\%E4\%B9\%9D\%E7\%AB\%A0\%E7\%AE\%97\%E8\%A1\%93_\%28\%E5\%9B\%9B\%E9\%83\%A8\%E5\%8F\%A2\%E5\%88\%8A\%E6\%9C\%AC\%29/\%E5\%8D\%B7\%E7\%AC\%AC\%E4\%B8\%89}{卷第三}, revision 2658136
\item \href{https://zh.wikisource.org/wiki/\%E4\%B9\%9D\%E7\%AB\%A0\%E7\%AE\%97\%E8\%A1\%93_\%28\%E5\%9B\%9B\%E9\%83\%A8\%E5\%8F\%A2\%E5\%88\%8A\%E6\%9C\%AC\%29/\%E5\%8D\%B7\%E7\%AC\%AC\%E5\%9B\%9B}{卷第四}, revision 2658137
\item \href{https://zh.wikisource.org/wiki/\%E4\%B9\%9D\%E7\%AB\%A0\%E7\%AE\%97\%E8\%A1\%93_\%28\%E5\%9B\%9B\%E9\%83\%A8\%E5\%8F\%A2\%E5\%88\%8A\%E6\%9C\%AC\%29/\%E5\%8D\%B7\%E7\%AC\%AC\%E4\%BA\%94}{卷第五}, revision 2658138
\item \href{https://zh.wikisource.org/wiki/\%E4\%B9\%9D\%E7\%AB\%A0\%E7\%AE\%97\%E8\%A1\%93_\%28\%E5\%9B\%9B\%E9\%83\%A8\%E5\%8F\%A2\%E5\%88\%8A\%E6\%9C\%AC\%29/\%E5\%8D\%B7\%E7\%AC\%AC\%E5\%85\%AD}{卷第六}, revision 2658139
\item \href{https://zh.wikisource.org/wiki/\%E4\%B9\%9D\%E7\%AB\%A0\%E7\%AE\%97\%E8\%A1\%93_\%28\%E5\%9B\%9B\%E9\%83\%A8\%E5\%8F\%A2\%E5\%88\%8A\%E6\%9C\%AC\%29/\%E5\%8D\%B7\%E7\%AC\%AC\%E4\%B8\%83}{卷第七}, revision 2658140
\item \href{https://zh.wikisource.org/wiki/\%E4\%B9\%9D\%E7\%AB\%A0\%E7\%AE\%97\%E8\%A1\%93_\%28\%E5\%9B\%9B\%E9\%83\%A8\%E5\%8F\%A2\%E5\%88\%8A\%E6\%9C\%AC\%29/\%E5\%8D\%B7\%E7\%AC\%AC\%E5\%85\%AB}{卷第八}, revision 2658141
\item \href{https://zh.wikisource.org/wiki/\%E4\%B9\%9D\%E7\%AB\%A0\%E7\%AE\%97\%E8\%A1\%93_\%28\%E5\%9B\%9B\%E9\%83\%A8\%E5\%8F\%A2\%E5\%88\%8A\%E6\%9C\%AC\%29/\%E5\%8D\%B7\%E7\%AC\%AC\%E4\%B9\%9D}{卷第九}, revision 2658142
\end{itemize}
\end{document}
